\documentclass[11pt]{article}
\usepackage[a4paper,margin=1.45cm]{geometry}
\usepackage{fontspec}
\usepackage{xeCJK}
\setmainfont{Noto Serif}
\setCJKmainfont{Noto Serif CJK JP}
\setCJKsansfont{Noto Sans CJK JP}
\usepackage{amsmath,amssymb,mathtools,array,booktabs,pdflscape,adjustbox,diagbox}
\setlength{\parindent}{1em}
\setlength{\parskip}{0pt}
\newcommand{\widebar}[1]{\overline{#1}}
\newcommand{\A}{\Delta}
\newcommand{\D}{\Delta}
\allowdisplaybreaks
\setlength{\emergencystretch}{12em}
\sloppy
\hbadness=10000
\vbadness=10000
\hfuzz=2pt
\vfuzz=10pt
\raggedbottom
\newcommand{\R}{\varrho}
\newcommand{\hata}{\widehat{a}}
\newcommand{\hatv}{\widehat{v}}
\newcommand{\modu}[1]{\;\operatorname{mod} #1}
\newcommand{\mcell}[1]{\ensuremath{\begin{gathered}#1\end{gathered}}}
\newcommand{\pair}[2]{(#1\,|\,#2)}
\newcommand{\eps}{\varepsilon}
\providecommand{\Hgrp}{\mathfrak{H}}
\providecommand{\GaloisRes}{\Phi}
\providecommand{\pol}[2]{\left(#1\frac{\partial}{\partial #2}\right)}
\providecommand{\Deltaop}{\Delta\!\left(\frac{\partial}{\partial x},\frac{\partial}{\partial y},\frac{\partial}{\partial z}\right)}
\providecommand{\Omegaop}{\Omega}
\makeatletter
\renewcommand\footnotesize{\@setfontsize\footnotesize{7.5}{8.7}\selectfont}
\makeatother
\newcommand{\bdelta}{\bar{\delta}}
\providecommand{\dd}{\mathrm d}
\providecommand{\Div}{\operatorname{Div}}
\providecommand{\G}{\mathfrak G}
\providecommand{\bd}{\bar{\delta}}
\providecommand{\p}{\partial}
\begin{document}

\begin{center}
{\Large\bfseries 1. 三元四次形式の形式系の形成について}\par
\vspace{1.5em}
\emph{Sitzungsberichte der Physikalisch-medizinischen Sozietät in Erlangen} 39 (1907), pp. 176--179
\end{center}

\vspace{3em}
\begin{center}
（著者の学位論文からの抜粋。）
\end{center}

三元四次形式の形式系については、Gordan、Maisano、Pascal の研究\footnote{P. Gordan, ``Über das volle Formensystem der ternären biquadratischen Form'' \(f=x_1^3x_2+x_2^3x_3+x_3^3x_1\) (Math. Annalen, vol. XVII, pp. 217--233, 1880); G. Maisano, 1. ``Sistemi completi dei primi cinque gradi della forma ternaria biquadratica e degl' invarianti, covarianti e contravarianti di sesto grado'' (Giorn. di Battaglini XIX). 2. ``Sui covarianti indipendenti di \(6^\circ\) grado nei coefficienti della forma biquadratica ternaria'' (Rend. Circ. Mat. di Palermo I, 1887); E. Pascal, ``Contributo alla teoria della forma ternaria biquadratica e delle sue varie decomposizioni in fattori'' (Memoria premiata dalla R. Accademia delle scienze fisiche e matematiche di Napoli, 1905).}がある。Gordan は、特殊な自己同型形式
\[
f=x_1^3x_2+x_2^3x_3+x_3^3x_1
\]
に対して、54 個の形成から成る完全な形式系を、二元の場合の形式系について彼が与えたものと同種の原理に基づいて構成した。

Maisano の研究では、一般の四次形式について、五次の位数まで（それを含む）\footnote{ここで「位数」とは係数における次元を意味し、「次数」とは変数における次元を意味する。}の形式が、さらに高位数のいくつかの不変式・共変式・反変式とともに、Gordan が \emph{Mathematische Annalen} 第 I 巻で三元三次形式に用いた方法に従って構成されている。Pascal は Maisano の結果を用い、主として四次形式の因子への分解の問題を扱っている\footnote{Maisano は第二の短いノートで、位数 6 かつ次数 6 の三つの共変式の間の線形従属を証明しようとしている。この完全とはいえない証明に対して、Pascal は、位数 6 の三つの共変式の線形独立性、および位数 9 の三つの不変式の線形独立性を証明したと考えている。しかし、三つの共変式の間にも、また三つの不変式の間にも、実際には線形関係が存在することが私の計算の過程で明らかになった。Pascal の記法では、その明示公式は
\[
0=20\Omega_1+6\Omega_2-3\Omega_3+4fC_1-12fC_2-4A\cdot\Delta,
\]
\[
0=4A^3-15AB+30C-90D+9E.
\]
である。}。

\textbf{私の研究の目的は、一般三元四次形式の形式系を構成することである。さしあたり、主要な基礎だけを与え、いわゆる「相対的完全系」\footnote{用語については Gordan--Kerschensteiner, \emph{Vorlesungen über Invariantentheorie}, vol. II, p. 227 を参照。}を確立する。}

三元形式の系を作る基本的な考え方は、二元の場合と同じである。最初の相対的完全系、すなわち二元の場合から移された形式の系から出発し、一定の法則に従って、法が次第に高くなる系へと進む。その過程は、ある法を基本形式と見なしたときにその法の系が有限かつ既知になるか、またはある法が不変式を因子にもつ形式に還元できるようになるまで続けられる。第一の場合には、相対的完全系とその法の系とのトランスヴェクションによって絶対的完全系が得られ、第二の場合には相対的完全系と絶対的完全系とが一致する。Hilbert による形式系の有限性の一般的証明により、この手続きは必ず終わらなければならない。

本場合では、第一の相対的完全系の法 \((abc)\) は、法
\[
\Delta=(abc)^2a_x^2b_x^2c_x^2\quad\text{および}\quad \nu=(abu)^4
\]
に帰着できる。ここから、法 \(\nu\) に関する相対的完全系は容易に得られる。そこで法の列として次の形式を選ぶ：
\[
\nu;\qquad \nu^{(\nu)}=(\nu\nu_1x)^4=s_x^4,
\]
\[
\nu^{(s)}=(ss'u)^4\ldots,
\]
さらに、法 \(s\) に関する相対的完全系の形成に現れる二つの二次形式、すなわち \(u_\sigma^2\) と \(t_x^2\) を、追加の法として付け加える。すると、\(s\) に続く法、すなわち \((ss'u)^4\) は、法 \((\varrho,t)\) に還元可能であることを示すことができる。しかし二つの二次形式の同時系は有限かつ既知であるから、上に述べた終結点に到達したことになる。法 \((\varrho,t)\) に関する相対的完全系として、331 個の形成が得られる。

用いた方法について若干の詳細を加える。

形式を作る基本過程である折り畳み過程は、三元の場合には次のように定義できる。

記号的積
\[
s_x^m t_x^n u_\sigma^\mu u_\tau^\nu
\]
において、因子の対
\[
\begin{array}{c|c|c|c|c}
 & s_xt_x & u_\sigma u_\tau & s_xu_\sigma\ \text{または}\ s_xu_\tau & t_xu_\tau\ \text{または}\ t_xu_\sigma\\
\text{それぞれを} & (stu) & (\sigma\tau x) & s_\sigma\ \text{または}\ s_\tau & t_\tau\ \text{または}\ t_\sigma\\
\text{で置き換える折り畳み} & \mathrm{I} & \mathrm{II} & \mathrm{III} & \mathrm{IV},
\end{array}
\]
を行うなら、そのようにして生じる形式は、もとの形式から折り畳みによって得られたものである\footnote{Gordan, Math. Annalen, vol. XVII, p. 219.}。

ここではなお常に \(s_\sigma=0;\ t_\tau=0\) とおくことができる。個々の折り畳み相互の関係については、次の定理が成り立つ。

\textbf{折り畳み I と II は基本折り畳みであり、合成の順序に依存せず、折り畳み III と IV はこれらから合成できる。換言すれば、与えられた式から折り畳みによって生じるすべての形式を作るには、折り畳み I と II だけを適用すればよい}\footnote{この定理には、\(n\) と \(\mu\)、または \(m\) と \(\nu\) がそれぞれ同時に零になる場合に例外がある。このとき「形式列」は対角項に縮退する。}。

「形式列」とは、初期形式と、それを自己の中へ折り畳むことによって生じるすべての形式とを合わせたものをいう。上の定理により、形式列
\[
s_x^m t_x^n u_\sigma^\mu u_\tau^\nu
\]
は長方形の図式に配列することができる。この図式で進んでいくと、次が得られる：
\begin{enumerate}
\item 一列右へ進むことにより、隣接する形式から折り畳み I によって生じるすべての形式。
\item 一行下へ進むことにより、その上に立つ形式から折り畳み II によって生じるすべての形式、したがって折り畳みによって生じるすべての形式。
\end{enumerate}

これらの仮定のもとで、還元子に関する定理は、
\begin{center}
「還元子とは可約な形式列である」
\end{center}
という定義のもと、最も一般的な形で次のように述べられる。

\begin{center}
\textbf{形式列の初期形式が可約であるのは、その一つの項が還元子との折り畳みによって生じているためであり、かつその形式列の終形式がこの同じ項から折り畳みによって生じているならば、形式列全体は可約である。}
\end{center}

さらに挙げるべき還元方法は次の二つである：
\begin{enumerate}
\item いわゆる「二重還元」、すなわち高次の形式の間の関係を得るために、一つの形式を二通りの方法で還元すること。
\item 「分解した形式による折り畳み」。これは一部では還元子による還元の性格をもち、一部では「二重還元」の性格をもつ。
\end{enumerate}

\clearpage

\begin{center}
{\Large\bfseries 2. 三元四次形式の形式系の形成について}\par
\vspace{1.5em}
Journal f. d. reine u. angew. Math. 134 (1908), pp. 23--90 および二つの表
\end{center}

\vspace{3em}

\begin{center}
\textbf{序論。}
\end{center}

三元四次形式の形式系については、Gordan、Maisano、Pascal の研究\footnote{序論および第 I 章からの短い抜粋は、\emph{Sitzungsberichte der phys.-med. Sozietät Erlangen 1907}, pp. 176--179 に発表された。}\footnote{P. Gordan, 三元四次形式の完全な形式系について
\[
f=x_1^3x_2+x_2^3x_3+x_3^3x_1.
\]
(\emph{Math. Annalen}, vol. XVII (1880), pp. 217--233.)
G. Maisano, 1) \emph{Sistemi completi dei primi cinque gradi della forma ternaria biquadratica e degl' invarianti, covarianti e contravarianti di sesto grado}. (Giorn. di Battaglini XIX.)
2) \emph{Sui covarianti indipendenti di $6^\circ$ grado nei coefficienti della forma biquadratica ternaria}. (Rend. Circ. Mat. di Palermo I, 1887.)
E. Pascal, \emph{Contributo alla teoria della forma ternaria biquadratica e delle sue varie decomposizioni in fattori}. (Memoria premiata dalla R. Accademia delle scienze fisiche e matematiche di Napoli. 1905.)}がある。Gordan 氏は、特殊な自己同型形式
\[
f=x_1^3x_2+x_2^3x_3+x_3^3x_1
\]
について、54 個の形成から成る完全な形式系を、二元領域の形式系について彼が与えたものと同種の原理に基づいて構成している。

Maisano 氏においては、一般の四次形式について、五次の位数まで（それを含む）\footnote{ここで「位数」とは係数における次元を、「次数」とは変数における次元を意味する。}の形式が、さらに高位数のいくつかの不変式、共変式、反変式とともに、Gordan 氏が \emph{Math. Annalen} 第 I 巻で三元三次形式に用いた方法に従って構成されている。Pascal 氏は Maisano の結果を用い、主として四次形式の因子への分解の問題を扱っている\footnote{第二の短いノートで、Maisano 氏は位数 6 かつ次数 6 の三つの共変式の線形従属を証明しようとしている。この完全とはいえない証明に対して、Pascal 氏は、位数 6 の三つの共変式の線形独立性、さらに位数 9 の三つの不変式の線形独立性を証明したと考えている。しかし、三つの共変式の間にも、また三つの不変式の間にも、実際には線形関係が存在することは、本論文の式 (13), § 11 および § 17 の注の式 (22) が示しており、そこではこれらの関係が明示的に与えられている。Pascal 氏からの連絡によれば、この矛盾は、彼の反変式 $p$（ここでは $\varrho$ と呼ぶ）の表式、同論文 p. 46 における数値的誤りによって説明される。}。

\emph{以下の研究の目的は、一般三元四次形式の形式系を構成することである。すなわち本論文では主要な基礎だけを与え、いわゆる「相対的完全系」\footnote{この名称については § 3a を参照。}を構成する。}

本論文は Gordan の研究と密接に結びついている。ただし、そこではごく一般的にしか与えられていなかった原理を、まず個々に展開する必要があった。折り畳み相互の関係に関する一つの定理を用い、さらに「形式列」（§ 1）を導入することによって、還元定理を明確に定式化し補完する（§ 3）。一方、特別な級数展開を再帰的に構成すること（§ 2）が、還元を実際に遂行するための計算手段を与える。

三元形式の系を形成する基本思想は、二元領域の場合と同じである。最初の相対的完全系、すなわち二元の場合から移された形式の系から出発して、一定の法則に従い、法がしだいに高くなる系へ進む。この過程は、ある法の系が、その法を基本形式としてとったときに、有限かつ既知になるまで、またはある法が不変式を因子にもつ形式へ還元できるようになるまで続く。第一の場合には、相対的完全系とその法の系とのトランスヴェクションによって絶対的完全系が生じ、第二の場合には相対的完全系と絶対的完全系とが一致する。Hilbert 氏によって一般に証明された形式系の有限性により、この手続きは必ず終結しなければならない。

本場合では、第一の相対的完全系の法 $(abc)$（§ 4）は法
\[
\Delta=(abc)^2a_x^2b_x^2c_x^2\quad\text{および}\quad \nu=(abu)^4
\]
（§ 5）に帰着され、ついで法 $\nu$ に関する相対的完全系が求められる（§ 6）。これらの結果は、一般四次形式については Gordan の研究にすでに与えられている。しかし、ここでの導出方法は、より高い次数の形式へいっそう容易に移すことができる。

法の列として、いま次の形式を選ぶ（§ 7）：
\[
\nu,\qquad \nu(\nu)=(\nu\nu_1x)^4=s_x^4,\qquad \nu(s)=(ss'u)^4,\ldots .
\]

法 $s$ に関する相対的完全系を形成する際には、系の中に現れる二つの二次形式 $u_\varrho^2$ と $t_x^2$ を法として付け加え（§§ 9, 10）、それに応じて法 $(s,\varrho,t)$ に関する相対的完全系を作る。ついで、法 $(ss'u)^4$ が法 $\varrho$ と $t$ に還元可能であることを示す（§ 17）。これによって、次の高い系は法 $(\varrho,t)$ に関する相対的完全系へ移る。ところが二つの二次形式の同時系は有限かつ既知であり、一般の場合には 20 個の形成から成る\footnote{Clebsch--Lindemann, \emph{Vorlesungen über Geometrie}. (vol. I, p. 288 以下) 二つの共変的形式の系。Gordan, \emph{Über Büschel von Kegelschnitten}. (\emph{Math. Ann.} vol. XIX, p. 530.) 二つの反変的形式の系。}ので、法 $(\varrho,t)$ に関する相対的完全系の構成によって上に述べた終結点に到達する。この系を $(\varrho,t)$ の系とトランスヴェクションして絶対的完全系を作ることは、後に残しておく。

法 $(\varrho,t)$ に関する相対的完全系として、331 個の形成が得られ、それらは付された表において、変数 $x$ と $u$ における次数に従って配列されている。

最後に、導入した記号の概観を与える：
\begin{align*}
f&=a_x^4,\\
\theta&=\theta_x^4u_\vartheta^2=(abu)^2a_x^2b_x^2,\\
K&=K_x^6u_k^3=(a\theta u)u_\vartheta^2a_x^3\theta_x^3,\\
j&=u_j^6=(a\theta u)^4u_\vartheta^2,\\
\Delta&=\Delta_x^6=a_\vartheta^2a_x^2\theta_x^4=(abc)^2a_x^2b_x^2c_x^2,\\
N&=N_x^8u_\eta=(a\Delta u)a_x^3\Delta_x^5,\\
\nu&=u_\nu^4=(abu)^4,\\
H&=H_x^2u_\eta^4=(\nu\nu_1x)^2u_\nu^2u_{\nu_1}^2,\\
L&=L_x^3u_l^6=(\nu\eta x)H_x^2u_\nu^3u_\eta^3,\\
g&=g_x^6=(\nu\eta x)^4H_x^2,\\
\sigma&=u_\sigma^6=H_\nu^2u_\nu^2u_\eta^4,\\
s&=s_x^4=(\nu\nu_1x)^4,\\
Z&=Z_x^4u_\zeta^2=(ss'u)^2s_x^2s_x'^2,\\
\varrho&=u_\varrho^2=\theta_\nu^4u_\vartheta^2,\\
t&=t_x^2=a_\eta^4H_x^2,\\
i&=a_\nu^4,\\
J&=s_\nu^4.
\end{align*}

\begin{center}
{\Large\bfseries 第 I 章。\quad 三元形式に関する一般定理。}
\end{center}

\begin{center}
\textbf{§ 1. \quad 折り畳み過程。\quad 形式列。}
\end{center}

形式を作る基本過程は折り畳み過程であり、三元領域では次のように定義できる。記号的積
\[
s_x^m t_x^n u_\sigma^\mu u_\tau^\nu .
\]
が与えられているとする。因子の対
\[
s_xt_x\qquad u_\sigma u_\tau\qquad s_xu_\sigma\ \text{または}\ s_xu_\tau
\qquad t_xu_\tau\ \text{または}\ t_xu_\sigma
\]
をそれぞれ
\[
(stu)\qquad (\sigma\tau x)\qquad s_\sigma\ \text{または}\ s_\tau
\qquad t_\tau\ \text{または}\ t_\sigma,
\]
すなわち折り畳み
\[
\text{I}\qquad\text{II}\qquad\text{III}\qquad\text{IV},
\]
によって置き換えると、得られる形式は元の形式から折り畳みによって生じたものである\footnote{Gordan, \emph{Math. Annalen} vol. XVII, p. 219.}。

このとき式 $s_x^m t_x^n u_\sigma^\mu u_\tau^\nu$ は、二つの形式 $S\cdot T$ の実際の積として見ることもできるし、複数の記号列をもつ一つの形式
\[
s_x^m t_x^n u_\sigma^\mu u_\tau^\nu=A_x^{m+n}u_x^{\mu+\nu}
\]
として見ることもできる。第一の場合には形式 $S$ と $T$ との折り畳みといい、第二の場合には「形式 $A$ の自己への折り畳み」という。簡潔のため、以下では常に（実際、後で考察されるすべての形式についてそうである）
\[
\tag{a}
s_\sigma^\lambda=0,\qquad t_\tau^\lambda=0
\]
が、零でない任意の $\lambda$ と任意の $x$ について、また任意の他の折り畳みと合わせて成り立つものと仮定する。これは、形式 $s_x^m t_y^n u_\sigma^\mu u_\tau^\nu$ がいわゆる「正規形式」である、すなわち $\Omega$ 過程を変数 $x,u$ についても、また変数 $y,v$ についても適用すると消える、ということを意味する。

したがって条件 (a) は何らの制限でもなく、常に達成できる\footnote{Gordan, \emph{Math. Annalen} vol. V, p. 104 を参照。}。

個々の折り畳みの関係については次の定理が成り立つ：

\emph{定理 I. 折り畳み I と II は基本折り畳みであり、合成の順序によらず、折り畳み III と IV はこれらから合成される。換言すれば、与えられた式から折り畳みによって生じるすべての形式を作るには、折り畳み I と II だけを適用すればよい。}

証明：同一性定理、または行列に対する積定理により、(a) を考慮すれば、
\[
(\widehat{st}\sigma x)=-t_\sigma,\qquad
(\widehat{\sigma\tau}su)=-s_\tau,
\]
\[
(\widehat{st}\tau x)=s_\tau,\qquad
(\widehat{\sigma\tau}tu)=t_\sigma,
\]
\[
(stu)(\sigma\tau x)=s_\tau+t_\sigma-u_x\,s_\tau t_\sigma .
\]
を得る。したがって折り畳み I と II を合成することで折り畳み III と IV が生じる。これは、折り畳み II を折り畳み I、すなわち因子 $(stu)u_\sigma u_\tau$ に適用する場合にも、折り畳み I を折り畳み II、すなわち因子 $(\sigma\tau x)s_xt_x$ に適用する場合にも成り立つ。

\emph{形式列}\footnote{Clebsch, \emph{Über eine Fundamentalaufgabe der Invariantentheorie} (Abh. der Gött. Ges. d. Wiss. vol. XVII): § 17 および § 18 の末尾を参照。「形式列」は、そこに導入された「本来的に還元された同値系」とは、より高い形式による順序づけの原理によって異なる。}とは、初期形式と、それを自己の中へ折り畳むことによって生じるすべての形式とを合わせたものをいう。また形式列をその初期形式によって表す。ここで「より高い形式」とは、自己への折り畳みをより多く含む任意の形式を意味する。ただし、同等の折り畳み $s_\tau$ と $t_\sigma$ の間では、一方をより高いものとして正規化しなければならない。形式列の形成法則から次が従う：

1) 形式列の各形式は初期形式の記号に関して線形である。

2) 形式列は、それぞれ二つの反変的記号列をもつ初期形式については一意に定まる。より多くの記号列をもつ初期形式については、同一性定理によって結びつけられる折り畳みのうちから特別なものを区別することによって、一意に正規化されなければならない。

定理 I により、形式列 $s_x^m t_x^n u_\sigma^\mu u_\tau^\nu$（$m>n;\ \mu>\nu$）は次の長方形の図式に従って配置できる：
\begingroup\small
\[
\begin{array}{c|c|c|c}
s_x^m t_x^n u_\sigma^\mu u_\tau^\nu
& (stu)
& (stu)^2
& \cdots\\[0.3em]
(\sigma\tau x)
& t_\sigma,\ s_\tau
& t_\sigma(stu),\ s_\tau(stu)
& \cdots\\[0.3em]
(\sigma\tau x)^2
& t_\sigma(\sigma\tau x),\ s_\tau(\sigma\tau x)
& t_\sigma^2,\ t_\sigma s_\tau,\ s_\tau^2
& \cdots\\[0.3em]
\vdots&\vdots&\vdots&\ddots\\[0.3em]
(\sigma\tau x)^\nu
& t_\sigma(\sigma\tau x)^{\nu-1},\ s_\tau(\sigma\tau x)^{\nu-1}
& t_\sigma^2(\sigma\tau x)^{\nu-2},\ t_\sigma s_\tau(\sigma\tau x)^{\nu-2},\ s_\tau^2(\sigma\tau x)^{\nu-2}
& \cdots\\[0.3em]
\vdots
& t_\sigma(\sigma\tau x)^\nu
& t_\sigma^2(\sigma\tau x)^{\nu-1},\ t_\sigma s_\tau(\sigma\tau x)^{\nu-1}
& \cdots
\end{array}
\]
\endgroup
および、対応して継続される下部
\footnote{ここで、および以下でも同様に、星印で示された下部は、図式の右上に同じ印で示された場所に置かれるべきである。}
\begingroup\small
\[
\begin{array}{c|c|c}
(stu)^n & \cdots & \cdots\\[0.3em]
t_\sigma(stu)^{n-1},\ s_\tau(stu)^{n-1}
& s_\tau(stu)^n & \cdots\\[0.3em]
t_\sigma^2(stu)^{n-2},\ t_\sigma s_\tau(stu)^{n-2},\ s_\tau^2(stu)^{n-2}
& t_\sigma s_\tau(stu)^{n-1},\ s_\tau^2(stu)^{n-1}
& s_\tau^2(stu)^n .
\end{array}
\]
\endgroup

ここで、

1) 一列右へ進むと、隣接する形式から折り畳み I によって生じるすべての形式が得られ、

2) 一行下へ進むと、その上にある形式から折り畳み II によって生じるすべての形式が得られ、

したがって折り畳みによって生じるすべての形式が得られる。

全形式列の任意の形式 $t_\sigma^\kappa s_\tau^\lambda(stu)^\mu$ または $t_\sigma^\kappa s_\tau^\lambda(\sigma\tau x)^\nu$ は新しい形式列の始まりとなる。その形式列は、長方形の図式のうち $t_\sigma^\kappa s_\tau^\lambda(stu)^\mu$ または $t_\sigma^\kappa s_\tau^\lambda(\sigma\tau x)^\nu$ を左上隅とし、因子 $t_\sigma^\kappa s_\tau^\lambda$ をもつすべての形式から成る。

補足。定理 I には例外がある。数 $n$ と $\mu$（それぞれ $m$ と $\nu$）が零となり、折り畳み I, II, IV がもはや存在しない一方で、折り畳み III（それぞれ IV）はなお存在する場合である。この場合には、図式の対角項を形式列と呼ぶ：
\[
\begin{array}{ccccc}
s_x^m u_\tau^\nu &&&& t_x^n u_\sigma^\mu\\
& s_\tau && t_\sigma &\\
& s_\tau^2 && t_\sigma^2 &\\
&&\ddots&&\\
&&s_\tau^\nu && t_\sigma^\mu .
\end{array}
\]
折り畳み I と II が存在しないことに応じて、この場合には形式列の極形式による級数展開は常に初期項に縮退する。しかし、還元子に関する定理は有効に残る。

\begin{center}
\textbf{§ 2. \quad 形式列の極形式による級数展開。}
\end{center}

$n$ 個の変数をもつ形式については、二種類の級数展開が明示的に立てられている\footnote{Gordan, \emph{Über Kombinanten}, §§ 2, 5 (\emph{Math. Ann.} vol. V) を参照。}：

1) 反変的変数の列を一つずつもつ形式 $s_x^m u_\tau^\nu$ について、$u_x$ の冪によって進む級数で、その係数は「正規形式」である。

2) 共変的変数の二列をもつ形式 $s_x^m t_x^n$ について、基本共変式の極形式（形式列 $s_x^m t_x^n$ の極形式）による展開であり、三元の場合には次のように書かれる\footnote{Study, \emph{Methoden zur Theorie der ternären Formen} (Teubner 1889) II. § 7 を参照。そこに与えられた導出と異なり、われわれの導出は数値係数 $C_{pr\varrho}$ を迅速に計算する方法を与える。特殊化 a), a') のもとでの Clebsch--Study 展開は
\[
(stu)^\mu u_\tau^{\nu-\kappa}v_\tau^\kappa s_x^m t_x^{\,n-\lambda}(tuv)^\lambda
=\sum_{p,r,\varrho} C_{pr\varrho}
\bigl[s_\tau^\varrho(stu)^{\mu+r-\varrho}\bigr]_{\widehat{uv}^{\lambda-r+\varrho}v^{\nu-\varrho},\,u^p,\,(vw=\widehat{xw})},
\]
となり、したがって、$u_x$ を因子にもつすべての項を省くべきことが分かっている場合に、われわれの展開で計算の途中ですでに生じる簡略化を許さない。}：
\[
\tag{I}
s_x^m t_x^{n-\lambda}t_y^\lambda
=\sum_{\varrho}
\frac{\binom{m}{\varrho}\binom{\lambda}{\varrho}}
{\binom{m+n-\varrho+1}{\varrho}}
\bigl[(st\widehat{xy})^\varrho s_x^{m-\varrho}t_x^{n-\varrho}\bigr]_{y^{\lambda-\varrho}}.
\]

この二つの展開を組み合わせ、反変的変数列を二つずつもつ形式
\[
s_x^m t_x^{n-\lambda}t_y^\lambda u_\sigma^\mu u_\tau^{\nu-\kappa}v_\tau^\kappa
\]
について、$u_x$ の冪に従う展開を作る。その係数は、形式列 $s_x^m t_x^n u_\sigma^\mu u_\tau^\nu$ の、変数 $x$ と $u$ に関して取られた合成極形式である。

記号（または変数）の二つの反変的列の場合に級数を立てれば十分である。というのも、記号（変数）の列を二つずつ一つの新しい列にまとめることによって、より多くの記号列（変数列）をもつ形式はこの場合に帰着されるからである。

形式列のすべての形式が記号列、または変数列を二つだけ含むようにするため、次の特殊化を行う：
\[
\begin{array}{ll}
\text{a) } \sigma=\widehat{st}, & \text{a') } y=\widehat{uv},\\[0.3em]
\text{b) } s=\widehat{\sigma\tau}, & \text{b') } v=\widehat{xy},
\end{array}
\]
そして特殊化 a), a') について展開を行う。

展開 I を直接転送すると、形式 $(stu)^\mu s_x^m t_x^{n-\lambda}t_y^\lambda$ については合成極形式が得られる。一方、形式 $(stu)^\mu u_\tau^{\nu-\kappa}v_\tau^\kappa s_x^m t_x^{n-\lambda}t_y^\lambda$ については、合成極形式の代わりに、評価のために次々と新しい補助変数を導入し、最後にだけ同一視しなければならない項が現れる。このため計算が不必要に複雑になる。

そこで間接的な道をとる。形式列のすべての形式の合成極形式、すなわち式
\[
\bigl[s_\tau^\varrho(stu)^{\mu-\varrho+r}\bigr]_{y^\lambda}^{v^\kappa}
\]
を、これらの極形式の初期項によって表し、その結果得られる再帰的な方程式系を反転することによって、求める極形式による展開を得る。

転送原理によれば、形式 $s_x^m t_x^n:[s_x^m t_x^n]_{y^\lambda}$ の $\lambda$ 番目の極形式（$\lambda\le n$）について次の展開が成り立つ（$\lambda>n$ の場合は $[s_y^m t_y^n]_{x^{m+n-\lambda}}$ の展開によって第一の場合に帰着される）\footnote{Gordan, \emph{Die Resultante binärer Formen} (chap. I, § 5). Rend. Circ. Mat. di Palermo XXII. 1906.}：
\[
\bigl[s_x^m t_x^n\bigr]_{y^\lambda}
=\sum_r (-1)^r\,
\frac{\binom{m}{r}\binom{\lambda}{r}}{\binom{m+n}{r}}\,
(st\widehat{xy})^r s_x^{m-r}t_x^{n-\lambda}t_y^{\lambda-r},
\]
同様に、$u_\sigma^\mu u_\tau^\nu:[u_\sigma^\mu u_\tau^\nu]_{v^\kappa}$ の $\kappa$ 番目の極形式について、
\[
\bigl[u_\sigma^\mu u_\tau^\nu\bigr]_{v^\kappa}
=\sum_\varrho (-1)^\varrho\,
\frac{\binom{\mu}{\varrho}\binom{\kappa}{\varrho}}{\binom{\mu+\nu}{\varrho}}\,
(\sigma\tau\widehat{uv})^\varrho u_\sigma^{\mu-\varrho}u_\tau^{\nu-\kappa}v_\tau^{\kappa-\varrho}.
\]

これらを掛け合わせ、特殊化 a) と a') を導入すると、
\[
\begin{aligned}
\bigl[s_x^m t_x^n(stu)^\mu u_\tau^\nu\bigr]_{\widehat{uv}^{\lambda}}^{v^\kappa}
={}&
\sum_{r,\varrho}c_{r\varrho}\,
(st\widehat{uv})^r(st\widehat{uv})^\varrho
s_x^{m-r}t_x^{n-\lambda}(tuv)^{\lambda-r}(stu)^{\mu-\varrho}
u_\tau^{\nu-\kappa}v_\tau^{\kappa-\varrho},
\end{aligned}
\]
したがって、$u_x$ の冪に従って展開し、第一項を区別し（$\sum'_{r,\varrho}$ は $r=0$, $\varrho=0$ の項を省くことを意味する）、(a) p. 26 を考慮すると、
\[
\tag{II}
\begin{aligned}
&s_x^m t_x^{n-\lambda}(tuv)^\lambda(stu)^\mu u_\tau^{\nu-\kappa}v_\tau^\kappa\\
&\quad =
\bigl[s_x^m t_x^n(stu)^\mu u_\tau^\nu\bigr]_{\widehat{uv}^{\lambda}}^{v^\kappa}\\
&\qquad
-\sum_{r,\varrho}' c_{r\varrho}\,s_\tau^\varrho(stu)^{\mu-\varrho+r}
u_\tau^{\nu-\kappa}v_\tau^{\kappa-\varrho}
s_x^{m-r}t_x^{n-\lambda}(tuv)^{\lambda-r+\varrho}v_x^r\\
&\qquad
+u_x\sum_{\varrho}\sum_{r=1}^{\lambda} c_{r\varrho}\,
s_\tau^\varrho(stu)^{\mu-\varrho+r-1}(stv)
u_\tau^{\nu-\kappa}v_\tau^{\kappa-\varrho}
s_x^{m-r}t_x^{n-\lambda}(tuv)^{\lambda-r+\varrho}v_x^{r-1}\\
&\qquad
\vdots\\
&\qquad
-(-1)^\lambda u_x^\lambda\sum_\varrho c_{\lambda\varrho}\,
s_\tau^\varrho(stu)^{\mu-\varrho}(stv)^\lambda
u_\tau^{\nu-\kappa}v_\tau^{\kappa-\varrho}
s_x^{m-\lambda}t_x^{n-\lambda}(tuv)^\varrho,
\end{aligned}
\]
ただし
\[
c_{r\varrho}=(-1)^{r+\varrho}\cdot
\frac{\binom{m}{r}\binom{\lambda}{r}}{\binom{m+n}{r}}\cdot
\frac{\binom{\mu}{\varrho}\binom{\kappa}{\varrho}}{\binom{\mu+\nu}{\varrho}}
\qquad
\begin{array}{l}
\text{ただし }\lambda\le n,\\
\kappa\le\nu.
\end{array}
\]
他の数の組合せ
\[\lambda,n;\kappa,\nu\]
についても同様である。式
\[s_x^m t_x^{n-\lambda}(tuv)^\lambda(stu)^\mu u_\tau^{\nu-\kappa}v_\tau^\kappa\]
は、したがって合成極形式へ帰着され、さらに恒等式
\[
(stv)u_\tau=(stu)v_\tau-s_\tau(tuv),
\]
を用いると、同様に作られた式の和へ帰着される。それらは形式列のより高い形式に対応する。

反復によって次の展開に至る：
\[
\tag{III}
s_x^m t_x^{n-\lambda}(tuv)^\lambda(stu)^\mu u_\tau^{\nu-\kappa}v_\tau^\kappa
=
\sum_{p,r,\varrho}u_x^p\,C_{pr\varrho}\cdot
\bigl[s_\tau^\varrho(stu)^{\mu-\varrho+r}\bigr]_{\widehat{uv}^{\lambda-r+\varrho}v^{\kappa-\varrho+p},\,v_x^{r-p}}.
\]
数値係数 $C_{pr\varrho}$ は、(II.) を反復して得られる方程式系の係数 $c_{r\varrho}$, $c'_{r\varrho}$ から一意かつ再帰的に計算される（例 p. 42 を参照）。他の特殊化についても同様の展開が得られる。

\begin{center}
\textbf{§ 3. \quad 還元定理。}
\end{center}

形式および形式系の還元に関する既知の定理を、いま §§ 1, 2 と関連づける。そのために、二元形式の理論から取ったいくつかの定義が必要である。

a) あらかじめ与えられた形式列を法として、任意のこれらの形式の積を折り畳むことによって生じるすべての形成が、系の形式の整有理関数として表され、ただし系の形式を法の系と折り畳んで生じる式から成る加法項を除いてよい、という性質をもつ形式系を、\emph{相対的完全系}と呼ぶ\footnote{Gordan--Kerschensteiner, \emph{Vorlesungen über Invariantentheorie}. vol. II, p. 227.}。

b) 形式が、不変式を因子にもつ形式、または「より高い形式」によって表されるとき、その形式を\emph{可約}と呼ぶ。すなわち、その形式が属する全形式列のより高い形式、または法の記号をより高い位数で含む形式によって表される場合である。

還元子に関する定理\footnote{Gordan, \emph{Math. Annalen} vol. XVII, p. 222.}は、最も一般的な形で次のように述べられる：

定義：\emph{還元子とは、可約な形式列である。}

\emph{定理 II. 形式列の初期形式が可約であるのは、その一つの項が還元子との折り畳みによって生じている（還元子を因子にもつ）ためであり、さらに形式列の終形式がこの同じ項から折り畳みによって生じているならば、形式列全体は可約である}\footnote{定理 II による簡略化は次の例で見られる。特殊形式 $f=x_1^3x_2+x_2^3x_3+x_3^3x_1$ に対して、$a_y^2a_x^2u_y^2$ は還元子である。したがって定理 II により、形式列
\[
\theta_\nu(\vartheta\nu x)\theta_x^3u_\vartheta u_\nu^2
=
a_\nu^2(abu)-a_\nu b_\nu(aba),
\]
の還元が従う。なぜなら $\theta_\nu^4=a_\nu^2b_\nu^2(abu)^2$ は項 $a_\nu^2(abu)$ から折り畳みによって生じているからである。Gordan, loc. cit. p. 231 では、これに対して形式
\[
\theta_\nu(\vartheta\nu x),\quad \theta_\nu(\vartheta\nu x)^2,\quad
\theta_\nu^2,\quad \theta_\nu^2(\vartheta\nu x),\quad
\theta_\nu^2(\vartheta\nu x)^2
\]
の還元が個別に行われている。}。

証明：形式列は初期項から自己への折り畳みによって生じる。したがって一方では還元子とのより高い折り畳みによって、他方では還元子の自己への折り畳みによって生じる。いずれの場合にも、§ 2 により、形式列のすべての形式は還元子の形式列上の極形式（トランスヴェクション）として表せる。

初期形式から全形式列へ結論することは、他の還元方法ではもはや成り立たない。それらは次のものである。

1) いわゆる「二重還元」。

これは、互いに独立な二つの還元子を因子にもつ式を二通りに還元することを意味し、それによって、還元先となったより高い形式の間に関係が生じる。

「二重還元」を体系的に適用することによって、一方ではすべての関係\footnote{例えば「二重還元」によって、形式 $a_\eta^4$ (§ 10) の還元が得られた。これは Maisano が不要に系に入れていた唯一の形式である。さらに、序論の注で述べた三つの共変式および三つの不変式の間の関係も得られた。}が得られなければならない。他方、さまざまな式に「二重還元」を適用しても新しい関係が得られない場合には、より高い形式の独立性の判定基準と計算の検算とを同時に得る。

2) \emph{分解可能形式による折り畳み。}

「分解可能形式」とは、通常の概念をわずかに一般化して、低い次数の形式の積および「より高い」形式によって表される形式を意味する。形式の積は、一回の折り畳みによって積に移る。同様に、非線形形式の積も、特殊化 a'), b') (§ 2) における二回の折り畳みによって、積定理
\[
a_yb_y a_xb_x=\frac12a_y^2b_x^2+\frac12b_y^2a_x^2-\frac12(ab\widehat{yx})^2
\]
に従って積に移る。
したがって、分解可能形式の第一極形式（トランスヴェクション）およびある第二極形式は、ふたたび分解可能形式である。よって次が従う。

分解可能形式との一回（または二回）の折り畳みによって生じた、分解可能形式上の一回（または二回）のトランスヴェクションの一項は、それ自身分解可能形式となる。

a) 分解可能形式の形式列における次に高い形式が分解するか可約になる場合、または一回の折り畳みで特殊化 $y=\widehat{uv}$ または $v=\widehat{xy}$ が現れる場合。

b) 他の一回の折り畳みがより高い形式へ導く場合（§ 12, B. $H_k$ および $H_k(k\eta x)$ 参照）。

そのようなトランスヴェクションの一項は、さらに次の場合にも分解し得る。

c) 他の一回の折り畳みが、分解可能形式との折り畳みとは独立に分解する低い形式へ導く場合、すなわち積および還元すべき形式によって表せる場合である。ただしその都度、該当する項の数値係数が恒等的に零でないことを計算によって確認しなければならない。これは低い形式の「二重還元」にほかならない（§ 23, C. $H_q(\theta H u)(\theta s u)$ 参照）。

分解可能形式は、関数行列式とのすべての一回の折り畳みによって生じる\footnote{Gordan, loc. cit. § 3 を参照。}ほか、ある高次の折り畳みによっても生じる。すなわち
\[
(st\widehat{yx})s_y=\frac12\,t\,s_y^2-\frac12\,s\,t_y^2+\frac12(st\widehat{yx})^2,
\]
\[
(st\widehat{yx})^2s_y^2=\frac13\,t\,s_y^3-\frac13\,s\,t_y^3+s_y(st\widehat{yx})^2-\frac13(st\widehat{yx})^3.
\]
双対的な形式
\[
(\sigma\tau\widehat{\nu u})v_\sigma
\quad\text{および}\quad
(\sigma\tau\widehat{\nu u})v_\sigma^2
\]
についても同様の公式が成り立つ。

この節の定理から、$S$ と $T$ の折り畳みから生じるすべての既約形式を構成するため、次の規則が得られる。

最も低い折り畳みから始め、あらかじめ定めた任意の道筋（例えば行ごと、または対角項に対して対称的）に従って、形式列 $ST$ の形成を進め、還元子を因子にもつ形式に到達するまで行う。その形式によって定められる形式列は、終形式が定理 II の条件を満たすならば、ただちに省く。すべての可約な形式列を除いた後、二重還元を用いて、残りのすべての可約形式およびすべての分解可能形式を取り除く。互いに独立に可約な形式列によって二重に覆われる全形式列の場所は、より高い形式の還元を生じさせる（§ 11, D 参照）。

\clearpage
\noindent\textbf{定理 III.}
\emph{二元領域から移された形式、すなわち、対応する二元形式の系の形式から Clebsch の転送原理による縁取りによって生じる形式は、法 $(abc)$ に関する相対的完全系を成す。}

\noindent\textbf{証明：}
二元の基本形式について、形式 $Q'_1,\ldots,Q'_n$ が完全系を成すことを表す二元関係
\[
Q'-F(Q')=0=\sum_\lambda\bigl\{(ab)c_x+(bc)a_x+(ca)b_x\bigr\}\varphi_\lambda^{*}
\]
には、縁取りによって三元関係
\[
Q-F(Q_i)=\sum_\lambda u_x\cdot(abc)\varphi_\lambda
\]
が対応する。これはまさに、法 $(abc)$ に関する相対的完全系の定義式である。四次形式の場合、移された形式の系は次の形式から成る：
\[
\begin{gathered}
f=a_x^4,\qquad
\theta=\theta_x^4u_\vartheta^2=(abu)^2a_x^2b_x^2,\qquad
K=K_x^6u_k^3=(a\theta u)u_\vartheta^2a_x^3\theta_x^3,\\
j=u_j^6=(a\theta u)^4u_\vartheta^2,\qquad
\nu=u_\nu^4=(abu)^4,
\end{gathered}
\]
これらのうち最初の四つは、法 $((abc),\nu)$ に関する相対的完全系を成す。

\begin{center}
\textbf{§ 5. \quad 法 $(abc)$ の法 $\A$ および $\nu$ への還元。}
\end{center}

単に括弧因子によって与えられる法 $(abc)$ は、定義 (§ 3, a) に従い、系の形式へ還元されなければならない。換言すれば、形式列 $(abc)$ は一意に正規化されなければならない（§ 1 参照）。

\[
a_s^2a_xu_x=\frac13 a_s^3u_x=\frac13 S\cdot u_x. \tag{2}
\]
\[
\left.
\begin{aligned}
(a\A u)^2&=a_s-\frac{S}{6}u_x^2,\\
(a\A u)^3&=0.
\end{aligned}
\right\} \tag{3}
\]
\[
\left.
\begin{aligned}
\theta(s)=(ss,x)^2&=2a_t+\frac13 S\cdot\theta-\frac13 Tu_x^2,\\
\theta(t)=(tt,x)^2&=-\frac13 S\cdot a_t+\frac23 T\cdot a_s+\frac1{18}S^2\cdot\theta-\frac1{18}u_x^2\cdot S\cdot T.
\end{aligned}
\right\} \tag{4}
\]

法の列：$(abc)$；$\A$；$s$；$t$（法 $t$ の系はただ一つの形式 $t$ から成る）。\footnote{Gordan--Kerschensteiner, loc. cit. p. 134.}

次の形式
\[
\A=(abc)^2a_x^2b_x^2c_x^2,\qquad \nu=(abu)^4
\]
が法として十分であること、すなわち形式
\[
\begin{gathered}
\A=(abc)^2a_x^2b_x^2c_x^2,\qquad
a_\nu=(abc)(bcu)^3a_x^3,\\
a_\nu^2=(abc)^2(bcu)^2u_x^2,\qquad
i=a_\nu^4=(abc)^4
\end{gathered}
\]
が形式列 $(abc)$ を定めることを示す。形式列 $a_\nu$ では、恒等式
\[
a_\nu^3=(abc)^3a_x(bcu)=\frac13 u_x\cdot(abc)^4=\frac13 i\cdot u_x
\]
により、形式 $a_\nu^3$ が欠ける。

法をより高い法へ還元するためには、定義に従い、最も一般的な折り畳み、またはその法とともに考慮されるすべての特殊な折り畳みを、より高い法との折り畳みへ還元しなければならない。

本場合に最も一般的な折り畳みは、式
\[
(abc)a_x^3b_y^3c_z^3,\qquad
(abc)a_x^2b_y^2c_z^2a_zb_xc_x
\]
（三次形式から移されたもの）、
\[
(abc)a_xb_yc_za_x^2b_x^2c_x^2
\]
（二次形式から移されたもの）、およびこれらの式の極化から生じる。

これらの式を $\A$ の極形式および形式列 $a_\nu$ の極形式、あるいはそれらの初期項によって計算するため、§ 2 と同じく間接的な道をとる。極項
\[
(abc)^2a_x^2b_y^2c_z^2,\qquad
a_\nu(\nu xy)(\nu yz)(\nu zx)a_xa_ya_z
=(abc)(bc\widehat{x}u)(bc\widehat{y}z)(bc\widehat{z}x)a_xa_ya_z\quad\text{等}
\]
を、記号因子 $(abc)$ をもつ式の線形結合として展開し、方程式系を反転することによって、求める関係、および変数の異なる組合せに従って取った極形式の間の関係を得る。評価には、行列式についての同一性定理および積定理を用いる。

$(abc)a_x^3b_y^3c_z^3$ に対する方程式系は次の通りである：
\begingroup\small
\[
\begin{aligned}
1)\quad
\sum_3 a_\nu(\nu yz)^3a_x^3
&=6(abc)a_x^3b_y^3c_z^3
-6\sum_3(abc)a_x^3b_y^2b_zc_z^2c_y^{*},\\[0.4em]
2)\quad
3(abc)^2a_x^2b_y^2c_z^2\cdot(xyz)
-\frac12\sum_3 a_\nu^2(\nu yz)^2\nu_x^2(xyz)
&=2\sum_3(abc)a_x^3b_y^2b_zc_z^2c_y\\
&\quad-4\sum_3(abc)a_xa_ya_zb_y^2b_zc_z^2c_x,\\[0.4em]
3)\quad
a_\nu(\nu xy)(\nu yz)(\nu zx)a_xa_ya_z
&=2\sum_3(abc)a_xa_ya_zb_y^2b_zc_z^2c_x,\\[0.4em]
4)\quad
\sum_3 a_\nu(\nu yz)^3u_x^3
-\frac32\sum_3 a_\nu^2(\nu yz)^2\nu_x^2(xyz)
+\frac16 i\cdot(xyz)^3
&=6\sum_3(abc)a_xa_ya_zb_y^2b_zc_z^2c_x.
\end{aligned}
\]
\endgroup

したがって
\[
(abc)a_x^3b_y^3c_z^3
\]
について、また同様の計算により
\[
(abc)a_x^2b_y^2c_z^2a_zb_xc_x,\qquad
(abc)a_xb_yc_za_z^2b_x^2c_x^2:
\]
について、次を得る：
\[
\begin{aligned}
(abc)a_x^3b_y^3c_z^3
&=\frac32(abc)^2a_x^2b_y^2c_z^2(xyz)
+\frac32 a_\nu(\nu xy)(\nu yz)(\nu zx)a_xa_ya_z\\
&\quad-\frac1{12}i\cdot(xyz)^3,\\[0.4em]
(\text{IV.})\quad
(abc)a_x^2b_y^2c_z^2a_zb_xc_x
&=\frac23(abc)^2a_x^2b_y^2c_z^2c_x\cdot(xyz)\\
&\quad+\frac13a_\nu(\nu xy)(\nu yz)(\nu zx)a_x^3
-\frac16 a_\nu^2(\nu xy)(\nu zx)a_x^2\cdot(xyz),\\[0.4em]
(abc)a_xb_yc_za_z^2b_x^2c_x^2
&=\frac16(abc)^2a_x^2b_z^2c_z^2\cdot(xyz).
\end{aligned}
\]

こうして、四つの形式 $f,\theta,K,j$ から成る、法 $(\A,\nu)$ に関する相対的完全系が得られた。したがって、二元領域から移されたものではない $f$ の $\theta$ 上のすべてのトランスヴェクション、および二元の場合に $(ab)^4$ へ還元可能なトランスヴェクション $(a\theta u)^2$ は、記号 $\A$ と $\nu$ によって表すことができる。\footnote{$\sum_3$ は、初期項から変数の循環置換によって得られる三項の和を指す。各方程式はその和の各項について別々にも成り立つ。}

以下で基本となる還元公式は、展開 IV の特殊化、またはより短く直接計算（記号の置換）によって得られる。すなわち、$a_\vartheta(a\theta u)^2$, $a_\vartheta^2((a\theta u)^2$, $(a\theta u)^2$ について、次の仮定から出発する：
\[
\begin{aligned}
a_\vartheta(a\theta u)^2
&=(abc)(bcu)(abu)^2-\frac13(abc)(bcu)\{(bcu)a_x-u_x(abc)\}^2,\\
a_\vartheta^2(a\theta u)^2
&=(abc)^2(abu)^2-\frac13(abc)^2\{(bcu)a_x-u_x(abc)\}^2,
\end{aligned}
\]
$((a\theta u)^2)^{*}$ については
\[
\begin{aligned}
(abu)a_y^3b_x^3&=\frac32u_\vartheta(\vartheta yx)\theta_y^2+\frac14(\nu yx)^3,\\
((abu)(acu))b_x^3&=\frac32u_\vartheta(\vartheta\widehat{a}ux)(\theta au)^2+\frac14(\nu\widehat{a}ux)^3.
\end{aligned}
\]
\[
\left.
\begin{aligned}
(a\theta u)^2&=\frac16\nu\cdot f-\frac23u_x\cdot a_\nu
+\frac23u_x^2\cdot a_\nu^2-\frac{i}{18}u_x^4,\\
a_\vartheta&=\frac13u_x\cdot\A,\\
a_\vartheta(a\theta u)&=0,\\
a_\vartheta^2&=\A,\\
(a\theta u)^3&=0,\\
a_\vartheta(a\theta u)^2&=-\frac56a_\nu+\frac76u_x\cdot a_\nu^2-\frac{i}{9}u_x^3,\\
a_\vartheta^2(a\theta u)&=0,\\
(a\theta u)^4&=j,\\
a_\vartheta(a\theta u)^3&=0,\\
a_\vartheta^2(a\theta u)^2&=\frac23a_\nu^2-\frac{i}{9}u_x^2.
\end{aligned}
\right\} \tag{1}
\]

これに、同じく法 $(abc)$ の還元から生じる公式
\[
a_\nu^3a_xu_\nu=\frac13 i\cdot u_x. \tag{2}
\]
を加える。

\emph{式 (1.) および (2.) と、基本形式 $\nu$ について同様に作られる公式から、以後のすべての還元公式は極化によって導かれる}。ただし、分解可能形式に関する関係を除く。式 (1.) から次が分かる。

形式列：
\[
\begin{array}{rcl}
a_\vartheta(a\theta u)^2 &\text{は}& \text{記号 }\nu\text{ のみへ導く},\\[0.2em]
a_\vartheta &\text{は}& \text{記号 }\A\text{ と }\nu\text{ へ導く},\\[0.2em]
(a\theta u)^2 &\text{は}& \text{記号 }j,\A\text{ と }\nu\text{ へ導く},\\[0.2em]
(a\theta u)^2\text{ の折り畳み I} &\text{は}& \text{記号 }j\text{ と }\nu\text{ へ導く},\\[0.2em]
(a\theta u)^2\text{ の折り畳み II} &\text{は}& \text{記号 }\A\text{ と }\nu\text{ へ導く}.
\end{array}
\]

したがって、次の置換が可能である：
\begin{enumerate}
\item 法 $(\A,\nu)$ に関して：$j$ との折り畳みを、$(a\theta u)^2$ または $(a\theta u)^3$ との対応する折り畳みに置き換える。
\item 法 $(\nu)$ に関して：$\A$ との折り畳みを、$a_\vartheta$ または $a_\vartheta(a\theta u)$ との対応する折り畳みに、または $(a\theta u)^2$ との特殊な折り畳みに置き換える（これは単なる一回の折り畳み I にしか導かない）。
\end{enumerate}
特に次が成り立つ：
\[
\begin{aligned}
(a\theta u)^2\theta_y^2\theta_x^2&=\frac13(jyx)^2\pmod{\nu},&
(a\theta u)^2u_y\theta_y&=-\frac16(jyx)^2\pmod{\nu},\\
(a\theta u)^2\nu_y^2&=\frac13(\A u\nu)^2\pmod{\nu},&
(a\theta u)(a\theta\nu)u_\vartheta\nu_\vartheta&=-\frac16(\A u\nu)^2\pmod{\nu},\\
(a\theta u)^2u_\nu^2\vartheta_\nu^2&=\frac13(\A u\nu)^2\A_\nu\pmod{\nu}.
\end{aligned}
\]

\begin{center}
\textbf{§ 6. \quad 法 $(\A,\nu)$ の法 $(\nu)$ への還元。}
\end{center}

前節の還元公式は、法 $\nu$ に関する相対的完全系を直接構成する手段を与える。移された形式の系には、形式 $\A$ および $(a\A u)=N$ だけを付け加えればよいことが分かる（これは三次形式の場合と同様であり、より一般にも成り立つ）。

法 $\A$ を還元するために、考慮すべきすべての特殊な折り畳み、すなわち法 $(\A,\nu)$ に関する相対的完全系と $\A$ の系との折り畳みを考え、係数における最も低い位数の形式、すなわち $f$ と $\A$ との折り畳みから出発する。

形式 $(a\A u)^2$, $(a\A u)^3$, $(a\A u)^4$ を還元するため、§ 5 に従い、記号 $\A$ を形式列の低い形式で置き換える。すなわち式
\[
a_\vartheta b_\vartheta(b\theta u)^2,\qquad
a_\vartheta(a\theta u)b_\vartheta(b\theta u)^2,\qquad
a_\vartheta(a\theta u)b_\vartheta(b\theta u)^3
\]
を、

1) 形式列 $a_\vartheta$（記号 $\A$ と $\nu$）の極形式に従って、

2) 形式列 $b_\vartheta(b\theta u)^2$（記号 $\nu$）の極形式に従って、

展開し、次の還元公式を得る（計算は下を参照）：
\begin{align}
\text{a)}\quad (a\A u)^2
&=\frac12(\vartheta\nu x)^2-\frac25 f\cdot a_\nu^2
-\frac45u_x\cdot\theta_\nu(\vartheta\nu x)^2 \notag\\
&\quad+u_x^2\left\{\frac16 i\cdot f-\frac1{40}S\right\},\notag\\
\text{b)}\quad (a\A u)^3
&=-\frac3{10}\theta_\nu(\vartheta\nu x),\tag{3}\\
\text{c)}\quad (a\A u)^4
&=\frac{21}{10}\theta_\nu^2-\frac3{10}\Pi
-\frac65u_x\cdot\theta_\nu^3+\frac1{10}u_x^2\cdot\theta .\notag
\end{align}

（同じ結果は、$a_\vartheta^2(b\theta u)^2$, $a_\vartheta^2(b\theta u)^3$, $a_\vartheta^2((a\theta u)^2(b\theta u)^2$ および $a_\vartheta^2((a\theta u)(b\theta u))^3$ を二重還元し、$a_\vartheta^2$ を消去することによっても得られる。）

式 (3.) から、$(a\A u)^2$ は法 $\nu$ に関する還元子である。これより次が従う。

1) $\A$ の系の還元：形式列 $(\A\A' u)^2=a_\vartheta^2((a\A u)^2$ は還元子 $(a\A u)^2$ を因子にもつ。

2) 法 $\A$ との折り畳みによって生じる系の還元：

形式列 $(\theta\A u)^2$ は還元子 $(a\A u)^2$ を因子にもつ。

形式 $(\theta\A u)$ と $\A_\vartheta$ については、
\[
(a\A u)^2(abu)=(\theta\A u)-u_x\cdot\A_\vartheta(\theta\A u),
\]
\[
(a\A u)(a\A b)=-\frac12\A_\vartheta+\frac12u_x\cdot\A_\vartheta^2.
\]
を得る。

しかし還元子 $(\theta\A u)$ から、法 $\A$ との折り畳みによって生じる系の還元が従う。

\bigskip

したがって、法 $\nu$ に関する相対的完全系は六つの形式
\[
f,\ \theta,\ K,\ j,\ \A,\ N
\]
から成る。

§ 2 の級数展開の例として、式 (3.)a の導出を詳しく与える。以後の計算では、最終公式 III を直接記すことにする。（零形式の極形式は計算中すでに省かれている。第一行は展開 II に従って記入した値を与え、第二行は方程式系の最後の式から始めて再帰的に求めた最終値を与える。）展開 II から次が従う：

\begin{align*}
\text{1)}\quad&m=3,\quad n=4,\quad \lambda=2,\quad \mu=0,\quad \nu=\chi=1,\\
a_\vartheta b_\vartheta(b\theta u)^2
&=[a_\vartheta]_{b^2}b+\frac67\{a_\vartheta b_\vartheta(a\theta u)(b\theta u)-u_xa_\vartheta b_\vartheta(a\theta b)(b\theta u)\}\\
&\quad-\frac17\{a_\vartheta(a\theta u)^2b_\vartheta-2u_xa_\vartheta(a\theta u)(a\theta b)b_\vartheta+u_x^2a_\vartheta(a\theta b)^2b_\vartheta\}\\
&=\frac23(a\Delta u)^2+\frac15[a_\vartheta(a\theta u)^2]b+\frac1{30}f[a_\vartheta^2(a\theta u)^2]-\frac25u_x[a_\vartheta(a\theta u)^2]b^2\\
&\quad-\frac1{15}u_x[a_\vartheta^2(a\theta u)^2]b+\frac15u_x^2[a_\vartheta(a\theta u)^2]b^3+\frac1{30}u_x^2[a_\vartheta^2(a\theta u)^2]b^2;\\[.6em]
\text{2)}\quad&m=2,\quad n=3,\quad \lambda=1,\quad \mu=1,\quad \nu=\chi=1,\\
a_\vartheta(a\theta u)b_\vartheta(b\theta u)
&=\frac25\{a_\vartheta(a\theta u)^2b_\vartheta-u_xa_\vartheta(a\theta u)(a\theta b)b_\vartheta\}+\frac12a_\vartheta^2(b\theta u)^2\\
&\quad-\frac15\{a_\vartheta^2(b\theta u)(a\theta u)-u_xa_\vartheta^2(a\theta b)(b\theta u)\}\\
&=\frac12(a\Delta u)^2+\frac25[a_\vartheta(a\theta u)^2]b+\frac1{15}[a_\vartheta^2(a\theta u)^2]f\\
&\quad-\frac25u_x[a_\vartheta(a\theta u)^2]b^2-\frac1{10}u_x[a_\vartheta^2(a\theta u)^2]b+\frac1{30}u_x^2[a_\vartheta^2(a\theta u)^2]b^2;\\[.6em]
\text{3)}\quad&m=2,\quad n=3,\quad \lambda=1,\quad \mu=1,\quad \nu=1,\quad \chi=2,\\
a_\vartheta(a\theta b)b_\vartheta(b\theta u)
&=\frac25\{a_\vartheta(a\theta b)(a\theta u)b_\vartheta-u_xa_\vartheta(a\theta b)^2b_\vartheta\}\\
&=\frac25[a_\vartheta(a\theta u)^2]b^2+\frac1{30}[a_\vartheta^2(a\theta u)^2]b-\frac25u_x[a_\vartheta(a\theta u)^2]b^3-\frac1{30}u_x[a_\vartheta^2(a\theta u)^2]b^2;\\[.6em]
\text{4)}\quad&m=1,\quad n=2,\quad \lambda=0,\quad \mu=2,\quad \nu=\chi=1,\\
a_\vartheta(a\theta u)^2b_\vartheta
&=[a_\vartheta(a\theta u)^2]b+\frac23a_\vartheta^2(a\theta u)(b\theta u)\\
&=[a_\vartheta(a\theta u)^2]b+\frac16[a_\vartheta^2(a\theta u)^2]f-\frac16u_x[a_\vartheta^2(a\theta u)^2]b;\\[.6em]
\text{5)}\quad&m=1,\quad n=2,\quad \lambda=0,\quad \mu=2,\quad \nu=1,\quad \chi=2,\\
a_\vartheta(a\theta u)(a\theta b)b_\vartheta
&=[a_\vartheta(a\theta u)^2]b^2+\frac13a_\vartheta^2(a\theta b)(b\theta u)\\
&=[a_\vartheta(a\theta u)^2]b^2+\frac1{12}[a_\vartheta^2(a\theta u)^2]b-\frac1{12}u_x[a_\vartheta^2(a\theta u)^2]b^2;\\[.6em]
\text{6)}\quad&m=1,\quad n=2,\quad \lambda=0,\quad \mu=2,\quad \nu=1,\quad \chi=3,\\
a_\vartheta(a\theta b)^2b_\vartheta&=[a_\vartheta(a\theta u)^2]b^3;\\[.6em]
\text{7)}\quad&m=2,\quad n=4,\quad \lambda=2,\quad \mu=\nu=\chi=0,\quad (\text{展開 I による}),\\
a_\vartheta^2(b\theta u)^2&=[a_\vartheta^2]_{ub^2}+\frac1{10}\{[a_\vartheta^2(a\theta u)^2]f-2u_x[a_\vartheta^2(a\theta u)^2]b+u_x^2[a_\vartheta^2(a\theta u)^2]b^2\};\\[.6em]
\text{8)}\quad&m=1,\quad n=3,\quad \lambda=1,\quad \mu=1,\quad \nu=\chi=0,\quad (\text{展開 I}),\\
a_\vartheta^2(a\theta u)(b\theta u)&=\frac14[a_\vartheta^2(a\theta u)^2]f-\frac14u_x[a_\vartheta^2(a\theta u)^2]b;\\[.6em]
\text{9)}\quad&m=1,\quad n=3,\quad \lambda=1,\quad \mu=\chi=1,\quad \nu=0,\quad (\text{展開 I}),\\
a_\vartheta^2(a\theta b)(b\theta u)&=\frac14[a_\vartheta^2(a\theta u)^2]b-\frac14u_x[a_\vartheta^2(a\theta u)^2]b^2.
\end{align*}

1) と 4) から次が従う：
\[
\begin{aligned}
(a\Delta u)^2
&=\frac65[a_\vartheta(a\theta u)^2]b+\frac15f[a_\vartheta^2(a\theta u)^2]+\frac35u_x[a_\vartheta(a\theta u)^2]b^2-\frac3{20}u_x[a_\vartheta^2(a\theta u)^2]b\\
&\quad-\frac3{10}u_x^2[a_\vartheta(a\theta u)^2]b^3-\frac1{20}u_x^2[a_\vartheta^2(a\theta u)^2]b^2,\\
(a\Delta u)^2
&=-a_\nu b_\nu+\frac35fa_\nu^2+\frac45u_xa_\nu^2b_\nu-\frac3{20}u_x^2a_\nu^2b_\nu^2-\frac1{12}u_x^2if.
\end{aligned}
\]

（記号 $\theta$ への変換については § 8 および § 9 を参照。）式 (3.)a は、補助変数 $w=x\widehat{u}b$ を導入すれば、展開 I の直接転送によってさらに短く計算できる。しかし式 (3.)b および (3.)c 以降では、ここで採った道のほうが簡単になる。(3.)b については、§ 9 により、因子 $u_x$ を含むすべての項を省くべきことがあらかじめ分かっている。(3.)b と (3.)c を計算するために次を得る：
\[
\begin{aligned}
(3.)\mathrm{b)}\quad 1)&\quad a_\vartheta(a\theta u)b_\vartheta(b\theta u)^2=\frac12(a\Delta u)^3+\frac45[a_\vartheta(a\theta u)^2]_{ub}b+\frac7{360}[a_\vartheta^2(a\theta u)^2]_{ub^2},\\
2)&\quad a_\vartheta(a\theta u)b_\vartheta(b\theta u)^2=[a_\vartheta(a\theta u)^2]_{ub}b+\frac{13}{36}[a_\vartheta^2(a\theta u)^2]_{ub^2};\\[.5em]
(3.)\mathrm{c)}\quad 1)&\quad a_\vartheta(a\theta u)b_\vartheta(b\theta u)^3=\frac12(a\Delta u)^4+\frac65[a_\vartheta(a\theta u)^2]_{ub^2}b+\frac{53}{120}[a_\vartheta^2(a\theta u)^2]_{ub^3}\\
&\qquad-\frac65u_x[a_\vartheta(a\theta u)^2]_{ub^3}b^2-\frac35u_x[a_\vartheta^2(a\theta u)^2]_{ub^2}b+\frac{19}{120}u_x^2[a_\vartheta^2(a\theta u)^2]_{ub^3}b^2,\\
2)&\quad a_\vartheta(a\theta u)b_\vartheta(b\theta u)^3=\frac34[a_\vartheta^2(a\theta u)^2]_{ub^2}.
\end{aligned}
\]

最後の注意：§ 5 によって可約な $f$ の $j$ 上へのトランスヴェクションも同様に計算される。すなわち
\[
\tag{4}
\begin{aligned}
g_\nu&=-\theta_\nu+u_x\left(\frac92\theta_\nu^2-\frac12H\right)-3u_x^2\theta_\nu^3+\frac12u_x^3\theta,\\
g_\nu^2&=\frac{21}{10}\theta_\nu^2-\frac3{10}H-2u_x\theta_\nu^3+\frac12u_x^2\theta,\\
g_\nu^3&=-\frac7{10}\theta_\nu^3+\frac12u_x^2\theta,\qquad g_\nu^4=\frac35\theta.
\end{aligned}
\]

\[
\tag{4}
g_\nu^3=-\frac7{10}\theta_\nu^3+\frac12u_x^2\theta,
\qquad
g_\nu^4=\frac35\theta.
\]

(3.) と (4.) から、二元領域からの転送により、
\[
(\theta\theta'u)^2=\frac13\,j\cdot f\pmod{\nu}.\footnotemark
\]
形式列 $(\theta\theta'u)^2$ の他の形式および形式 $\theta'_\nu$ は、記号 $\nu$ に還元可能である。したがって、次の置換が可能である：
\[
\bmod\ \nu:\quad \text{ } f\cdot j\text{ との折り畳みを}
\text{ }(\theta\theta'u)^2\text{ との対応する折り畳みで置き換える}.
\]
さらに
\[
(\vartheta\vartheta'x)^2=\frac43 f\cdot\A,
\qquad
\theta_{g\nu}(\vartheta\vartheta'x)=-N.
\]
\footnotetext{Gordan--Kerschensteiner, loc. cit. p. 181.}



\begin{center}
{\Large\bfseries 第 III 章。\quad 法 $(s,\varrho,t)$ に関する相対的完全系。}
\end{center}

\begin{center}
\textbf{\S\ 7.\quad 系の形成の概観。}
\end{center}

法 $\nu$ に関する相対的完全系から、次に高い法をもつ相対的完全系へ移るには、\S\ 3 の定義に従って、この高い法に関する $\nu$ の系を作り、それを法 $\nu$ に関する相対的完全系の形式と折り畳まなければならない。ところが $\nu=u_\nu^4$ は、形式 $f$ の四次反変式として $f$ と双対的に対応している。したがって $\nu$ を基本形式と見れば、法 $\nu(\nu)=(\nu\nu_1x)^4=s_x^4$ に関する六つの形式から成る相対的完全系が得られる。

したがって、法 $s$ に関する相対的完全系（系 III）を形成するためには、次をトランスヴェクションすればよい：
\[
\begin{array}{ll}
\text{系 I:} & f,\ \theta,\ K,\ j,\ \A,\ N\quad \text{を}\vphantom{\dfrac11}\\[0.3em]
\text{系 II:} & \nu,\ H=(\nu\nu_1x)^2u_\nu^2u_{\nu_1}^2,\ L=(\nu\eta x)H_x^2u_\nu^3u_\eta^3,\\
& g=(\nu\eta x)^4H_x^2,\ \sigma=H_\eta^2u_\nu^2u_\eta^4,\ (\nu\sigma x)
\end{array}
\]
とである。後で（式 (13)）示すように、形式 $g$ は還元可能であり、従って系 II は五つの形式だけから成る。

計算の途中で、二次形式
\[
\varrho=u_\varrho^2=\theta_\nu^4u_\vartheta^2,
\qquad
 t=t_x^2=a_\eta^4H_x^2
\]
を法として付け加える。そのため法 $(s,\varrho,t)$ に関する相対的完全系が得られる。ただし法 $(\varrho,t)$ は、法 $s$ より高い法とみなす。

個々の形式の配列は、係数における位数に従って行う。折り畳み
\[
\theta_\nu(K_\nu,\theta_\nu,\ldots)
\]
を折り畳み
\[
H_y(H_y,L_y,\ldots)
\]
より高いものとして正規化する。これにより、\S\ 1 および \S\ 3b に従って、同じ位数をもつ個々の形式の順序は一意に定まる。

同じ位数の形式の形成は表によって実行する：

I. 係数における所定の位数を得るため、系 I と系 II のどの形式を互いに折り畳むべきかの指定。

II. 形式列の図式に従った既約形式の列挙。

III. 還元可能または分解可能な形式列あるいは形式の還元、すなわち全形式列が中断する箇所の還元。これにより、すべての既約な形成が、挙げられたものによって尽くされていることが示される。

IV. 帰結。すなわち 1) 新しく生じた還元子の指定、2) 同じ系の形式との積において、他方の系の形式との折り畳みに入らない形式の指定。

現れるものは次の場合だけであることが分かる：

1) 系 I の一つの形式と系 II の一つの形式との折り畳み。

2) $\theta$、または $H$ の冪の折り畳み、あるいは一方の系の二つの形式と第二の系の一つの形式との折り畳み。

折り畳みに関して次の図式を得る：
\[
\begin{array}{r@{\ }c@{\ }l@{\qquad}c@{\qquad}r@{\ }c@{\ }l}
\multicolumn{3}{c}{\text{系 I}} & \text{と折り畳む} & \multicolumn{3}{c}{\text{系 II}}\\[0.25em]
1. & \text{位数} & f & & 2. & \text{位数} & \nu\\
2. & ,, & \theta & & 4. & ,, & H\\
3. & ,, & K,j,\A & & 6. & ,, & L,\sigma\\
4. & ,, & N,\theta^2,f\cdot j & & 8. & ,, & (\nu\sigma x),H^2\\
5. & ,, & \theta\cdot K & & 10. & ,, & H\cdot L\\
6. & ,, & \theta^3,\A\cdot j & & 12. & ,, & H^3.
\end{array}
\]

計算の検査には、「二重還元」のほか、次の二つの特殊形式を用いる：

\[
\begin{array}{ll}
\text{I.}&\begin{array}{rlrlrl}
f&=\displaystyle\sum_3x_1^3x_2,& u_\nu^2&=\frac{i}{6}u_x^2,& s&=\frac{i}{3}f,\\
a_\eta^2&=-\frac19 i\theta,& H_\nu^2&=\frac{i}{6}\theta,& \sigma&=-\frac{i}{6}j,\\
g&=-\frac23i\Delta,& \varrho&=0,& t&=0,
\end{array}\\[1.2em]
\text{II.}&\begin{array}{rlrlrl}
f&=\displaystyle\sum_3x_1^4,& s&=\frac43if,& a_\eta^2&=\frac29i\theta,\\
\Delta_\nu^2&=\frac1{15}i\theta,& \sigma&=\frac43ij,& g&=\frac43i\Delta,\\
\varrho&=0,& t&=0.
\end{array}
\end{array}
\]

後注：基本形式 $\nu$ について、式 (1), (2), (3), (4) に対応するものは次である：
\[
\tag{5}
\begin{array}{c|c|c}
(\nu\eta x)^2=A & H_\nu(\nu\eta x)=0 & H_\nu^2=\sigma\\[0.4em]
(\nu\eta x)^3=0 & H_\nu(\nu\eta x)^2=B & H_\nu^2(\nu\eta x)=0\\[0.4em]
(\nu\eta x)^4=g & H_\nu(\nu\eta x)^3=0 & H_\nu^2(\nu\eta x)^2=C,
\end{array}
\]
\[
A=\frac16s\nu-\frac23u_xs_\nu+\frac23u_x^2s_\nu^2-\frac{J}{18}u_x^4,
\quad B=-\frac56s_\nu+\frac76u_xs_\nu^2-\frac{J}{9}u_x^3,
\quad C=\frac23s_\nu^2-\frac{J}{9}u_x^2.
\]
\[
\tag{6}
s_\nu^3u_xs_\nu=\frac13u_xs_\nu^4=\frac13Ju_x.
\]

\[
\tag{7}
\begin{aligned}
\text{a)}\quad (\nu\sigma x)^2
&=\frac12(Hsu)^2-\frac25\nu\cdot s_\nu^2-
\frac45u_x\,s_\eta(Hsu)^2
+u_x^2\left\{\frac16J\cdot\nu-\frac1{40}(ss'u)^4\right\},\\
\text{b)}\quad (\nu\sigma x)^3&=-\frac3{10}s_\eta(Hsu),\\
\text{c)}\quad (\nu\sigma x)^4&=\frac{21}{10}s_\eta^2-\frac3{10}Z-
\frac65u_xs_\eta^3+\frac1{10}u_x^2s_\eta^4.
\end{aligned}
\]
\[
\tag{8}
\begin{aligned}
\text{a)}\quad g_\nu&=-s_\eta+u_x\left(\frac92s_\eta^2-\frac12Z\right)-3u_x^2s_\eta^3+\frac12u_x^3s_\eta^4,\\
\text{b)}\quad g_\nu^2&=\frac{21}{10}s_\eta^2-\frac3{10}Z-2u_xs_\eta^3+\frac12u_x^2s_\eta^4,\\
\text{c)}\quad g_\nu^3&=-\frac7{10}s_\eta^3+\frac12u_xs_\eta^4,\\
\text{d)}\quad g_\nu^4&=\frac35s_\eta^4.
\end{aligned}
\]

\begin{center}
\textbf{\S\ 8.\quad 三次の形式（系 III）。}
\end{center}

\noindent\textbf{$f$ と $\nu$ の折り畳み。}

\noindent\emph{既約形式：}
\[
\begin{array}{cccc}
a_\nu&\cdot&\cdot&\cdot\\
\cdot&a_\nu^2&\cdot&\cdot\\
\cdot&\cdot&\cdot&\cdot\\
\cdot&\cdot&\cdot&a_\nu^4=i.
\end{array}
\]

\noindent\emph{還元：}
\[
a_\nu^3=\frac13 i u_x\qquad\text{(式 (2)).}
\]
\noindent\emph{帰結：} 1) $a_\nu^3$ は還元子である。2) $f$ は、反傾形式 $(j)$ との積においてのみ、$\nu$ との折り畳みに入る。$f$ に関する $\nu$ についても同じことが成り立つ。

\begin{center}
\textbf{\S\ 9.\quad 四次の形式（系 III）。}
\end{center}

\noindent\textbf{$\theta$ と $\nu$ の折り畳み。}

\noindent\emph{既約形式：}
\[
\begin{array}{cccc}
(\theta\nu x)&\theta_\nu&\cdot&\cdot\\
(\theta\nu x)^2&\theta_\nu(\theta\nu x)&\theta_\nu^2&\cdot\\
\cdot&\theta_\nu(\theta\nu x)^2&\cdot&\theta_\nu^3.
\end{array}
\]

\noindent\emph{還元：} 還元子 $a_\nu^3$ から、式 $a_\nu^3(abu)$, $a_\nu^3b_\nu$, $a_\nu^3b_\nu(abu)$ の二重還元を、次の設定に従って行う：

\[
\begin{aligned}
\theta_\nu^2(\vartheta\nu x)
&=a_\nu^2(\widehat{a}b\nu x)(abu)-\frac13(\nu\nu_1x)^3=a_\nu b_\nu(\widehat{a}b\nu x)(abu)+\frac16(\nu\nu_1x)^3\\
&=a_\nu^3(abu)-a_\nu^2b_\nu(abu)=2a_\nu^2b_\nu(abu)=\frac23a_\nu^3(abu),\\[.4em]
\theta_\nu^2(\vartheta\nu x)^2
&=a_\nu^2(\widehat{a}b\nu x)^2-\frac13(\nu\nu_1x)^4=a_\nu b_\nu(\widehat{a}b\nu x)^2+\frac16(\nu\nu_1x)^4\\
&=if-2a_\nu^3b_\nu+a_\nu^2b_\nu^2-\frac13s=2a_\nu^3b_\nu-2a_\nu^2b_\nu^2+\frac16s=\frac23if-\frac23a_\nu^3b_\nu-\frac16s,\\[.4em]
\theta_\nu^3(\vartheta\nu x)&=a_\nu^2b_\nu(\widehat{a}b\nu x)(abu)=a_\nu^3b_\nu(abu).
\end{aligned}
\]

得られる式は次である：
\[
\tag{9}
\begin{array}{c|c|c}
\theta_\nu^2(\theta\nu x)=0
& \theta_\nu^3(\theta\nu x)=0
& \theta_\nu^4u_x^2=u_\varrho^2=\varrho\quad(\text{付加された法, \S\ 7}),\\[0.6em]
\theta_\nu^2(\theta\nu x)^2=\dfrac49 i f-\dfrac16s
& &
\end{array}
\]

\noindent\emph{帰結：} 1) $\theta_\nu^2(\theta\nu x)^2$ は還元子である。2) $\nu$ は、反傾形式 $(g)$ との積においてのみ、$\theta$ との折り畳みに入る。

\begin{center}
\textbf{\S\ 10.\quad 五次の形式（系 III）。}
\end{center}

\[
\text{次の折り畳み： }\left\{\begin{array}{l}
K,j,\A \text{ と }\nu,\\
f\text{ と }H.
\end{array}\right.
\]

\noindent\emph{既約形式：}
\[
\begin{array}{c@{\qquad}c@{\qquad}c@{\qquad}c}
\text{A)} & \begin{array}{cccc}
\cdot&K_\nu&\cdot&\cdot\\
(k\nu x)^2&\cdot&K_\nu^2&\cdot\\
(k\nu x)^3&K_\nu(k\nu x)^2&\cdot&\cdot\\
\cdot&K_\nu(k\nu x)^3&\cdot&\cdot
\end{array}
&
\text{B)}\ \begin{array}{c}(j\nu x)\\(j\nu x)^2\\(j\nu x)^3\\ \cdot\end{array}
&
\text{C)}\ \begin{array}{c}\A_\nu\\\cdot\\\cdot\\\cdot\end{array}
\end{array}
\]
\[
\text{D)}\qquad
\begin{array}{cccc}
(aHu)&(aHu)^2&\cdot&\cdot\\
a_\eta&a_\eta(aHu)&a_\eta(aHu)^2&\cdot\\
\cdot&a_\eta^2&\cdot&\cdot
\end{array}
\]

\noindent\emph{還元：}

\noindent A) 形式列 $K_\nu^3$: 還元子 $a_\nu^3$。

形式 $K_\nu^2(k\nu x)^2$: 還元子 $\theta_\nu^2(\vartheta\nu x)^2$。
\[
(k\nu x),\quad K_\nu(k\nu x),\quad K_\nu^2(k\nu x)
\]
は \S\ 3 による分解可能形式である。

\noindent B) 形式
\[
\begin{aligned}
(j\nu x)^4&=(\hata\theta\nu x)^4
=a_\nu^4\theta+\theta_\nu^4 f-4a_\nu^3\theta_\nu
-4a_\nu\{a_\nu\theta_x-(\hata\theta\nu x)\}^3\\
&\quad+6a_\nu^2\{a_\nu\theta_x-(\hata\theta\nu x)\}^2
=\R f+\frac53 i\theta-6a_\nu^2(\hata\theta\nu x)^2
+4a_\nu(\hata\theta\nu x)^3,
\end{aligned}
\]
従って
\[
(j\nu x)^4=\R f+\frac53 i\theta\modu{(\A,\nu)},\qquad (\text{\S\ 5 の末尾参照}).
\]

\noindent C) 形式 $\A_\nu^4$: 還元子 $\theta_\nu^4$。

形式 $\A_\nu^2$ および $\A_\nu^3$: 形式 $\A$ を形式列の低い形式（\S\ 5 の末尾）で置き換えることにより、すなわち式
\[
a_\vartheta a_\nu^3,
\qquad a_\vartheta a_\nu^3\theta_\nu,
\qquad a_\vartheta\theta_\nu^4
\]
の二重還元により、還元子 $a_\nu^3$ または $\theta_\nu^4$ をもつ。$\A_\nu^3$ の二通りの計算は、高次形式 $a_\eta^3$ の還元を引き起こす。

\noindent\emph{還元公式：}
\[
\tag{10}
\begin{aligned}
\text{a)}\quad \A_\nu^2
&=\frac35a_\eta^2-\frac3{10}(asu)^2+\frac13 i\theta
+\frac15u_x^2(2a_\R^2-t),\\
\text{b)}\quad \A_\nu^3
&=-\frac3{10}a_\R+\frac35u_x\left(\frac{13}{10}a_\R^2-\frac25t\right),\\
\text{c)}\quad \A_\nu^4&=a_\R^2-\frac25t.
\end{aligned}
\]

\noindent\emph{公式の導出：}
\[
\begin{aligned}
\text{a)}\quad
 a_\vartheta a_\nu^3
&=\frac13 i\theta
=[a_\vartheta]_{\nu^3}+\frac67[a_\vartheta^2]_{\nu^2}
+\frac65[a_\vartheta(a\theta u)^2]_{\nu^2}\, u x^2
+\frac{11}{30}[a_\vartheta^2(a\theta u)^2]_\nu\,\nu x^2\\
&\quad-\frac67u_x[a_\vartheta^2]_{\nu^3}
-\frac16u_x[a_\vartheta^2(a\theta u)^2]_{\nu^2}\,\nu x^2,\\
\frac13 i\theta
&=\A_\nu^2-\frac23u_x\A_\nu^3-a_\nu a_{\nu_1}(\nu\nu_1x)^2
+\frac25a_\nu^2(\nu\nu_1x)^2
+\frac15u_x^2a_\nu^2a_{\nu_1}(\nu\nu_1x)^2;
\end{aligned}
\]
\[
\begin{aligned}
\text{b)}\quad
 a_\vartheta a_\nu^3\theta_\nu
&=0=[a_\vartheta]_{\nu^4}+\frac9{14}[a_\vartheta^2]_{\nu^3}
+\frac35[a_\vartheta(a\theta u)^2]_{\nu^2}\,\nu x^2
+\frac{17}{120}[a_\vartheta^2(a\theta u)^2]_{\nu}\,\nu x^2\\
&\quad-\frac9{14}u_x[a_\vartheta^2]_{\nu^4}
-\frac1{24}u_x[a_\vartheta^2(a\theta u)^2]_{\nu^2}\,\nu x^2,\\
0&=\frac56\A_\nu^3-\frac12u_x\A_\nu^4-\frac14a_\eta^3+\frac1{20}u_xa_\eta^4;
\end{aligned}
\]
\[
\begin{aligned}
\text{b')}\quad
 a_\vartheta\theta_\nu^4
&=a_\R=a_\vartheta\theta_\nu\{a_\nu-(a\theta\nu x)\}^3\\
&=-\frac32[a_\vartheta]_{\nu^3}+\frac35[a_\vartheta(a\theta u)^2]_{\nu^2}\,\nu x^2
-\frac{37}{120}[a_\vartheta^2(a\theta u)^2]_{\nu}\,\nu x^2\\
&\quad+\frac32u_x[a_\vartheta^2]_{\nu^4}
+\frac{49}{120}u_x[a_\vartheta^2(a\theta u)^2]_{\nu^2}\,\nu x^2,\\
a_\R&=-\frac32\A_\nu^3+\frac32u_x\A_\nu^4-\frac{11}{20}a_\eta^3+\frac7{20}u_xa_\eta^4;
\end{aligned}
\]
\[
\begin{aligned}
\text{c)}\quad
 a_\vartheta^2\theta_\nu^4
&=a_\R^2=[a_\vartheta^2]_{\nu^4}+\frac35[a_\vartheta^2(a\theta u)^2]_{\nu^2}\,\nu x^2,\\
a_\R^2&=\A_\nu^4+\frac25a_\eta^4.
\end{aligned}
\]

\noindent D) $\A_\nu^3$ の二重還元 (C) による $a_\eta^3$ の還元。

残りの還元は、\S\ 9 の計算と双対的に反対の計算によって得られる。

\noindent\emph{還元公式：}
\[
\tag{11}
\begin{array}{rl|rl}
a_\eta^2(aHu)&=-\dfrac16(asu)^3,&
 a_\eta^2(aHu)^2&=\dfrac49 i\nu-\dfrac16(asu)^4,\\[0.7em]
a_\eta^3&=-a_\R+\dfrac35u_xa_\R^2+\dfrac15u_xt,
& a_\eta^3(aHu)&=0,\\[0.7em]
&&a_\eta^4&=t\quad(\text{法として付け加えた}).
\end{array}
\]

\noindent\emph{帰結：}

1) $(j\nu x)^4$, $\A_\nu^2$, $a_\eta^2(aHu)$ は還元子である。

2) $\nu$ は、$K,j,\A$ との折り畳みにおいて、反傾形式との積にだけ入る。$H$ に関する $f$、および $\nu$ に関する $j$ についても同じことが成り立つ。なぜなら $\theta_\nu(\theta\cdot j,\nu,x^3)$ は \S\ 11 B により還元可能になるからである。

\begin{center}
\textbf{\S\ 11.\quad 六次の形式（系 III）。}
\end{center}
\[
\text{次の折り畳み： }\left\{\begin{array}{l}
\theta^2,
 f\cdot j,
 N \text{ と }\nu,\\
\theta\text{ と }H.
\end{array}\right.
\]

\noindent\emph{既約形式：}
\[
\begin{array}{c@{\quad}c@{\quad}c}
\text{A)}\ (\vartheta^2\nu x)^3,\quad \theta_\nu(\vartheta^2\nu x)^3
&
\text{B)}\ a_\nu(j\nu x),\quad a_\nu^2(j\nu x)
&
\text{C)}\ N_\nu
\end{array}
\]
\[
\text{D)}\quad
\begin{array}{c|c|c|c}
 & (\theta Hu)&(\theta Hu)^2&\cdot\\
(\vartheta\eta x)&\theta_\eta&H_\vartheta(\theta Hu),\ \theta_\eta(\theta Hu)&\theta_\eta(\theta Hu)^2\\
(\vartheta\eta x)^2&H_\vartheta(\vartheta\eta x),\ \theta_\eta(\vartheta\eta x)&\theta_\eta^2&\theta_\eta^2(\theta Hu)\\
\cdot&\theta_\eta(\vartheta\eta x)^2&\cdot&\cdot
\end{array}
\]

\noindent\emph{還元：}

\noindent A) $(\vartheta^2\nu x)^4$ は式 (3)a に従って分解する：
\[
(a\A\vartheta x)^4=\frac12(\vartheta^2\nu x)^4-\frac35 f\cdot a_\nu^2(\nu\vartheta x)^2
=\frac13\A^2+f\A_\vartheta^2.
\]

\noindent B) $a_\nu(j\nu x)^2$ は式 (9)a に従って分解する。ここでは \S\ 6 の末尾に従い $f\cdot j$ を $(\theta\theta'u)^2$ で置き換え、$\theta_\vartheta'$ の還元可能性に注意する：
\[
\frac13a_\nu(j\nu x)^2=(\theta\theta'\nu x)^2\theta_\nu
=\theta_\nu^3\theta-\theta_\nu^2\theta_\nu'
=\theta_\nu^3\theta+\theta_\nu^2(\widehat{j}\nu\theta'u)
=\theta_\nu^3\theta-\frac23\theta_\nu^3\theta.
\]
形式 $a_\nu(j\nu x)^3$ および $a_\nu^2(j\nu x)^2$：還元子 $a_j$ および $(j\nu x)^4$：
\[
a_\nu(j\nu x)^3=a_j(j\nu x)^3-(j\nu x)^3(j\nu\widehat{a}u),
\]
\[
a_\nu^2(j\nu x)^2=a_j^2(j\nu x)^2-2a_j(j\nu x)^2(j\nu\widehat{a}u)+(j\nu x)^2(j\nu\widehat{a}u)^2.
\]

\noindent C) 形式列 $N_\nu^2$；還元子 $\A_\nu^2$。$(n\nu x)$ および $N_\nu(n\nu x)$ は、\S\ 3 に従って関数行列式に関するトランスヴェクションとして分解する。

\noindent D) \emph{還元の概観：}

a) 形式列 $\theta_\eta^2H_\vartheta$: 還元子 $a_\eta^2(aHu)$。（高次記号をもつ形式へ導く。）

b) 形式列 $\theta_\eta H_\vartheta$: 二重還元による還元子 $a_\eta^2(aHu)$。（高次記号をもつ形式、および全形式列の高次形式、すなわち形式列 $\theta_\eta^2$ へ導く。）$\theta_\eta H_\vartheta$ の代わりに形式列
\[
\theta_\eta(\vartheta\eta x)(\theta Hu)=\theta_\eta H_\vartheta+\theta_\eta^2
\]
を考え、その中で記号 $\theta$ を低い形式 $(abu)$ で置き換え、こうして形式列
\[
a_\eta^2(aHu)(abu)=\theta_\eta H_\vartheta+\frac32\theta_\eta^2\modu{s}
\]
を、終端形式
\[
a_\eta^3b_\eta(bHu)(abu)=\theta_\eta^3H_\vartheta+\frac12\theta_\eta^4
\]
まで二重還元する。

c) 形式列 $\theta_\eta^3$: 形式列 $\theta_\eta H_\vartheta$ と $\theta_\eta^2H_\vartheta$ の両方に属する形式、すなわち形式列 $\theta_\eta^2H_\vartheta$ の形式を二重還元する。一方では高次記号をもつ形式へ、他方では高次記号をもつ形式と高次形式列 $\theta_\eta^3$ へ還元する。

d) 形式 $\theta_\eta^2(\vartheta\eta x)$, $\theta_\eta^2(\vartheta\eta x)^2$, $\theta_\eta^3(\vartheta\eta x)$: \S\ 9 の計算に類似して、還元子 $a$ による二重還元。

e) 形式 $\theta_\eta^2(\theta Hu)^2$ および $\theta_\eta^2(\vartheta\eta x)^2$: 還元子 $\A_\nu^2$ および $(a\A u)^4$ を用いる式 $\A_\nu^2(a\A u)^4$ および $(a\A\nu x)^4$ の二重還元。

f) 形式 $H_\vartheta$ および $H_\vartheta^2$: 分解可能形式。$H_\vartheta^2$ については、式 $\A_\nu^2(a\A u)^2$ の二重還元、または系の形式による積 $(a_\nu^2)^2$, $a_\nu^2b_{\nu_1}$ の直接計算から得られる。（$(a_\nu)^2$ の類似計算は分解可能形式間の関係を与える。）

g) 形式列 $\theta\cdot H$ の形式が二つの還元可能性を許す（すなわち二つの還元可能な形式列に属する）ことによる高次形式の還元。すなわち：
\begin{quote}\small
$g$: $\theta_\eta^2(\vartheta\eta x)^2$ の二重還元による。還元 d) と e) を許す。\par
$s_\vartheta^2(\theta su)$: $\theta_\eta^3(\vartheta\eta x)$ の二重還元による。還元 c) と d) を許す。\par
$s_\vartheta^2(\theta su)^2$: $\theta_\eta^2H_\vartheta^2$ の二重還元による。還元 a) を許し、還元子 $\theta_\nu^2(\vartheta\nu x)^2$ を因子にもつ。\par
$s_\nu^2$: $\theta_\eta^3H_\vartheta$ の二重還元による。還元 a) と c) を許す。
\end{quote}
残りの形式の独立性の証明は、形式列 $\A_\nu^2(a\A u)^2=\theta_\eta^2$ の二重還元によって得られる。

\noindent\emph{還元の実行：}

還元 a), b), c), d) の存在はアプリオリに明らかである。なぜなら、示された形成法則に従えば、二重還元の際に生じる関係は互いに独立だからである。還元 e), f), g) については、恒等式が生じないことを計算により示さなければならない。

形式列 $\theta\cdot H$ の対角項に対して対称的に、形式 $\theta_\eta$ まで進め、計算の一部を示しながら、還元 a), b), c), d) から得られる公式も与える。これらは、一部は還元 e), f), g) に、一部は後に用いられるからである。

\noindent 1) $H_\vartheta$ および $H_\vartheta^2$、f) による。設定\footnote{等号の代わりに記号 $\sim$ を用いることは、$u_x$ と、折り畳み III および IV によって生じる形式より高次の形式との積を省略していることを意味する。}
\[
\begin{aligned}
a_\nu b_{\nu_1}^2
&=\nu a_\nu b_\nu^2-(aHu)(bHu)b_\eta
\sim \nu a_\nu b_\nu^2-f a_\eta(aHu)^2-(abu)(aHu)b_\eta+(abu)^2b_\eta,\\
a_\nu^2b_{\nu_1}^2
&=b_\nu^2\{a_\nu u_\nu-(a\nu\nu_1u)\}^2
\sim \nu\{a_\nu^2b_\nu^2-2a_\nu b_\nu(aHu)(bHu)
-\frac16(asu)^2(bsu)^2\}\\
&\sim \nu a_\nu^2b_\nu^2-2a^2(aHu)(abu)-\frac16(asu)^2(bsu)^2+(\vartheta\eta x)^2(\theta Hu)^2
\end{aligned}
\]
において、$\theta_\eta H_\vartheta$ の値を代入すると、次の公式が生じる：
\[
\tag{12}
\begin{aligned}
\text{a)}\quad H_\vartheta&\sim 2a_\nu a_\R^2+2\nu b_\nu(\vartheta\eta x)^2+2f a_\eta(aHu)^2-2\theta_\eta,\\
\text{b)}\quad H_\vartheta^2&\sim (a^2)^2+2\theta_\eta^2+\frac23(\theta su)^2
+\frac1{12}s\nu-\frac19if\nu-\frac16f(asu)^4.
\end{aligned}
\]

\noindent 2) $\theta_\eta H_\vartheta$、b) による：
\[
\begin{aligned}
a_\eta^2(aHu)(abu)&\sim [a_\eta^2(aHu)]_{bu}+\frac12 f[a_\eta^2(aHu)^2]
=a_\eta b_\eta(aHu)(abu)\\
&\quad+a_\eta(a\widehat{b}\eta x)(abu)(aHu)
\sim \theta_\eta(\vartheta\eta x)(\theta Hu)+\frac12\theta_\eta^2+\frac14(\nu\eta x)^2.
\end{aligned}
\]
式 (5) および (11) から次が従う：
\[
\theta_\eta H_\vartheta\sim -\frac32\theta_\eta^2-\frac14(\theta su)^2-\frac1{12}s\nu+\frac29if\nu.
\]

\noindent 3) $\theta_\eta H_\vartheta(\theta Hu)$ は、$\theta_\eta H_\vartheta$ と同様に b) に従って計算される：
\[
\theta_\eta H_\vartheta(\theta Hu)\sim-\theta_\eta^2(\theta Hu)-\frac16(\theta su)^3.
\]

\noindent 4) $\theta_\eta H_\vartheta(\vartheta\eta x)$ は b) による。$\theta_\eta^2(\vartheta\eta x)$ は d) による（\S\ 9 参照）：
\[
\theta_\eta H_\vartheta(\vartheta\eta x)\sim-\theta_\eta^2(\vartheta\eta x)-\frac16s_\vartheta(\theta su)
\sim-\frac13(\vartheta\R x)-\frac16s_\vartheta(\theta su),
\]
\[
\theta_\eta^2(\vartheta\eta x)\sim\frac23a_\eta^3(abu)\sim-\frac13(\vartheta\R x).
\]

\noindent 5)
\[
\begin{array}{ccc}
\theta_\eta H_\vartheta^2\text{ は b) による},&
\theta_\eta^2H_\vartheta\text{ は a) による},&
\theta_\eta^2H_\vartheta\text{ は b) による、従って }\theta_\eta^3\text{ は c) による},\\
\text{ }a_\eta^2(aHb)(bHu)\text{ から};&
\text{ }a_\eta^3(Hab)(abu)\text{ から};&
\text{ }a_\eta^3(bHu)(abu)\text{ から}.
\end{array}
\]
\[
\theta_\eta H_\vartheta^2\sim-\frac32\theta_\eta^2H_\vartheta-\frac14s_\vartheta(\theta su)^2+\frac16s_\nu-\frac49ia_\nu
\sim-\theta_\R-\frac16s_\nu-\frac19ia_\nu,
\]
\[
\theta_\eta^2H_\vartheta\sim\frac23\theta_\R-\frac16s_\vartheta(\theta su)^2+\frac29s_\nu-\frac29ia_\nu,
\]
\[
\theta_\eta^3H_\vartheta\sim\frac23\theta_\R-\frac23\theta_\eta^3+\frac5{36}s_\nu,
\qquad
\theta_\eta^3\sim\frac14s_\vartheta(\theta su)^2-\frac18s_\nu+\frac13ia_\nu.
\]

\noindent 6) $\theta_\eta^2(\theta Hu)^2$、e) による：
\[
\A_\nu^2(a\A u)^2=[(a\A u)^4]_{\nu^2}
=-\frac{12}{5}\theta_\eta^2(\theta Hu)^2-\frac3{10}\sigma-\frac1{20}(\theta su)^4+\nu\R
\quad(\text{式 (3)c による}),
\]
\[
\A_\nu^2(a\A u)^4=[\A_\nu^2]_{aa}^4=-\frac35\theta_\eta^2(\theta Hu)^2+\frac15\sigma-\frac3{10}(\theta su)^4+\frac13ij
\quad(\text{式 (10)a による}),
\]
\[
\theta_\eta^2(\theta Hu)^2=-\frac16\sigma+\frac1{12}(\theta su)^4-\frac19ij+\frac13\nu\R.
\]

\noindent 7) $\theta_\eta^2(\vartheta\eta x)^2$、d) および e) による。これから高次形式 $g$ の還元が得られる。

\S\ 9 と同様の計算により、次が従う：
\[
\begin{aligned}
\theta_\eta^2(\vartheta\eta x)^2
&=\frac12 i f-\frac12a_\eta^2b_\eta-\frac1{12}g
=\frac23tf-\frac23a_\eta^3b_\eta-\frac16g\\
&=-\frac16g-\frac13(\vartheta\R x)^2+\frac4{15}fa_\R^2+\frac8{15}ft
\quad(\text{式 (11) による}).
\end{aligned}
\]
上と対応する計算により、次が従う：
\[
((a\A\nu x)^4)=[(a\A u)^4]_{\nu x^4}
=\frac{21}{10}\theta_\eta^2(\vartheta\eta x)^2+\frac7{10}s_\vartheta^2-\frac3{10}g
\quad(\text{式 (3)c による}),
\]
\[
\begin{aligned}
(a\A\nu x)^4&=-\frac13i\A+6[\A_\nu^2]a^2-4[\A_\nu^3]a+\A_\nu^4f\\
&= -\frac{36}{5}\theta_\eta^2(\vartheta\eta x)^2-\frac95s_\vartheta^2-\frac35g-\frac35(\vartheta\R x)^2\\
&\quad+\frac53i\A+\frac{37}{25}fa_\R^2+\frac{74}{25}ft\quad(\text{式 (10) による}).
\end{aligned}
\]
したがって
\[
\begin{aligned}
\frac{93}{10}\theta_\eta^2(\vartheta\eta x)^2
&=-\frac3{10}g-\frac52s_\vartheta^2-\frac35(\vartheta\R x)^2\\
&\quad+\frac53i\A+\frac{37}{25}fa_\R^2+\frac{74}{25}ft,
\end{aligned}
\]
となり、g) に従って $\theta_\eta^2(\vartheta\eta x)^2$ の二つの値を比較すると：
\[
\tag{13}
0=g-2s_\vartheta^2+2(\vartheta\R x)^2+\frac43i\A-\frac45fa_\R^2-\frac85ft.
\]

\noindent 8) $\theta_\eta^2H_\vartheta(\theta Hu)$ は a) および b) による。これから $\theta_\eta^3(\theta Hu)$ は c) による（因子 $u_x$ をもつ形式は消える）：
\[
\theta_\eta^2H_\vartheta(\theta Hu)=-\frac16s_\vartheta(\theta su)^3-\frac13(\nu\R x),
\]
\[
\theta_\eta^2H_\vartheta(\theta Hu)=-\frac13\theta_\eta^3(\theta Hu)-\frac13(\nu\R x),
\]
\[
\theta_\eta^3(\theta Hu)=\frac12s_\vartheta(\theta su)^3.
\]

\noindent 9) $\theta_\eta^2H_\vartheta(\vartheta\eta x)$ は a) および b) による。これから $\theta_\eta^3(\vartheta\eta x)$ は c) による。また $\theta_\eta^3(\vartheta\eta x)$ は d) による。これから高次形式 $s_\vartheta^2(\theta su)$ の還元が g) に従って得られる（因子 $u_x$ をもつ形式は消える）：
\[
\theta_\eta^2H_\vartheta(\vartheta\eta x)=\frac13(atu),
\]
\[
\theta_\eta^2H_\vartheta(\vartheta\eta x)=\frac13(atu)+\frac13\theta_\eta^3(\vartheta\eta x)-\frac16s_\vartheta^2(\theta su),
\]
\[
\theta_\eta^3(\vartheta\eta x)=\frac12s_\vartheta^2(\theta su),
\]
\[
\theta_\eta^3(\vartheta\eta x)=\frac25\theta_\R(\vartheta\R x)+\frac15(atu),
\]
\[
\tag{14}
s_\vartheta^2(\theta su)=\frac45\theta_\R(\vartheta\R x)+\frac25(atu).
\]

\noindent 10) $\theta_\eta^2H_\vartheta^2$ は a) によるものと、還元子 $\theta_\nu^2(\vartheta\nu x)^2$ によるもの；$\theta_\eta^3H_\vartheta$ は a) と c) による；$\theta_\eta^4$ は c) による。ここから高次形式 $s_\vartheta^2(\theta su)^2$ および $s_\nu^2$ の還元が g) に従って得られる：
\[
\begin{aligned}
\text{a) によって}\quad
\theta_\eta^2H_\vartheta^2&=[a_\eta^2(aHu)^2]b^2+\frac13[a_\eta^4]_{bu^2}-\frac13H_\nu^2(\nu\eta x)^2\\
&=\frac49ia_\nu^2-\frac16s_\vartheta^2(\theta su)^2-\frac5{18}s_\nu^2+\frac13(atu)^2+\frac1{54}Ju_x^2.
\end{aligned}
\]
\[
\begin{aligned}
\theta_\nu^2(\vartheta\nu x)^2\text{ によって}:
\quad\theta_\eta^2H_\vartheta^2
&=[\theta_\nu^2(\vartheta\nu x)^2]_{\nu_1}+\frac13[\theta_\nu^4]_{\nu_1}\hat\nu_1x^2-\frac13s_\vartheta^2(\theta su)^2\\
&=\frac49ia_\nu^2-\frac16s_\nu^2-\frac13s_\vartheta^2(\theta su)^2+\frac13(\nu\R x)^2.
\end{aligned}
\]
\[
\begin{aligned}
\text{a) によって}\quad
\theta_\eta^3H_\vartheta&=-[a_\eta^2(aHu)]_{bu}b^2-\frac12[a_\eta^2(aHu)^2]b^2-\frac13[a_\eta^3]b+\frac1{12}[a_\eta^4]_{bu^2}\\
&\quad+\frac12u_x[a_\eta^2(aHu)^2]b^3\\
&=-\frac29ia_\nu^2+\frac1{12}s_\nu^2+\frac4{15}\theta_\R^2-\frac7{90}(\nu\R x)^2+\frac1{30}(atu)^2+\frac2{27}i^2u_x^2-\frac1{36}Ju_x^2.
\end{aligned}
\]
\[
\begin{aligned}
\text{c) によって}\quad
\theta_\eta^3H_\vartheta:
 a_\eta^3(Hab)(bHu)&=\frac12(atu)^2
=\theta_\eta^2H_\vartheta(\theta Hu)(\vartheta\eta x)+\frac14H_\nu^2(\nu\eta x)^2\\
&\quad+\frac12\theta_\eta^2H_\vartheta^2
=\frac32\theta_\eta^2H_\vartheta^2+\theta_\eta^3H_\vartheta-u_x[a_\eta^2(aHu)^2]b^3\\
&\quad+\frac14H_\nu^2(\nu\eta x)^2,
\end{aligned}
\]
従って
\[
\begin{aligned}
\theta_\eta^3H_\vartheta&=-\frac23ia_\nu^2+\frac5{12}s_\nu^2+\frac12(\nu\R x)^2-\frac12(atu)^2\\
&\quad+\frac4{27}i^2u_x^2-\frac1{12}Ju_x^2.
\end{aligned}
\]
c) によって
\[
\begin{aligned}
\theta_\eta^4&=\frac49ia_\nu^2-\frac16s_\nu^2+\frac{16}{15}\theta_\R^2\\
&\quad-\frac{14}{45}(\nu\R x)^2+\frac2{15}(atu)^2.
\end{aligned}
\]
g) により、$\theta_\eta^2H_\vartheta^2$ の二つの値から：
\[
\tag{15}
\begin{aligned}
s_\vartheta^2(\theta su)^2
&=\frac23s_\nu^2-2(atu)^2+2(\nu\R x)^2-\frac19Ju_x^2\\
&=\frac89ia_\nu^2+\frac8{15}\theta_\R^2-\frac{14}{15}(atu)^2+\frac{38}{45}(\nu\R x)^2-\frac4{27}i^2u_x^2.
\end{aligned}
\]
g) により、$\theta_\eta^3H_\vartheta$ の二つの値から：
\[
\tag{16}
s_\nu^2=\frac43ia_\nu^2+\frac45\theta_\R^2+\frac85(atu)^2-\frac{26}{15}(\nu\R x)^2+\frac16Ju_x^2-\frac29i^2u_x^2.
\]
式 (16) は法 $(ss'u)^4$ の還元の基礎を与え、式 (13) は $s$ の系の還元の基礎を与える（\S\ 17）。

\noindent\emph{帰結：}

1) $a_\nu(j\nu x)^3$, $\theta_\eta H_\vartheta$, $\theta_\eta^2(\vartheta\eta x)$, $\theta_\eta^2(\theta Hu)^2$ は還元子であり、$a_\nu(j\nu x)^2$, $H_\vartheta$, $H_\vartheta^2$ は分解可能形式である。

2) 系 II の形式 $j(\nu)=g=(\nu\eta x)^4H_x^2$ は還元可能である。

3) 唯一の反傾形式 $g$ が還元可能になるので、$\nu$ は、系 I との折り畳みにおいて、冪にも、同じ系の形式との積にも、まったく入らない。

$\theta$ は、反変式と見たときにのみ、積または冪として折り畳みに入る。対応して $H$ は、共変式と見たときにのみ入る（\S\ 13 B 参照）。




\begin{center}
\textbf{\S\ 12.\quad 七次の形式。}
\end{center}
\[
\text{次の折り畳み：}\left\{\begin{array}{l}
\theta\cdot K\text{ と }\nu,\\
K,j,\A\text{ と }H,\\
f\text{ と }L,\sigma.
\end{array}\right.
\]

\noindent\emph{既約形式：}
\[
\text{B)}\quad
\begin{array}{c|c|c|c|c}
\cdot&\cdot&(KHu)^2&\cdot&\cdot\\
\cdot&K_\eta&\cdot&K_\eta(KHu)^2&\cdot\\
(k\eta x)^2&\cdot&K_\eta^2&\cdot&\cdot\\
(k\eta x)^3&H_k(k\eta x)^2,\ K_\eta(k\eta x)^2&\cdot&\cdot&\cdot\\
\cdot&K_\eta(k\eta x)^3&\cdot&\cdot&\cdot
\end{array}
\]
\[
\text{C)}\quad
\begin{array}{cc}
(j\eta x)&H_j\\
\cdot&H_j(j\eta x)\\
\cdot&H_j^2
\end{array}
\qquad
\text{D)}\quad
\begin{array}{c}
(\A Hu)\\\A_\eta
\end{array}
\]
\[
\text{E)}\quad
\begin{array}{cccc}
a_i&(aLu)^2&(aLu)^3&\cdot\\
\cdot&\cdot&a_i(aLu)^2&a_i(aLu)^3\\
\cdot&a_i^2&\cdot&\cdot
\end{array}
\]

\noindent\emph{還元：}

\noindent A) $(\vartheta\cdot k,\nu,x)^4$ は、分解可能形式 $(\vartheta^2\nu x)^4$ に関する一回のトランスヴェクションとして分解する。

\noindent B) 形式列 $H_kK_\eta$：還元子 $H_\vartheta\theta_\eta$。

形式列 $K_\eta^2(KHu)$：還元子 $a_\eta^2(aHu)$。

形式列 $K_\eta^2(k\eta x)$：還元子 $\theta_\eta^2(\vartheta\eta x)$。

ここから生じる形式列 $K_\eta^3$ の二重還元は、より高い形式列 $(j\eta x)^3$ の還元をもたらす。

$(KHu)$、$(k\eta x)$、$K_\eta(k\eta x)$ は \S\ 3 による分解可能形式である。

$H_k(KHu)$、$K_\eta(KHu)$ は、分解可能形式 $K_\nu(k\nu x)$ および関数行列式
\[
K_\nu^2=\theta_\nu^2(a\theta u)\equiv a_\nu^2(a\theta u)\modu{\nu^2}
\]
との一回の折り畳みによって生じるものとして分解する。
$K_\nu^2$ の二重表示に対応して、分解可能形式の間の関係がある。すなわち
\[
H_k(KHu)=2K_\nu(\nu\nu_1k)=2\theta_{\nu_1}(\nu\nu_1a\theta)
=2\theta_\nu^2a_\nu-a_\nu^2\theta_\nu+f\theta_\eta(\theta Hu)^2-f\nu\theta_\vartheta^3
\modu{\nu^3},
\]
\[
\begin{aligned}
K_\eta(KHu)&=-K_\nu^2(\nu\nu_1x)
=-a_\nu^2(a\theta u)(\nu\nu_1x)\\
&=a_\eta(aHu)^2\theta+a_\nu^2\theta_\nu
=-\theta_\nu(\theta Hu)^2f-\theta_\nu^2a_\nu.
\end{aligned}
\]
従って
\[
\tag{17}
0=a_\nu^2\theta_\nu+\theta_\nu^2a_{\nu_1}+f\theta_\eta(\theta Hu)^2+\theta a_\eta(aHu)^2\modu{\nu^3}.
\]
形式 $H_k$、$H_k(k\eta x)$、$H_k^2$、$H_k^2(k\eta x)$ は、分解可能形式 $H_\vartheta$ および $H_\vartheta^2$ との一回の折り畳みによって生じるものとして分解する（\S\ 3. 2), a), b) 参照）。

実際、
\[
H_k=H_\vartheta(a\theta u)\sim [H_\vartheta]_{ua}-fH_\vartheta(\theta Hu)
\]
である（2)a の特殊化 $y=uv$）。
\[
H_k(k\eta x)=H_\vartheta a_\eta-H_\vartheta\theta_\eta f,
\qquad
H_\vartheta a_\eta=\frac54[H_\vartheta]a-\frac14H_\vartheta a_\vartheta
\quad(\S\ 3. 2)b),
\]
\[
H_\vartheta^2(a\theta u)=[H_\vartheta^2]_{ua},
\qquad
H_\vartheta^2a_\eta=[H_\vartheta^2]a=H_k^2(k\eta x)+H_\vartheta^2\theta_\eta f.
\]

\noindent C) 形式列 $(j\eta x)^3$：B) による形式列 $K_\eta^3$ の二重還元によって還元可能である。次の Ansatz による：
\[
0=a_\eta^3(a\theta u)=\theta_\eta^3(a\theta u)
=\theta_\eta+(a\theta\eta x)^3(a\theta u)-\theta_\eta^3(a\theta u)
=\frac12(j\eta x)^3\modu{(\nu^3,s,\R,t,i,J)}.
\]
同様の Ansatz は、形式列の終形式まで成り立つ：
\[
0=H_\vartheta^2(a\theta\eta x)a_\eta^3
=H_\vartheta^2(a\theta\eta x)\theta_\eta^3
-\frac12H_\vartheta^2(j\eta x)^4\modu{(\nu^3,s,\R,t,i,J)}.
\]
形式 $(j\eta x)^2$：1) 分解可能形式 $H_\vartheta^2$ に対する関係式の二重偏極によって分解する（(12.) b)）；

2) 積 $a_\nu\theta_\nu^3$ の直接計算によっても分解する。これにより高次形式 $(\A Hu)^2$ も分解する。すなわち
\[
\tag{18}
\left.\begin{array}{ll}
\text{a)}&\dfrac13(\A Hu)^2+\dfrac13(j\eta x)^2=\theta_\nu^2a_{\nu_1}^2,\\[0.6em]
\text{b)}&(j\eta x)^2=\dfrac43\theta_\nu^3a_{\nu_1},
\end{array}\right\}\modu{(\nu^3,s,\R,t,i,J)}.
\]
これは次の Ansatz による：
\[
H_\vartheta^2(a\theta u)^2-\frac13(\A Hu)^2
=2\theta_\eta^2(a\theta u)(aHu)+\theta_\nu^2a_{\nu_1}^2
=(a\theta u)(aHu)\{a_\eta-(a\theta\eta x)^2\}+\theta_\nu^2a_{\nu_1}^2,
\]
\[
\begin{aligned}
a_{\nu_1}\theta_\nu^3
&=-(a\widehat{\nu}\nu_1u)\theta_\nu^2(\theta\widehat{\nu}\nu_1u)
-\frac12(u\nu\nu_1u)\theta_\nu\theta_{\nu_1}(\vartheta\nu\nu_1u)\\
&=-\frac32\theta_\eta^2(\theta Hu)(aHu)
=-\frac32\theta_\eta^2(\theta Hu)(a\theta u)\\
&=-\frac32\{a_\eta-(a\theta\eta x)^2\}\{(aHu)-(a\theta u)\}(a\theta u).
\end{aligned}
\]

\noindent D) 形式列 $\A_\eta(\A Hu)$：還元子 $\A_\nu^2$。

$(\A Hu)^2$ は C) による分解可能形式である。

\noindent E) 形式列 $a_i^2(aLu)$：還元子 $a_\eta^2(aHu)$。形式 $(aLu)$ および $a_i(aLu)$ は、\S\ 3 により関数行列式に関するトランスヴェクションとして分解する。

\noindent F) 形式列 $a_\sigma^3$：還元子 $a_\R^3$。

形式 $a_\sigma^2$ および $a_\sigma^3$：対応する式
\[
H_\nu a_\nu^3\quad\text{および}\quad H_\nu a_\eta^3,
\qquad
H_\nu a_\nu^3a_\nu\quad\text{および}\quad H_\nu a_\eta^4
\]
の二重還元により、還元子は $a_\nu^3$ および $a_\eta^3$ である（\S\ 10. C), $\A_\nu^2$ および $\A_\nu^3$ の還元を参照）。

これから、(16.) を考慮して、高次形式 $a_\nu s_\nu(asu)^2$, $a_\nu^2s_\nu(asu)^2$ および記号 $\R,t(a_\nu,a_\eta^2)$ における形式の還元が得られる：
\[
\tag{19}
\left.\begin{array}{rcl}
a_\nu s_\nu(asu)^2&=&\dfrac13 i\theta_\nu^2-\dfrac1{18}iH,\\[0.5em]
a_\nu^2s_\nu(asu)^2&=&0,
\end{array}\right\}\modu{(\R,t)}.
\]
形式 $a_\sigma$ は、(12.)b の偏極によって、またはより短く積 $a_\nu^2\theta_\nu^3$ を次の Ansatz に従って直接展開することによって分解する：
\[
a_\nu^2\theta_{\nu_1}^3=(a_\nu^2\theta_{\nu_1})a_{\nu_1}-(a\theta\nu_1x)^2
=-2a_\eta(aHu)(a\theta H)(a\theta\eta x)+(a\theta\nu_1x)^2a_\nu^2\theta_{\nu_1}
\]
さらに $H_j(j\eta x)^2$ の還元 (C) を考慮すれば、
\[
\tag{20}a_\sigma=2a_\nu^2\theta_{\nu_1}^3\modu{(s,\R,t,i)}.
\]

\noindent\emph{帰結：}

1) $(j\eta x)^3$, $\A_\eta(\A Hu)$, $a_\sigma^2$ は還元子であり、$(j\eta x)^2$, $(\A Hu)^2$, $a_\sigma$ は分解可能形式である。

2) $f$ は、反傾形式との積としてのみ、$L$ との折り畳みに入る。従って $f$、また $K$ と $N$ も、$\sigma$ したがって $(\nu\sigma x)$ との折り畳みにまったく入らない。$j$ は反傾形式との積としてのみ入る。$\A$ は $H$ および $L$ との折り畳みに積としてまったく入らない（これは $\A_\eta(\A Hu)$ と $(\A Hu)^2$ から従う）。同様に $\sigma$、従って $(\nu\sigma x)$ は、系 I のいかなる形式との積にも折り畳みとして入らない。\S\ 11 の帰結を考慮すると、系 II の積として残るのは $H$ の冪と $H$ と $L$ の積だけである。

\begin{center}
\textbf{\S\ 13.\quad 八次の形式（系 III）。}
\end{center}
\[
\text{次の折り畳み：}\left\{\begin{array}{l}
\A\cdot j\text{ と }\nu,\\
\theta^2,\\
 f\cdot j,\\
 N\text{ と }H,\\
\theta\text{ と }L,\sigma.
\end{array}\right.
\]

\noindent\emph{既約形式：}
\[
\begin{array}{c@{\quad}l@{\qquad}c@{\quad}l}
\text{A)}&\A_\nu(j\nu x)&
\text{B)}&\begin{array}{c}(\vartheta^2\eta x)^3\\(\vartheta^2\eta x)^4\end{array}
\quad H_\vartheta(\vartheta^2\eta x)^3,
\ \theta_\eta(\vartheta^2\eta x)^3,\\[1em]
\text{C)}&(aHu)H_j,
\quad a_\eta(aHu)H_j&
\text{D)}&H_n,
\ N_\eta,
\end{array}
\]
\[
\text{E)}\quad
\begin{array}{c|c|c|c}
\cdot&(\theta Lu)^2&(\theta Lu)^3&\cdot\\
\theta_i&\cdot&L_\vartheta(\theta Lu)^2,
\ \theta_i(\theta Lu)^2&\theta_i(\theta Lu)^3\\
(\vartheta lx)^2&\theta_i^2&\cdot&\cdot\\
\cdot&\theta_i(\vartheta lx)^2&\cdot&\cdot
\end{array}
\]

\noindent\emph{還元：}

\noindent A) $\A_\nu(j\nu x)^3$：還元子 $a_\nu(j\nu x)^3$。

$\A_\nu(j\nu x)^2$ は、\S\ 3. 2)b により分解可能形式 $a_\nu(j\nu x)^2$ によって還元可能であり、還元子 $\theta_{\vartheta j}$ を考慮する。実際、\S\ 11 B) により
\[
\frac3{10}\A_\nu(j\nu x)^2+\frac35a_\nu(j\nu x)a_\vartheta(j\nu x)
+\frac1{10}a_\nu(j\nu x)^2
=\frac35\theta_\nu^3\theta_{\vartheta j}+\frac25\theta_\nu^3\theta_{\vartheta j}\theta_{\vartheta j},
\]
（還元子 $a_j$ および $\theta_{\vartheta j}$）。

\noindent B) 形式 $H_\vartheta(\vartheta'\eta x)^2$, $H_\vartheta^2(\vartheta'\eta x)$, $H_\vartheta^2(\vartheta'\eta x)^2$ は、\S\ 3. 2)b) および関係式
\[
H_\vartheta(\vartheta'\eta x)^2=2H_\vartheta a_\eta^2f-2H_\vartheta a_\eta b_\eta
\]
により、分解可能形式 $H_\vartheta$ と $H_\vartheta^2$、それぞれ $H_\vartheta a_\eta$ と $H_\vartheta^2a_\eta$ との逐次的な一回の折り畳みによって生じるものとして分解する。

\noindent C) 形式列 $a_\eta(j\eta x)^2$：還元子 $a_j$ および $(j\eta x)^3$。

形式 $a_\eta(j\eta x)$ は分解する。還元子 $a_j$ と、\S\ 3, 2)a（特殊化 $y=uv$）による分解可能形式 $(j\eta x)^2$ との一回の折り畳みである。

形式 $a_\eta H_j$ は分解する。1) 分解可能形式 $a_\nu(j\nu x)^2$ との一回の折り畳みによる。

2) 分解可能形式 $(j\eta x)^2$ との特殊な二回の折り畳みによる。これから分解可能形式間の関係が出る。

$a_\nu(j\nu x)^2$ から：
\[
0=\theta_\nu'\theta_\nu+2a_\eta H_j\modu{(s,\R,t,i,J)}.
\]
$(j\eta x)^2$ から：
\[
0=\theta_\nu'\theta_\nu-\frac32a_\eta H_j\modu{(s,\R,t,i,J)}
\]
（反変式との積は省略する）。これは次の Ansatz による：
\[
0=\frac12a_\nu(abu)^2\theta_\nu^3
+\frac12u_x(ubu)\theta_{\nu_1}^3(\theta bu)
-\frac34(\widehat{j}\eta bu)^2+\frac34H_j(\widehat{j}\eta bu)
-\frac18fH_j^2
\]
および $a_\nu(ubu)(\theta bu)\theta_{\nu_1}^3$ における記号の交換による。

\noindent D) 形式列 $N_\eta(NHu)$：還元子 $\A_\nu^2$。

$(NHu)$, $(n\eta x)$, $N_\eta(n\eta x)$, $H_\eta(NHu)$ は分解可能形式である（\S\ 12 B 参照）。$(NHu)^2$ は形式 $(\A Hu)^2$ との一回の折り畳みによって分解する。

\noindent E) 形式列 $\theta_iL_\vartheta$, $\theta_i^2(\theta Lu)^2$, $\theta_i^2(\vartheta lx)$：

還元子：$\theta_\eta H_\vartheta$, $\theta_\eta^2(\theta Hu)^2$, $\theta_\eta^2(\vartheta\eta x)$。

形式 $(\theta Lu)$, $(\vartheta lx)$, $L_\vartheta$, $L_\vartheta(\theta Lu)$, $\theta_i(\theta Lu)$, $L_\vartheta(\vartheta lx)$, $\theta_i(\vartheta lx)$, $L_\vartheta^2$, $L_\vartheta^2(\theta Lu)$, $\theta_i^2(\theta Lu)$ は分解する。（\S\ 12 B の分解可能形式に双対的に反対である。）

\noindent F) 形式列 $\theta_\sigma(\vartheta\sigma x)$：還元子 $a_\sigma^2$。

形式 $(\vartheta\sigma x)$ および $(\vartheta\sigma x)^2$ は、分解可能形式 $a_\sigma$ との一回の折り畳みによって生じるものとして分解する。二回の折り畳みによって生じる $\theta_\sigma$ も、次の形式列
\[
a_{\nu_2}^2(a\theta u)\theta_{\nu_1}^3\equiv H_j^2(j\eta x)\equiv0\modu{u_x(s,\R,t,i,J)}:
\]
により分解する：
\[
\theta_\sigma=a_\sigma(abu)^2=a_\nu^2(abu)^2\theta_\nu^3.
\]

\noindent\emph{帰結：}

$\theta^3$ または $\theta^2K$ は $H$ との折り畳みに入らない。$\theta$ は、反変式と見たときにのみ、冪または積として $L$ との折り畳みに入るが、$\sigma$ および $(\nu\sigma x)$ との折り畳みにはまったく入らない。

$N$ は系 II のいかなる形式との積にも折り畳みに入らず、また系 II の形式の冪または積とも同様に入らない。直前の諸段落の帰結を考慮すると、系 I の積として残るのは、$f\cdot j$, $\A\cdot j$, $\theta$ の冪、および $\theta\cdot K$ の積だけである。

\begin{center}
\textbf{\S\ 14.\quad 九次の形式（系 III）。}
\end{center}
\[
\text{次の折り畳み：}\left\{\begin{array}{l}
\theta\cdot K\text{ と }H,\\
K,j,\A\text{ と }L,\\
j,\A\text{ と }\sigma,\\
f\text{ と }H^2.
\end{array}\right.
\]

\noindent\emph{既約形式：}
\[
\begin{array}{c@{\quad}l@{\qquad}c@{\quad}l}
\text{A)}&(\vartheta\cdot k,\eta,x)^4,
\ H_k(\vartheta\cdot k,\eta,x)^4&
\text{B)}&(KLu)^3,
\ K_i(KLu)^3,
\ K_i^2,
\ (klx)^3,
\ K_i(Klx)^3,\\[0.7em]
\text{C)}&L_j,
\ L_j^2&
\text{D)}&\A_i,\\[0.7em]
\text{E)}&(j\sigma x)&
\text{G)}&(aH^2u)^3,
\ (aH^2u)^4,
\ a_\eta(aH^2u)^3.
\end{array}
\]

\noindent\emph{還元：}

A) 形式 $H_\vartheta(k\eta x)^3$, $H_\vartheta^2(k\eta x)^2$, $H_\vartheta^2(k\eta x)^3$ は、分解可能形式との逐次的な一回の折り畳みによって生じるものとして分解する。

B) 形式列 $K_iL_k$, $K_i^2(KLu)$, $K_i^2(klx)$：

還元子：$\theta_\eta H_\vartheta$, $a_\eta^2(aHu)$, $\theta_\eta^2(\vartheta\eta x)$。

形式 $L_i$, $K_i(KLu)$, $K_i(klx)$, $(KLu)^2$, $(klx)^2$ は、関数行列式に関するトランスヴェクションとして分解する（\S\ 3）。

形式 $L_k$, $L_k(KLu)$, $L_k(KLu)^2$, $L_k(klx)$, $L_k(klx)^2$, $L_k^2$, $L_k^2(KLu)$, $L_k^2(klx)$, $L_k^3$, $K_i(KLu)^2$, $K_i(klx)^2$ は、次の分解可能形式との一回の折り畳みによって生じるものとして分解する：
\[
L_\vartheta,
\ H_k(KHu),
\ L_\vartheta(\vartheta lx),
\ H_k^2,
\ L_\vartheta^2,
\ L_\vartheta^2(\theta Lu),
\ K_\eta(KHu),
\ \theta_i(\vartheta lx)
\]
これは \S\ 3. 2), a), b) による。

C) 形式列 $(jlx)^3$：還元子 $(j\eta x)^3$。

形式 $(jlx)^2$ および $L_j(jlx)$：分解可能形式 $(j\eta x)^2$ との一回の折り畳みによって生じる。

D) 形式列 $\A_i(\A Lu)$：還元子 $\A_\nu^2$。

形式 $(\A Lu)^2$, $(\A Lu)^3$：分解可能形式 $(\A Hu)^2$ との一回の折り畳み。

E) 形式列 $(j\sigma x)^2$：$j$ を形式列の低次形式で置き換えることにより、還元子は $a_\sigma^2$ である（\S\ 5）：
\[
a_\sigma^2(a\theta u)^2=\frac13(j\sigma x)^2,
\qquad
a_\sigma^2(a\theta\sigma x)^2=\frac13(j\sigma x)^4.
\]
F) 形式列 $\A_\sigma^2$：還元子 $a_\sigma^2$。

形式 $\A_\sigma$：$\A$ を形式列の低次形式で置き換えることにより、還元子は $a_\sigma^2$ である：
\[
a_\sigma a_\R^2=\frac23\A_\sigma\modu{\nu}.
\]

\noindent\emph{帰結：} $\sigma$ および $(\nu\sigma x)$ との折り畳みに入るのは形式 $j$ だけである。

$N$ は $L$ との折り畳みに入らない。なぜなら、$\A_i$ に対応する形式 $N_i$ は分解するからである。

\begin{center}
\textbf{\S\ 15.\quad 十次の形式（系 III）。}
\end{center}
\[
\text{次の折り畳み：}\left\{\begin{array}{l}
\theta^2,\\
\ f\cdot j\text{ と }L,\\
\theta\text{ と }H^2.
\end{array}\right.
\]

\noindent\emph{既約形式：}
\[
\begin{array}{c@{\quad}l@{\qquad}c@{\quad}l}
\text{A)}&(\vartheta^2lx)^4,
\ \theta_i(\vartheta^2lx)^4&
\text{B)}&(aLu)^2L_j,
\ a_i(aLu)^2L_j,\\[0.8em]
\text{C)}&(\theta H^2u)^3,
\ (\theta H^2u)^4,
\ H_\vartheta(\theta H^2u),
\ \theta_\eta(\theta H^2u)^3.
\end{array}
\]

\noindent\emph{還元：}

A) および B)：残りの形式は、分解可能形式との一回の折り畳みによって分解する。

C) 残りの形式は、\S\ 13 B) に従い、双対的に反対の考察によって分解する。

\noindent\emph{帰結：}

$\theta^2K$ は $L$ との折り畳みに入らない。対応して、$H^2L$ は $K$ との折り畳みに入らない。$H^3$ および $H^2L$ は $\theta$ との折り畳みに入らない。これにより \S\ 7 の折り畳み図式が証明された。

\begin{center}
\textbf{\S\ 16.\quad 11、12、13、15 次の形式（系 III）。}
\end{center}

\noindent\emph{11 次。}
\[
\text{次の折り畳み：}\left\{\begin{array}{l}
\theta\cdot K\text{ と }L,\\
K,j\text{ と }H^2,\\
j\text{ と }(\nu\sigma x),\\
f\text{ と }H\cdot L.
\end{array}\right.
\]
\noindent\emph{既約形式：}
\[
\begin{array}{c@{\quad}l@{\qquad}c@{\quad}l}
\text{A)}&(\vartheta\cdot k,l,x)^5,
\ \theta_i(\vartheta\cdot k,l,x)^5&
\text{B)}&(KH^2u)^4,
\ K_\eta(KH^2u)^4,\\[0.8em]
\text{C)}&(\nu\sigma j),&
\text{D)}&(a_\eta H\cdot L,u)^4.
\end{array}
\]

\noindent\emph{12 次。}
\[
\text{次の折り畳み：}\left\{\begin{array}{l}
\theta^3\text{ と }L,\\
\theta\text{ と }H\cdot L.
\end{array}\right.
\]
\noindent\emph{既約形式：}
\[
\text{E)}\ (\vartheta^3lx)^6,
\qquad
\text{F)}\ (\theta,H\cdot L,u)^4,
\qquad L_\vartheta(\theta,H\cdot L,u^4).
\]

\noindent\emph{13 次。}
\[
K\text{ と }H\cdot L\text{ の折り畳み。}
\]
\noindent\emph{既約形式：}
\[
\text{G)}\ (K,H\cdot L,u)^5,
\qquad K_\eta(K,H\cdot L,u)^5.
\]

\noindent\emph{15 次。}
\[
K\text{ と }H^3\text{ の折り畳み。}
\]
\noindent\emph{既約形式：}
\[
\text{H)}\ (KH^3u)^6.
\]

\noindent\emph{還元：}

形式 $H_j^2$, $H_j^2$ は、次の関係式から還元子 $L_j^3$ をもつ：
\[
(\nu lx)L_x^3=-\frac12H^2\modu{(\sigma,s)}.
\]
残りの形式の分解または還元可能性は、以前の各節の考察から従う。

\noindent\emph{帰結：}

\S\ 7 の折り畳み図式および \S\S\ 8--15 の帰結により、系 III の既約形式（117 個）はこれで尽くされた。

\clearpage

\begin{center}
{\Large\bfseries 第 IV 章。\quad 法 $(\varrho,t)$ に関する相対的完全系。}
\end{center}

\begin{center}
\textbf{\S\ 17.\quad 法 $(ss'u)^4$ および $s$ の系の還元。}
\end{center}

次に高い相対的完全系を形成するには、\S\ 7 と同様に、次の法
\[
\nu(s)=(ss'u)^4
\]
に関する $s$ の系を作り、これを法 $s$ に関する相対的完全系の形式と折り畳む必要がある。高い法として法 $(\varrho,t)$ を導入すると、次が示される：
\[
\begin{array}{ll}
\text{A)}&\text{法 }(ss'u)^4\text{ は還元可能になる，}\vphantom{\dfrac11}\\
\text{B)}&s\text{ の系はただ一つの形式 }s\text{ から成る。}
\end{array}
\]

A) 法 $(ss'u)^4$ の還元は、\S\ 11 の式 (16.) で得られた還元子 $s_\nu^2$ から直ちに従う。すなわち
\[
 s_\nu^2=\frac43 i a_\nu^2+\frac45\theta_\varrho^2+\frac85(atu)^2-
 \frac{26}{15}(\nu\varrho x)^2+\frac16J\cdot u_x^2-
 \frac29i^2\cdot u_x^2
\]
から、$x$ を $\nu_1$ に偏極して、
\[
\tag{21}
\begin{aligned}
(ss'u)^4&=-\frac23J\cdot\nu+6s_\nu^2s_{\nu_1}^2\\
&=\frac{24}{5}\theta_\nu^2\theta_\varrho^2+
\frac{48}{5}a_\nu^2(atu)^2-
\frac{52}{5}H_\varrho^2+
\frac43i\cdot(asu)^4+\frac13J\cdot\nu-
\frac49i^2\cdot\nu,
\end{aligned}
\]
従って法 $(ss'u)^4$ の還元が得られる。\footnote{式 (21.) から、序論の注で述べられた 9 次の三つの不変式の間の関係が、式 (10.)c を用いて従う。}

B)
\[
Z=(ss'u)^2=\theta(s)
\]
の還元、従って $s$ の系の還元は、\S\ 7 の式 (8.)b により、関係式
\[
 s_\eta^2=s_\nu^2(\nu\nu_1x)^2-\frac13Z
\]
を考慮して、形式 $Z$ をより低い形式 $g_\nu^2$ で置き換えることによって得られる。
形式 $g_\nu^2$ は、\S\ 11 の式 (13.) によって、還元子 $s_\nu^2$ を因子にもつ。式 (8.) および (16.) から
\[
 g_\nu^2=-Z+\frac{14}{5}ia_\eta^2+\frac{14}{15}i(asu)^2
 \modu{(\varrho,t)},
\]
また式 (13.), (14.), (15.), (16.) から
\[
\begin{aligned}
g_\nu^2&=2[s_\vartheta^2]_{\nu^2}-\frac43 i\A_\nu^2
      =2s_\vartheta^2s_\nu^2-\frac65[s_\vartheta^2(\theta su)^2]\widehat{\nu x}^{\,2}
        -\frac43i\A_\nu^2\\
&=\frac45i\cdot a_\eta^2-\frac25i\cdot(asu)^2+\frac13J\cdot\theta
\modu{(\varrho,t)}.
\end{aligned}
\]
従って
\[
\tag{23}
 Z=2i\cdot a_\eta^2+\frac43i\cdot(asu)^2-\frac13J\cdot\theta
 \modu{(\varrho,t)}.
\]

したがって、法 $((ss'u)^4,\varrho,t)$ に関する相対的完全系は、法 $(\varrho,t)$ に関する相対的完全系に移る。そしてここから、二つの二次形式の既知の系に関するトランスヴェクションによって、絶対的完全系が生じる。法 $(\varrho,t)$ に関する相対的完全系を形成するには、式 (23.) によって $s$ の系が形式 $s$ に還元されるので、法 $(s,\varrho,t)$ に関する相対的完全系を $s$ および $s$ の冪に関してトランスヴェクトしなければならない。$s$ の冪は系 I との折り畳みにのみ入ることが分かる。

予告した関係式は次の通りである：
\[
\tag{22}
\begin{aligned}
(ass')^4&=\frac{24}{5}\theta_\nu^2\theta_\varrho^2a_\nu^2a_\varrho^2+
\frac{48}{5}a_\nu^2b_\nu^2(tab)^2-
\frac{52}{5}H_\varrho^2a_\nu^4+
\frac53J\cdot i-\frac49i^3,\\
(ass')^4&=\frac{24}{5}a_\varrho^2a_{\varrho'}^2-
\frac{12}{5}t_\varrho^2+\frac53J\cdot i-\frac49i^3.
\end{aligned}
\]

同じ位数かつ同じ次数の形式のうち、「高次形式」とは、法 $(s,\varrho,t)$ に関する相対的完全系の高次形式から $s$ との折り畳みによって生じたものをいう。ただし、$s$ との折り畳みは元の形式の偏極にすぎないものと見る。これにより、個々の形式の配列は一意に定まる（\S\ 1、形式列 2 参照）。例えば
\[
(\theta_\eta(\theta Hu),s,u)^2\text{ は }s_\eta((\theta Hu)^2,s,u)\text{ より高い，}
\]
\[
(\theta_\eta(\theta Hu)^2,s,u)\text{ は }(\theta_\eta(\theta Hu),s,u)^2\text{ より高い。}
\]

得られる系を四つの部分系に分ける：
\[
\begin{array}{ll}
\text{系 }s_I:&s\text{ および }s\text{ の冪と系 I との折り畳みによって生じる；}\vphantom{\dfrac11}\\
\text{系 }s_{II}:&s\text{ と系 II との折り畳みによって生じる；}\vphantom{\dfrac11}\\
\text{系 }s_{III}:&s\text{ と系 III の個々の形式}\\
&\text{（積を除く）との折り畳み、および系 I と II から}\\
&\text{それぞれ一つずつ取った形式の積との折り畳みによって生じる；}\vphantom{\dfrac11}\\
\text{系 }s_{IV}:&s\text{ と、系 I と III、II と III、III と III からそれぞれ}\\
&\text{一つずつ取った形式の積との折り畳みによって生じる。}
\end{array}
\]
注意：以後、すべての公式は法 $(\varrho,t)$ で理解する。式 (27.) 以降は、これを毎回明記せずに、法 $(\varrho,t,i,J,\text{高次形式})$ で理解する。

還元のため、以前に得た公式をまとめる：
\[
\tag{24}
\begin{aligned}
s_\nu^2&=\frac43ia_\nu^2+\frac16J\cdot u_x^2-\frac29i^2\cdot u_x^2,
& s_\nu^3&=\frac13J\cdot u_x,
& s_\nu^4&=J,\\[0.3em]
g&=2s_\vartheta^2-\frac43i\A,
& s_\vartheta^2(\theta su)&=0,\\[0.3em]
s_\vartheta^2(\theta su)^2&=\frac89i\cdot a_\nu^3-\frac4{27}i^2\cdot u_x^2,\\[0.3em]
Z&=2i\cdot a_\eta^2+\frac43i(asu)^2-\frac13J\cdot\theta,\\[0.3em]
g_\nu&=-s_\eta+u_x\left\{2ia_\eta^2-\frac23i(asu)^2+\frac23J\cdot\theta\right\},\\[0.3em]
g_\nu^2&=\frac45ia_\eta^2-\frac25i(asu)^2+\frac13J\cdot\theta,
&g_\nu^3&=0,
&g_\nu^4&=0.
\end{aligned}
\]

\noindent\emph{帰結：}
\[
 s_\nu^3,
\quad s_\vartheta^2(\theta su),
\quad s_\vartheta^2s_\nu,
\quad s_\vartheta^2\theta_\nu
\]
は還元子である。



\clearpage
\begin{center}
\textbf{\S\ 18.\quad 系 $s_I$.}
\end{center}

\[
\text{次の折り畳み：}\quad
\left\{\begin{array}{l}
s\text{ と } f;\ \theta;\ K,j,\A;\ f\cdot j,N;\ \A\cdot j,\\
s^3\text{ と }\theta,K.
\end{array}\right.
\]

\noindent\emph{既約形式。}

\noindent 5次：
\[
(asu),\qquad (asu)^2,\qquad (asu)^3,\qquad (asu)^4.
\]

\noindent 6次：
\[
\begin{array}{cccc}
(\theta su)&(\theta su)^2&(\theta su)^3&(\theta su)^4\\
s_\vartheta&s_\vartheta(\theta su)&s_\vartheta(\theta su)^2&s_\vartheta(\theta su)^3\\
\cdot&s_\vartheta^2&\cdot&\cdot
\end{array}
\]

\noindent 7次：
\[
\begin{array}{c@{\qquad}c@{\qquad}c@{\qquad}c}
\cdot&(Ksu)^2&(Ksu)^3&(Ksu)^4\\
\text{A)}\ s_k&s_k(Ksu)&s_k(Ksu)^2&s_k(Ksu)^3\\
\cdot&s_k^2&\cdot&\cdot
\end{array}
\qquad
\begin{array}{l}
\text{B)}\ s_j,\\[0.4em]
\text{C)}\ (\A su),\ (\A su)^2.
\end{array}
\]

\noindent 8次：
\[
\begin{array}{ll}
\text{A)}&s_j(asu),\ s_j(asu)^2,\ s_j(asu)^3,\\[0.4em]
\text{B)}&(Nsu)^2,\quad s_n,\quad s_n(Nsu).
\end{array}
\]

\noindent 10次：
\[
\text{A)}\ s_j(\A su),\qquad \text{B)}\ s_\vartheta(\theta s'u)^4.
\]

\noindent 11次：
\[
(Ks^2u)^6,\qquad s_k(Ks^2u)^5,
\qquad s_k(Ks^2u)^6.
\]

\noindent\emph{還元：}

形式列
\[
 s_\vartheta^2(\theta su),\quad s_k^2(Ksu),\quad s_j^2,\quad (\A su)^3,
 \quad (Nsu)^3
\]
は、還元子 $s_\vartheta^2(\theta su)$ を因子にもつ。$s_j^2$ と $(\A su)^3$ は、$j$ と $\A$ を形式列の低次形式へ戻すことによって得られる：
\[
\begin{aligned}
s_\vartheta^2(\theta su)(a\theta u)^3&=-\frac12s_j^2,\\
s_\vartheta^2(\theta su)(a\theta u)^2&=\frac13(\A su)^3,\\
s_\vartheta^2(\theta su)^2(a\theta u)^2&=\frac13(\A su)^4+\frac13s_j^2.
\end{aligned}
\]
形式 $s_\vartheta^2s_{\vartheta'}$ と $s_\vartheta^2s_{\vartheta'}^2$：還元子 $s_\vartheta^2(\theta su)$ と $\theta_{\vartheta'}$：
\[
\begin{aligned}
s_\vartheta^2s_{\vartheta'}&=s_\vartheta^2\theta_{\vartheta'}-s_\vartheta^2(\theta s\widehat{\vartheta'}x),\\
s_\vartheta^2s_{\vartheta'}^2&=s_\vartheta^2s_{\vartheta'}\theta_{\vartheta'}-s_\vartheta^2s_{\vartheta'}(\theta s\widehat{\vartheta'}x).
\end{aligned}
\]
形式列 $s_j(\A su)^2$：還元子 $\A_j$ と $(\A su)^3$。

\noindent\emph{帰結：}

$s$ は $f,j,\A,N$ との折り畳みにおいて冪としては入らない。$f,j,\A,N$ は、反傾形式との積においてだけ折り畳みに入る。ただし、$f$ との積ではそれらの形式はなお $x$ に関して二次でありうるし、$\A$ との積ではなお $x$ に関して一次でありうる。従って、次の形式の積を考えればよい：
\[
 f\cdot(0,\lambda),\quad f\cdot(1,\lambda),\quad f\cdot(2,\lambda),\quad
 j\cdot(\mu,0),\quad N,\quad
 \A\cdot(0,\lambda),\quad \A\cdot(1,\lambda)\quad(\lambda>0),
\]
ここで $(\mu,\lambda)$ は、$x$ に関して次数 $\mu$、$u$ に関して次数 $\lambda$ の形式を表す。

$\theta$ も $K$ も、$s$ または $s^2$ との折り畳みにおいて、任意の形式との冪または積には入らない。実際、関数行列式 $K$ を含む積 $K\cdot\varphi_x$（$\varphi\ne f$ または $\theta$）については、分解可能形式間の関係式
\[
 K\cdot\varphi_x=(a\theta u)\varphi_x=f\cdot(\varphi\theta u)-\theta\cdot(\varphi au)
\]
が成り立つ。従って、上に掲げた系 $s_I$ の折り畳みの図式が証明された。

\begin{center}
\textbf{\S\ 19.\quad 系 $s_{II}$.}
\end{center}

$s$ と $\nu;\ H;\ L;\ H^2$ との折り畳み。

\noindent\emph{既約形式：}
\[
\begin{array}{c@{\qquad}c@{\qquad}c@{\qquad}c}
\text{6次} & \text{8次} & \text{10次} & \text{12次}\\[0.3em]
s_\nu & (Hsu),\ (Hsu)^2 & (Lsu)^2,\ (Lsu)^3 & (H^2su)^3\\
\cdot & s_\eta & s_l & \cdot\\
\cdot & s_\nu^4=J & \cdot & \cdot
\end{array}
\]

\noindent\emph{還元：}

形式列 $s_\eta(Hsu)$ と $s_l(Lsu)$：還元子 $s_\nu^2$。

形式列 $s_\sigma$：$H$ から来る還元子 $s_\nu^2$、すなわち $s_\nu^2=\dfrac23s_\sigma$。

$(H^2su)^4$ は式 (7.)a により分解する（\S\ 11 A 参照）。従って $(H\cdot L,s,u)^4$ も分解する。

\noindent\emph{帰結：}

$s^2$ は系 II との折り畳みに入らない。形式 $\nu$ は反傾形式との積においてだけ入る。形式 $H$ は、共変形式として見れば、任意の $\lambda$ について形式 $(1,\lambda)$、$(2,\lambda)$、$(3,\lambda)$ との積において折り畳みに入る。形式 $\sigma$ と $(\nu\sigma x)$ は全く入らず、$L$ は関数行列式としては折り畳みの積に入らない。系 III の共変形式に関する完全な超過は還元可能になる。系 I と II の積として残るのは
\[
 f\cdot\nu,\qquad \A\cdot\nu.
\]

\begin{center}
\textbf{\S\ 20.\quad 七次の形式（系 $s_{III}$）。}
\end{center}

$s$ と $f\cdot\nu$、$a_\nu$、$a_\nu^2$ との折り畳み。

\noindent\emph{既約形式：}
\[
\begin{gathered}
s_\nu(asu),\quad a_\nu(asu),\quad s_\nu(asu)^2,\quad a_\nu(asu)^2,
\quad s_\nu(asu)^3,\quad a_\nu(asu)^3,\\
a_\nu s_\nu,
\quad a_\nu s_\nu(asu),
\quad a_\nu^2(asu),
\quad a_\nu^2(asu)^2.
\end{gathered}
\]

\noindent\emph{還元：}

還元子 $s_\nu^2$ と $s_\eta(Hsu)$ との直接の折り畳みから生じる自明な還元は述べない。また、関数行列式を含むトランスヴェクションの分解も述べない。

形式 $a_\nu s_\nu(asu)^2$ と $a_\nu^2s_\nu(asu)^2$ については式 (19.)、\S\ 12 を見よ。さらに
\[
 a_\nu s_\nu(asu)^3=a_\nu(a\widehat{\nu}_1\nu_2u)^3(\nu_1\nu_2\nu)=0.
\]
形式列 $a_\nu^2s_\nu$：$a_\nu^2$ を低次形式へ戻すこと、すなわち式
\[
 s_\vartheta^2(a\theta s)(a\theta u),
 \qquad s_\vartheta^2(a\theta s)(a\theta u)^2
\]
の二重還元によって得られる還元子 $s_\vartheta^2(\theta su)$。Ansatz は
\[
\begin{aligned}
s_\vartheta^2(a\theta s)(a\theta u)
&=[(a\theta u)^2]s^3+\frac15[a_\vartheta(a\theta u)^2]s^2+\frac1{60}[a_\vartheta^2(a\theta u)^2]s\\
&=\frac12a_\nu^2s_\nu-\frac23a_\nu^2s_\nu^2,\\
s_\vartheta^2(a\theta s)(a\theta u)^2
&=\frac45[a_\vartheta(a\theta u)^2]_{us}s^2+\frac{23}{180}[a_\vartheta^2(a\theta u)^2]_{us}s\\
&=-\frac12a_\nu^2s_\nu(asu).
\end{aligned}
\]
ここで $a_\nu^2b_\nu(abu)=0$ である。

得られる公式は
\[
\tag{25}
\begin{array}{rcl@{\qquad}rcl}
a_\nu^2s_\nu&=&\dfrac49i\,a_\nu^2b_\nu,
& a_\nu s_\nu(asu)^2&=&\dfrac13i\theta_\nu^2-\dfrac1{18}iH,\\[0.7em]
a_\nu s_\nu(asu)^3&=&0,
& a_\nu^2s_\nu(asu)&=&0,\\[0.7em]
a_\nu^2s_\nu(asu)^2&=&0.
\end{array}
\]

\noindent\emph{帰結：}

1) $a_\nu s_\nu(asu)^2$ と $a_\nu^2s_\nu$ は還元子である。

2) \S\ 19 で既に還元可能と認められた形式 $(\A su)^3$、$(\A su)^4$ の値が得られる。同様に、対応する形式列 $(agu)^2$ についても、$\nu$ との折り畳みに際して還元子 $g_\nu$ により二重還元が生じる：
\[
\tag{26}
\begin{aligned}
\text{a)}\quad (\A su)^3&=-\frac65a_\nu^2(asu)+3a_\nu s_\nu(asu)+2i\theta_\nu(\theta\nu x),\\
\text{b)}\quad (\A su)^4&=-\frac25a_\nu^2(asu)^2+\frac43i\theta_\nu^2+\frac19iH.
\end{aligned}
\]
式 (24.) と (26.) から、法 $(i,J)$ で（p. 71 参照）、
\[
\tag{27}
\begin{aligned}
\text{a)}\quad (agu)^2&=\frac23(\A su)^2-\frac43a_\nu s_\nu+\frac35s\cdot a_\nu^2,\\
\text{b)}\quad (agu)^3&=2a_\nu s_\nu(asu)-a_\nu^2(asu)^2,\\
\text{c)}\quad (agu)^4&=-a_\nu^2(asu)^2.
\end{aligned}
\]

\begin{center}
\textbf{\S\ 21.\quad 八次の形式（系 $s_{III}$）。}
\end{center}

$s$ と四次の形式（系 III）との折り畳み（\S\ 9）。

\noindent\emph{既約形式：}
\[
\begin{array}{lll}
(\vartheta\nu x)s_\nu,& ((\vartheta\nu x),s,u)^2,\\[-0.1em]
((\vartheta\nu x)^2,s,u),& ((\vartheta\nu x)^2,s,u)^2,\ \theta s_\nu,\\[-0.1em]
(\vartheta\nu x)^2s_\nu,& ((\vartheta\nu x)^2,s,u)s_\nu,
\end{array}
\]
\[
\begin{array}{lll}
((\vartheta\nu x),s,u)^3,\ (\theta_\nu su)^2,&
((\vartheta\nu x),s,u)^4,\ (\theta_\nu su)^3,\\[-0.1em]
((\vartheta\nu x)^2,s,u)^3,\ (\theta_\nu(\vartheta\nu x),s,u)^2,\ (\theta_\nu^2su),&
((\vartheta\nu x),s,u)^3,\ (\theta_\nu^2su)^2,\ (\theta_\nu(\vartheta\nu x),s,u)^4.
\end{array}
\]

\noindent\emph{還元：}

$((\vartheta\nu x),s,u)s_\nu$ は、\S\ 3 により関数行列式に関するトランスヴェクションとして分解する。

形式列 $((\vartheta\nu x),s,u)^2s_\nu$：還元子 $s_\nu^2$ と $s_\vartheta^2\theta_\nu$：
\[
(\widehat\vartheta\nu su)s_\nu(\theta su)=s_\nu s_\vartheta(\theta su)=-\frac12(\widehat\vartheta\nu su)^2(\theta su),
\]
\[
\theta_\nu(\widehat\vartheta\nu su)s_\nu=\theta_\nu s_\nu s_\vartheta=-\frac12\theta_\nu(\widehat\vartheta\nu su)^2.
\]
形式列 $\theta_\nu s_\nu(\theta su)$：次の式の二重還元による還元子 $a_\nu s_\nu(asu)^2$：
\[
a_\nu s_\nu(asu)^2(abu)\quad\text{および}\quad a_\nu s_\nu(asu)^2(abu)(sbu);
\]
$\theta_\nu s_\nu(\theta su)^3$ は $(agu)^3(abu)(bgu)^3$ の二重還元による。

形式列 $\theta_\nu(\vartheta\nu x)s_\nu$：$\theta_\nu(\vartheta\nu x)s_\nu=a_\nu^2(abu)s_\nu$ からの還元子 $a_\nu^2s_\nu$。

形式列 $(\theta_\nu(\vartheta\nu x)^2,s,u)$：形式列 $\theta_\nu(\widehat\vartheta\nu su)s_\nu$ の二重還元。この列は同時に $((\vartheta\nu x)su)^2s_\nu$ と $\theta_\nu(\vartheta\nu x)s_\nu$ の形式列に属する。

形式 $(\theta_\nu(\vartheta\nu x),s,u)$ は、関数行列式 $\theta_\nu(\vartheta\nu x)=a_\nu^2(abu)$ に関する一回のトランスヴェクションとして分解する。

形式 $((\vartheta\nu x)^2,s,u)^4$：式 (27.) から、かつ $(j\eta x)^4$ の還元（\S\ 12 C）と同様に、式 $(a\A u)^2(\A su)^4$、または $(\theta gu)^4$ の二重還元。

二重還元により次の公式が得られる：
\[
\tag{28}
\begin{aligned}
0&=\theta_\nu s_\nu(\theta su)-\frac16(\widehat\vartheta\nu su)^2(\theta su)-\frac1{12}(Hsu),\\
0&=\theta_\nu s(\theta su)^2-\frac14(\widehat\vartheta\nu su)^2(\theta su)^2-\frac1{24}(Hsu)^2,\\
0&=(\widehat\vartheta\nu su)^2(\theta su)^2+2\theta_\nu(\widehat\vartheta\nu su)(\theta su)^2
  +3\theta_\nu^2(\theta su)^2-\frac13H(su)^2,\\
0&=\theta_\nu s_\nu(\theta su)^3+\frac12\theta_\nu(\widehat\vartheta\nu su)(\theta su)^3,\\
0&=\theta_\nu s_\nu s_\vartheta-\frac14s_\eta,
\quad 0=\theta_\nu s_\nu s_\vartheta(\theta su),
\quad 0=\theta_\nu s_\nu s_\vartheta(\theta su)^2.
\end{aligned}
\]
最後の群は $((\theta_\nu(\vartheta\nu x)^2,s,u)^2,\ldots)$ に対応する。

\noindent\emph{帰結：}

1) $\theta(\vartheta\nu x)s_\nu$、$(\widehat\vartheta\nu su)(\theta su)s_\nu$、$\theta_\nu(\widehat\vartheta\nu su)^2$、$(\widehat\vartheta\nu su)^2(\theta su)^2$ は還元子である。

2) 四次の形式 $(5,\lambda)$、$(6,\lambda)$ などは $s^2$ との折り畳みに入らない。

3) 形式列 $(\theta gu)^2$ について、式 (27.) と本節の還元を合わせると次の公式が得られる：
\[
\tag{29}
\begin{aligned}
\text{a)}\quad (\theta gu)^2&=\frac43\theta_\nu s_\nu+\frac23s_\nu s_\vartheta-\frac13s\theta_\nu^2-\frac19sH-\frac13f\,a_\nu^2(asu)^2+\frac13a_\nu^2(asu)^2,\\
\text{b)}\quad (\theta gu)^3&=-\frac13(\widehat\vartheta\nu su)^2(\theta su)-\theta_\nu(\widehat\vartheta\nu su)(\theta su)-\theta_\nu^2(\theta su),\\
\text{c)}\quad (\theta gu)^4&=2\theta_\nu^2(\theta su)^2-\frac13(Hsu)^2,\\
g\vartheta(\theta gu)&=(\widehat\vartheta\nu su)(\vartheta\nu x)s_\nu-\theta_\nu(\vartheta\nu x)(\widehat\vartheta\nu su),\\
g\vartheta(\theta gu)^2&=\frac23s\theta_\nu^3,
\quad g\vartheta(\theta gu)^3=\frac13\theta_\nu^3(\theta su),
\quad g\vartheta(\theta gu)^4=0.
\end{aligned}
\]

\begin{center}
\textbf{\S\ 22.\quad 九次の形式（系 $s_{III}$）。}
\end{center}

$s$ と五次の形式（系 III）および積 $\A\cdot\nu$ との折り畳み（\S\ 10）。

\noindent\emph{既約形式：}
\[
\begin{array}{ll}
\text{A)}& (k\nu x)^2s_\nu,
((k\nu x)^3,s,u),
(k\nu x)^3s_\nu,
((k\nu x)^2,s,u)^2,
K_\nu s_\nu,\\
& ((k\nu x)^3,s,u)^2,
((k\nu x)^2,s,u)s_\nu,
(K_\nu(k\nu x)^3,s,u),\\
& (K_\nu su)^2,\ (K_\nu su)^3,\ (K_\nu su)^4,
((k\nu x)^2,s,u)^3,
(K_\nu^2su)^4,\\[0.3em]
\text{B)}& (j\nu x)s_\nu,
((j\nu x)^2,s,u),
(j\nu x)^3,s,u,
((j\nu x)^2,s,u)^2,\\[0.3em]
\text{C)}& s_\nu(\A su),\quad (\A_\nu su),\\[0.3em]
\text{D)}& (aHu)s_\eta,
(a_\eta su),
((aHu),s,u)^2,
((aHu)^2,s,u),\\
& (aHu)^2s_\eta,
(a_\eta su)^2,
(a_\eta^2su),
((aHu),s,u)^3,
((aHu)^2,s,u)^2,\\
& ((aHu),s,u)^4,
(a_\eta su)^3,
(a_\eta(aHu),s,u)^2,
(a_\eta(aHu)^2,s,u),
(a_\eta(aHu),s,u)^3.
\end{array}
\]

\noindent\emph{還元：}

A) 還元は、\S\ 21 の還元から一回の折り畳みによって直接従う。形式列 $K_\nu s_\nu(Ksu)^2$ は還元子 $a_\nu s_\nu(asu)^2$ と $\theta_\nu s_\nu(\theta su)^2$ を因子にもつ。従って高次形式 $K_\nu^2(Ksu)^2$、$K_\nu^2(Ksu)^3$ は Ansatz
\[
0=a_\nu s_\nu(asu)^2(a\theta u)=K_\nu s_\nu(Ksu)^2-\theta_\nu s_\nu(\theta su)^2(a\theta u)
\]
に従って分解する。

B) 形式列 $(j\nu x)^2s_\nu$：
\[
(j\nu x)^2s_\nu=3a_\nu^2s_\nu(a\theta u)^2
\]
から得られる還元子 $a_\nu^2s_\nu$。
$((j\nu x)^3,s,u)^2$：
\[
((j\nu x)^3,s,u)^2=-6a_\nu^2s_\nu(a\theta u)^2(\theta su)
\]
から得られる還元子 $a_\nu^2s_\nu$。

C) 形式 $\A_\nu(\A su)^2$ と $s_\nu(\A su)^2$：式 (27.)a から得られる還元子 $g_\nu$。

形式 $\A_\nu s_\nu$：
1) $a_g a_\nu^2s_\nu=\frac23\A_\nu s_\nu$ から得られる還元子 $a_\nu^2s_\nu$；
2) 式 (27.)a により、式 $(\A s\widehat\nu x)^2$ の二重還元を通じて分解する。

従って高次形式 $a_\eta s_\eta$ が分解する。

D) 形式列 $(aHu)^2s_\eta(asu)$：形式列 $[a_\nu s_\nu(asu)^2]_{\nu_1x^2}$ の二重還元による還元子 $a_\nu s_\nu(asu)^2$；$(aHu)^2s_\eta(asu)^2$ は $a_\nu s_\nu(asu)^2$ と $a_\nu s_\nu(asu)^3$ から還元される。従って $(a_\eta,s,u)^4$ が還元される。

形式列 $a_\eta(aHu)s_\eta$：式 $(\A s\widehat\nu x)^3$ の二重還元による還元子 $a_\nu^2s_\nu$、$a_\eta^2s_\eta$。ただし $a_\eta^2s_\eta$ は、形式 $a_\eta(aHu)s_\eta$ の還元可能項 $a_\nu^2s_\nu(\nu\nu_1x)$ からの折り畳みでは生じない。

形式列 $(a_\eta^2su)^2$：
\[
a_\nu^2s_\nu s_{\nu_1}=-\frac12(a_\eta^2(Hsu))^2
\]
からの還元子 $a_\nu^2s_\nu$。
形式 $(a_\eta(aHu),s,u)$ は関数行列式に関するトランスヴェクションとして分解する。

得られる還元公式は
\[
\tag{30}
\begin{aligned}
\text{a)}\quad 0&=(aHu)^2(asu)s_\eta-a_\eta(asu)(Hsu)^2-a_\eta(aHu)(asu)(Hsu),\\
\text{b)}\quad 0&=(aHu)^2(asu)^2s_\eta-a_\eta(aHu)(asu)^2(Hsu),\\
\text{c)}\quad 0&=a_\eta(asu)^2(Hsu)^2,\\
0&=a_\eta s_\eta+\frac14a_\nu^2g-\frac14s\cdot a_\eta^2,\\
0&=a_\eta(aHu)s_\eta-\frac12a_\eta^2(Hsu),
\quad 0=a_\eta^2(Hsu)^2,\ldots,
\quad 0=a_\eta^2s_\eta.
\end{aligned}
\]

\noindent\emph{帰結：}

1) $(j\nu x)^2s_\nu$、$\A_\nu s_\nu$、$\A_\nu(\A su)^2$、従って $N_\nu(Nsu)^2$、$(aHu)^2(asu)s_\eta$、$a_\eta(aHu)s_\eta$、$a_\eta^2(Hsu)^2$ は還元子である。$a_\eta s_\eta$ は分解可能形式であり、すべての一回および二重の折り畳みで分解可能形式へ移る。

2) 五次の形式 $(5,\lambda)$、$(6,\lambda)$ などは $s^2$ との折り畳みに入らない。

\begin{center}\textbf{\S\ 23.\quad 十次の形式（系 $s_{III}$）。}\end{center}
$s$ と六次の形式（系 III）との折り畳み（\S\ 11）。

\noindent\emph{既約形式：}
\[
\begin{array}{rl}
\text{A)}& (\vartheta^2\nu x)^3s_\nu,
((\vartheta^2\nu x)^3,s,u),
((\vartheta^2\nu x)^3,s,u)^2,
(\vartheta^2\nu x)^3s_\nu s_\vartheta,
\theta_\nu(\vartheta^2\nu x)^3s_\vartheta,\\[0.3em]
\text{B)}& a_\nu(j\nu x)s_\nu,
(a_\nu(j\nu x),s,u)^2,
(a_\nu(j\nu x),s,u)^3,
(a_\nu(j\nu x),s,u)^4,
(a_\nu^2(j\nu x),s,u)^2,
(a_\nu^2(j\nu x),s,u)^3,\\[0.3em]
\text{C)}& ((\vartheta\eta x)^2,s,u),
((\vartheta\eta x),s,u)^2,
(\theta Hu)s_\eta,
(\theta_\eta su),
((\theta Hu),s,u)^2,
((\theta Hu)^2,s,u),\\
& ((\theta\eta x),s,u)^3,
(\theta_\eta su)^2,
(\theta_\eta^2su),
((\theta Hu),s,u)^3,
((\theta Hu)^2,s,u)^2,
(H_\vartheta(\theta Hu),s,u)^2,\\
& (\theta_\eta(\theta Hu),s,u)^2,
(\theta_\eta(\theta Hu)^2,s,u),
((\theta Hu),s,u)^4,
(\theta_\eta su)^4,
(\theta_\eta(\theta Hu),s,u)^3,
(\theta_\eta^2(\theta Hu),s,u)^2.
\end{array}
\]
\noindent\emph{還元：}
A) 形式 $((\vartheta^2\nu x)^3,s,u)^2$、$((\vartheta^2\nu x)^3,s,u)^3$ は分解する：
\[
(\vartheta\nu x)(\widehat{\vartheta\nu su})(\widehat{\vartheta\nu su})=(\vartheta\nu x)s_\vartheta s_\nu-(\vartheta\nu x)s_\nu s_\vartheta=-2\theta(\vartheta\nu x)s_\nu s_\vartheta+(\vartheta\nu x)s_\vartheta^2.
\]
残りの形式は、還元子を因子にもつか、あるいは分解された形式との一回の折り畳みである。

B) 還元子 $a_\nu^2s_\nu$ と $(\widehat{j\nu su})s_\nu$。

C) 1. 形式列 $(\theta Hu)^2s_\eta$：a) 還元子 $(aHu)^2(asu)s_\eta$；b) 形式列 $(\theta gu)^3g_\nu$ と $g_\nu(agu)^3(bgu)^3(abu)$ の二重還元。従って高次形式 $(\theta,s,u)^3$、$(H_\vartheta(\theta Hu),s,u)^3$、$(a_\nu b_{\nu_1}^2,s,u)^4$［後者は $(H_\vartheta su)^4$ を置き換える形式］の還元が得られる。

2. $(\vartheta\eta x,s,u)^4$：還元子 $(a_\eta su)^4$。

3. $((\vartheta\eta x)^2,s,u)^3$：$(\widehat{\vartheta\eta su})^2(Hsu)$ からの還元子 $s_\vartheta^2s_\nu$。形式 $(\vartheta\eta x)s_\eta$、$(\vartheta\eta x)^2s_\eta$、$((\vartheta\eta x)^2,s,u)^2$、$\theta_\eta s_\eta$ は、分解された形式 $a_\eta s_\eta$ との折り畳みによって分解する。

4. 形式列 $H_\vartheta(\theta Hu)s_\eta$、$\theta_\eta(\theta Hu)s_\eta$、$H_\vartheta(\vartheta\eta x)s_\eta$、$\theta_\eta(\vartheta\eta x)s_\eta$：還元子 $a_\eta(aHu)s_\eta$ と $a_\eta^2s_\eta$。

形式列 $(\theta_\eta(\vartheta\eta x),s,u)^2$：還元子 $a_\eta^2(Hsu)^2$［ただし、折り畳みによって項 $a_\eta^2(\theta su)(Hsu)$ から生じる $(\theta_\eta^2(\theta Hu),s,u)^2$ は除く］。

5. $(H_\vartheta(\vartheta\eta x),s,u)$、$(H_\vartheta(\vartheta\eta x),s,u)^2$：$a_\eta(aHu)s_\eta(abu)$ と $g_\vartheta(\theta gu)g_\nu$ の二重還元。

形式列 $(H_\vartheta(\vartheta\eta x),s,u)^3$：還元子 $((\vartheta\nu x)^2,s,u)^2s_\nu$。

6. $(H_\vartheta(\theta Hu),s,u)$ は、\S\ 3. 2), c) に従って、分解された形式 $(\theta_\nu(\vartheta\nu x),s,u)$ との一回の折り畳みによって分解する。すなわち、低次の分解された形式 $(H_\vartheta su)^2$ の二重還元による：\footnote{記号 $\equiv$ は、より低次数の形式の積を省略したことを意味する。}
\begin{align*}
0&=\theta_\nu(\vartheta\nu_1)(\theta su)+\frac23\theta_\nu^2(\vartheta\nu_1x)(\theta su)+\theta_\nu(\vartheta\nu x)(\theta su)s_{\nu_1}\\
&\equiv -\frac12H_\vartheta(\theta Hu)(\theta su)+\theta_\eta(\theta su)(Hsu)+\frac12H_\vartheta(\theta su)(Hsu)\\
&\equiv -\frac12H_\vartheta(\theta Hu)(\theta su)+\theta_\eta(\theta su)(Hsu)+\frac12[H_\vartheta]_{sx}^2-\frac12\frac35H_\vartheta(\theta Hu)(\theta su)\\
&\equiv -\frac35H_\vartheta(\theta Hu)(\theta su).
\end{align*}
1. と 5. に従う二重還元により、次の公式が得られる：
\begin{equation}
\tag{31}
\begin{aligned}
0&=\theta_\eta(\theta su)(Hsu)^2+\frac14H_\vartheta(\theta Hu)(\theta su)^2\\
&\quad-\frac34\theta_\eta(\theta Hu)(\theta su)(Hsu)-\frac54\theta_\eta(\theta Hu)^2(\theta su)^3\\
&\quad-(asu)^3b_\eta(\theta Hu)^2-a_\nu(asu)^3b_\nu^2,\\
0&=H_\vartheta(\theta Hu)(\theta su)^3-7\theta_\eta(\theta Hu)(\theta su)^2(Hsu),\\
0&=a_\nu(asu)^2b_{\nu_1}^{2}(bsu)^2-\theta_\eta(\theta su)^2(Hsu)^2\\
&\quad+6\theta_\eta(\theta Hu)(\theta su)^2(Hsu),\\
0&\equiv (H_\vartheta(\vartheta\eta x),s,u),
&0&\equiv H_\vartheta(\widehat{\vartheta\eta su})(Hsu)-4\theta_\eta^2(Hsu).
\end{aligned}
\end{equation}
\noindent\emph{帰結：}
\begin{enumerate}
\item $(\theta Hu)^2s_\eta$、$H_\vartheta(\theta Hu)s_\eta$、$\theta_\eta(\theta Hu)s_\eta$、$H_\vartheta(\vartheta\eta x)s_\eta$、$\theta_\eta(\vartheta\eta x)s_\eta$、$(H_\vartheta(\vartheta\eta x),s,u)$、$(\theta_\eta(\vartheta\eta x),s,u)^2$ は還元子である。
\item $s^2$ は、6次の形式 $(5,\lambda)$ などとの折り畳みに入らない。
\end{enumerate}

\begin{center}\textbf{\S\ 24.\quad 十一次の形式（系 $s_{III}$）。}\end{center}
$s$ と七次の形式（系 III）との折り畳み（\S\ 12）。

\noindent\emph{既約形式：}
\[
\begin{array}{rl}
\text{A)}& ((k\eta x)^3,s,u),\quad (K_\eta(k\eta x)^3,s,u),\quad (K_\eta su)^2,
\quad ((KHu)^2su)^2,\\
& ((KHu)^2,s,u)^3,
\quad ((KHu)^2,s,u)^4,
\quad (K_\eta(KHu)^2,s,u)^3,\\[0.3em]
\text{B)}& ((j\eta x),s,u)^2,
\quad (H_jsu),
\quad ((j\eta x),s,u)^3,\\[0.3em]
\text{C)}& (\D_\eta su),\\[0.3em]
\text{D)}& (aLu)^2s_\eta,
\quad (a_\eta su)^2,
\quad ((aLu)^2,s,u)^2,
\quad ((aLu)^3,s,u)^2,
\quad ((aLu)^3,s,u)^3,
\quad (a_l(aLu)^2,s,u)^2.
\end{array}
\]
\noindent\emph{還元：}

A) 記号 $\theta-H$ における対応する形式との一回の折り畳みによる還元または分解。

$(K_\eta(KHu)^2,s,u)$ と $(K_\eta(KHu)^2,s,u)^2$：$K_\nu^2(Ksu)$、$K_\nu^2(Ksu)^2$ との折り畳みによって、\S\ 3. 2), a) に従って分解する。

$(K_\eta(KHu),s,u)^4\equiv(a_\nu^2\theta_{\nu_1},s,u)^4$：分解された形式 $K_\nu^2(Ksu)^3$ から、\S\ 3. 2), b) に従って分解する。

B) 形式列 $(j\eta x)s_\eta$：還元子 $(j\nu x)^2s_\nu$。

形式 $H_j(\widehat{j\eta su})(Hsu)$：還元子 $(j\nu x)(\widehat{j\nu su})^2$。

形式 $H_j(j\eta x)(Hsu)$：分解された形式 $((j\eta x)^2,s,u)^2$ を還元子 $(j\nu x)^2s_\nu$（式 (18.)a）で還元することにより分解する：
\[
0=-2(j\nu x)^2s_\nu s_{\nu_1}=[(j\eta x)^2]_{su^2}+H_j(j\eta x)(Hsu).
\]
C) 還元子 $\D_\nu s_\nu$ と $\D_\nu(\D su)^2$。

D) 記号 $f-H$ における対応する形式との一回の折り畳みによる還元および分解。

E) (30.)c から、分解された形式列 $(a_\sigma su)^2$ の還元が得られる：
\begin{align*}
0&=a_\eta(Hsu)^2(as\widehat{\nu x})^2
=a_\eta(Hsu)^2(a\widehat{\nu\eta}u)^2+2a_\eta(a\widehat{\nu\eta}u)(\widehat{\eta\sigma u})(Hsu)^2
=\frac13a_\sigma(asu)^2,\\
0&=a_\eta a_\nu(Hsu)^2(as\widehat{\nu x})(asu)
=a_\eta(a\widehat{\nu\eta}u)^2(Hsu)^2(asu)+a_\eta(a\widehat{\nu\eta}u)(\widehat{\eta\sigma u})(Hsu)^2(asu)
=\frac13a_\sigma(asu)^3.
\end{align*}
\noindent\emph{帰結：} $s^2$ は七次の形式との折り畳みに入らない。

\begin{center}\textbf{\S\ 25.\quad 十二次、十三次、十四次、十五次の形式（系 $s_{III}$）。}\end{center}
$s$ と八、九、十、十一次の形式（系 III）との折り畳み（\S\S\ 13--16）。

\noindent\emph{既約形式：}

\noindent 12次。
\[
\begin{array}{rl}
\text{A)}& ((\vartheta^2\eta x)^4,s,u),\quad \theta_\eta(\vartheta^2\eta x)^3s_\vartheta,\\
\text{B)}& ((aHu)H_j,s,u)^2,
\quad ((aH\cdot u)H_j,s,u)^3,\\
\text{C)}& (\theta_\nu su)^2,
\quad ((\theta Lu)^2,s,u)^2,
\quad ((\theta Lu)^3,s,u),
\quad ((\theta Lu)^2,s,u)^3,
\quad (\theta_l(\theta Lu)^2,s,u)^2.
\end{array}
\]
\noindent\emph{還元：} 還元子または分解された形式との直接の折り畳みから生じる還元は、もはや述べない。

B) $((aHu)H_j,s,u)$ と $(a_\eta(aHu)H_j,s,u)$ は、関数行列式に関するトランスヴェクションとして分解する。

$(a_\eta(aHu)H_j,s,u)^2$ の分解：分解された形式 $[\theta_\nu\theta_{\nu_1}^3]_{su^3}$（\S\ 13. C）の二重還元：
\begin{align*}
0&\equiv[\theta_{\nu_1}^3\theta_\nu]_{su^3}
\equiv-\frac34\theta_\nu(\theta su)^2(\theta\theta'u)\theta_{\nu_1}^3\\
&\equiv-\frac98(\theta\theta'\eta su)(\theta\theta'u)(\theta Hu)(\theta'Hu)\theta_\nu'(\theta su)
\equiv-\frac38a_\eta(aHu)H_j(asu)^2.
\end{align*}

C) 公式 (31.) から、分解された形式
\[
 [L_\vartheta(\theta Lu)]_{su^3}\quad\hbox{および}\quad [\theta_l(\theta Lu)]_{s^2u^3}
\]
の還元、すなわち積 $[\theta_\nu^2H]_{su^3}$ と $[a_\nu^2(bHu)^2]_{su^3}$ の還元が従う：
\begin{equation}
\tag{32}
\begin{aligned}
0&\equiv\theta_\nu^2(\theta su)^2(Hsu),\\
0&\equiv [L_\vartheta(\theta Lu)]_{su^4}-7[\theta_l(\theta Lu)]_{su^4},\\
0&\equiv a_\nu^2(asu)^2(bHu)^2(bsu)\\
&\quad+6\theta_l(\theta Lu)^2(\theta su)(Lsu),\\
0&\equiv\theta_\nu^3(\theta su)(Hsu)^2\\
&\quad\hbox{から }0\equiv\theta_l^2(\theta Lu)(\theta su)(Lsu)^2\\
&\quad(\hbox{還元子 }a_\eta^2(Hsu)^2).
\end{aligned}
\end{equation}

\noindent 13次。
\[
\text{A)}\ ((KLu)^3,s,u)^2,\ ((KLu)^3,s,u)^3,\qquad
\text{B)}\ (L_jsu)^2,
\qquad
\text{C)}\ ((aH^2u)^3,s,u)^2.
\]
\noindent 14次。
\[
\text{A)}\ ((aLu)^2L_j,s,u)^2,
\qquad
\text{B)}\ ((\theta H^2u)^3,s,u)^2.
\]
\noindent 15次。
\[
 ((KH^2u)^4,s,u)^2.
\]
\noindent\emph{還元：} $(K_l(KLu)^3,s,u)^2$、$(H_\vartheta(\theta H^2u)^3,s,u)$、$((K,H\cdot L,u)^5,s,u)$ は、\S\ 3. 2), c) に従って分解する。すなわち、分解された形式 $(K_l(KLu)^2,s,u)^3$、$(H_\vartheta(\theta H'u)^2,s,u)^2$、$((K,H\cdot L,u)^4,s,u)^2$ を同時に考慮する。

\noindent\emph{帰結：} $s^2$ は系 III の個々の形式との折り畳みに入らない。

\begin{center}\textbf{\S\ 26.\quad 系 $s_{IV}$。}\end{center}
\noindent 1) \emph{$s$ と系 I--III との折り畳み。}

\noindent\emph{還元：} 最後の諸節の還元（$(a_\nu^2s_\nu(aHu)^2(asu)s_\eta$ など）により、形式 $(0,\lambda)$、$(1,\lambda)$、$(2,\lambda)$ は I と III の折り畳みを同時には許さない。従って \S\ 18 により、形式 $f,\D,N$ は系 $s_{IV}$ を形成するための折り畳みの積には入らない。

形式 $j$ は折り畳みに入らない。なぜなら、同じ還元によって、$j$ に反変な折り畳み
\[
(K_\nu(k\nu x)^3,s,u),\qquad ((\vartheta^2\eta x)^4,s,u)\quad\hbox{など}
\]
は、$j$ に共変な折り畳み
\[
K_\nu(k\nu x)^2s_k,
\qquad (\vartheta^2\eta x)^3s_\vartheta\quad\hbox{など}
\]
に対応するからである。

\noindent 2) \emph{$s$ と系 II--III との折り畳み}、すなわち $s$ と $H\cdot(1,\lambda)$、$H\cdot(2,\lambda)$、$H\cdot(3,\lambda)$ との折り畳み。

\noindent\emph{既約形式：}
\[
\begin{array}{rl}
\text{A)}&(a_\nu H,s,u)^4,
((aHu)^2H,s,u)^3,
((aHu)^2H,s,u)^4,\\
&((aLu)^2H,s,u)^4,
((aLu)^3H,s,u)^3,
((aH^2u)^3H,s,u)^4,\\[0.3em]
\text{B)}&(a_\nu^2H,s,u)^3,
(a_\eta(aHu)^2H,s,u)^3,
(a_l(aLu)^2H,s,u)^4,\\[0.3em]
\text{C)}&(\theta_\nu H,s,u)^4,
((\theta Hu)^2H,s,u)^3,
((\theta Hu)^2H,s,u)^4,\\
&((\theta Lu)^2H,s,u)^4,
((\theta Lu)^3H,s,u)^3,
((\theta H^2u)^3H,s,u)^4,\\[0.3em]
\text{D)}&(\theta_\eta(\theta Hu)^2H,s,u)^3,
(\theta_l(\theta Lu)^2H,s,u)^4,\\[0.3em]
\text{E)}&((KLu)^3H,s,u)^4,
((KH^2u)^4H,s,u)^4,\\[0.3em]
\text{F)}&((j\nu x)^2H,s,u)^3,
((j\nu x)^2H,s,u)^4,
((j\eta x)H,s,u)^4,\\
&(H_jH,s,u)^3,
(L_jH,s,u)^4.
\end{array}
\]
\noindent\emph{還元：} $\nu$ は $j$ と同様に折り畳みに入らない。

B) および D)。$(a_\nu^2H,s,u)^4$ と $(\theta_\nu^2H,s,u)^4$ は、$(gHu)^2=-\frac13s\sigma$ を考慮して、式 (27.)c) と (29.)c) により分解する：
\[
(a_\nu^2H,s,u)^4=\frac13(asu)^4\sigma,
\qquad
(\theta_\nu^2H,s,u)^4=-\frac16(\theta su)^4\sigma;
\]
従って $(a_\eta(aHu)H,s,u)^4$、$(\theta_\eta(\theta Hu)H,s,u)^4$ が分解する。

$(\theta_\nu^2H,s,u)^3$、$(\theta_\nu^3H,s,u)^3$、従って $(\theta_\eta^2(\theta Hu)H,s,u)^4$ は、\S\ 25、12次の形式 C) に従って分解する。

\noindent 3) \emph{$s$ と系 III--III との折り畳み}、すなわち $s$ と系 III の形式
\[
(3,\lambda)(3,\lambda),\quad (3,\lambda)(2,\lambda),\quad (3,\lambda)(1,\lambda),\quad
(2,\lambda)(2,\lambda),\quad (2,\lambda)(1,\lambda),\quad (1,\lambda)(1,\lambda)
\]
との折り畳み。
\noindent\emph{既約形式：}
\[
\begin{array}{rl}
\text{A)}&(a_\nu^2a_\nu^2,s,u)^3,
(a_\nu^2a_\nu^2,s,u)^4,
(a^2a_\eta(aHu)^2,s,u)^3,\\[0.3em]
\text{B)}&(a_\nu H_j,s,u)^4,
((aHu)^2H_j,s,u)^3,
((aLu)^2H_j,s,u)^4,
((aLu)^3H_j,s,u)^2,
((aH^2u)^3H_j,s,u)^3.
\end{array}
\]
\noindent\emph{還元：}

記号 $\nu$ と $f$ における同じ位数に対して、積 II--III は積 III--III より高いものと見なす。

還元は、積の間の関係（シジジー）から部分的に生じ、また積を対応する分解形式で置き換え、これらの形式と $s$ との折り畳みを還元することから部分的に生じる。反変形式を含む積は、積へ移るので常に省略される。

A) $f$ が二次、記号 $\nu$ は任意の位数。\S\ 11、式 (12.) および同様の計算により、
\begin{align*}
1)\quad 0&=a_\nu b_{\nu_1}+\frac12f(aHu)^2-\frac14\theta H+\frac12u_xH_\vartheta-\frac14u_x^2H_\vartheta^2,\\
2)\quad 0&=a_\nu b_{\nu_1}^2+f a_\eta(aHu)^2-\theta_\eta-\frac12H_\vartheta,\\
3)\quad 0&=a_\nu^2b_{\nu_1}^2-H_\vartheta^2+2\theta_\eta^2.
\end{align*}
従って
\begin{align*}
4)\quad 0&=a_\nu^2b_{\nu_1}^2(j\nu_1x),\\
5)\quad L_\vartheta(\theta Lu)&=-a_\nu b_\eta(bHu)^2+\frac34a_\nu^2(bHu)^2+\frac34\theta_\nu^2H,\\
6)\quad L_\vartheta(\theta Lu)&=-6a_\nu b_\eta(bHu)^2+2a_\nu^2(bHu)^2-\frac12\theta_\nu^2H,\\
7)\quad 0&=a_\nu(bHu)^2+f(aLu)^3+\theta_\nu H,\\
8)\quad 0&=a_\nu(bLu)^3+\frac34(\theta Hu)^2H,\\
9)\quad 0&=(aHu)^2(bHu)^2+2(\theta Hu)^2H,\\
10)\quad 0&=a_\eta(aHu)^2b_\eta(bHu)^2,\\
11)\quad 0&\equiv [a_\nu^2b_\eta(bHu)]_{\widehat{su^4}}.
\end{align*}
\[
(H_\vartheta su)^4,
\qquad (L_\vartheta(\theta Lu),s,u)^3,
\qquad (L_\vartheta(\theta Lu),s,u)^4
\]
の還元（公式 (31.) と (32.)）は、公式 1) から 11) と合わせて、記号 $f$ に関して二次、$\nu$ に関して任意の積との $s$ のすべての折り畳みを還元する。

B) 記号 $f,\theta,j$ のうち二つずつの積、$\nu$ は任意の位数。式 (17.)、(18.)、(20.) および直接計算により、次の関係式が成り立つ：
\begin{align*}
1)\quad 0&=a_\nu\theta_{\nu_1}+\frac14\theta(aHu)^2+\frac14f(\theta Hu)^2,\\
2)\quad 0&=\theta_\nu\theta_{\nu_1}+\frac12\theta(\theta Hu)^2.
\end{align*}
これより
\begin{align*}
3)\quad 0&=3a_\nu(\theta Hu)^2+\theta(aLu)^3+2f(\theta Lu)^3,\\
4)\quad 0&=3\theta_\nu(aHu)^2+f(\theta Lu)^3+2\theta(aLu)^3,\\
5)\quad 0&=a_\nu(\theta Lu)^3,
&0&=\theta_\nu(aLu)^3,
&0&=(aHu)^2(\theta Hu)^2,\\
6)\quad 0&=a_\nu\theta_\nu^2+\theta_\nu a_\nu^2+f\theta_\eta(\theta Hu)^2+\theta a_\eta(aHu)^2,\\
7)\quad 0&=\theta_\nu\theta_\nu^2+\theta\theta_\eta(\theta Hu)^2+\frac16fH_j,\\
8)\quad 0&=a_\nu(j\nu x)^2+fH_j,\\
9)\quad 0&=(aHu)^2(j\nu x)^2-2a_\nu H_j,\\
10)\quad 0&=\theta_\nu(j\nu x)^2,
&0&=\theta H_j,\\
11)\quad 0&=a_\nu\theta_\nu^3-\frac34(j\eta x)^2,\\
12)\quad 0&=a_\nu^2\theta_\nu^2-\frac13(j\eta x)^2-\frac13(\D Hu)^2,\\
13)\quad 0&=a_\nu^2\theta_\nu^3-\frac12a_\sigma,\\
14)\quad 0&=\theta_\nu\theta_\nu^3,
&0&=a_\eta H_j.
\end{align*}
公式は、積が偏極可能になるところまで進められている。すなわち、折り畳みによって新しい積が一つだけ生じるのは、a) その積自身への折り畳みが還元可能になり、b) 折り畳みによって生じるその他の積が還元可能になるか、または反変形式に至るからである（\S\ 3. 2), a), b) 参照）。

公式 1) から 5) は、$a_\nu\theta_\nu$ から折り畳みによって生じる積との $s$ のすべての折り畳みの還元を与える。公式 6) から 10) は、$a_\nu\theta_\nu^2$ から折り畳みによって生じる積の還元を与える。\S\ 24 A) により、公式 6) から 10) を考慮すると、$a_\nu^2\theta_\nu$ から生じる積との $s$ の折り畳みは分解する。

すなわち
\begin{align*}
0&\equiv a_\nu^2(asu)^2\theta_{\nu_1}(\theta su)^2,
&0&\equiv a_\nu^2(asu)\theta_{\nu_1}(\theta su)^3-K_\eta(KHu)^2(Ksu)^3,\\
0&\equiv a_\nu^2(asu)^2\theta_{\nu_1}(\theta su),
&0&\equiv a_\nu^2(asu)\theta_{\nu_1}(\theta su)^2,\\
0&\equiv a_\nu^2(asu)\theta_{\nu_1}(\theta su).
\end{align*}
還元子
\[
 ((j\eta x)^2,s,u)^3\ [\hbox{由来 }(\widehat{\vartheta\eta su})^2(Hsu),\ \S\ 23.\ C3],
\]
\[
(H_j(j\eta x),s,u)^2\ [\S\ 24,\ B],
\quad ((\D Hu)^2,s,u)^2\ [\S\ 24,\ C],
\quad (a_\sigma su)^2\ [\S\ 24,\ E]
\]
は、公式 11) から 14) と合わせて、$a_\nu\theta_{\nu_1}^3$、$a_\nu^2\theta_{\nu_1}^2$、$a_\nu^2\theta_{\nu_1}^3$ から生じるすべての積の還元を与える。

以上の計算は次の\emph{結論}にまとめられる：

法 $(\R,t)$ に関する相対的完全系は、以下の 331 個の形式および部分系からなり、それらを続いて表としても整理して掲げる。
\clearpage
\begin{landscape}
\thispagestyle{plain}
\begin{center}\bfseries 表 I（90頁参照）\end{center}
\vspace{0.25em}
\begingroup
\setlength{\tabcolsep}{2pt}
\renewcommand{\arraystretch}{1.35}
\tiny
\begin{center}
\begin{adjustbox}{max width=0.98\linewidth,max totalheight=0.72\textheight,center}
\begin{tabular}{|c|c|c|c|c|c|c|c|c|c|c|c|c|c|c|c|c|c|c|c|c|c|c|}
\hline
\diagbox{$u$}{$x$} & 0 & 1 & 2 & 3 & 4 & 5 & 6 & 7 & 8 & 9 & 10 & 11 & 12 & 13 & 14 & 15 & 16 & 17 & 18 & 19 & 20 & 21 \\ \hline
0 & \mcell{i} &  &  &  & \mcell{f} &  & \mcell{\A} &  & \mcell{K_{\nu}(k\nu x)^3} &  & \mcell{K_{\eta}(k\eta x)^3} &  &  &  & \mcell{(\vartheta^3\eta x)^4} & \mcell{H_k(k\vartheta,\eta x)^4} &  & \mcell{\theta_l(\vartheta k,lx)^5} &  &  &  & \mcell{(\vartheta^2lx)^6} \\ \hline
1 &  & \mcell{u_x} &  &  &  & \mcell{\theta_\nu(\vartheta\nu x)^2} &  & \mcell{\theta_\eta(\vartheta\eta x)^2} & \mcell{N} & \mcell{(k\nu x)^3} & \mcell{\theta_\nu(\vartheta^2\nu x)^3} & \mcell{(k\eta x)^3} & \mcell{H_{\vartheta}(\vartheta^2\eta x)^3\\ \theta_\eta(\vartheta^2\eta x)^3} &  & \mcell{\theta_l(\vartheta^2lx)^4} &  & \mcell{(\vartheta k,\eta,x)^4} &  & \mcell{(\vartheta k,l,x)^5} &  &  &  \\ \hline
2 &  &  & \mcell{a_\nu^2} &  & \mcell{\theta\\ a_\eta^2} &  & \mcell{(\vartheta\nu x)^2} & \mcell{K_\nu(k\nu x)^2} & \mcell{(\vartheta\eta x)^2} & \mcell{H_k(k\eta x)^2\\K_\eta(k\eta x)^2} &  & \mcell{(\vartheta^2\nu x)^3\\ K_l(klx)^3} &  & \mcell{(\vartheta^2\eta x)^3} &  & \mcell{(\vartheta^2lx)^4} &  &  &  &  &  &  \\ \hline
3 &  & \mcell{\theta_\nu^3} &  & \mcell{a_\nu} & \mcell{\theta_\nu(\vartheta\nu x)} & \mcell{\A_\nu\\a_\eta} & \mcell{K\\H_{\vartheta}(\vartheta\eta x)\\\theta_\eta(\vartheta\eta x)} & \mcell{\A_\eta} & \mcell{(k\nu x)^2\\\theta_l(\vartheta lx)^2} &  & \mcell{(k\eta x)^2} &  & \mcell{(klx)^3} &  &  &  &  &  &  &  &  &  \\ \hline
4 & \mcell{\nu} &  & \mcell{H\\\theta_\nu^2} & \mcell{(j\nu x)^3\\a_\eta(aHu)} & \mcell{\theta_\eta^2} & \mcell{(\vartheta\nu x)\\a_l^2} &  & \mcell{N_\nu\\(\vartheta\eta x)} &  & \mcell{H_\eta\\N_\eta\\(\vartheta lx)^2} &  &  &  &  &  &  &  &  &  &  &  &  \\ \hline
5 &  & \mcell{a_\eta(aHu)^2} & \mcell{\theta_\eta^2(\theta Hu)} & \mcell{\theta_\nu} & \mcell{\frac{K_\nu^2}{(aHu)}} & \mcell{\theta_\eta} & \mcell{K_\eta^2\\(\A Hu)\\a_l} &  & \mcell{\A_l} &  &  &  &  &  &  &  &  &  &  &  &  &  \\ \hline
6 & \mcell{j\\\sigma} &  & \mcell{(j\nu x)^2\\(aHu)^2} & \mcell{L\\a_\nu^2(j\nu x)\\H_{\vartheta}(\theta Hu)\\\theta_\eta(\theta Hu)} & \mcell{\frac{K_\nu}{\theta_l^2}} &  &  & \mcell{K_\eta} &  &  &  &  &  &  &  &  &  &  &  &  &  &  \\ \hline
7 &  & \mcell{\theta_\eta(\theta Hu)^2} & \mcell{H_j(j\eta x)\\a_l(aLu)^2} &  & \mcell{a_\nu(j\nu x)\\(\theta Hu)} &  & \mcell{\A_\nu(j\nu x)\\\theta_l} & \mcell{K_l^2} &  &  &  &  &  &  &  &  &  &  &  &  &  &  \\ \hline
8 & \mcell{H_j^2\\a_l(aLu)^3} & \mcell{(\nu\sigma x)\\(j\nu x)} & \mcell{(\theta Hu)^2} & \mcell{K_\eta(KHu)^2\\(j\eta x)\\(aLu)^2} &  &  &  &  &  &  &  &  &  &  &  &  &  &  &  &  &  &  \\ \hline
9 &  & \mcell{H_j\\(aLu)^3} & \mcell{a_\eta(aHu)H_j\\L_{\vartheta}(\theta Lu)^2\\\theta_l(\theta Lu)^2} &  & \mcell{(KHu)} &  &  &  &  &  &  &  &  &  &  &  &  &  &  &  &  &  \\ \hline
10 & \mcell{\theta_l(\theta Lu)^3} & \mcell{L_j^2\\(j\sigma x)\\a_\eta(aH^2u)^3} &  & \mcell{H_j(aHu)\\(\theta Lu)^2} &  &  &  &  &  &  &  &  &  &  &  &  &  &  &  &  &  &  \\ \hline
11 &  & \mcell{(\theta Lu)^3} & \mcell{K_l(KLu)^3\\L_j\\(aH^3u)^3} &  &  &  &  &  &  &  &  &  &  &  &  &  &  &  &  &  &  &  \\ \hline
12 & \mcell{(aH^2u)^4} & \mcell{a_\eta(aLu)^2L_j\\H_{\vartheta}(\theta H^2u)^3\\\theta_\eta(\theta H^2u)^3} &  & \mcell{(KLu)^3} &  &  &  &  &  &  &  &  &  &  &  &  &  &  &  &  &  &  \\ \hline
13 & \mcell{(\nu\sigma j)} &  & \mcell{(aLu)^2L_j\\(\theta H^3u)^3} &  &  &  &  &  &  &  &  &  &  &  &  &  &  &  &  &  &  &  \\ \hline
14 & \mcell{(\theta H^2u)^4} & \mcell{K_\eta(KH^2u)^4\\(a_\eta H\!\cdot L,u)^4} &  &  &  &  &  &  &  &  &  &  &  &  &  &  &  &  &  &  &  &  \\ \hline
15 & \mcell{L_9(\theta,H\!\cdot L,u)^4} &  & \mcell{(KH^2u)^4} &  &  &  &  &  &  &  &  &  &  &  &  &  &  &  &  &  &  &  \\ \hline
16 & \mcell{(\theta,H\!\cdot L,u)^5} &  &  &  &  &  &  &  &  &  &  &  &  &  &  &  &  &  &  &  &  &  \\ \hline
17 & \mcell{K_\eta(K,H\!\cdot L,u)^5} &  &  &  &  &  &  &  &  &  &  &  &  &  &  &  &  &  &  &  &  &  \\ \hline
18 &  & \mcell{(K,H\!\cdot L,u)^5} &  &  &  &  &  &  &  &  &  &  &  &  &  &  &  &  &  &  &  &  \\ \hline
19 &  &  &  &  &  &  &  &  &  &  &  &  &  &  &  &  &  &  &  &  &  &  \\ \hline
20 &  &  &  &  &  &  &  &  &  &  &  &  &  &  &  &  &  &  &  &  &  &  \\ \hline
21 & \mcell{(KH^3u)^6} &  &  &  &  &  &  &  &  &  &  &  &  &  &  &  &  &  &  &  &  &  \\ \hline
\end{tabular}
\end{adjustbox}
\end{center}
\vspace{0.2em}
\begin{center}\footnotesize （最上段の横列の数字は $x$ に関する次数を、最初の縦列の数字は $u$ に関する次数を表す。）\end{center}
\endgroup
\end{landscape}
\clearpage
\begin{landscape}
\thispagestyle{plain}
\begin{center}\bfseries 表 II（90頁参照）\end{center}
\vspace{0.25em}
\begingroup
\setlength{\tabcolsep}{2pt}
\renewcommand{\arraystretch}{1.35}
\tiny
\begin{center}
\begin{adjustbox}{max width=0.98\linewidth,max totalheight=0.72\textheight,center}
\begin{tabular}{|c|c|c|c|c|c|c|c|c|c|c|c|c|c|c|c|c|c|}
\hline
\diagbox{$u$}{$x$} & 0 & 1 & 2 & 3 & 4 & 5 & 6 & 7 & 8 & 9 & 10 & 11 & 12 & 13 & 14 & 15 & 16 \\ \hline
0 & \mcell{J} &  &  &  & \mcell{s} &  & \mcell{s_j^3} &  &  &  &  & \mcell{s_\nu} & \mcell{(k\nu x)^3s_\nu} & \mcell{(\vartheta^3\nu x)^2s_{\vartheta}\\(\vartheta^2\nu x)^2s_{\vartheta}} &  & \mcell{\theta_\eta(\vartheta^2\eta x)^2s_\eta} &  \\ \hline
1 &  &  &  &  &  &  & \mcell{(asu)} & \mcell{s_\eta} & \mcell{s_l^2\\(\A su)} & \mcell{s_\eta(Nsu)\\(\vartheta\nu x)^2s,u} & \mcell{((k\nu x)^3,s,u)_s\\(K_\nu(k\nu x)^3,s,u)} &  & \mcell{K_\eta(k\eta x)^3,s,u} &  & \mcell{(\vartheta^2\nu x)^2s_\nu} &  & \mcell{(\vartheta^2\eta x)^4,s,u} \\ \hline
2 &  &  &  &  & \mcell{(asu)^2} & \mcell{s_\theta(\theta su)} & \mcell{(\A su)^2/a_\nu s_\nu} & \mcell{(\vartheta\nu x)^2,s,u\\(\theta_l(\vartheta\nu x)^2,s,u)} &  & \mcell{\theta_\eta(\vartheta^2\eta x)^2,s,u} &  & \mcell{(k\eta x)^2s_\nu\\((k\nu x)^2,s,u)} &  & \mcell{(k\eta x)^3,s,u} &  &  &  \\ \hline
3 &  &  & \mcell{(asu)^3} & \mcell{s_\vartheta(\theta su)^2\\s_\nu} & \mcell{a_\nu s_\nu(asu)\\a_l^2(asu)} & \mcell{s_\eta} & \mcell{(\theta su)\\a_l^2su} &  & \mcell{(Nsu)^2\\(\vartheta\eta x)s_\nu\\((\vartheta\nu x)^2,s,u)} & \mcell{((k\nu x)^2,s,u)^2} & \mcell{(\vartheta^2\eta x)^2,s,u} &  &  &  &  &  &  \\ \hline
4 & \mcell{(asu)^4} & \mcell{s_\vartheta(\theta su)^3} & \mcell{a_\nu^2(asu)^2} & \mcell{(\theta_\nu^2su)} & \mcell{(\theta su)^2} & \mcell{s_k(Ksu)^3\\a_\nu(asu)\\s_\nu(asu)} & \mcell{((\vartheta\nu x)^2,s,u)^2} & \mcell{s_\nu(\A su)\\(\A su)\\(aHu)s_\eta\\(a_\nu su)} &  & \mcell{(\A_\eta su)} &  &  &  &  &  &  &  \\ \hline
5 &  &  & \mcell{(\theta su)^3} & \mcell{s_k(Ksu)^3,s_j\\s_\vartheta(\theta s^\prime u)^4\\s_\nu(asu)^2\\a_\nu(asu)^2} & \mcell{(Hsu)\\((\vartheta\nu x)^2,s,u)\\(\theta_\nu(\vartheta\nu x)^2,s,u)\\(\theta_\nu^2su)} & \mcell{((j\nu x)^2,s,u)\\(aHu)^2s_\nu\\(a_\eta su)^2} & \mcell{(Ksu)^2\\s_\eta\\(\theta_\eta^2su)} &  & \mcell{((k\nu x)^2,s,u)^2\\K_\nu s_\nu} &  &  &  &  &  &  &  &  \\ \hline
6 & \mcell{(\theta su)^4} & \mcell{s_\nu(asu)^3\\a_\nu(asu)^3} & \mcell{(Hsu)^2\\(\theta_\nu(\vartheta\nu x),s,u)^3\\(\theta_\nu^2su)^2} & \mcell{a_\eta s_\nu^2\\(a_\eta(aHu),s,u)^2\\(a_\eta(aH^2u),s,u)} & \mcell{(Ksu)^3} & \mcell{s_\eta(asu)\\((\vartheta\nu x),s,u)^3\\(\theta_\eta su)} & \mcell{((k\nu x)^2,s,u)^3} & \mcell{s_\nu(\A su)\\a_\nu(j\nu x)s_\nu\\(\theta H u)s_\eta\\(\theta_\eta su)} &  &  &  &  &  &  &  &  &  \\ \hline
7 & \mcell{\theta_\nu(\vartheta\nu x),s,u)^4} & \mcell{(a_\eta(aHu),s,u)^3} & \mcell{(Ksu)^4\\(\theta_\eta^2(\theta Hu),s,u)^2\\(a_\eta^2-a_l^2,s,u)^3} & \mcell{s_j(asu)^3\\s_k(K^2u)^3\\((j\nu x),s,u)^2\\(\theta_\eta su)^2} & \mcell{((j\nu x),s,u)\\((j\eta x),s,u)\\(aHu),s,u)\\(aH^2u),s,u)} & \mcell{((\vartheta\eta x),s,u)^3} & \mcell{(aLu)^3s_l/(asu)^3} &  &  &  &  &  &  &  &  &  &  \\ \hline
8 & \mcell{(a^2-a_l^2,s,u)^4} & \mcell{s_j(asu)^3\\s_k(K^2u)^3\\((j\nu x),s,u)^4\\(\theta_\eta su)^3} & \mcell{((j\nu x)^2,s,u)^2\\(aHu),s,u)^3\\((aHu)^2,s,u)^3} & \mcell{(Lsu)^2\\(a_l^2(j\nu x),s,u)^3\\H_\vartheta(\theta Hu),s,u)^3\\\theta_l(\theta Hu),s,u)^3} & \mcell{(asu)^3} & \mcell{(K,su)^3} &  & \mcell{(K_\vartheta su)^2} &  &  &  &  &  &  &  &  &  \\ \hline
9 & \mcell{K_j^2su^4\\((aHu),s,u)^4} & \mcell{(Lsu)^4\\(a_l^2(j\nu x),s,u)^4\\(\theta_\nu(\theta Hu),s,u)^2} & \mcell{(K_\nu^3u)^6\\(a_\eta(aLu)^2,s,u)^3\\(a_l^2H,s,u)^3} & \mcell{(K,su)^3} & \mcell{(a_\nu(j\nu x),s,u)^3\\(\theta Hu),s,u\\(\theta H u)^2,s,u} &  & \mcell{(\theta_lsu)^3} &  &  &  &  &  &  &  &  &  &  \\ \hline
10 &  &  & \mcell{(K,su)^4\\(a_l^2a_\eta(aHu)^2,s,u)^3} & \mcell{((j\eta x),s,u)^3\\(H_j,su)\\((aLu)^2,s,u)^2\\((aLu)^3,s,u)} &  &  &  &  &  &  &  &  &  &  &  &  &  \\ \hline
11 & \mcell{(a_\nu(j\nu x),s,u)^4\\((\theta Hu),s,u)^4} & \mcell{(K_l(KH u)^2,s,u)^3\\((j\eta x),s,u)^3\\((aLu)^2,s,u)^4\\(a_\eta H,s,u)^3} & \mcell{(H^2su)^3\\\theta_l(\theta Lu)^2,s,u)^3} &  & \mcell{((KHu)^2,s,u)^3} &  &  &  &  &  &  &  &  &  &  &  &  \\ \hline
12 &  & \mcell{(a_\eta(aHu)^2H,s,u)^3} & \mcell{((KH u)^3,s,u)^3} & \mcell{((aHu)H_j,s,u)^2\\((\theta Lu)^2,s,u)^3} &  &  &  &  &  &  &  &  &  &  &  &  &  \\ \hline
13 & \mcell{((KH u)^3,s,u)^4} & \mcell{((aHu)H_j,s,u)^3\\((\theta Lu)^2,s,u)^4\\(\nu\cdot H,s,u)^3} & \mcell{(L_jsu)^2\\((aH^3u)^3,s,u)^3\\((j\eta x)H,s,u)^3\\((aHu)^2H,s,u)^3} &  &  &  &  &  &  &  &  &  &  &  &  &  &  \\ \hline
14 & \mcell{((j\nu x)^2H,s,u)^4\\((aHu)^2H,s,u)^4} & \mcell{(\theta_\eta(\theta Hu)^2H,s,u)^3} &  & \mcell{((KLu)^3,s,u)^2} &  &  &  &  &  &  &  &  &  &  &  &  &  \\ \hline
15 & \mcell{(a_l(aLu)^3H,s,u)^4} & \mcell{((KLu)^3,s,u)^3} & \mcell{((aLu)^2L_j,s,u)^2\\((\theta Hu)^3H,s,u)^3\\((\theta Hu)^2H,s,u)^2} &  &  &  &  &  &  &  &  &  &  &  &  &  &  \\ \hline
16 & \mcell{((\theta Hu)^2H,s,u)^4\\(\sigma\cdot H_j,s,u)^4} & \mcell{((j\eta x)H,s,u)^4\\(H_jH,s,u)^3\\((aLu)^2H,s,u)^4\\((aLu)^3H,s,u)^3} &  &  &  &  &  &  &  &  &  &  &  &  &  &  &  \\ \hline
17 & \mcell{(\theta(\theta Lu)^2H,s,u)^4} &  & \mcell{((KH^3u)^3,s,u)^3} &  &  &  &  &  &  &  &  &  &  &  &  &  &  \\ \hline
18 &  & \mcell{((\theta Lu)^2H,s,u)^4\\((\theta Lu)^3H,s,u)^3\\(aHu)^2H,s,u)^4} &  &  &  &  &  &  &  &  &  &  &  &  &  &  &  \\ \hline
19 & \mcell{((aH^3u)^3H,s,u)^4\\(L_jH,s,u)^5} &  &  &  &  &  &  &  &  &  &  &  &  &  &  &  &  \\ \hline
20 &  & \mcell{((KLu)^3H,s,u)^5} & \mcell{((aLu)^3H,s,u)^5} &  &  &  &  &  &  &  &  &  &  &  &  &  &  \\ \hline
21 & \mcell{((\theta H^3u)^3H,s,u)^4\\((aLu)^2H_j,s,u)^4} &  &  &  &  &  &  &  &  &  &  &  &  &  &  &  &  \\ \hline
22 &  &  &  &  &  &  &  &  &  &  &  &  &  &  &  &  &  \\ \hline
23 & \mcell{((KH^3u)^4H,s,u)^4} & \mcell{((aH^3u)^4H,s,u)^3} &  &  &  &  &  &  &  &  &  &  &  &  &  &  &  \\ \hline
\end{tabular}
\end{adjustbox}
\end{center}
\vspace{0.2em}
\begin{center}\footnotesize （最上段の横列の数字は $x$ に関する次数を、最初の縦列の数字は $u$ に関する次数を表す。）\end{center}
\endgroup
\end{landscape}




\clearpage

\providecommand{\pair}[2]{(#1\,|\,#2)}
\providecommand{\eps}{\varepsilon}
\providecommand{\rhoidx}{\varrho}

\begin{center}
{\Large\bfseries 3. $n$ 変数形式の不変式論について\footnote{1909 年、ザルツブルク自然科学者集会で行われた講演。}}\par
\vspace{1.2em}
Jahresbericht der DMV 19 (1910), pp. 101--104
\end{center}

\vspace{1em}
$n$ 変数形式の射影的不変式論では、主問題、すなわち形式系の有限性は解決されている。しかし、さらに別の一連の問題、すなわち形式間の関連を調べる方法に関わる問題は扱われていない。私はそのような問題について得た結果を手短に述べたいが、明瞭さのため、まずそれらの問題設定を、それが知られている三元の領域で概説しておく。

まず念頭にあるのは、記号的方法の二つの基本定理である。第一は、すべての不変な形成が記号的に表示できるという定理であり、第二は、不変式の間に存在するすべての関係が、少数の記号的恒等式を逐次適用することによって得られるという定理である。したがって、記号的方法は、不変式の領域を離れることなくすべての関係を与える。\footnote{Study: \emph{Methoden zur Theorie der ternären Formen}. II. \S\ 6. (Leipzig, Teubner, 1889.)} その上に、形式を生成する過程、すなわち記号的な折り畳み過程と非記号的な微分過程、級数展開の理論などが置かれる。これに関連して、私の学位論文で証明した定理\footnote{\emph{Über die Bildung des Formensystems der ternären biquadratischen Form}. \S\ 1. (Crelle's Journal Bd. 134. 1908.)}も挙げておきたい。すなわち、すべての折り畳みは、二つの可能な「共変的」折り畳みを逐次適用することによって生成できる、という定理である。ここで、記号的積 $s_x t_x u_\sigma u_\tau$ に対して私がいう二つの折り畳みとは、$(stu)$ と $(\sigma\tau x)$ のことである。この定理の上に、二元領域の場合と同様に、還元子の理論および形式系の還元の理論を築くことができることは、私が学位論文で示したとおりである。

さて $n$ 元領域に移ると、すでに記号的方法の第一定理は、制限された範囲でしか成り立たない。よく知られているように、四元形式では、点座標および平面座標 $x$ と $u$ のほかに、二つの点座標列から作られる小行列式である直線座標 $p_{ik}$ が現れる。これと同様に、$n$ 元領域では変数列 $p_\rho$ があり、その個々の要素は $\rho$ 個の列 $x$ から作られる小行列式である。また記号列 $S^\rho$ があり、その要素は、$u$ と共変的な $\rho$ 個の列 $s$ から作られる小行列式のように振る舞う。行列式による記号的表示は、知られているように、記号列 $S^\rho$ を列 $s$ に分解し、それらを別々の行列式に分配した場合にしか成り立たない。このとき実際の意味をもつのは、$S^\rho$ に依存する行列式の集合体だけである。しかし、どのような行列式の集合体がこれを実現するのかは見通せず、そのため三元領域で述べた他のすべての定理が失われてしまう。

そこで私は、列 $S^\rho$ において明示的な記号的表示を得るための簡単な原理を示し、それによって残りの定理も回復したい。これは、対応行列に関するよく知られた定理の応用である。「各変数列 $p_\rho$ には、一対一に別の列 $q_{n-\rho}$ が対応する。その個々の要素は $n-\rho$ 個の列 $u$ から作られる小行列式であり、これらの列は方程式 $(u\mid x)=0$ によって $\rho$ 個の列 $x$ と結ばれている。」これに対応して、各記号列 $S^\rho$ には別の $S_{n-\rho}$ が対応し、それは $x$ と共変的な $n-\rho$ 個の列から構成されているかのように振る舞う。

この定理には、応用に適した第二の定式化もある。「各行列積 $(v^{(1)}\cdots v^{(\rho)}\mid y^{(1)}\cdots y^{(\rho)})$ には、\footnote{すなわち、$\rho$ 個の列 $v$ の行列と $\rho$ 個の列 $y$ の行列から作られる、対応する行列式の積の和である。}第二の行列の列が第一の行列の列に反変的であるとき、一つの行列式が対応する。また逆に、各行列式は、あらかじめ与えられた列数 $\rho$ をもつ行列積に変換できる。」したがって、行列式と行列積とは完全に同等のものとして現れ、記号的方法の第一定理は次の形をとる。「すべての不変な形成は、行列積によって記号的に表示できる。」二つの記号列 $S^\sigma$, $T^\tau$ からは、変数列を用いて、これらの記号列を明示的に含む行列積、したがって形式 $(S^\sigma\mid p_\sigma)(T^\tau\mid p_\tau)$ の不変な形成を表す行列積を、非常に容易に作ることができる。すなわち次のものである。
\begin{equation}
 \pair{q_{n-\sigma-\lambda}S^\sigma}{T_{n-\tau}p_{\tau-\lambda}},
 \qquad (\lambda=1,2,\ldots,\tau\ \text{または}\ n-\sigma).
\tag{1}
\end{equation}

より重要なのはその逆である。すなわち、すべての不変な形成が行列積によって明示的に表示できるということであり、したがって上の過程は二元および三元の折り畳み過程の拡張を意味する。この逆を証明するには、記号的方法の第二定理を $n$ 元領域へ移す必要がある。すなわち、すべての列 $S^\rho$, $p_\rho$ を分解不能なものと見なす場合の、独立で分解不能な恒等式の全体を立てなければならない。これらの恒等式は、列 $x$ と $u$ だけを含む既知の恒等式\footnote{E. Pascal in \emph{Memorie d. R. Acc. d. Lincei}. (4) 5 (1888), p. 375.}から、行列式の積定理を行列に対する積定理の系で置き換え、さらにこれらの積定理そのものによって、いくつかの行列積を一つにまとめることによって得られる。このようにして二つの双対的な公式に到達するが、そのうち一つだけを記す。
\begin{equation}
\begin{aligned}
 \pair{S_1^{\rho_1}S_2^{\rho_2}\cdots S_k^{\rho_k}}{x p_{\rho-1}}
 &=\sum_{i=1}^{k}\eps\,\pair{S_1^{\rho_1}\cdots S_{i-1}^{\rho_{i-1}}S_{i+1}^{\rho_{i+1}}\cdots S_k^{\rho_k}q_{n-\rho_i+1}}{S_{n-\rho_i}x},\\
 \rho&=\sum_{i=1}^{k}\rho_i<n,\qquad \eps=\pm1.
\end{aligned}
\tag{2}
\end{equation}

これらの公式に含まれる唯一の特殊化、すなわち列 $x$ が入るということは避けられない。完全に一般の記号列および変数列の場合には、「分解恒等式」に到達する。これは、列 $p_\rho$ が個々の列 $x$ に分解される恒等式である。分解恒等式の最も簡単な特殊例は、Clebsch がその「不変式論の基本問題」において、級数展開における基本共変式を導くためにすでに用いていた。非明示的な行列式表示から行列積への移行を媒介するこれらの恒等式によって、行列表現により望ましい明示的な記号的表示が実際に得られることを示すことができる。

しかし折り畳み過程については、三元領域の場合とまったく同様の定理が成り立ち、それに基づいて還元定理も同様に築くことができる。(1) において数 $\lambda$ を折り畳みの欠損数と呼び、$\sigma=\tau$ である折り畳みを共変的と呼ぶならば、すべての折り畳みは欠損数 1 の共変的折り畳みを逐次適用することによって生成できる。この定理の証明は、本質的に、最も一般の形式を「正規形」、すなわちすべての不変な微分演算の適用のもとで恒等的に消える形式で置き換えることができる、という事実に依拠する。四元領域ではこの後者の事実を Mertens\footnote{Über invariante Gebilde quaternärer Formen. Wiener Berichte. Bd. 98. 1889.} が証明した。最も一般の微分演算、すなわちよく知られているように記号列を微分記号で置き換えることによって折り畳み過程から生じる微分演算を我々が知っていることにより、分解恒等式を用いて Mertens の証明をほとんど字句どおりに $n$ 元領域へ移すことができる。

\clearpage
\begin{center}
{\Large\bfseries 4. $n$ 変数形式の不変式論について}\par
\vspace{1.2em}
Journal für die reine und angewandte Mathematik 139 (1911), pp. 118--154
\end{center}

\section*{序論.}
$n$ 変数形式の射影的不変式論では、主問題、すなわち形式系の有限性の証明に付随する問題は、二元および三元の領域を越えると、ほとんど扱われていない。それは、形式の相互依存性およびこの形式間の関連を把握するための方法に関する問題である。これらの方法は本質的に、二元および三元の領域で完全に証明されているように、Study が「記号的方法の基本定理」と呼んだ二つの定理、すなわち形式の記号的表示可能性に関する定理と記号的恒等式に関する定理に基づいている。\footnote{E. Study, \emph{Methoden zur Theorie der ternären Formen} (Leipzig, Teubner, 1889). II. \S\ 6. これらの基本定理が特殊な変換群の不変式論にも現れることについては、Study, \emph{Das Apollonische Problem} (Math. Ann. Bd. 49 [1897]) および \emph{Zur Differentialgeometrie der analytischen Kurven} (Transactions of the American Math. Society, Bd. 1 [1909]) を参照。} 後者の定理は、整な不変形成の間に存在するすべての関係が、少数の記号的恒等式を逐次適用することによって得られることを述べる。したがって、記号的方法、しかもそれだけが、不変式の領域を離れることなく形式間の関連についての知識を与えるのである。

Gordan はそのエルランゲン・プログラム\footnote{P. Gordan, \emph{Über das Formensystem binärer Formen}. S. 46. (Leipzig, Teubner, 1875).}において、すでに $n$ 変数への移行が困難になる理由を指摘していた。それは、記号的方法の第一定理が、もはやその一般的な形では有効でないという点にある。$n$ 元領域では、よく知られているように、高い段階の変数列 $p_\rho$\footnote{記法については \S\ 1 を参照。}が現れ、同様に高い段階の記号列 $S^\rho$ も現れる。もともと Clebsch\footnote{A. Clebsch, über eine Fundamentalaufgabe der Invariantentheorie. (Abh. der Ges. d. Wiss. zu Göttingen, Bd. XVII, 1872).} のように任意個の列 $p_\rho$ をもつ形式から出発する場合にも、F. Deruyts と A. Capelli\footnote{Cf. A. Capelli, \emph{Lezioni sulla teoria delle forme algebriche} (Napoli 1902), \S\ 22--24, およびそこに引用された文献。} に従って任意個の共変的な列 $x$ だけを含む形式から出発する場合にも、これらの列に到達する。列 $p_\rho$, $S^\rho$ を明示的に含む記号的表示\footnote{「明示的表示」のより詳しい定義は \S\ 2。}は存在しない。$p_\rho$, $S^\rho$ を個々の列 $x$ と、それに反変的な $s$ とに分解し、それらを別々の行列式に分配しなければならない。確かに、微分条件（公式 (19.) 参照）によって、与えられた行列式の集合体が列 $p_\rho$, $S^\rho$ だけに依存する性質をもつことを確かめることはできる。しかしこの方法では、不変な形成の全体およびそれらを生成する過程についての見通しは得られない。\footnote{Study による、計算上きわめて有用な記号法の展開（F. Engel, Ber. d. K. Sächs. Ges. d. Wiss., math.-phys. Kl., 1900, p. 126 によって伝えられ、四元領域については E. Study, \emph{Geometrie der Dynamen}, p. 135 による）もまた、ここでいう意味での明示的表示ではない。たとえば、それは形式を生成する非記号的な微分過程にはほとんど導かないであろう。}

そこで \S\S\ 2 および 6 で、「対応行列に関する定理」\footnote{Grassmann の 1862 年の \emph{Ausdehnungslehre}, no. 112 を別にすれば、最初は A. Brill, über zwei Berührungsprobleme, \S\ 2 (Math. Ann. Bd. 4. 1871)。}を用い、行列式による表示を「行列積」による表示で置き換えることによって、第一基本定理がその一般性において回復されることを示す。\S\ 2 では、列 $S^\rho,p_\rho$ を明示的に含む各行列積が基礎形式の不変な形成であることを示す。\S\ 6 で与えられる逆、すなわち各不変な形成が行列積によって明示的に表示できるということは、第二基本定理の移行（\S\S\ 3--5）を通して初めて成り立つ。

この定理は、変数列 $x$ だけを含む形式については知られている。\footnote{E. Pascal, Giorn. di mat. 26 (1888), pp. 33, 102; Rendic. d. Accad. d. Linc. (4) 4 (1888), p. 119; ib. Mem. (4) 5 (1888), p. 375.} すべての列 $S^\rho,p_\rho$ を分解不能と見なす場合の、分解不能な恒等式の全体を立てるという一般問題は、Study\footnote{``Methoden \dots'', p. 204, note 17).} によって定式化された。同時に彼は、現れる恒等式の数に、$n$ 元領域における方法の困難さの尺度を見ている。いまや恒等式の全体を一つの公式 (36.) にまとめることができるので、この困難は少なくとも最小限に抑えられる。しかし応用では、(36.) のある際立った特殊例、たとえば (20.), (21.) がしばしば現れる。それらは、新しい列を代入することなしに明示的表示を許すという点で特徴づけられる。また公式 (31.) は「分解恒等式」と呼ばれる。そこでは列 $p_\rho$ が列 $y$ に分解されているように現れるからである。

二つの基本定理の上に、さらなる理論を容易に築くことができる。第一定理は、形式を生成する過程、すなわち記号的折り畳み過程と非記号的微分過程\footnote{論文 ``Invariante Gebilde von Nullsystemen'' (Sitzungsber. d. Ak. d. Wiss. Wien, Bd. 97, Abt. IIa, 1888) において、Mertens は別の道筋により、直接には一般化できない四元の折り畳み過程に到達した。そこに現れる 6 個の列 $S^2$ の交代不変式は、ここで生じるものとは、置換 (11.) を一度適用することによって、偶不変式の集合体だけ異なる。特殊な場合については、四元の折り畳み過程は P. Gordan, \emph{Das simultane System von zwei quadratischen quaternären Formen} (Math. Ann. Bd. 56. 1903) にも見られる。}を直接に与える。一方、分解恒等式は、これらの過程のうち同値なものをすべて消去する（\S\ 6）。微分過程を知ることにより、さらに「正規形」の概念を $n$ 元領域へ移すことができ、また Mertens\footnote{F. Mertens, Über invariante Gebilde quaternärer Formen. (Sitzber. d. Ak. d. Wiss. Wien. Bd. 98. Abt. IIa. 1889.)} が四元領域で与えた証明手続きによって、不変な形成の系を同値な正規形の系で置き換えることができるという定理が得られる（\S\ 7）。この後者の仮定のもとで、私は学位論文\footnote{Über die Bildung des Formensystems der ternären biquadratischen Form. (Dieses Journal, Bd. 134. 1908.)}において、三元の折り畳み過程が二つの基本的な折り畳みを逐次適用することによって生成できることを示した。この定理と級数展開の理論に基づき、さらに「形式列」という概念を導入して、形式系の主要な還元定理を展開した。完全に同様に、$n$ 元領域では、折り畳み過程は $n-1$ 個の基本的な折り畳みから生成でき（\S\ 8）、還元定理を立てることができる（\S\ 9）。

\paragraph{後記。} ザルツブルクでの報告の際に、ウィーンの E. Müller 教授は、この論文の基礎にある行列積の計算が、原理的には Grassmann の拡張論の計算法と同一であることを私に指摘した。\footnote{Hermann Graßmanns gesammelte mathematische und physikalische Werke. Ersten Bandes zweiter Teil: Die Ausdehnungslehre von 1862. In Gemeinschaft mit H. Graßmann d. Jüngeren herausgegeben v. Fr. Engel (Leipzig, Teubner 1896).} 実際、Grassmann の「組合せ積」$[S^{\rho_1}S^{\rho_2}\cdots S^{\rho_i}]$ は、これらの列によって与えられた行列と見なすことができる（\S\ 2 参照）。具体的には、「進行積」はその列自身の行列、「退行積」は対応する列 $S_{n-\rho_1},S_{n-\rho_2},\ldots,S_{n-\rho_i}$ の行列であり、「混合積」に対応する行列は、いくつかの列 $S$ が置換 (9.) によって新しい列 $P,Q$ にまとめられるときに生じる。我々の「対応する列」は Grassmann の「補元」に対応し、我々の「行列積」は「段階数 0 の混合積」に対応する。

この対応のもとで、1862 年の \emph{Ausdehnungslehre} no. 113 に与えられた展開は、我々の公式 (32.) に双対的な一公式と同一になる。特殊な場合 $\rho=1$ では、(32.) は (17.) へ移り、双対的公式は (16.) へ移る。Grassmann の展開のこれら特殊な場合から、Müller\footnote{E. Müller, Beiträge zur Graßmannschen Ausdehnungslehre. 1. Mitteilung: Einige allgemeine Sätze. (Sitzungsber. d. Ak. d. Wiss. Wien; Bd. 118. Abt. IIa. Juli 1909), Satz 16 und 18.} は最近、上に挙げた恒等式 (20.), (21.) を導いた。

Grassmann--Müller の導出と異なり、我々の導出は同時に恒等式の完全性の証明を与える。ここで考察している問題にとって本質的なのはまさにこの点であると思われる。さらに我々の導出は、(20.), (21.) から出発して、これまで注意されていなかったと思われる最も一般の分解不能な恒等式 (36.) へ進む。

\section*{\S\ 1. 記法と定義。既知の結果のまとめ。}
通常どおり、変数列 $x$ によって要素 $x_1,x_2,\ldots,x_n$ の列を表し、同様に変数列 $p_\rho$（$\rho=1,2,\ldots,n-1$）によって $\binom n\rho$ 個の要素 $p_{i_1i_2\ldots i_\rho}$ の列を表す。ここで $p_{i_1i_2\ldots i_\rho}$ は、$\rho$ 個の列 $x^{(1)},x^{(2)},\ldots,x^{(\rho)}$ の行列から作られる $\binom n\rho$ 個の行列式を表す。
\[
\left|\begin{array}{cccc}
 x_{i_1}^{(1)}&x_{i_2}^{(1)}&\cdots&x_{i_\rho}^{(1)}\\
 x_{i_1}^{(2)}&x_{i_2}^{(2)}&\cdots&x_{i_\rho}^{(2)}\\
 \vdots&\vdots&&\vdots\\
 x_{i_1}^{(\rho)}&x_{i_2}^{(\rho)}&\cdots&x_{i_\rho}^{(\rho)}
\end{array}\right|,
\qquad (i_1<i_2<\cdots<i_\rho=1,2,\ldots,n).
\]
対応行列に関する定理\footnote{たとえば Gordan--Kerschensteiner, \emph{Vorlesungen über Invariantentheorie}, Bd. I, p. 99 を参照。}によれば、各列 $p_\rho$ には、列の各要素について成り立つ次の関係によって、第二の列 $q_{n-\rho}$ が一対一に対応する。
\begin{equation}
 p_{i_1i_2\ldots i_\rho}=(-1)^{\frac{\rho(\rho+1)}{2}+i_1+i_2+\cdots+i_\rho}q_{j_1j_2\ldots j_{n-\rho}},
\tag{1}
\end{equation}
ここで $i$ と $j$ はそれぞれ整列しており、合わせて 1 から $n$ までの数の完全な置換をなす。列 $q_{n-\rho}$ は、$x$ に反変的な $n-\rho$ 個の列 $u^{(1)},u^{(2)},\ldots,u^{(n-\rho)}$ の行列から作られる、次数 $n-\rho$ の $\binom n\rho$ 個の行列式 $q_{j_1j_2\ldots j_{n-\rho}}$ から成る。これらの列は、$\rho(n-\rho)$ 個の条件方程式
\begin{equation}
 (u^{(k)}\mid x^{(l)})=\sum_{\alpha=1}^{n}u_\alpha^{(k)}x_\alpha^{(l)}=0;
 \qquad (k=1,2,\ldots,n-\rho;\ l=1,2,\ldots,\rho).
\tag{2}
\end{equation}
を満たす。以下では、Clebsch\footnote{a. a. O. \S\ 5.} に従って、(1.) に本来現れる比例係数が 1 になるように、これらを正規化しておく。

Clebsch\footnote{a. a. O. \S\ 12.} によれば（また F. Deruyts と A. Capelli の後の研究、序論参照、によっても）、最も一般の形式は、$n-1$ 個の列 $p_\rho$ の各々を高々一つずつ含むものに還元でき、したがって次のように記号的に表すことができる。
\begin{equation}
 (S^1\mid p_1)^{m_1}(S^2\mid p_2)^{m_2}\cdots(S^{n-1}\mid p_{n-1})^{m_{n-1}},
\tag{3}
\end{equation}
ここで (2.) と同様に
\[
 (S^\rho\mid p_\rho)=\sum S^{i_1i_2\ldots i_\rho}p_{i_1i_2\ldots i_\rho}
\]
と置いた。列 $p_\rho$ の要素の間に存在する非線形の恒等的関係のために、記号列 $S^\rho$ は、それらがまったく同じ恒等式を満たすように正規化できる。したがって、それらは、行 $s^{(1)},s^{(2)},\ldots,s^{(\rho)}$ が $x$ に反変的である $\rho$ 行の行列の行列式のように振る舞う。\footnote{Clebsch, a. a. O. \S\ 4.} こうして、各列 $S^\rho$ には、次の関係を通じて別の $S_{n-\rho}$ が一対一に対応する。
\begin{equation}
 S_{i_1i_2\ldots i_{n-\rho}}=(-1)^{\frac{(n-\rho)(n-\rho+1)}2+i_1+i_2+\cdots+i_{n-\rho}}S^{j_1j_2\ldots j_\rho},
\tag{4}
\end{equation}
ここでも $i$ と $j$ はそれぞれ整列しており、合わせて 1 から $n$ までの数をちょうど尽くす。列 $S_{n-\rho}$ の要素 $S_{i_1i_2\ldots i_{n-\rho}}$ は、$x$ と共変的な $(n-\rho)$ 個の列 $\sigma^{(1)},\sigma^{(2)},\ldots,\sigma^{(n-\rho)}$ から作られる行列式のように振る舞い、列 $s$ と $\sigma$ は $\rho(n-\rho)$ 個の条件
\begin{equation}
 (s^{(k)}\mid \sigma^{(l)})=0,\qquad(k=1,2,\ldots,\rho;\ l=1,2,\ldots,n-\rho).
\tag{5}
\end{equation}
によって結ばれる。関係 (1.) および (4.) の結果として、(3.) の各因子は、同値な四重の記号的表示をもつ。
\begin{equation}
 (S^\rho\mid p_\rho)=(S^\rho q_{n-\rho})=(S_{n-\rho}p_\rho)=(-1)^{\rho(n-\rho)}(q_{n-\rho}\mid S_{n-\rho}),
\tag{6}
\end{equation}
ここで、以下つねにそうであるように、$(S^\rho q_{n-\rho})$ および $(S_{n-\rho}p_\rho)$ は、それぞれ行列式 $(s^{(1)}\cdots s^{(\rho)}u^{(1)}\cdots u^{(n-\rho)})$ および $(\sigma^{(1)}\cdots\sigma^{(n-\rho)}x^{(1)}\cdots x^{(\rho)})$ の略記である。Laplace 展開によって、これらはそれぞれ列 $S^\rho,q_{n-\rho}$ および $S_{n-\rho},p_\rho$ のみに依存することが分かり、上の記法はそのことを示すためのものである。式 $(S^\rho\mid p_\rho)$ および $(q_{n-\rho}\mid S_{n-\rho})$ は、前述の議論によれば、列 $s$ と $x$ の行列、ならびに列 $u$ と $\sigma$ の行列の積を意味する。ここで行列積とは、対応する行列式の積の和、すなわち積行列の行列式の値を意味する。

記号列（それぞれ変数列）に属する特性数 $\rho(n-\rho)$ を、Clebsch に従ってその列の重みと呼ぶ。\footnote{a. a. O. \S\ 11.} 列 $s$ の数を与える数 $\rho$ を上重み添字と呼び、列 $\sigma$ の数を与える数 $n-\rho$ を下重み添字と呼ぶ。関係 (4.) から、一つの記号列は、同じ関係の中で一度は上重み添字をもって、また一度は下重み添字をもって現れ得ることが従う。対応する列 $p_\rho$ と $q_{n-\rho}$ の出現についても同じことが成り立つ。最後に、一般に、対応する行列が列 $x$ から成る高い段階の変数列を $p$ で表し、その行列が列 $u$ から成るものを $q$ で表すことを記しておく。$x$ と共変的な記号列は小ギリシア文字 $\sigma,\tau,\alpha,\beta$ で表し、$u$ と共変的な記号列は小ラテン文字 $s,t,a,b$ で表す。その他の記号列はすべて、重み添字を付した大ラテン文字 $S^\rho$, $T^\tau$, $A^\rho$ または $S_{n-\sigma}$, $T_{n-\tau}$, $A_{n-\rho}$ で表す。上または下の重み添字に対応して、たとえば (4.) のように、列の個々の要素も上付きまたは下付きの添字で書く。
\clearpage
\section*{\S\ 2. 行列積による表示.}
以下では常に、「行列積」とは、同じ数の行をもつ二つの行列について、対応する小行列式の積を総和したものを意味する。ただし第二の行列の列は第一の行列の列に反変的であるものとする。各行列の列は、1. 個別に与えられていてもよく、2. (6.) のように一つの列 $S^\rho$ にまとめられていてもよく、3. 異なる列 $S^\rho,S^\sigma\ldots$ に結合されていてもよい。行列積は (6.) で用いた記号 $(\ \mid\ )$ で表す。(6.) の表示は、対応行列に関する定理の第二の定式化の特別な応用である。「各 $\rho$ 行の行列積 $(v^{(1)}v^{(2)}\cdots v^{(\rho)}\mid y^{(1)}\cdots y^{(\rho)})$ には一つの行列式が対応し、逆に各行列式には、あらかじめ定めた行数 $\rho$ をもつ行列積を対応させることができる（$\rho=1,2,\ldots,n-1$）。」\footnote{$\rho>n$ では、行列積はよく知られているように消える。$\rho=n$ では二つの行列式の積になる。} 実際、与えられた行列積から行列式へ、またその逆へ移るには、列 $y$ を一つの列 $p_\rho$ にまとめ（または列 $v$ を一つの列 $q_\rho$ にまとめ）、(1.) で与えられた置換を実行すればよい。ただし、行列式の値はその個々の要素から直接与えられるのではなく、ラプラス展開によって初めて得られる。このように行列式と行列積が同値であることがわかるので、形式の記号的表示可能性に関する定理（第一基本定理）は、次のようにも表される。

\emph{定理 I.} 「すべての不変な形成は行列積によって記号的に表示でき、また逆にすべての行列積は不変な形成である。」\footnote{もちろん、指定された基礎形式の不変な形成とは、記号列 $S^\rho\ldots$ に依存する行列積の集合体だけである。}

行列積による表示は、行列式表示によって生じる非明示的表示の本質的困難を克服させる。\footnote{序論を参照。} ここで「明示的表示」とは、列 $S^\rho,T^\tau,\ldots$ そのもの、またはそれらから一意に導かれる列 $R^\rho$ だけが入り、個々の成分列 $s,t,\ldots$ はもはや現れない表示をいう。\S\ 6 では、記号列 $S^\rho,T^\tau,\ldots$ のみに依存するすべての不変な形成について、これらの記号列において明示的な表示を得る。そしてこれにより、生成過程、折り畳み過程、微分過程を $n$ 元領域へ移すことができる。明示的表示が可能であることは、行列式表示に現れる行列式に正確に対応しうるのが最も単純な行列積だけであり、したがって他のすべての行列積は、異なる行列式を一つの式にまとめたものになる、という事実から明らかである。実際、記号列および変数列を個々の列 $s,u,x$ に分解すれば、通常の形の第一基本定理によって、行列式の集合体が必然的に現れなければならない。逆に、指定された基礎形式の不変な形成を表す行列式の集合体がすべて明示的な行列積にまとめられることは、第二基本定理を用いて証明される（\S\ 6）。

二つの記号列 $S^\sigma,T^\tau$ から、変数列を用いて、基礎形式 $(S^\sigma\mid p_\sigma)(T^\tau\mid p_\tau)$ の不変な形成であってすべての列を明示的に含むもの（折り畳み過程）を得るには、この節の初めに与えた定義に従って第三の種類の行列積を作ればよい。
\begin{equation}
 \pair{q_{n-\sigma-\lambda}S^\sigma}{T_{n-\tau}p_{\tau-\lambda}},
 \qquad (\lambda=1,2,\ldots,\tau,n-\sigma).
\tag{7}
\end{equation}
展開すると (7.) は次を与える。
\begin{equation}
\begin{aligned}
&\pair{v^{(1)}\cdots v^{(n-\sigma-\lambda)}s^{(1)}\cdots s^{(\sigma)}}
       {\alpha^{(1)}\cdots\alpha^{(n-\tau)}y^{(1)}\cdots y^{(\tau-\lambda)}} \\
&\qquad =\sum_{i,k,l}\eps\,q_{k_1\ldots k_{n-\sigma-\lambda}}
 S^{k_{n-\sigma-\lambda+1}\ldots k_{n-\lambda}}
 T_{l_1\ldots l_{n-\tau}}
 p^{l_{n-\tau+1}\ldots l_{n-\lambda}},
\end{aligned}
\tag{8}
\end{equation}
すなわち、列 $S,T,q,p$ のみに依存する式である。ここで $\eps=\pm1$\footnote{この記法は以下全体で保持する。}であり、和は、各列 $S\ldots$ の添字が正しい順序にあり、各組合せ $i_1,i_2,\ldots,i_{n-\lambda}$ に対して、$k$ も $l$ もそれぞれ数 $i$ のすべての可能な順列を走るように理解する。(7.) に入る数 $\lambda$ は、$\sigma>\tau$ の場合に折り畳みの欠損数と呼ぶ。この最後の制限の理由は \S\ 6 で現れる。\footnote{数 $\lambda$ は、常に最後に示された二つの数の小さい方まで走る。}

(7.) に入る二つの行列の行列式は、新しい列 $Q^{n-\lambda},P_{n-\lambda}$ の要素と見なすことができ、それらに対応する列は $Q_\lambda,P^\lambda$ である。(8.) によれば、これらの列 $Q,P$ は、もとの $S$ と $q$、それぞれ $T$ と $p$ によって表せる。これを次の等式で示す。
\begin{equation}
\begin{aligned}
 q_{n-\sigma-\lambda}S^\sigma&=Q^{n-\lambda}\sim Q_\lambda, &\quad\text{または}\quad Q_\lambda&=[q_{n-\sigma-\lambda}S^\sigma]_\lambda,\\
 T_{n-\tau}p_{\tau-\lambda}&=P_{n-\lambda}\sim P^\lambda, &\quad\text{または}\quad P^\lambda&=[T_{n-\tau}p_{\tau-\lambda}]^\lambda.
\end{aligned}
\tag{9}
\end{equation}
(4.) を考慮すると、(6.) と同様に、行列積 (7.) について次の同値な四重の記号的表示を得る。
\begin{equation}
\pair{q_{n-\sigma-\lambda}S^\sigma}{T_{n-\tau}p_{\tau-\lambda}}
=(q_{n-\sigma-\lambda}S^\sigma P^\lambda)=(Q_\lambda T_{n-\tau}p_{\tau-\lambda})=(-1)^{\lambda(n-\lambda)}\pair{P^\lambda}{Q_\lambda}.
\tag{10}
\end{equation}
さらに、(7.) に対して新しい列を用いる別の置換を、次のように置くことによって行うことができる。
\begin{equation}
 \pair{q_{n-\sigma-\lambda}S^\sigma}{T_{n-\tau}p_{\tau-\lambda}}
 =\pair{R^{\sigma+\lambda}}{p_{\sigma+\lambda}}(R^{\tau-\lambda}p_{\tau-\lambda}).
\tag{11}
\end{equation}
ここでも (8.) は、列 $R$ の要素の積が、列 $S$ と $T$ によって一意に表されることを示す。置換 (9.) と (11.) は、特に (7.) によって他の列とのさらなる折り畳みを行うときに用いられる。

行列表現により、列 $p_\rho$ の要素間の二次関係（そこから、知られているように、他のすべての関係が従う）を不変的に述べることができる。したがって、\S\ 1 によれば、基礎形式の係数について線形な、列 $S^\rho$ が満たす関係も述べられる。それらは次の恒等式である。
\begin{equation}
 \pair{q_{n-\rho-\lambda}S^\rho}{S_{n-\rho}p_{\rho-\lambda}}=0,
 \qquad (\lambda=1,2,\ldots,\rho,n-\rho),
\tag{12}
\end{equation}
ここで $p$ と $q$ は任意量と見なす。\S\ 5 で、欠損数 $\lambda$ が奇数である恒等式は、より高い欠損の恒等式から従うことを見る。関係 (12.) は関係 (5.) の直接の帰結である。実際、
\[
 \pair{q_{n-\rho-\lambda}S^\rho}{S_{n-\rho}p_{\rho-\lambda}}
 =(v^{(1)}\cdots v^{(n-\rho-\lambda)}s^{(1)}\cdots s^{(\rho)}\mid \sigma^{(1)}\cdots\sigma^{(n-\rho)}y^{(1)}\cdots y^{(\rho-\lambda)}),
\]
であり、積定理を適用すると、(5.) によって消える $(n-\lambda)$ 行の行列式が得られる。恒等式 (12.) は、すべての折り畳みの型を見つけるためには、各変数列について線形な形式の積を折り畳めば足りることを述べる。あるいはまた、\S\ 6 により、正規化された形式 $(S^\rho\mid p_\rho)^m$ は係数について線形な不変形成をもたない、ということでもある。\footnote{Grassmann の著作、第一巻第二部 p. 510 にある Study の補注（量 $S^\rho$ が単純であるための条件）を参照。}

条件 (12.) は、列 $p_\rho$ を特徴づけるのに十分でもある。というのも、$p_\rho$ の個々の要素がある行列の行列式であるという性質は不変な性質であり、したがって不変的に表現されなければならないからである。しかし、定理 VI によれば、条件 (12.) は $p_\rho$ に対する可能な最大の不変的制限、すなわち列 $p_\rho$ を係数として作られる線形形式が不変形成をまったくもたないという制限を与える。

\section*{\S\ 3. 記号的恒等式.}
ここで、記号的方法の第二基本定理を移さなければならない。すなわち、すべての列 $S^\rho,p_\rho$ を分解不能と見なす場合の、既約で分解不能な恒等式の全体を設定するのである。さらに次の取り決めを置く。記号列と変数列の意味に対応して、変数列 $p_\rho$ を $p_{\rho-1}y$ に特殊化し、その結果の式に系の恒等式の一つを適用して生じる恒等式は合成的なものと見なす。一方、$k$ 個の記号列を含む恒等式は、上の過程によってより少ない記号列の恒等式から生じる場合でも、分解不能と数える。列 $S^\rho,p_\rho$ は必ずしも恒等式に明示的に入る必要はない。これらを分解不能な量として解釈するには、現れる行列積がこれらの列だけに依存することを示せば足りる。しかしこの場合にも、部分的には置換 (11.) を適用することで、常に明示的表示が得られることを示す。一般の恒等式は、行列表現の原理によって、Pascal が与えた特殊な恒等式から得られる。これらは列 $x$ と $u$ だけを含み、Study が三元領域で与えた恒等式の直接の移行である。\footnote{問題の定式化と文献については序論を参照。}
\begin{align}
\sum_{i=1}^n(-1)^{i+1}\pair{a^{(i)}}{x}(a^{(1)}\cdots a^{(i-1)}a^{(i+1)}\cdots a^{(n)}u)&=(-1)^{n+1}\pair{u}{x}(a^{(1)}\cdots a^{(n)}),\tag{13}\\
\sum\pm\pair{a^{(1)}}{x^{(1)}}\pair{a^{(2)}}{x^{(2)}}\cdots\pair{a^{(n)}}{x^{(n)}}&=(a^{(1)}\cdots a^{(n)})(x^{(1)}\cdots x^{(n)}),\tag{14}\\
\sum_{i=1}^n(-1)^{i+1}\pair{u}{a^{(i)}}(a^{(1)}\cdots a^{(i-1)}a^{(i+1)}\cdots a^{(n)}x)&=(-1)^{n+1}\pair{u}{x}(a^{(1)}\cdots a^{(n)}).\tag{15}
\end{align}
(13.) と (15.) からは、置換 $(a^{(i)}\mid x)=(a^{(i)}q_{n-1})$ および $(u\mid a^{(i)})=(p_{n-1}a^{(i)})$ によってさらに二つの恒等式が生じるが、我々の解釈ではそれらは本質的に異なるものとは見なさない。(14.) は行列式の積定理を表す。我々は、(13.) と (14.)（また双対的に (15.) と (14.)）を、適切に形成された行列の積定理の系で置き換える方法を示す。$\rho$ 行の行列と $(\rho-1)$ 行の行列についてそれぞれ作られる積定理は、
\[
\pair{a^{(1)}a^{(2)}\cdots a^{(\rho)}}{xp_{\rho-1}}
=(a^{(1)}a^{(2)}\cdots a^{(\rho)}\mid xy^{(2)}\cdots y^{(\rho)})
=\sum\pm\pair{a^{(1)}}{x}\pair{a^{(2)}}{y^{(2)}}\cdots\pair{a^{(\rho)}}{y^{(\rho)}},
\]
\[
\sum\pm\pair{a^{(2)}}{y^{(2)}}\cdots\pair{a^{(\rho)}}{y^{(\rho)}}
=(a^{(2)}\cdots a^{(\rho)}\mid y^{(2)}\cdots y^{(\rho)})
=(a^{(2)}\cdots a^{(\rho)}p_{\rho-1})=(a^{(2)}\cdots a^{(\rho)}q_{n-\rho+1});
\]
となる。ここから、$\rho=2,3,\ldots,n$ について次の積定理の系が従う。
\begin{equation}
\begin{aligned}
\pair{a^{(1)}a^{(2)}\cdots a^{(\rho)}}{xp_{\rho-1}}
 &=\sum_{i=1}^{\rho}(-1)^{i+1}\pair{a^{(i)}}{x}(a^{(1)}\cdots a^{(i-1)}a^{(i+1)}\cdots a^{(\rho)}\mid p_{\rho-1})\\
 &=\sum_{i=1}^{\rho}(-1)^{i+1}\pair{a^{(i)}}{x}(a^{(1)}\cdots a^{(i-1)}a^{(i+1)}\cdots a^{(\rho)}q_{n-\rho+1}).
\end{aligned}
\tag{16}
\end{equation}
$\rho=n$ では (16.) は (13.) に移る。さらに $p_{n-1}=y^{(2)}p_{n-2}$, $p_{n-2}=y^{(3)}p_{n-3}$ などと特殊化し、恒等式 (16.) を式 $(a^{(1)}\cdots a^{(i-1)}a^{(i+1)}\cdots a^{(n)}\mid y^{(2)}p_{n-2})$, $(a^{(1)}\cdots a^{(i-1)}a^{(i+1)}\cdots a^{(h-1)}a^{(h+1)}\cdots a^{(n)}\mid y^{(3)}p_{n-3})$ などに適用すると、(13.) から (14.) へ移る。したがって恒等式 (14.) は、我々の定義では、恒等式 (16.) から合成されたものである。同値に、系 (16.) は恒等式 (13.) と (14.) に同値である。系 (16.) は一般の恒等式に要求される条件の一つを満たす。すなわち列 $p_{\rho-1}$ を分解不能な量として含み、同時に明示的な形で含む。この条件だけを満たす系はこれだけである。というのも、列 $a,x,p$ からはこれ以上の行列式または行列積を作れず、またそれらの間に (16.) から独立な関係は存在しえないからである。もし存在すれば、$(a^{(1)}\cdots a^{(\rho)}\mid xp_{\rho-1})$ を消去することで、$\rho$ 個（$\rho<n$）の式 $(a^{(i)}\mid x)(a^{(1)}\cdots a^{(i-1)}a^{(i+1)}a^{(\rho)}q_{n-\rho+1})$ の間に関係が生じるが、(13.) によれば少なくとも $n+1$ 個のそのような式が関係に入らなければならないからである。同様の考察から、(15.) と (14.) に同値な系を得る。
\begin{equation}
\begin{aligned}
\pair{u q_{\rho-1}}{a^{(1)}\cdots a^{(\rho)}}
 &=\sum_{i=1}^{\rho}(-1)^{i+1}\pair{u}{a^{(i)}}(q_{\rho-1}\mid a^{(1)}\cdots a^{(i-1)}a^{(i+1)}\cdots a^{(\rho)})\\
 &=\sum_{i=1}^{\rho}(-1)^{i+1}\pair{u}{a^{(i)}}(p_{n-\rho+1}a^{(1)}\cdots a^{(i-1)}a^{(i+1)}\cdots a^{(\rho)}).
\end{aligned}
\tag{17}
\end{equation}
(16.) から任意の列 $A^{\rho_1},B^{\rho_2},\ldots,x,p$ の間の恒等式へ移るため、次を観察する。これらの恒等式は、$A^{\rho_1}=a^{(1)}a^{(2)}\cdots a^{(\rho_1)}$, $B^{\rho_2}=b^{(1)}b^{(2)}\cdots b^{(\rho_2)}$ などと分解すれば直ちに (16.) と同一になるべきであり、また個々の式が (16.) に入る型である限り、最終的な式へまとめることも (16.) で定められた操作によってしか行えない。したがって、さらに
\begin{equation}
 \rho=\rho_1+\rho_2+\cdots+\rho_k\le n,
\tag{18}
\end{equation}
と置けば、
\begin{align*}
&\pair{a^{(1)}\cdots a^{(\rho_1)}b^{(1)}\cdots b^{(\rho_2)}\cdots k^{(1)}\cdots k^{(\rho_k)}}{xp_{\rho-1}}
  =\pair{A^{\rho_1}B^{\rho_2}\cdots K^{\rho_k}}{xp_{\rho-1}} \\
&\quad=\sum_{i=1}^{\rho_1}(-1)^{i+1}\pair{a^{(i)}}{x}(a^{(1)}\cdots a^{(i-1)}a^{(i+1)}\cdots a^{(\rho_1)}B^{\rho_2}\cdots K^{\rho_k}q_{n-\rho+1})\\
&\qquad+(-1)^{\rho_1\rho_2}\sum_{i=1}^{\rho_2}(-1)^{i+1}\pair{b^{(i)}}{x}(b^{(1)}\cdots b^{(i-1)}b^{(i+1)}\cdots b^{(\rho_2)}A^{\rho_1}\cdots K^{\rho_k}q_{n-\rho+1})\\
&\qquad+\cdots\\
&\qquad+(-1)^{(\rho_1+\cdots+\rho_{k-1})\rho_k}\sum_{i=1}^{\rho_k}(-1)^{i+1}\pair{k^{(i)}}{x}(k^{(1)}\cdots k^{(i-1)}k^{(i+1)}\cdots k^{(\rho_k)}A^{\rho_1}\cdots K^{\rho_{k-1}}q_{n-\rho+1}).
\end{align*}
各個別の和（ただしまだ和の一部分ではない）は、実際に列 $A^{\rho_1},B^{\rho_2},\ldots$ を分解不能量として含む。なぜなら、それは必要十分条件\footnote{Capelli, a. a. O. \S\ 23 は十分条件 $\bigl(a^{(k)}\mid \frac{\partial}{\partial a^{(i)}}\bigr)=0$, $(k>i,\ i=1,2,\ldots,\rho_1-1)$ を与える。これらは、Wellstein, Von den Differentialgleichungen der projektiven Invarianten, Math. Ann. 67 (1909) \S\ 3 に従う Poisson 括弧過程を適用することで、最も簡単に必要十分条件を含意する。}
\begin{equation}
 \left(a^{(i+1)}\middle|\frac{\partial}{\partial a^{(i)}}\right)=0;\ \ldots\ ;
 \left(k^{(j+1)}\middle|\frac{\partial}{\partial k^{(j)}}\right)=0.
 \quad (i=1,2,\ldots,\rho_1-1;\ j=1,2,\ldots,\rho_k-1)
\tag{19}
\end{equation}
を満たすからである。さらに、それがすべての列を明示的に含むことを示す。実際、$B^{\rho_2}\cdots K^{\rho_k}q_{n-\rho+1}=q_{n-\rho_1+1}$ と置くと、(16.) と (10.) により、
\begin{align*}
\sum(-1)^{i+1}\pair{a^{(i)}}{x}(a^{(1)}\cdots a^{(i-1)}a^{(i+1)}\cdots a^{(\rho_1)}q_{n-\rho+1})
&=\pair{A^{\rho_1}}{xp_{\rho_1-1}}\\
&=(A_{n-\rho_1}xp_{\rho_1-1})=(-1)^{(\rho_1+1)(n+1)}\pair{q_{n-\rho_1+1}}{A_{n-\rho_1}x}\\
&=(-1)^{(\rho_1+1)(n+1)}\pair{B^{\rho_2}\cdots K^{\rho_k}q_{n-\rho+1}}{A_{n-\rho_1}x}
\end{align*}
を得る。残りの和を同様に計算し、ついで、列はすでに重み添字によって互いに区別されているので、$B^{\rho_2}=A^{\rho_2}$, $C^{\rho_3}=A^{\rho_3},\ldots$ と置くと、求める恒等式を得る。
\begin{equation}
\begin{aligned}
\pair{A^{\rho_1}A^{\rho_2}\cdots A^{\rho_k}}{xp_{\rho-1}}
&=\sum_{i=1}^{k}(-1)^{\delta_i}\pair{A^{\rho_1}\cdots A^{\rho_{i-1}}A^{\rho_{i+1}}\cdots A^{\rho_k}q_{n-\rho+1}}{A_{n-\rho_i}x},\\
\delta_i&=(\rho_1+\rho_2+\cdots+\rho_{i-1})\rho_i+(\rho_i+1)(n+1),
\end{aligned}
\tag{20}
\end{equation}
また (17.) から双対的な恒等式
\begin{equation}
\begin{aligned}
\pair{u q_{\sigma-1}}{A_{\sigma_1}A_{\sigma_2}\cdots A_{\sigma_k}}
&=\sum_{i=1}^{k}(-1)^{\eps_i}\pair{u A^{n-\sigma_i}}{p_{n-\sigma+1}A_{\sigma_1}\cdots A_{\sigma_{i-1}}A_{\sigma_{i+1}}\cdots A_{\sigma_k}},\\
\eps_i&=(\sigma_1+\sigma_2+\cdots+\sigma_{i-1})\sigma_i+(\sigma_i+1)(n+\sigma)
\end{aligned}
\tag{21}
\end{equation}
を得る。
\emph{定理 II.} (20.) と (21.) により、列 $A^{\rho_i},x,p$、それぞれ $A_{\sigma_i},u,q$ の間の分解不能な恒等式の全体が尽くされる。同時に、置換 (11.) を行わずに明示的表示を許す分解不能な恒等式の全体も尽くされる。

定理の第一部は (20.) と (21.) へ至る考察から従う。第二部は、二つの記号列の間の恒等式
\[
\pair{q_\rho A^\sigma}{B_{\rho+\sigma-\tau}p_\tau}
=\eps_1\pair{q_\rho B^{n-\rho-\sigma+\tau}}{A_{n-\sigma}p_\tau}
+\eps_2\pair{q_\rho q_{n-\tau}}{B_{\rho+\sigma-\tau}A_{n-\sigma}},
\]
が不可能であることから従う。ただし $\rho\le\tau\le\sigma$ で、どの列も重み $1\cdot(n-1)$ をもたないとする。（$\rho,\tau,\sigma$ に関する他の仮定は、列の名称を入れ替えることでこれに帰着できる。）この不可能性は、たとえば特殊化 $q_\rho=lM^{\rho-1}$, $q_{n-\tau}=lN^{n-\tau-1}$ のもとで (8.) を展開することにより従う。ここで $l_1=1$, $M^{2,3,\ldots,\rho}=1$, $A^{\rho+1,\ldots,\rho+\sigma}=1$ とし、列 $l,M,A$ の他のすべての要素を零にし、$N$ と $B$ は任意とする。任意の列の間の恒等式は、列 $p_\tau$ を $y^{(1)}y^{(2)}\cdots y^{(\tau)}$ に分解すると、(20.) から合成された恒等式へ移る。したがって、(20.) が二つの記号列だけをもつ恒等式から生成できるならば、第二部は証明される。実際、
\begin{equation}
\pair{A^{\rho_1}A^{\rho_2}}{xp_{\rho-1}}
=\eps\pair{A^{\rho_1}q_{n-\rho+1}}{A_{n-\rho_2}x}
+\pair{A^{\rho_1}}{xp_{\rho_1-1}},
\quad \text{ここで }p_{\rho_1-1}\sim A^{\rho_2}q_{n-\rho+1},
\tag{20a}
\end{equation}
から、$A^{\rho_1}=B^{\sigma_1}B^{\sigma_2}$ と特殊化すること、すなわち
\[
\pair{B^{\sigma_1}B^{\sigma_2}}{xp_{\rho_1-1}}
=\eps\pair{B^{\sigma_1}q_{n-\rho_1+1}}{B_{n-\sigma_2}x}
+\pair{B^{\sigma_1}}{xp_{\sigma_1-1}},
\quad \text{ここで } p_{\sigma_1-1}\sim B^{\sigma_2}q_{n-\rho_1+1},
\]
によって、三つの記号列をもつ恒等式 (20.) が得られる。さらに $B^{\sigma_1}=C^{\tau_1}C^{\tau_2}$ などと特殊化していけば、一般の恒等式 (20.) が得られる。したがって、二つの記号列と二つの任意の変数列の間の恒等式について明示的表示が存在しないことは、任意個の記号列の場合にもそのような表示が存在しないことを含意する。ただし、変数列の中に列 $x$ または $u$ が含まれない場合である。

\section*{\S\ 4. 分解恒等式.}
ここで、任意の記号列と変数列の間の恒等式の最も一般の場合を考える。その最終式では、置換 (11.) を行わずに、列 $p_\tau$ が個々の列 $y^{(1)}\cdots y^{(\tau)}$ に分解された形で現れる。このような恒等式を「分解恒等式」と呼ぶ。前節末の注意に従い、まず記号列が二つだけの場合についてそれらを決定する。すなわち式 $\pair{q_\rho A^\sigma}{B_{\rho+\sigma-\tau}p_\tau}$ を展開する。次を置く。
\begin{equation}
\begin{aligned}
q_{\rho-\alpha}^{(i_1\ldots i_\alpha)}&\sim p_{n-\rho}\,y^{(i_1)}\cdots y^{(i_\alpha)},\qquad
p_{\tau-\alpha}^{(i_1\ldots i_\alpha)}=y^{(i_{\alpha+1})}\cdots y^{(i_\tau)},\qquad
p_\tau=y^{(1)}\cdots y^{(\tau)}.
\end{aligned}
\tag{22}
\end{equation}
数 $\rho,\sigma,\tau$ に関しては四つの場合を区別しなければならない。
\[
\begin{array}{llll}
1)&\rho+\sigma\le n,& \tau\le\rho,& \tau\le\sigma;\\[2pt]
2)&\rho+\sigma\le n,& \tau\ge\rho,& \tau\le\sigma;\\[2pt]
3)&\rho+\sigma\le n,& \tau\le\rho,& \tau>\sigma;\\[2pt]
4)&\rho+\sigma\le n,& \tau\ge\rho,& \tau\ge\sigma.
\end{array}
\]
第一の場合を取り上げ、他の場合はこれから簡単に特徴づける。(20.), (10.), (9.) により、列 $y^{(1)}y^{(2)}\cdots y^{(\tau)}B_{\rho+\sigma-\tau}$ を $y^{(1)}$ と $y^{(2)}\cdots y^{(\tau)}B_{\rho+\sigma-\tau}$ に分解すると、
\begin{equation}
\begin{aligned}
\pair{q_\rho A^\sigma}{y^{(1)}y^{(2)}\cdots y^{(\tau)}B_{\rho+\sigma-\tau}}
&=\pair{q_\rho^{(1)}A^\sigma}{y^{(2)}\cdots y^{(\tau)}B_{\rho+\sigma-\tau}}\\
&\quad+(-1)^\rho\pair{q_\rho[A_{n-\sigma}y^{(1)}]^{\sigma-1}}{y^{(2)}\cdots y^{(\tau)}B_{\rho+\sigma-\tau}}
\end{aligned}
\tag{23}
\end{equation}
を得る。(23.) の最後の項は、もとの式とまったく同じ形をしている。ただし $\sigma$ は $\sigma-1$ に、$\tau$ は $\tau-1$ に置き換わっている。したがってそれは再び方程式 (23.) を許す。よって、$\tau\le\sigma$ であるから、操作 (23.) を $\tau$ 回適用でき、(10.) を考慮して次の関係を得る。
\begin{equation}
\begin{aligned}
\pair{q_\rho A^\sigma}{p_\tau B_{\rho+\sigma-\tau}}
&=(-1)^{\rho\sigma+(\sigma+\tau)(n+1)}
  \pair{q_\rho B^{n-\rho-\sigma+\tau}}{A_{n-\sigma}p_\tau}\\
&\quad+\sum_{i_1=1}^{\tau}(-1)^{\rho(i_1-1)}
 \pair{q_{\rho-1}^{(i_1)}[A_{n-\sigma}y^{(1)}\cdots y^{(i_1-1)}]^{\sigma-i_1+1}}
 {y^{(i_1+1)}\cdots y^{(\tau)}B_{\rho+\sigma-\tau}}.
\end{aligned}
\tag{24}
\end{equation}
和記号の下の各項もまたもとの式の型をしている。$\rho$ は $\rho-1$ に、$\sigma$ は $\sigma-i_1+1$ に、$\tau$ は $\tau-i_1$ に置き換わるので、条件 1) はさらに強く満たされる。$\tau\le\rho$ であるから、操作 (24.) を $\tau$ 回適用でき、求める分解恒等式を得る。
\begin{equation}
\begin{aligned}
\pair{q_\rho A^\sigma}{p_\tau B_{\rho+\sigma-\tau}}
&=\eps_{i_0}\pair{q_\rho B^{n-\rho-\sigma+\tau}}{A_{n-\sigma}p_\tau}\\
&\quad+\sum_{i_1=1}^{\tau}\eps_{i_1}
 \pair{q_{\rho-1}^{(i_1)}B^{n-\rho-\sigma+\tau}}{A_{n-\sigma}p_{\tau-1}^{(i_1)}}\\
&\quad+\sum_{i_2=i_1+1}^{\tau}\eps_{i_2}
 \pair{q_{\rho-2}^{(i_1i_2)}B^{n-\rho-\sigma+\tau}}{A_{n-\sigma}p_{\tau-2}^{(i_1i_2)}}+\cdots\\
&\quad+\eps_{i_\tau}\pair{q_{\rho-\tau}^{(i_1\ldots i_\tau)}B^{n-\rho-\sigma+\tau}}{A_{n-\sigma}}\\
&=\sum_{\alpha=0}^{\tau}\sum_{i_\alpha=i_{\alpha-1}+1}^{\tau}\eps_{i_\alpha}
 \pair{q_{\rho-\alpha}^{(i_1\ldots i_\alpha)}B^{n-\rho-\sigma+\tau}}{A_{n-\sigma}p_{\tau-\alpha}^{(i_1\ldots i_\alpha)}}.
\end{aligned}
\tag{25}
\end{equation}
符号因子 $\eps$ については (24.) から
\begin{equation}
\begin{aligned}
\eps_{i_\alpha}&=(-1)^{\delta_{i_\alpha}},\qquad
\delta_{i_\alpha}=\sum_{\beta=1}^{\alpha}(\rho-\beta+1)(i_{\beta-1}+i_\beta+1)\\
&\quad +(\rho-\alpha)(\sigma+i_\alpha+\alpha)+(\sigma+\tau+\alpha)(n+1)
\end{aligned}
\tag{26}
\end{equation}
を得る。ここで $i_0=0$ と置く。 (25.) の和記号は次のように理解する。$i_1<i_2<\cdots<i_{\alpha-1}$ である各組合せ $i_1i_2\ldots i_{\alpha-1}$ に対し、$i_{\alpha-1}<i_\alpha\le\tau$ を満たすすべての値 $i_\alpha$ を付け加える。(26.) を考慮すると、(25.) の個々の和が (19.) に類似する微分方程式
\begin{equation}
 \left(\frac{\partial}{\partial y^{(i)}}\middle|y^{(i+1)}\right)=0,
 \qquad (i=1,2,\ldots,\tau-1)
\tag{27}
\end{equation}
を満たすことは容易に確かめられる。したがってそれらは列 $p_\tau$ のみに依存する。よって分解恒等式は列 $A,B,q,p$ の間の分解不能な恒等式である。その明示的表示には次節で戻る。

場合 2) では、操作 (23.) は同じく $\tau$ 回適用できるが、操作 (24.) は $\rho$ 回目の適用後に止まる。したがって (25.) の最終式では、第一の和記号を $\sum_{\alpha=0}^{\rho}$ に置き換えなければならない。場合 3) と 4) では、操作 (23.) は $\sigma$ 回だけ行える。そのとき (24.) の代わりに次の関係をもつ。
\begin{equation}
\begin{aligned}
\pair{q_\rho A^\sigma}{p_\tau B_{\rho+\sigma-\tau}}
&=(-1)^{\rho\sigma}
 \pair{q_\rho^{(\sigma+1\ldots\tau)}B^{n-\rho+\tau-\sigma}}{A_{n-\sigma}p_{\tau-\sigma}^{(\sigma+1\ldots\tau)}}\\
&\quad+\sum_{i_1=1}^{\sigma}(-1)^{\rho(i_1-1)}
 \pair{q_{\rho-1}^{(i_1)}[A_{n-\sigma}y^{(1)}\cdots y^{(i_1-1)}]^{\sigma-i_1+1}}
 {y^{(i_1+1)}\cdots y^{(\tau)}B_{\rho+\sigma-\tau}};
\end{aligned}
\tag{28}
\end{equation}
(28.) を $\tau-\sigma$ 回繰り返す。このとき和記号は、和記号下の項に (23.) を $(\sigma-i_1+1)$ 回適用することからわかるように、$\sum_{i_\beta=i_{\beta-1}+1}^{\sigma+\beta-1}$ になる。これにより、(25.) の個別和のうち添字 $\alpha=\tau-\sigma$ に対応するすべての項が、対応する符号とともに得られる。すると場合 3) と 4) は場合 1) と 2) に移る。なぜなら $\sigma,\tau$ に対応する数は $\sigma'=\sigma-i_{\tau-\sigma}+\tau-\sigma$, $\tau'=\tau-i_{\tau-\sigma}=\sigma'$ となり、手続きが続くにつれて $\tau'<\sigma'$ となるからである。したがって (25.) の最終式では、第一の和記号を $\sum_{\alpha=\tau-\sigma}^{\tau,\rho}$ に置き換えなければならない。

同様に (21.) からは、列 $q_\rho$ が個々の列 $v^{(1)},\ldots,v^{(\rho)}$ に分解される双対的な分解恒等式が得られる。もし
\begin{equation}
 \dot p_{\tau-\beta}^{(i_1\ldots i_\beta)}\sim v^{(i_1)}\cdots v^{(i_\beta)}q_{n-\tau};\qquad
 \dot q_{\rho-\beta}^{(i_1\ldots i_\beta)}=v^{(i_{\beta+1})}\cdots v^{(i_\rho)};\qquad
 q_\rho=v^{(1)}\cdots v^{(\rho)},
\tag{29}
\end{equation}
と置けば、
\begin{equation}
 \pair{q_\rho A^\sigma}{p_\tau B_{\rho+\sigma-\tau}}
 =\sum_{\beta=\max(0,\tau-\sigma)}^{\min(\tau,\rho)}
   \sum_{i_\beta=i_{\beta-1}+1}^{\rho}\eps
\pair{\dot q_{\rho-\beta}^{(i_1\ldots i_\beta)}B^{n-\rho-\sigma+\tau}}{A_{n-\sigma}\dot p_{\tau-\beta}^{(i_1\ldots i_\beta)}} ,
\tag{30}
\end{equation}
となる。ここで和の添字 $\beta$ は、下限の大きい方から上限の小さい方まで走る。(25.) と (30.) における個別和のすべての項は、一つの形式の極であり、それらは極化過程
\[
\left(\frac{\partial}{\partial p_{\tau-\alpha}}\middle|p_{\tau-\alpha}^{(i_1\ldots i_\alpha)}\right),\qquad
\left(q_{\rho-\alpha}^{(i_1\ldots i_\alpha)}\middle|\frac{\partial}{\partial q_{\rho-\alpha}}\right)
\]
によって作られる。これを表すため、こうした極化操作を定数倍して集めたものを $\Pi$ と書く。すると (25.) と (30.) の代わりに、
\begin{equation}
 \pair{q_\rho A^\sigma}{p_\tau B_{\rho+\sigma-\tau}}
 =\sum_{\alpha=\max(0,\tau-\sigma)}^{\min(\tau,\rho)}
 \Pi\pair{q_{\rho-\alpha}B^{n-\rho-\sigma+\tau}}{A_{n-\sigma}p_{\tau-\alpha}}
\tag{31}
\end{equation}
を得る。

\section*{\S\ 5. 明示的な形での分解恒等式.}
(25) の場合 4) で $\tau$ を $\sigma+\rho$ に特殊化すると、
\begin{equation}
 \pair{q_\rho A^\sigma}{y^{(1)}\cdots y^{(\rho+\sigma)}}
 =\sum_{1\le i_1<\cdots<i_\rho\le \rho+\sigma}
 \eps_{i_\rho}\pair{q_\rho}{y^{(i_1)}\cdots y^{(i_\rho)}}
 \pair{A^\sigma}{y^{(i_{\rho+1})}\cdots y^{(i_{\rho+\sigma})}},
\tag{32}
\end{equation}
を得る。ここで $\eps_{i_\rho}=(-1)^{\delta_{i_\rho}}$ かつ $\delta_{i_\rho}=\rho(\rho+1)/2+i_1+\cdots+i_\rho$ である。これは (16.) と (17.) より一般的な行列の積定理の形であり、積行列の行列式
\[
 \sum\pm(v^{(1)}y^{(1)})\cdots(v^{(\rho)}y^{(\rho)})(a^{(1)}y^{(\rho+1)})\cdots(a^{(\sigma)}y^{(\rho+\sigma)})
\]
をラプラス展開することで従う。したがって、より一般の積定理は恒等式 (20.) と (21.) からそれぞれ合成され、特殊な形 (16.), (17.) の選択もこれによって説明される。

関係 (32.) は、分解恒等式を明示的に表示する手段を与える。このため、(31.) に現れる形式を基礎形式と見なす。すなわち置換 (11.) を行う。この置換は、列 $A,B$ における明示的表示を何も変えない。なぜなら、(8.) により、新しく導入される列 $R$ は、成分列 $a^{(1)}a^{(2)}\ldots$, $b^{(1)}b^{(2)}\ldots$ ではなく、列 $A,B$ そのものによって表せるからである。そこで
\begin{equation}
 \pair{q_{\rho-\alpha}B^{n-\rho-\sigma+\tau}}{A_{n-\sigma}p_{\tau-\alpha}}
 =\pair{R_{\rho-\alpha}p_{n-\rho+\alpha}}{R^{\tau-\alpha}p_{\tau-\alpha}}
\tag{33}
\end{equation}
と置く。すると (25.) の添字 $\alpha$ に対応する個別和について、
\begin{equation}
\begin{aligned}
&\sum \eps_{i_\alpha}\pair{R_{\rho-\alpha}p_{n-\rho}y^{(i_1)}\cdots y^{(i_\alpha)}}{R^{\tau-\alpha}y^{(i_{\alpha+1})}\cdots y^{(i_\tau)}}\\
&\quad=\sum \eps_{i_\alpha}\pair{[R_{\rho-\alpha}p_{n-\rho}]^\alpha y^{(i_1)}\cdots y^{(i_\alpha)}}{R^{\tau-\alpha}y^{(i_{\alpha+1})}\cdots y^{(i_\tau)}}\\
&\quad=\eps\pair{[R_{\rho-\alpha}p_{n-\rho}]^\alpha R^{\tau-\alpha}}{p_\tau}
 =\eps\pair{R^{\tau-\alpha}q_{n-\tau}}{R_{\rho-\alpha}p_{n-\rho}}
\end{aligned}
\tag{34}
\end{equation}
を得る。このように実際に明示的な最終式が得られ、また $q_\rho$ と $p_\tau$ が現れる対称的な形は、(30.) も同じ最終式へ導くことを示す。したがってこの最終式はもとの分解に依存しない。(25.) と (30.) の間に差があるのは、列 $y$ と $v$ が独立の意味をもっている間だけであり、そのとき分解恒等式は我々の定義により分解可能な恒等式へ移る。より多くの記号列をもつ恒等式へ進むには、\S\ 3（末尾）に従って、列 $q_\rho$ を $q_{\rho_1}C^{\sigma_1}$ に、または $p_\tau$ を $p_{\tau_1}C_{\tau-\tau_1}$ に分解すればよい。得られる式
\begin{equation}
 \pair{q_{\rho_1}C^{\sigma_1}}{[R^{\tau-\alpha}q_{n-\tau}]_\lambda R_{\rho+\sigma_1-\alpha-\lambda}},
 \qquad
 \pair{[R_{\rho-\alpha}p_{n-\rho}]^\lambda R^{\tau-\alpha}}{p_{\tau_1}C_{\tau-\tau_1}}
\tag{35}
\end{equation}
は、二つの記号列の場合に現れる型とまったく同じであり、したがって (33.) に対応する置換の後に再び明示的表示を許す。これを繰り返すことにより、任意個の列について明示的表示を得る。また、(25.) および (30.) の最初と最後の和は、それぞれ (33.) を適用することなく、公式 (32.) だけで明示的表示を与えるか、あるいはすべての列を明示的に含む単一項に完全に還元されることも注意しておく。(32.) から、$q_\rho$ を $q_{\rho_1}C^{\sigma_1}$ などに分解して得られる関係は、最も一般的な形の積定理と見なすことができる。

要約すれば、次を示す。

\emph{定理 III:} 恒等式
\begin{equation}
 \pair{q_\rho A^\sigma}{p_\tau B_{\rho+\sigma-\tau}}
 =\sum_{\alpha=\max(0,\tau-\sigma)}^{\min(\tau,\rho)}
 \eps\pair{R^{\tau-\alpha}q_{n-\tau}}{R_{\rho-\alpha}p_{n-\rho}}
\tag{36}
\end{equation}
と、ここで列 $R$ は (33) によって決定されるものとし、さらに $q_\rho,p_\tau$ を $q_{\rho_1}C^{\sigma_1}$ などへ分解してそこから得られる恒等式を合わせると、分解不能な恒等式の全体が尽くされる。

実際、得られる恒等式は、その種類のうち最も一般的である。なぜなら、個々の列は、(20.) と (21.) の場合のように重みや数に関する制限を受けないからである。また、これらの列の間にはそれ以上の恒等式は存在しない。列 $p_\tau$ を $y^{(1)}\cdots y^{(\tau)}$ に分解すると、最も一般の恒等式は (20.) から、したがって \S\ 3（末尾）により (20a.) から合成される。そして (16.) の後に論じた合成は、まさに \S\ 4 で行われた合成である。任意個の列への移行は、(20a.) の議論が示すように、$q_\rho$ を $q_{\rho_1}C^{\sigma_1}$ などに分解して生じる式へ、常に同じ操作を適用することで与えられる。また、(20a.) ではなく (20.) から直接移行しても新しいものは得られない。これは計算によって容易に確かめられる。一方では (25.) において列 $q_\rho$ を一度特殊化し、他方では \S\ 4 の方法で (20.) から移行を行う。どちらの場合にも同じ和が生じ、それらは一般の積定理（上を見よ）の適用により、(36.) から得られる恒等式へ移る。恒等式 (20.), (21.) は特殊な場合 $\tau=1$, $\rho=1$ に対応する。そのとき個別和は二つだけ現れ、したがって前の注意に従って、置換 (33.) なしに、以前に得たものと一致する明示的表示が得られる。

最後に、分解恒等式の二つの特殊例を簡単に述べる。(31) で $\tau=\sigma$, $\rho=n-\sigma$ と置き、列 $p_\tau$ を $p'_\sigma\sim q'_{n-\sigma}$ と書く。すると
\begin{equation}
\begin{aligned}
\pair{A^\sigma}{p_\sigma}\pair{B^\sigma}{p'_\sigma}
-\pair{A^\sigma}{p'_\sigma}\pair{B^\sigma}{p_\sigma}
=\sum_{\alpha=1}^{\min(\sigma,n-\sigma)}
\Pi\pair{q_{n-\sigma-\alpha}B^\sigma}{A_{n-\sigma}p_{\sigma-\alpha}}.
\end{aligned}
\tag{37}
\end{equation}
これは二つの共変的列 $p_\sigma,p'_\sigma$ をもつ形式を基本共変式へ展開する Clebsch の公式である。\footnote{Clebsch, loc. cit., pp. 25--32. Clebsch は、右辺の個々の項が列 $A,B$ のみに依存することを証明している。明示的表示は、彼の式 $\Phi_h$ に (32.) を適用することで得られる。} 形式 $\pair{A^\sigma}{p_\sigma}^m\pair{B^\sigma}{p'_\sigma}^n$ に属する基本共変式は次のものである。
\begin{equation}
 \left\{\sum_{\alpha=1}^{\min(\sigma,n-\sigma)}
 \eps\pair{q_{n-\sigma-\alpha}B^\sigma}{A_{n-\sigma}p_{\sigma-\alpha}}
 \right\}^{\rho}
 \pair{A^\sigma}{p_\sigma}^{m-\rho}\pair{B^\sigma}{p'_\sigma}^{n-\rho},
 \qquad (\rho=0,1,\ldots,m,n).
\tag{38}
\end{equation}
これらは、簡単に言えば、基礎形式から「欠損数 $\rho$ の共変的折り畳み」によって生じる（\S\ 2 参照）。

第二に、$\tau=\sigma-\lambda$, $\rho=n-\sigma-\lambda$, $B^\sigma=A^\sigma$ と置く。すると
\begin{equation}
\begin{aligned}
\pair{q_{n-\sigma-\lambda}A^\sigma}{A_{n-\sigma}p_{\sigma-\lambda}}
&=(-1)^\lambda\pair{q_{n-\sigma-\lambda}A^\sigma}{A_{n-\sigma}p_{\sigma-\lambda}}\\
&\quad+(-1)^{(n-\sigma)(\sigma-\lambda)}
 \sum_{\alpha=1}^{\min(\sigma-\lambda,n-\sigma-\lambda)}
 \Pi\pair{q_{n-\sigma-\lambda-\alpha}A^\sigma}{A_{n-\sigma}p_{\sigma-\lambda}}.
\end{aligned}
\tag{39}
\end{equation}
したがって、\S\ 2 で述べたように、奇数の欠損 $\lambda$ に対して記号列を特徴づける条件 (12.) は、より高い欠損の条件から従う。線形形式 $\pair{A^\sigma}{p_\sigma}=\pair{B^\sigma}{p_\sigma}$ については、(39.) は、係数について二次の既約な不変形成が偶数欠損のものに還元されることを含意する。これは、二元および三元領域で線形形式が不変形成をもたないという既知の事実と一致する。

一般の恒等式は、個々の列が条件 (12.) を満たすという仮定のもとで初めに導かれたことに注意すべきである。しかし、これらの個々の列は恒等式に線形にしか入らないので、恒等式は条件 (12.) を満たさないより一般の列についても成り立つ。

\section*{\S\ 6. 不変形成の明示的表示.\\折り畳み過程と微分過程.}
前節までの結果により、\S\ 2 で述べた記号的表示可能性の定理（第一基本定理）を、明示的表示可能性を加えて完成できる。すなわち次を示す。

\emph{定理 IV:} 「任意の基礎形式のすべての不変形成は行列積によって明示的に表示できる。すなわち、変数列 $p_\rho$ のほかには、もとの形式の列 $S^\sigma,\ldots,T^\tau$、または置換 (9.) あるいは (11.) によってそれらから導かれる列 $P,Q,R$ だけが最終式に入る。」

証明のため、ある不変形成が、変数列のほかに列 $S^\sigma,\ldots,T^\tau$ に依存するとする。また、これらの列が、行列式による表示を可能にする程度まで個々の列 $S^\sigma=S^{\sigma_1}\cdots S^{\sigma_i},\ldots,T^\tau=T^{\tau_1}\cdots T^{\tau_k}$ に分解されているとする。列を、任意ではあるが一度限り固定された順序 $S^\sigma,\ldots,T^\tau$ に並べる。最初の全体列 $S^\sigma=S^{\sigma_1}\cdots S^{\sigma_i}$ に依存し、個々の成分列だけには依存しない行列式の集合体を取り出す。この依存性は、\S\ 3--\S\ 5 により一般恒等式の帰結である。ここで変数列 $p_\rho$ を列 $y,v$ に、または後に現れる列 $T$ の一つを個々の列 $t$、すなわち $\tau$ に分解すると、最も一般の恒等式は (20.) と (21.) から合成される。しかしこれらの恒等式は常に、列 $S^\sigma=S^{\sigma_1}\cdots S^{\sigma_i}$ を明示的に含む式へ導く。すなわち (20.) と (21.) の左辺である。というのも、(19.) により、右辺の全和だけが、(20a.) の一項に対応する一部分ではなく、$S^\sigma$ のみに依存する性質をもつことが確認されるからである。手続きが続くにつれて、一度明示的に入ったすべての列は明示的なまま残る。恒等式の適用が要求しうるのは、せいぜい複数の列を一つにまとめること、すなわち置換 (9.) だけである。最後に、一つの列 $T^\tau$ だけを明示的な形にしなければならない場合には、二つの場合がある。まだ自由な変数列、すなわち置換 (9.) によって他の列と結びついていない列が残っているかもしれない。その列を成分列 $y$ または $v$ に分解すれば手続きは完了し、(31.) のように、すべての列を明示的に含む一つまたは複数の形式の極が得られる。自由な変数列が残っていなければ、個々の列 $t$ または $\tau$ を含むが列 $T^\tau$ に依存する式全体の中に、それらが生じた仕方から、分解恒等式の項を認めることができる。しかし \S\ 5 により、これらは置換 (11.) を適用した後、常にすべての列において明示的な表示をもつ。これで定理は完全な一般性において証明される。二つの記号列だけが現れる場合には、置換 (11.) は決して現れないことも注意しておく。なぜなら、$\sigma,n-\sigma,\tau,n-\tau<n$ なので、変数列がまったく入らないのは二つの記号列が一つの行列式に結合され、したがって明示的に現れる場合だけであり、しかし変数列が入るやいなや、(34.) と (33.) は置換 (11.) が不要になることを示すからである。

今述べたことから、すべての不変形成を生成するには、変数列を付け加えたのち、記号列を可能なすべての行列積に結合すれば足りることが従う。最も一般の行列積は今や
\begin{equation}
 \pair{P^{\mu_1}\cdots P^{\mu_i}}{Q_{\nu_1}\cdots Q_{\nu_k}},\qquad \sum\mu=\sum\nu,
\tag{40}
\end{equation}
の形をしている。ここで $P$ と $Q$ は、もとの形式の列でもよく、もとの列から置換 (9.) または (11.) の列を経て生じたものでもよく、最後には変数列でもよい。しかし式 (40.) は、変数列の助けを借りて、二つの記号列を一度に一つの行列積に結合していくことで生成できる。実際、
\[
 \pair{q_{n-\mu-\lambda}P^\mu}{Q_{\nu_1}p_{n-\mu-\nu_1-\lambda}}=\pair{R_\lambda Q_{\nu_1}}{p_{n-\mu-\nu_1-\lambda}}
\]
から、列 $Q_{\nu_2}$ と結合することにより、
\[
 \pair{[R_\lambda Q_{\nu_1}]^{n-\lambda-\nu_1}}{Q_{\nu_2}p_{n-\mu-\nu_1-\nu_2-\lambda}}
 =\pair{q_{n-\mu-\lambda}P^\mu}{Q_{\nu_1}Q_{\nu_2}p_{n-\mu-\nu_1-\nu_2-\lambda}}
\]
という式が得られる。したがってこの手続きを続ければ式 (40.) が得られる。$P$ と $Q$ がもとの形式の列でない場合にも、(9.) と (11.) は、今示したことと合わせて、それらも二つの記号列を一度に結合する列によって生じたことを意味する。したがって、

\emph{定理 V:} すべての不変形成を生成する過程、すなわち折り畳み過程は、変数列の助けを借りて、二つの記号列を一度に一つの行列積に結合していくことから成る。

二つの記号列 $S^\sigma,T^\tau$ について、$\sigma\ge\tau$ とすると、次の行列積が作れる。
\begin{align}
&\pair{q_{n-\sigma-\lambda}S^\sigma}{T_{n-\tau}p_{\tau-\lambda}}, &&(\lambda=1,2,\ldots,\min(\tau,n-\sigma)),\tag{41}\\
&\pair{q_{n-\tau-\mu}T^\tau}{S_{n-\sigma}p_{\sigma-\mu}}, &&(\mu=\sigma-\tau+\lambda),\tag{42}\\
&\pair{q_{n-\tau-\nu}T^\tau}{S_{n-\sigma}p_{\sigma-\nu}}, &&(\nu=1,2,\ldots,\sigma-\tau).
\tag{43}
\end{align}
しかし (31.) から、(42.) と (43.) について次の値が得られる。
\begin{align}
&\pair{q_{n-\sigma-\lambda}S^\sigma}{T_{n-\tau}p_{\tau-\lambda}}
 +\sum_{\alpha}\Pi\pair{q_{n-\sigma-\lambda-\alpha}S^\sigma}{T_{n-\tau}p_{\tau-\lambda-\alpha}},\tag{42a}\\
&\Pi(q_{n-\sigma}S^\sigma)(T_{n-\tau}p_\tau)
 +\sum_{\beta}\Pi\pair{q_{n-\sigma-\beta}S^\sigma}{T_{n-\tau}p_{\tau-\beta}}.
\tag{43a}
\end{align}
したがって、形式 (42.) は (41.) とその極によって線形に表され、形式 (43.) は基礎形式と (41.) の形式の極によって表される。同様に (31.) は、形式 (41.) が (42.) とその極に線形に依存することを与える。したがって Clebsch の定義\footnote{Clebsch, loc. cit., \S\ 7.}によれば、形式 (42.) は形式 (41.) と同値であり、(43.) は基礎形式へ戻る。したがって生成過程は (41.) または (42.) によって同じ権利で与えられる。我々は、数 $\lambda$ に関してより簡単な過程 (41.) を折り畳み過程と定義する。数 $\lambda$、すなわち折り畳みの欠損数（\S\ 2 参照）は、行列式積から欠けている列の数を与える。より正確には、その折り畳みにおいて最大の列数を含む行列積において欠けている列数である。$\sigma\ge\tau$ のもとで、$\lambda$ は二つの値 $\tau,n-\sigma$ の小さい方までしか走らないので、
\begin{equation}
 \lambda\le\frac n2,\ n\text{ は偶数};\qquad
 \lambda\le\frac{n-1}{2},\ n\text{ は奇数}.
\tag{44}
\end{equation}
となる。(42.) と (43.) の (41.) への還元は、さらなる折り畳みによって列 $p,q$ が記号列に移った場合にも有効である。実際、(36.) により
\[
 \pair{A^{n-\tau-\varkappa}T^\tau}{S_{n-\sigma}B_{\sigma-\varkappa}}
 =\sum_{\alpha=\max(0,\varkappa-\sigma+\tau)}^{\min(\tau,n-\sigma)}
 \eps\pair{R^{\tau-\alpha}B^{n-\sigma+\varkappa}}{A_{\tau+\varkappa}R_{n-\sigma-\alpha}}
\]
をもつ。ここで列 $R$ は (33.) に従って $S$ と $T$ から生じたものである。したがって、記号列だけを含むこれらの形式は、$\sigma-\tau\gtreqless 2\varkappa$ に応じて、$\pair{q_\tau A^{n-\tau-\varkappa}}{B_{\sigma-\varkappa}p_{n-\sigma}}$ または $\pair{q_{\tau-\alpha}B^{n-\sigma+\varkappa}}{A_{\tau+\varkappa}p_{n-\sigma-\alpha}}$ を、基礎形式および形式 (41.) と折り畳むことによって得られる。

記号的折り畳み過程からは、通常どおり、記号列 $S^\sigma,T_{n-\tau}$ を微分記号の列 $\frac{\partial}{\partial p_\sigma}$, $\frac{\partial}{\partial q_{n-\tau}}$ で置き換えるだけで、不変な微分過程が得られる。よく知られているように、積 $\frac{\partial}{\partial p_{i_1\ldots i_\sigma}}\frac{\partial}{\partial q_{j_1\ldots j_{n-\tau}}}$ は、そのとき実際には $\frac{\partial^2}{\partial p_{i_1\ldots i_\sigma}\partial q_{j_1\ldots j_{n-\tau}}}$ を意味する。列 $\frac{\partial}{\partial p_\sigma}$ と $\frac{\partial}{\partial q_{n-\tau}}$ は列 $S^\sigma,T_{n-\tau}$ と共変的であるから、それらと同じ恒等式を満たし、とくに (42.) と (42a.) の間、および (43.) と (43a.) の間に存在する恒等式も満たす。要約すれば、次を得る。

\emph{定理 VI:} 任意の基礎形式のすべての不変形成は、それらから、記号的折り畳み過程 (41.) を反復適用することによって、または不変な微分過程
\begin{equation}
 \left(q_{n-\sigma-\lambda}\frac{\partial}{\partial p_\sigma}\middle|\frac{\partial}{\partial q_{n-\tau}}p_{\tau-\lambda}\right),
 \qquad(\sigma\ge\tau,\ \lambda=1,2,\ldots,\min(\tau,n-\sigma))
\tag{45}
\end{equation}
によって得られる。
なお、各折り畳み (41.) において形式の総重量、すなわち個々の変数列の重みの和（\S\ 1 参照）は、正の数 $2(\lambda^2+\lambda(\sigma-\tau))$ だけ減少する。これは、一般の有限性証明とは独立に、形式列の有限性、すなわち、与えられた基礎形式の係数に線形な不変形成全体の有限性を含意する。\footnote{Clebsch, loc. cit., \S\ 11 と 12、基本共変式の系の有限性を参照。}

さらに、定理 VI は条件 (12.) を満たさない形式についてもなお成り立つことを注意しておく。(3) の各因子 $\pair{S^\rho}{p_\rho}^{m_\rho}$ は、$m_\rho$ 個の線形因子の積 $\pair{A^\rho}{p_\rho}\pair{B^\rho}{p_\rho}\cdots\pair{M^\rho}{p_\rho}$ で置き換えることができる。その場合、不変形成は線形形式の系の不変形成へ移り、後でそれらを互いに等しいものとすればよい。しかし各単独の線形形式については、条件 (12.) は、微分記号を導入すると二階の微分商だけを含むものであり、常に満たされる。

\section*{\S\ 7. 一般形式の正規形による展開.}
以下（\S\ 8）では、折り畳み過程 (41.) の主要な性質、すなわちそれが基本的折り畳みから合成できることを導く。そのために、四元領域で Mertens\footnote{F. Mertens, \emph{"Uber invariante Gebilde quatern"arer Formen}（序論参照）。以下「Mertens」と引用する。} が与えた重要な定理を $n$ 元領域へ移す必要がある。この定理は、任意の有限個の形式の系を、同値な正規形の系で置き換えることができるというものである。「正規形」とは、二元および三元の定義\footnote{Study, loc. cit., p. 54.}を一般化して、係数に線形な不変形成が、基礎形式自身を除いて恒等的に消える形式をいう。したがって定理 VI（\S\ 6）により、正規形 $\varphi$ は微分方程式系
\begin{equation}
 \left(q_{n-\sigma-\lambda}\frac{\partial}{\partial p_\sigma}\middle|\frac{\partial}{\partial q_{n-\tau}}p_{\tau-\lambda}\right)\varphi=0,
 \qquad(\sigma\ge\tau,\ \lambda=1,2,\ldots,\min(\tau,n-\sigma))
\tag{46}
\end{equation}
を満たす。ここで列 $q_{n-\sigma-\lambda},p_{\tau-\lambda}$ は任意量と見なす。四元領域では、$\sigma=\tau=2$ について (39.) を考慮すれば、これらの微分方程式は、Mertens が任意の列 $p,q$ を導入せずに与えた十個の微分方程式（loc. cit., 公式 28）と完全に一致する。それらを知ることと定理 VI とにより、彼の証明をほとんど字句どおり $n$ 変数へ移すことができる。付け加えるべきことはただ一つ、構成された形式が正規形であることの検証であり、これは分解恒等式 (30.) によって得られる。四元領域ではこの検証はさらなる変形なしに従っていた。したがって、ここでは証明の短い概説を与えれば十分である。簡単のため、以下で必要となる場合、すなわち一つの基礎形式 $f$ とその形式列（\S\ 6 末尾参照）の場合を扱う。任意個の基礎形式の係数に任意次数で現れる不変形成についても、証明は同じである。

証明の基本的な考えは、変換された係数を導入することにより、$n-1$ 個の列 $p_\rho$ をもつ形式を、$x$ と共変的な $n$ 個の列 $\xi$ をもつ形式に置き換えることである。これらの列は変換量の列である。求める正規形は、これらの形式から不変な微分過程によって直接得られる。一般形式の正規形による展開は、$\rho$ 個（$\rho<n$）の共変的列 $\xi$ をもつ形式の級数展開によって媒介される。不変な微分過程は、元の列 $p_\rho$ における形式へ戻す。

変換量の列を $\xi^{(1)},\ldots,\xi^{(n)}$ とし、
\begin{equation}
\begin{aligned}
x_i&=\xi_i^{(1)}x'_1+\xi_i^{(2)}x'_2+\cdots+\xi_i^{(n)}x'_n,\\
p_{i_1\ldots i_\rho}&=\sum_j(\xi_{i_1}^{(j_1)}\cdots \xi_{i_\rho}^{(j_\rho)})p'_{j_1\ldots j_\rho}
\end{aligned}
\tag{47}
\end{equation}
であるとする。形式 $f,\varphi,\ldots$ の変換された係数を $(f),(\varphi),\ldots$ と書き、
\[
 f=\pair{S^1}{p_1}^{m_1}\pair{S^2}{p_2}^{m_2}\cdots\pair{S^{n-1}}{p_{n-1}}^{m_{n-1}}
\]
とする。すると
\begin{equation}
\begin{aligned}
(f)&=\pair{S^1}{\xi^{(1)}}^{m_{11}}\cdots\pair{S^1}{\xi^{(n)}}^{m_{1n}}
\pair{S^2}{\xi^{(1)}\xi^{(2)}}^{m_{21}}\cdots
\pair{S^{n-1}}{\xi^{(2)}\cdots\xi^{(n)}}^{m_{n-1,n}}\\
&\quad\cdot\pair{s^{(1)}}{\xi^{(1)}}^{k_1}\pair{s^{(2)}}{\xi^{(2)}}^{k_2}\cdots\pair{s^{(n)}}{\xi^{(n)}}^{k_n}
\end{aligned}
\tag{48}
\end{equation}
である。ここで $m_{11}+\cdots+m_{1n}=m_1$ などである。各係数 $(f)$ について、もし
\begin{equation}
 k_1\ge k_2\ge k_3\ge\cdots\ge k_n
\tag{49}
\end{equation}
ならば、列 $\xi$ をちょうど尽くす微分過程
\begin{equation}
\left(\frac{\partial}{\partial\xi^{(1)}}\middle|p_1\right)^{k_1-k_2}
\left(\frac{\partial}{\partial\xi^{(1)}}\frac{\partial}{\partial\xi^{(2)}}\middle|p_2\right)^{k_2-k_3}
\cdots
\left(\frac{\partial}{\partial\xi^{(1)}}\cdots\frac{\partial}{\partial\xi^{(n-1)}}\middle|p_{n-1}\right)^{k_{n-1}-k_n}
\left(\frac{\partial}{\partial\xi^{(1)}}\cdots\frac{\partial}{\partial\xi^{(n)}}\right)^{k_n}
\tag{50}
\end{equation}
によって、形式
\begin{equation}
 \varphi=\bigl(s^{(1)}\mid p_1\bigr)^{k_1-k_2}
 \bigl(s^{(1)}s^{(2)}\mid p_2\bigr)^{k_2-k_3}\cdots
 \bigl(s^{(1)}\cdots s^{(n-1)}\mid p_{n-1}\bigr)^{k_{n-1}-k_n}
 \bigl(s^{(1)}s^{(2)}\cdots s^{(n)}\bigr)^{k_n}
\tag{51}
\end{equation}
が得られる。これらの形式 $\varphi$ が求める正規形である。実際、それらは正規形を定義する次の性質をもつ。

1. これらは基礎形式 $f$ の係数に線形な不変形成である。これは、それらが $f$ から不変な微分過程 $\left(\frac{\partial}{\partial p_\rho}\middle|\xi^{(i_1)}\xi^{(i_2)}\cdots\xi^{(i_\rho)}\right)$ と (50.) によって生じることから従う。したがって（\S\ 6 末尾により）それらは $f$ の有限な形式列の極によって線形に表される。これらの極は、$\varphi$ に入る変数 $p_\rho$ を、変数列 $p_\rho$ をもつ線形形式に分解された形式の係数と見なすことで得られる。
\begin{equation}
 H=(q_{n-1}\mid p_{n-1})^{k_1-k_2}(q_{n-2}\mid p_{n-2})^{k_2-k_3}\cdots(q_1\mid p_1)^{k_{n-1}-k_n}.
\tag{52}
\end{equation}
形式列 $f$ と $H$ の同時不変式は、必要な極をすべて与える。なぜなら、それらは $\varphi$ 自身と同じ次元を個々の列 $p_\rho$ においてもたなければならないからである。

2. それらは微分方程式 (46.) を満たす。微分過程 (45.) を $\varphi$ に適用すると、
\[
 \varphi'=\bigl(q_{n-\sigma-\lambda}s^{(1)}\cdots s^{(\sigma)}\mid [s^{(1)}\cdots s^{(\tau)}]_{n-\tau}p_{\tau-\lambda}\bigr),\qquad (\sigma>\tau)
\]
となる。列 $s$ から作られる列 $q_\sigma=s^{(1)}\cdots s^{(\sigma)}$, $p_{n-\tau}=[s^{(1)}\cdots s^{(\tau)}]_{n-\tau}$ を、そこに現れる $q_\rho,p_\tau$ と同一視すると、分解恒等式 (30.) から
\[
 \varphi'=\sum_{\beta=\sigma-\tau+\lambda}^{n-\tau,\sigma}\sum_{i_\beta}\eps\,
 \pair{q_{\sigma-\beta}^{(i_1\ldots i_\beta)}q_{n-\tau+\lambda}}{p_{\sigma+\lambda}p_{n-\tau-\beta}^{(i_1\ldots i_\beta)}}
\]
が従う。ここで
\[
 p_{n-\tau-\beta}^{(i_1\ldots i_\beta)}\sim s^{(1)}\cdots s^{(\tau)}s^{(i_1)}\cdots s^{(i_\beta)},
\]
であり、$i_1,\ldots,i_\beta$ は 1 から $\sigma$ までの数から選んだ任意の相異なる $\beta$ 個の数を表す。$\beta>\sigma-\tau$ であるから、これらの数 $i_1,\ldots,i_\beta$ の中には 1 から $\tau$ の間のものが必ずある。したがって $p_{n-\tau-\beta}^{(i_1\ldots i_\beta)}$ は恒等的に消え、ゆえに各行列積も、したがって $\varphi'$ も消える。

正規形の系が形式列 $f$ と同値であることを示すには、形式列のすべての形式が、1. で定めた意味で、正規形の極に展開できることを示せばよい。$\Pi$ を極化操作
\begin{equation}
 \left(\frac{\partial}{\partial\xi^{(i)}}\middle|\xi^{(k)}\right),
 \qquad (i\ne k;\ i=1,2,\ldots,\rho;\ k=1,2,\ldots,\rho)
\tag{53}
\end{equation}
の集合体とし、$\Pi^\sigma$ を、$i=\sigma$ かつ $k<\sigma$ である特殊な操作 (53.) の集合体とする。列 $\xi^{(1)},\ldots,\xi^{(\rho)}$ の形式 $F$ が最後の列 $\xi^{(\rho)}$ に関して次数 $k_\rho$ をもつとき、展開\footnote{F. Mertens, \emph{"Uber eine Formel der Determinantentheorie}. Sitzungsberichte der Akademie der Wissenschaften, Wien, Math.-Naturw. Kl., vol. 91, 2. Abteilung (1885), formula 20).}
\begin{equation}
 F=\sum_{\alpha=0}^{k_\rho}\Pi F_\alpha,
 \qquad(\rho\le n)
\tag{54}
\end{equation}
が成り立つ。ここで $F_\alpha$ は最後の列 $\xi^{(\rho)}$ において次数 $\alpha$ をもち、$F$ から不変な微分操作
\begin{equation}
 \left(\frac{\partial}{\partial\xi^{(1)}}\cdots\frac{\partial}{\partial\xi^{(\rho)}}\middle|\xi^{(1)}\cdots\xi^{(\rho)}\right)^\alpha
\tag{55}
\end{equation}
および $\Pi^\rho$ によって生じる。$F_\alpha$ を構成する際に、過程 (55.) を
\begin{equation}
 \left(\frac{\partial}{\partial\xi^{(1)}}\cdots\frac{\partial}{\partial\xi^{(\rho)}}\middle|p_\rho\right)^\alpha
\tag{56}
\end{equation}
で置き換えると、$\rho-1$ 個の列 $\xi^{(1)},\ldots,\xi^{(\rho-1)}$ だけをもつ形式 $F_\alpha(p_\rho)$ が得られ、これに再び (54.) を適用できる。この手続きを続けると、形式 $G(p_\rho,p_{\rho-1},\ldots,p_1)$ に至る。これらから (54.) に従う $F$ の展開を得るには、個々の列 $p_\rho$ を $\xi^{(1)},\ldots,\xi^{(\rho)}$ で置き換え、得られた形式に極化過程 $\Pi$ (53.) を適用する。これにより（(48.) 参照）変換された係数 $(G)$ の和が得られる。$\rho=n$ の場合には、過程 (55.) は第一段階で残る。$c$ を数値定数とし、簡単のため
\begin{equation}
 (\xi^{(1)}\xi^{(2)}\cdots\xi^{(n)})\sim\Delta;
 \qquad
 \left(\frac{\partial}{\partial\xi^{(1)}}\frac{\partial}{\partial\xi^{(2)}}\cdots\frac{\partial}{\partial\xi^{(n)}}\right)=\nabla
\tag{57}
\end{equation}
と置くと、$\rho=n$ の場合には
\begin{equation}
 F=\sum c\Delta^\alpha(G),
\tag{58}
\end{equation}
となる。一方、$\rho<n$ では右辺には $\alpha=0$、したがって $\sum c(G)$ だけが現れる。

(58.) を変換された係数 $(f)$ (48.) に適用し、極化過程 $\Pi^\sigma$ が $\rho\ge\sigma$ であるすべての操作 (56.) と可換であること、また変換された係数の極は再び変換された係数を与えることを用いると、
\[
 G=\left(\frac{\partial}{\partial\xi^{(1)}}\middle|p_1\right)^{\beta_1}
 \left(\frac{\partial}{\partial\xi^{(1)}}\frac{\partial}{\partial\xi^{(2)}}\middle|p_2\right)^{\beta_2}
 \cdots
 \left(\frac{\partial}{\partial\xi^{(1)}}\cdots\frac{\partial}{\partial\xi^{(n-1)}}\middle|p_{n-1}\right)^{\beta_{n-1}}
 \nabla^{\beta_n}\Pi(f)=\sum c\varphi
\]
を得る。さらに (58.) から
\begin{equation}
 (f)=\sum c\Delta^\alpha(\varphi)\qquad\text{（Mertens, 公式 23）}
\tag{59}
\end{equation}
が従う。(59.) から $f$ 自身の展開へ移るため、再び変数 $p_\rho$ を、$f$ に対応する次数 $m_1,\ldots,m_{n-1}$ をもつ形式 $H$ (52.) の係数と見なし、$m=m_1+\cdots+m_{n-1}$ と置く。すると不変式の定義により
\[
 f\cdot\Delta^m=\sum(f)(H)=\sum c\Delta^\alpha(\varphi)(H)
\]
であり、列 $\xi$ をちょうど尽くす過程 $\nabla^m$ を適用すると
\begin{equation}
 f=\sum c\Phi,
\tag{60}
\end{equation}
を得る。ここで形式 $\Phi$ は、変数に関する不変な微分過程によって $\varphi$ と $H$ から導かれ、したがって形式列 $\varphi$ と $H$ の同時不変式である。\footnote{直接の結論の代わりに、Mertens, \S\ 7 は、形式列 $H$ を形成するために本質的に用いられる一連のさらなる微分過程を導入する。これは我々の定理 VI によって行われる。} しかし $\varphi$ は正規形であるから、形式列 $\varphi$ は $\varphi$ 自身に還元される。形式列 $H$ は可能なすべての不変な変数結合を含む。したがって $\varphi$ と $H$ の同時不変式は、1. で示した意味での $\varphi$ の極である。言い換えれば、(60.) は $f$ の正規形の極による展開を与える。任意の形式列 $f$ の形式 $\theta$ についても同じ展開を確認する必要がある。$(\theta)$ から (50.) によって導かれる正規形を $\Gamma$ とし、
\begin{equation}
 (\theta)=\sum c\Delta^\beta(\Gamma)
\tag{61}
\end{equation}
とする。形式 $\theta$ は、各変換係数について有効な関係
\[
 (\theta)\Delta^\sigma=\sum c(f)\qquad\text{（Mertens, 公式 9）}
\]
によって特徴づけられる。しかし、因子 $\Delta^\sigma$ を含む $n$ 個の列 $\xi$ の任意の形式 $F$ は、第二の形式 $\nabla^\sigma\Pi F$ も含む（Mertens, 公式 20）。したがって
\[
 (\theta)=\sum c\nabla^\sigma(f),\qquad\text{したがって}\qquad \Gamma=\sum c\varphi,
\]
となる。(61.) から (59.) に対応する関係
\[
 (\theta)=\sum c\Delta^\beta(\varphi)
\]
を得る。これは $\theta$ について、正規形 $\varphi$ の極による展開 (60.) を含意する。したがって、

\emph{定理 VII.} 各形式列 $f$ は、同値な正規形の系で置き換えることができる。

\section*{\S\ 8. 折り畳み過程の基本的折り畳みへの還元.}
定理 VII に続いて、\S\ 7 の冒頭で述べた、折り畳み過程の基本的折り畳みへの還元を実行する。我々は（\S\ 5 参照）、二つの共変的列 $S^\rho,T^\rho$ の折り畳みをすべて「共変的折り畳み」と呼び、特に二つの列 $S^\rho,T^\rho$ の欠損数 1 の共変的折り畳みを「型 $\rho$ の基本的折り畳み」と呼ぶ。すると、定理 VI を補完して次を証明する。

\emph{定理 VIII:} 一般の折り畳み過程 (41.) は、型 $\rho$ の $(n-1)$ 個の基本的折り畳みから合成できる。ただし $\rho=1,2,\ldots,n-1$ である。言い換えれば、すべての不変形成を生成するには、$n-1$ 個の基本的折り畳みを逐次適用すれば足りる。

二元領域では、折り畳み過程 (41.) は唯一可能な基本的折り畳みと直接一致する。三元領域では、私は学位論文の \S\ 1 で定理 VIII を証明した（序論参照）。そこでは (44.) のため、一般の共変的折り畳みはまだ基本的折り畳みと同一である。$n$ 変数についての定理 VIII が、単に共変的折り畳みへのはるかに弱い還元ではなく、基本的折り畳みへの還元を与えることは注目に値する。証明は原理的には三元の場合の証明に倣う。すなわち、最も一般の折り畳みを、記号的恒等式を用いて基本的折り畳みから構成する。生成するのは次の三つである。
\begin{enumerate}
\item 欠損数 1 の任意の折り畳みを、欠損数 1 の共変的折り畳み、すなわち基本的折り畳みから生成する（三元の場合の類似）。
\item 欠損数 $\lambda$ の任意の折り畳みを、欠損数 $\lambda$ の共変的折り畳みと欠損数 1 の任意の折り畳みから生成する。
\item 欠損数 $\lambda$ の共変的折り畳みを、欠損数 $\lambda-1$ の共変的折り畳みと欠損数 1 の任意の折り畳みから生成する。
\end{enumerate}
本質的なこととして、定理 VII により、折り畳まれる二つの形式 $S$ と $T$ は正規形であると仮定する。したがって (46.) により、
\begin{equation}
\begin{aligned}
\pair{q_{n-\sigma_1-\lambda}S^{\sigma_1}}{S_{n-\sigma_2}p_{\sigma_2-\lambda}}&=0,
&& (\sigma_1\ge\sigma_2,\ \lambda=1,2,\ldots,\sigma_2,n-\sigma_1),\\
\pair{q_{n-\tau_1-\lambda}T^{\tau_1}}{T_{n-\tau_2}p_{\tau_2-\lambda}}&=0,
&& (\tau_1\ge\tau_2,\ \lambda=1,2,\ldots,\tau_2,n-\tau_1)
\end{aligned}
\tag{62}
\end{equation}
が成り立つ。さらに \S\ 2（末尾）により、形式 $S$ と $T$ がすべての変数列に関して線形であると仮定すれば十分である。すなわち、構成される形式は形式列
\[
 f=\pair{S^1}{p_1}\pair{S^2}{p_2}\cdots\pair{S^{n-1}}{p_{n-1}}
   \pair{T^1}{p_1}\pair{T^2}{p_2}\cdots\pair{T^{n-1}}{p_{n-1}}
\]
に属する。
1) $\sigma>\tau$ である $\pair{q_{n-\sigma-1}S^\sigma}{T_{n-\tau}p_{\tau-1}}$ を、型 $\tau+\alpha$, $\alpha=0,1,\ldots,\sigma-\tau$ の基本的折り畳みから生成する。型 $\tau$ の基本的折り畳み
\[
 \pair{q_{n-\tau-1}S^\tau}{T_{n-\tau}p_{\tau-1}}
\]
から、置換 $w\sim T_{n-\tau}p_{\tau-1}$ により、上の重み添字（\S\ 1 参照）$\tau+1$ をもつ三つの列をもつ形式列 $f$ の形式
\begin{equation}
 \pair{wS^\tau}{p_{\tau+1}}\pair{S^{\tau+1}}{p_{\tau+1}}\pair{T^{\tau+1}}{p_{\tau+1}}
\tag{63}
\end{equation}
を得る。(63.) に型 $\tau+1$ の基本的折り畳みを適用すると
\[
 \pair{q_{n-\tau-2}wS^\tau}{S_{n-\tau-1}p_\tau};
\]
が得られる。(31.) から、置換 $q_{n-\tau-2}w=q'_{n-\tau-1}$ により、これは値
\[
 \eps\pair{q_{n-\tau-2}wS^{\tau+1}}{S^\tau p_\tau}
 +\sum_{\alpha=1}^{\tau,n-\tau-1}\Pi\pair{q'_{n-\tau-1-\alpha}S^{\tau+1}}{S_{n-\tau}p_{\tau-\alpha}}
\]
をもつ。和記号下の式は (62.) によって恒等的に消える。したがって我々は
\[
 \pair{wS^{\tau+1}}{p_{\tau+2}}\pair{S^{\tau+2}}{p_{\tau+2}}\pair{T^{\tau+2}}{p_{\tau+2}}
\]
という形式に到達した。これは (63.) で $\tau$ を $\tau+1$ に変えたものである。この過程を $\sigma-\tau$ 回繰り返すと、求める折り畳み
\[
 \pair{wS^\sigma}{p_{\sigma+1}}=\eps\pair{q_{n-\sigma-1}S^\sigma}{T_{n-\tau}p_{\tau-1}}
\]
に到達する。今行った過程を公式
\begin{equation}
 T^\tau S^\tau,
\quad S^{\tau+1},\quad S^{\tau+2},\ldots,S^\sigma;
\tag{64}
\end{equation}
で示す。この過程では $S$ の正規形性だけが用いられ、$T$ のそれは用いられないことがわかる。(64.) に双対な過程も、$T$ の正規形性だけを用いて、逆順に行える。すなわち
\begin{equation}
 S^\sigma T^\sigma,
\quad T^{\sigma-1},\quad T^{\sigma-2},\ldots,T^\tau
\tag{65}
\end{equation}
であり、関係式
\begin{align*}
\pair{q_{n-\sigma-1}S^\sigma}{T_{1-\sigma}p_{\sigma-1}}&=\eps\pair{T_{n-\sigma}y}{p_{\sigma-1}},\quad y\sim q_{n-\sigma-1}S^\sigma,\\
\pair{q_{n-\sigma}T^{\sigma-1}}{T_{n-\sigma}yp_{\sigma-2}}&=\eps\pair{T_{n-\sigma+1}y}{p_{\sigma-2}},\\
&\vdots\\
\pair{q_{n-\tau-1}T^\tau}{T_{n-\tau-1}yp_{\tau-1}}&=\eps\pair{q_{n-\sigma-1}S^\sigma}{T_{n-\tau}p_{\tau-1}}
\end{align*}
に従う。\S\ 6（末尾）で与えた総重量の減少との関係も指摘しておく。欠損数 1 の折り畳みでは、この減少は $2(1+(\sigma-\tau))$ であり、基本的折り畳みでは $2\cdot1$ である。したがって、基本的折り畳みからの合成がそもそも可能であるならば、ちょうど $\sigma-\tau+1$ 個のそのような折り畳みが必要であり、上で実際にそれを行った。

2) 一般の折り畳みを共変的折り畳みと基本的折り畳みから生成する。一般の折り畳み (41) の重量減少は $2(\lambda^2+\lambda(\sigma-\tau))=2[\lambda^2+(\sigma-\tau)(1+(\lambda-1))]$ である。これは、それを欠損数 $\lambda$ の共変的折り畳み（重量減少 $2\lambda^2$）と、重み添字の差が $\lambda-1$ である $\sigma-\tau$ 個の欠損数 1 の折り畳み（重量減少 $2\{1+(\lambda-1)\}$）に分解できる可能性を示す。$\lambda=1$ の場合、この分解は 1) で与えたものと一致する。これが実際に実現できることを示す。

欠損数 $\lambda$ の共変的折り畳み
\[
 \pair{q_{n-\tau-\lambda}S^\tau}{T_{n-\tau}p_{\tau-\lambda}}
\]
から、置換 $R^\lambda\sim T_{n-\tau}p_{\tau-\lambda}$ により、上の重み添字 $\tau+\lambda$ と $\tau+1$ の二つの列をもつ形式
\begin{equation}
 \pair{R^\lambda S^\tau}{p_{\tau+\lambda}}\pair{S^{\tau+1}}{p_{\tau+1}}
\tag{66}
\end{equation}
を得る。より低い重み添字 $\tau+1$ をもつ列は正規形に属する。したがって (66.) は、(65.) の順序で $\lambda$ 個の基本的折り畳みから合成された欠損数 1 の折り畳みを許す。
\[
 (R^\lambda S^\tau)S^{\tau+\lambda},\quad S^{\tau+\lambda-1},\quad S^{\tau+\lambda-2},\ldots,S^{\tau+1}.
\]
(31.) と (62.) を考慮すると、これは
\[
 \pair{q_{n-\tau-\lambda-1}R^\lambda S^\tau}{S_{n-\tau-1}p_\tau}=\eps\pair{R^\lambda S^{\tau+1}}{p_{\tau+\lambda+1}}
\]
を与える。すなわち、1) の (63.) と同様に、(66.) で $\tau$ を $\tau+1$ に置き換えた形式である。この過程を $\sigma-\tau$ 回繰り返すと、求める折り畳み
\[
 \pair{R^\lambda S^\sigma}{p_{\sigma+\lambda}}=\eps\pair{q_{n-\sigma-\lambda}S^\sigma}{T_{n-\tau}p_{\tau-\lambda}}
\]
に至る。全過程で用いたのは $S$ の正規形性だけである。$T$ を正規形と見ると、順序を完全に逆にした双対過程が、次の公式に従って得られる。
\begin{align*}
\pair{q_{n-\sigma-\lambda}S^\sigma}{T_{n-\sigma}p_{\sigma-\lambda}}&=\eps\pair{T_{n-\sigma}R_\lambda}{p_{\sigma-\lambda}},\quad R_\lambda\sim q_{n-\sigma-\lambda}S^\sigma,\\
\pair{q_{n-\sigma}T^{\sigma-1}}{T_{n-\sigma}R_\lambda p_{\sigma-\lambda-1}}&=\eps\pair{T_{n-\sigma+1}R_\lambda}{p_{\sigma-\lambda-1}},\\
&\vdots\\
\pair{q_{n-\tau-1}T^\tau}{T_{n-\tau-1}R_\lambda p_{\tau-\lambda}}&=\eps\pair{q_{n-\sigma-\lambda}S^\sigma}{T_{n-\tau}p_{\tau-\lambda}}.
\end{align*}
3) 欠損数 $\lambda$ の共変的折り畳みを、欠損数 $\lambda-1$ の共変的折り畳みと基本的折り畳みから生成する。欠損数 $\lambda$ の共変的折り畳みの重量減少は $2\lambda^2=2((\lambda-1)^2+2\lambda-1)$ である。したがってそれは、欠損数 $\lambda-1$ の共変的折り畳みと $2\lambda-1$ 個の基本的折り畳みに分解できる可能性を示す。これはまず $n=2\lambda$ と置くと最も容易に見いだされる。欠損数 $\lambda$ の可能な共変的折り畳みは $\pair{S^\lambda}{T_\lambda}=\pair{S^\lambda}{T_{n/2}}$ だけである。これは同じ列の欠損数 $\lambda-1$ の折り畳み $\pair{uS^\lambda}{T_\lambda x}$ よりも、列 $u$ と列 $x$ を一つずつ少なく含む。逆に、欠損数 $\lambda$ の折り畳みを、列 $u$ と列 $x$ の消滅を実現し（重量減少 $2(n-1)=2(2\lambda-1)$）、同時に新しい記号列 $R^\lambda$ を生成する基本的折り畳みと、欠損数 $\lambda-1$ の共変的折り畳みとから合成する。これを行う最も単純な折り畳みは、次の欠損数 1 の折り畳みである。
\[
 \pair{sT^\lambda}{\sigma p_\lambda}=\eps\pair{S^{2\lambda-1}}{R_{\lambda-1}p_\lambda}=\eps\pair{R^\lambda}{p_\lambda},
\]
ここで
\[
 R^{\lambda+1}=sT^\lambda,
 \quad s=S^1,
 \quad \sigma=S_1,
 \quad R^\lambda\sim\sigma R_{\lambda-1}
\]
である。(65.) と (64.) により、これは $2\lambda-1$ 個の基本的折り畳み
\begin{equation}
 T^\lambda S^\lambda,S^{\lambda-1},\ldots,S^1;
 \qquad R^{\lambda+1}S^{\lambda+1},S^{\lambda+2},\ldots,S^{2\lambda-1}
\tag{67}
\end{equation}
から合成される。(21.) と (62.) を考慮した欠損数 $\lambda-1$ の折り畳みは
\[
 \pair{uS^\lambda}{\sigma R_{\lambda-1}x}=\eps\pair{R^{\lambda+1}}{S_\lambda x}
 =\eps\pair{sT^\lambda}{S^\lambda x}=\eps\pair{S^\lambda}{T_\lambda}
\]
を与える。二つの列 $S^\rho,T^\rho$ と任意の $n$ についての一般の場合は、この特殊な場合から直接得られる。(67.) の代わりに、次の $2\lambda-1$ 個の基本的折り畳みをもつ。
\begin{equation}
 T^\rho S^\rho,S^{\rho-1},\ldots,S^{\rho-\lambda+1};
 \qquad R^{\rho+1}S^{\rho+1},S^{\rho+2},\ldots,S^{\rho+\lambda-1}.
\tag{68}
\end{equation}
折り畳みの前半は形式
\[
 \pair{q_{n-\rho-1}T^\rho}{S_{n-\rho+\lambda-1}p_{\rho-\lambda}}
\]
を与える。ここから置換
\[
 S_{n-\rho+\lambda-1}p_{\rho-\lambda}\sim w,
 \qquad T^\rho w=R^{\rho+1}
\]
と後半の折り畳みの適用により、
\[
 \pair{q_{n-\rho-\lambda}S^{\rho+\lambda-1}}{R_{n-\rho-1}p_\rho}
\]
を得る。さらに
\[
 q_{n-\rho-\lambda}S^{\rho+\lambda-1}\sim y
\]
と置換すると、欠損数 $\lambda-1$ の共変的折り畳みにより
\begin{equation}
 \pair{q_{n-\rho-\lambda+1}S^\rho}{yR_{n-\rho-1}p_{\rho-\lambda+1}}
\tag{69}
\end{equation}
を得る。これが求める欠損数 $\lambda$ の折り畳みであることを示す。
\[
 R_{n-\rho-1}p_{\rho-\lambda+1}=R_{n-\lambda}
\]
と置き、$y$ を分解すると
\[
 \pair{q_{n-\rho-\lambda+1}R^\lambda}{S_{n-\rho}y}
 =\eps\pair{q_{n-\rho-\lambda}S^{\rho+\lambda-1}}{S_{n-\rho}p'_{\rho-1}}=0
\]
となることを観察する。したがって (20.) により、(69.) の値は
\[
 \eps\pair{S^\rho R^\lambda}{p_{\rho+\lambda-1}y}
 =\eps\pair{q_{n-\rho-\lambda}S^{\rho+\lambda-1}}{[S^\rho R^\lambda]_{n-\rho-\lambda}p_{\rho+\lambda-1}}
\]
である。この折り畳みは型 (43.) であり、したがって
\[
 \pair{q_{n-\rho-\lambda}S^\rho}{R_{n-\rho-1}p_{\rho-\lambda+1}}
 =\eps\pair{wT^\rho}{R_\lambda p_{\rho-\lambda+1}},
 \qquad R_\lambda\sim q_{n-\rho-\lambda}S^\rho
\]
に同値である。ここで $R_\lambda$ はすでに望む最終形になっている。この式は特殊な場合の形式 $\pair{sT^\lambda}{S_\lambda x}$ に対応する。今、$w$ と $R_\lambda$ を分解すると
\[
 \pair{wR^{n-\lambda}}{T_{n-\rho}p_{\rho-\lambda+1}}
 =\eps\pair{q'_{n-1}R^{n-\lambda}}{S_{n-\rho+\lambda-1}p_{\rho-\lambda}}
 =\eps\pair{q''_{n-\sigma-1}S^\rho}{S_{n-\rho+\lambda-1}p_{\rho-\lambda}}=0
\]
となることから、(21.) により値
\[
 \pair{wq_{n-\rho+\lambda-1}}{T_{n-\rho}R_\lambda}
 =\eps\pair{q_{n-\rho+\lambda-1}[T_{n-\rho}R_\lambda]^{n-\lambda}}{S_{n-\rho+\lambda-1}p_{\rho-\lambda}}
 \quad\text{これは}\quad
 \pair{q_{n-\rho-\lambda}S^\rho}{T_{n-\rho}p_{\rho-\lambda}}
 \text{に同値である}
\]
を得る。全過程で用いたのは $S$ の正規形性だけであり、$T$ のそれではない。したがって、(69.) を生成するために用いた欠損数 $\lambda-1$ の共変的折り畳みは、$2\lambda-3$ 個の基本的折り畳みと欠損数 $\lambda-2$ の共変的折り畳みに置き換えることができる。後者もさらに基本的折り畳みに分解でき、以下同様である。同様に、$T$ の正規形性を用い、双対的過程を適用することによっても生成が得られる。すなわち基本的折り畳み
\[
 S^\rho T^\rho,\quad T^{\rho-1},\ldots,T^{\rho-\lambda+1};
 \qquad R^{\rho-1}T^{\rho-1},\quad T^{\rho-2},\ldots,T^{\rho-\lambda+1},
\]
と、(69.) に双対な折り畳みである。そのとき前のものに同値な折り畳み
\[
\pair{q_{n-\rho-\lambda}T^\rho}{S_{n-\rho}p_{\rho-\lambda}}
\]
が得られる。

2) で得た分解を加えると、最も一般の折り畳みが確かに基本的折り畳みから生成されることがわかった。残るのは、この生成が順序の置換を除いて一意であること、すなわち各折り畳みに対して、型 $\rho$ の基本的折り畳みの数 $\alpha_\rho$ が一つだけ対応することを示すことである。ここで $\rho=1,2,\ldots,n-1$ である。

初期形式 $f$ の変数列 $p_\rho$ における次数を $m_\rho$ とし、折り畳みによって作られた形式の次数を $l_\rho$ とする。型 $\rho$ の各基本的折り畳みは、次数 $m_{\rho-1},m_\rho,m_{\rho+1}$ をそれぞれ $m_{\rho-1}+1,m_\rho-2,m_{\rho+1}+1$ に変える。したがって $\alpha$ は方程式系
\begin{equation}
 m_\rho-l_\rho=-\alpha_{\rho-1}+2\alpha_\rho-\alpha_{\rho+1},
 \qquad (\rho=1,2,\ldots,n-1)
\tag{70}
\end{equation}
を満たす。ここで $\alpha_0=\alpha_n=0$ と置く。この方程式系の行列式の値は $n$ である。\footnote{$\Delta_n$ が $n$ 行の行列式を表すなら、漸化式 $\Delta_n=2\Delta_{n-1}-\Delta_{n-2}$ が成り立つ。すなわち $\Delta_n-\Delta_{n-1}=\Delta_{n-1}-\Delta_{n-2}=\cdots=\Delta_2-\Delta_1=1$ である。なぜなら $\Delta_2=3$, $\Delta_1=2$ だからである。} したがって $\alpha$ は一意に定まり、しかも定理 VIII により正の整数として定まる。これにより差 $m_\rho-l_\rho$ についての条件が得られる。

この節の結果は関係 (62.) に基づく。しかし、もし関係 (62.) が成り立たないなら、最終式は初期式と決して同値ではない。\S\ 6 で与えた同値性の定義は、次の定式化も許す。\footnote{次節末の考察を参照。}「二つの形式が同値であるとは、それらの差がより高い折り畳みの形式、すなわちより大きい重量減少を示す形式である場合、かつその場合に限る。」このより大きい重量減少は、折り畳みに入る二つの列 $q_{n-\sigma-i},p_{\tau-\lambda}$ が実際に変数列であるときには常に起こる。これは (41.), (42.)、この節で同値として現れた折り畳み、また定理 VII の同値な形式において起こる。しかし、これらの列が置換によって記号列から導かれているときには、必ずしも起こるとは限らない。この節で、(62.) によって互いに表すことができた形式がその場合である。消える形式でさえ、部分的にはより高い総重量をもつことが容易にわかる。

\section*{\S\ 9. 形式列. 還元定理.}
ここで、\S\ 6（末尾）で導入した形式列を、やや厳密に定義する。「与えられた初期形式の係数に線形な不変形成の全体、すなわち、その初期形式をそれ自身で折り畳むことによって生じる形式の全体を、より高い形式に従って配列したもの」である。

より高い形式を説明するため、定理 VII により初期形式が二つの正規形 $S\cdot T$ の積として与えられているとする。より高い形式とは、それ自身においてより高い折り畳みをもつ形式、そして同じ折り畳みをもつ形式の中では特別に正規化された形式をいう。より正確には、形式 $A$ が型 $\rho$ の基本的折り畳みを $\alpha_\rho$ 回行って生じ、形式 $B$ がそのような折り畳みを $\beta_\rho$ 回行って生じたとする。このとき $B$ が $A$ より高いとは、次の場合である。
\begin{enumerate}
\item 不等式系 $\beta_\rho\ge\alpha_\rho$, $\rho=1,2,\ldots,n-1$ において、少なくとも一つの $\rho$ について不等号が真に成り立つ場合。\footnote{不等式系の中で $>$ の符号と $<$ の符号が部分的に混在するなら、形式は比較できない。なぜなら、ある形式から折り畳みによって生じる形式は常に基本的折り畳みの正の値系をもち、したがってこの場合には正または零の値 $\beta_\rho-\alpha_\rho$ をもつからである。}
\item すべての $\rho$ について等号が成り立つが、より少ない記号列が一つの行列積に結合されている場合。この正規化は \S\ 8 の還元のために必要である。したがって、たとえば $\pair{S^{n-1}}{T_{n-1}}$ は $\pair{S^{n-1}}{[S^1T^\rho]_{n-\rho-1}p_\rho}$ より高い。
\item すべての $\rho$ について等号が成り立ち、一つの行列積に結合される記号列の数も同じである場合には、任意ではあるが一度限り固定された正規化によって $B$ を $A$ より高くする。この正規化は、値系 $\alpha$ に属する形式が有限個であるため常に可能である。たとえば一つの形式 $S$ を優先する（記号列 $S$ の数が多い、列 $S$ の上の重み添字の和が大きい、など）ことによって実現できる。
\end{enumerate}
形式列は、$(n-1)$ 次元の格子点の多様体として配列されていると考えると便利である。その $n-1$ 個の方向は $n-1$ 個の基本的折り畳みに対応する。各形式には一つ、かつただ一つの値系 $\alpha$、すなわち正の整数格子点が対応し、同じ値系 $\alpha$ に属する異なる形式は、上の正規化によってその順序が一意に定まる。形式列の任意の形式は、新しい形式列の初期形式を与え、その形式列は全形式列の一部として含まれる。実際、それは新しい初期形式から $n-1$ 個の正方向へ進むことによって得られる形式からなる部分として含まれる。

定理 VIII と級数展開の一般理論により、形式列については、二元および三元領域の還元定理（学位論文、\S\ 3 参照）がそのまま成り立つ。主定理を簡単に導く。まず定義を述べる。

「与えられた形式系において、ある有限個の形式が一つの加群にまとめられているとする。ある形式が、因子として不変式をもつ形式によって、または『より高い形式』によって表せるなら、その形式を可約と呼ぶ。すなわち、その形式が属する全形式列のより高い形式によって、またはその加群の記号をより高い次数で含む形式によって表せる場合である。可約な形式列を還元子と呼ぶ。」すると、

\emph{定理 IX:} ある形式列の初期形式が、還元子との折り畳みによって生じたために可約であるなら、この初期形式から生じる形式列全体は可約である。

証明のため、形式 $f$ と第二の形式 $g$ を折り畳んで生じる形式 $h$ は、(45.) により、いくつかの共変的な変数列 $p_\rho,p'_\rho$ を含む形式 $f$ と見なせることを観察する。このような形式は、級数展開の一般理論により、各対の共変的変数列 $p_\rho,p'_\rho$ に順に級数展開を適用することで、$f$ の基本共変式の極として表される。(38.) によれば、現れる基本共変式は共変的折り畳みによって生成される形式である。しかし、級数展開を適用する順序はなお任意であるから、とくに一般の折り畳みを基本的折り畳みから生成する順序に対応するものを選べる。得られる基本共変式の系は、形式列 $f$ と同一になる。欠損数の高い共変的折り畳みによって生じる形式は、基本的折り畳みによって生じる形式の中に配列される。しかし、級数展開の任意の順序に対しても形式列 $f$ の極が得られ、それに第二の基本共変式の系が対応することがある。形式列は不変形成の全体を尽くすので、第二の系の各形式は形式列の極によって線形に表され、しかも同じか、またはより大きい重量減少を示す形式の極によって表される。後者は、極が（\S\ 7 参照）形式 $h$ に対応する次数をもつ形式列 $H$ (52.) と形式列 $f$ との同時不変式と見なされることから従う。$H$ のそれ自身における各折り畳みは重量減少に対応し、置換 (11.) が行われた後、この減少は記号列へ、したがって $f$ の形式へ移される。\footnote{この考察は、同値性概念の第二の定式化（\S\ 8 末尾）の基礎にもなっている。三元領域における基本共変式の異なる系についてのさらなる説明は、Study, loc. cit., II, \S\ 7, 定理 7) および 8) にある。我々の定理 VII により、同じ結果は $n$ 変数についても成り立つ。}

形式列 $h$ の任意の形式は、$h$ をそれ自身で折り畳むことによって生じた。すなわち一方では $g$ と $f$ のより高い折り畳みによって、他方では $f$ または $g$ をそれ自身で折り畳むことによって生じた。どちらの場合にも、級数展開は、初期形式 $h$ について上で示したのと同じように、形式列 $f$ の極へ導く。特に $f$ が還元子であるなら、還元子の極による展開が得られる。したがって形式列 $h$ は可約となる。

完全な形式系の還元へのこの定理の応用は、二元および三元領域の場合と同様に従う。

\clearpage
% ---- Continuation appended 2026-06-02: Paper 05 complete and Paper 06 through Section 6, Japanese ----
\providecommand{\Kfield}{\mathfrak{K}}
\providecommand{\Ssys}{\mathfrak{S}}
\providecommand{\Mfield}{\mathfrak{M}}
\providecommand{\tq}{\tau}
\providecommand{\Lsys}{\mathfrak{L}}
\providecommand{\Jdom}{\mathfrak{J}}
\providecommand{\ds}{\displaystyle}

\begin{center}
{\Large\bfseries 5. 有理関数体}\par
\vspace{1.2em}
\emph{Jahresbericht der DMV} 22 (1913), pp. 316--319
\end{center}

\vspace{1.5em}

以下の諸問題は、もともと E. Fischer との会話に由来する。なお、それらの問題のいくつかは、一つの不定元という特殊な場合について、しかも係数体に関するより一般的な仮定のもとで、すでに E. Steinitz によって提起され解決されている。\footnote{E. Steinitz: \emph{Algebraische Theorie der Körper}, とくに \S\ 24. Crelle's Journal 137, 1910.}

「有理関数体」とは、与えられた数体からの係数をもつ $n$ 個の不定元の有理関数を元とする体をいう。この係数体は、とくにすべての複素数を含んでもよい。このような体の例としては、$n$ 個の量の対称関数の体、またより一般に、ある群の置換を許す $n$ 個の不定元の有理関数全体（Lagrange の \emph{Gattungsbereiche}）、さらに不変式体がある。

最初に挙げた二つの体と最後の体との間には、特徴的な相違がある。前者は $n$ 個の代数的に独立な関数、すなわち基本対称関数を含むのに対し、代数的に独立な不変式の個数は常に不定元の個数より小さい。したがって、一般に、第二種のこのような体を、いくつかの不定元を数で置き換えることによって第一種の体と一対一に対応させられることが重要である。そこで以下では、第一種の体、すなわち $n$ 個の代数的に独立な関数をもつ体に限って考える。この対応により、結果は一般にも成り立つ。

問題は三つの基底概念、すなわち有理基底、最小基底、整性基底のまわりに集まる。

1. 「有理基底」とは、その体の有限個の関数であって、体のすべての関数が、与えられた係数体からの係数を用いて、それら有限個の関数の有理的結合として表されるものをいう。簡単な考察によって、任意の有理関数体について有理基底の存在が示される。一つの不定元の場合には、この考察は E. Steinitz の前掲書にも見られる。例えば Lagrange の定理によれば、基本対称関数と群に属する一つの関数は、先に述べた Lagrange の \emph{Gattungsbereiche} の有理基底をなす。

2. 任意の有理基底は、第一種の体に対して、少なくとも $n$ 個の関数を含まなければならない。ちょうど $n$ 個の関数を含むならば、それらは代数的に独立でなければならず、そのような基底を「最小基底」と呼ぶ。最小基底の存在問題は、Lüroth, Castelnuovo, Enriques の定理によって部分的に答えられる。\footnote{J. Lüroth: Beweis eines Satzes über rationale Kurven. \emph{Math. Ann.} 9, 1875. --- G. Castelnuovo: Sulla razionalità delle involuzioni piane. \emph{Math. Ann.} 44, 1893. --- F. Enriques: Sopra una involuzione non razionale dello spazio. \emph{Rend. Acc. Linc.} vol. 21, 1912年1月21日.} それらによれば、一つの不定元の体には常に最小基底が存在する。すなわち、その体の Lüroth 関数であり、これは一次分数変換を除いて一意に定まる（Steinitz \S\ 24 参照）。二つの不定元の体については、係数体の代数拡大を許すならば、そして一般にはその場合に限って、常に最小基底が存在する。これに対して三つ以上の不定元については、一般には最小基底はまったく存在しない。最小基底をもつ特殊な体を特徴づける試みは、まだ行われていない。

私は Lagrange の \emph{Gattungsbereiche} における最小基底の問題をさらに追究した。この場合、それは次の意味をもつ。「群 $G$ に属する Lagrange の \emph{Gattungsbereich} が最小基底をもつならば、指定された群 $G$ をもつ方程式全体を、パラメータ表示によって有理的に構成できる。しかも、最小基底の係数体を含む任意の数体についてそうであり、特にこの係数体が有理数体であるならば、任意の数体についてそうである。」

最も簡単な例は対称群である。この場合、基本対称関数は有理数係数をもつ最小基底をなす。そこからパラメータ表示として、単に不定係数をもつ方程式が得られる。これから任意の数体上の任意に多くの affektlos な方程式へ移る方法は、Hilbert の既約性定理が示している。

さて、Lüroth と Castelnuovo の定理から、3 次および 4 次方程式に現れるすべての群について最小基底の存在が直接に従う。同時に、この最小基底は有理数係数をもつことも分かる。したがって、任意の数体について、指定された群をもつ 3 次および 4 次方程式全体を有理的に構成できる。巡回群および二面体群については、$n$ 乗根を係数体とする最小基底が存在する。さらにいくつかのメタ巡回群についても同様である。最小基底がそもそもある種の群を特徴づけるかどうかは、未決のままにしておかなければならない。

3. 最後の基底概念に到達するため、体の中に含まれる多項式全体を考える。「整性基底」とは、これらの多項式の有限個であって、体の任意の多項式が、与えられた係数体からの係数を用いて、それら有限個の多項式の整有理結合として表されるものをいう。整性基底の存在問題は、例えば不変式系の有限性を特殊な場合として含むが、Hilbert の「数学の問題」\footnote{D. Hilbert: Mathematische Probleme. パリ講演, 1900. 問題14. \emph{Göttinger Nachr.} 1900.}では、やや異なる形、すなわち相対的に整な関数の問題として提出されている。

当面、私は整性基底が存在する体の一つの類だけを示すことができる。すなわち、その体が $n$ 個の多項式の系を含み、それらの最高次元項の同次終結式が零でないならば、整性基底が存在する。それは、線形形式
\[
  u_0=u_1x_1+u_2x_2+\cdots+u_nx_n\footnotemark
\]
に対する、その体で既約な方程式に現れるすべての $u$ の係数によって与えられる。
\footnotetext{この既約方程式は、一つの不定元の場合には、Steinitz による Lüroth の定理の証明に見られる。}
とくに終結式の値が係数体の単元に等しいならば、表示の整性も保証される。この類の体の例は Lagrange の \emph{Gattungsbereiche} である。なぜなら基本対称関数の終結式は値 $\pm 1$ をもつからである。さらに別の例（整性なし）は、ただ一つの代数的に独立な多項式だけを含むすべての体によって与えられる。この場合、整性基底は、すべての多項式を含む最小の部分体の Lüroth 関数と同一である。したがって、よく知られているように、$n$ 変数の二次形式のすべての整有理な射影不変式は、判別式の整関数である。

いま述べた条件は整性基底の存在に十分であるにすぎず、決して必要ではない。任意の $n$ 個の多項式の終結式が消えるにもかかわらず、あの既約方程式の係数によって与えられる整性基底をもつ体は、容易に構成できる。最も一般の場合にも、あの方程式の係数が整性基底を与えるのではないか、という問題が生じる。

\clearpage

\begin{center}
{\Large\bfseries 6. 有理関数の体と系}\par
\vspace{1.2em}
\emph{Mathematische Annalen} 76 (1915), pp. 161--196
\end{center}

\vspace{1.5em}

本論文は、有理関数および整有理関数の任意の系に関する基底問題を扱う。ここで用いられる体論の方法により、有理関数からなる体、すなわち有理関数体\footnote{ウィーン自然科学者会議の機会における予備的報告「Rationale Funktionenkörper」, \emph{Jahresber. d. D. Math.-Ver.} 22 (1913) を参照。そこでは、問題設定と関数体に関する結果の概観が与えられている。}についてこれらの問題を解くことが本質的な部分として現れ、任意の系への結果の一般化はその帰結として得られる。

一般の系に対する基底問題のうち、これまで知られていたのは Hilbert の定理（\emph{Math. Ann.} 36）によって保証される加群基底の存在だけであった。以下では、任意の系に対する有理基底の存在（\S\ 7）によって、有理的表示可能性の問題が完全に解決される。体の有理基底はすでに \S\ 4 で得られる。この有理基底の存在により、常に抽象的に定義された体または系から出発することができ、したがって、特殊な有理基底の選択だけに由来し、系そのものには由来しない困難、例えば特殊な分母や基底関数の基本点の出現を避けることができる。

体については、最小基底、すなわち代数的に独立な関数からなる有理基底の問題も扱われる（\S\ 6）。さらに体論の方法は、特別な有理基底、すなわちインボリューション基底\footnote{この名称は、幾何学的解釈において、有理関数体が線形空間内の最も一般的なインボリューションを表すことによる。}（\S\ 5）へ導く。これは多項式からなる整性領域にとって特に重要となる（\S\ 8）。その場合、固定された分母（整性領域によって定まる関数の冪）をもつ表示を与える。これは、例えば不変式の典型表示という特殊な場合に知られているものである。

有理的表示可能性から、可能な限り、整有理的表示可能性、すなわち狭い意味での有限性の問題への帰結を引き出す。

これまで知られていた有限性定理はすべて、Hilbert の加群基底定理が、表示
\[
  F=A_1f_1+A_2f_2+\cdots+A_kf_k
\]
から第二の表示
\[
  F=B_1f_1+B_2f_2+\cdots+B_kf_k
\]
が従い、ここで $B_i$ がその系に属し、かつ $F$ より低い次数をもつようなすべての系に対して有限性を保証する、という事実に基づいている。

これに対して、有理基底に関する定理と体論の方法は、まったく異なる種類の仮定のもとで有限性を推論することを可能にする。

こうして Hilbert の問題（数学の問題 14）への部分的な答えとして、相対的に整な関数の一類（第一種の相対的に整な領域）が得られる。これらは有限な整性領域であり、その整性基底は上に述べたインボリューション基底によって与えられる（\S\ 10）。さらに、Hilbert による Kronecker の代数的整数の基本系に関する定理の証明（\emph{Complete invariant systems}, \S\ 2; \emph{Math. Ann.} 42）と類似に、すべての正則多項式系、すなわち基本点をもたないものへ写すことのできるすべての系は整性基底をもつことが示される（\S\ 12）。一方、整性基底をもたない非正則系は容易に指定できる。この事実の簡単な例として、次の定理を挙げられる。「一つの不定元に関する任意の無限個の多項式の系において、その系の任意の多項式が有限個の多項式の整有理結合として表されるような有限個の多項式が存在する。」さらなる例は \S\ 13 にある。

最後に、末尾の二つの節（\S\S\ 14, 15）は、整性に関して基底定理を精密化する。

この研究の本質的な道具は \S\ 3 の転移原理である。それによれば、ここで提起されるすべての基底問題について、不定元の個数が系の代数的に独立な関数の個数と一致するような系だけを考えれば十分である。

本研究の発端は、E. Fischer 氏との会話、とくに Lagrange の \emph{Gattungsbereiche} の最小基底に関する同氏の問いであった。これに関わる研究は、指定された群をもつ方程式の構成へ導いたが（前掲の「有理関数体」第2部を参照）、それは後の公表に委ねる。

本論文が最初の報告から本質的に拡張されている点は、そこでは体または特殊な整性領域についてだけ述べられた基底定理が、ここでは任意の系へ移されていることである。この転移は、任意の系を含む最小の体という概念によって簡単に行われる。これは整性領域の商体の一般化である（\S\ 7）。私がこの概念を形成するに至ったのは、K. Hentzelt 氏が、体の有理基底を知った後、Hilbert の加群基底を経由する迂回によって任意の系の有理基底の存在を証明できた、という事実による。また、正則系の整性基底に関する定理へ私を導いたのは、整性領域に対して私が用いた推論が、同じ仮定を満たす任意の系にも有効であるという K. Hentzelt の指摘であった。さらに、いくつかの個別の小さな指摘も Hentzelt 氏に負っている。

最後に、一つの不定元の場合、体の有理基底と最小基底は E. Steinitz の \emph{Algebraische Theorie der Körper}, \S\ 24（Crelle's Journal 137）に見られることを述べておく。そこでは係数体は最も一般的な意味での体として仮定されている。これらのより一般的な仮定のもとで、ここに与える定理がどこまで保たれるかという問題は、以下では扱わない。いずれにせよ証明は修正を必要とする。例えば、関数行列式の理論からの手段は標数 $p$ の体では機能しないからである。

\section*{\S\ 1. 体 $\Kfield_{n\rho}$ と系 $\Ssys_{n\rho}$.}

以下では、$n$ 個の不定元の有理関数の系、とくに有理関数体を調べる。

$f(x_1,\ldots,x_n)$, $g(x_1,\ldots,x_n)$, \ldots、または略して $f(x)$, $g(x)$, \ldots は、常に $x_1,\ldots,x_n$ の有理関数を意味するものとする。これらは既約表示、すなわち分子と分母が互いに素である表示にあると仮定する。整有理関数（多項式）は大文字 $F(x)$, $G(x)$, \ldots で表す。係数領域 $\Omega$ は何らかの数体と仮定され、したがってとくにすべての複素数を含んでもよい。$\Omega$ はさらに有限個のパラメータを含んでもよい。

$\Omega$ からの量、したがって零次のすべての関数は、常に系に属するものと数える。

これらの取り決めのもとで、有理関数からなる体、すなわち有理関数体は抽象的に定義できる。

\emph{定義 I: 有理関数の系は、次の条件を満たすとき体と呼ばれる。}
\begin{enumerate}
\item \emph{$f(x)$ とともに、$\Omega$ からの任意の量 $c$ に対して $c\cdot f(x)$ も含む。}
\item \emph{$f(x)$ と $g(x)$ とともに、常に $f(x)+g(x)$, $f(x)\cdot g(x)$ および、$g(x)\ne0$ のとき商 $f(x):g(x)$ も含む。}\footnote{ここで $f(x)$ と $g(x)$ は零次、すなわち $\Omega$ の量であってもよい。}
\end{enumerate}
特殊な体とは、有限個の有理関数
\[
  f_1(x),\ f_2(x),\ldots, f_k(x)
\]
を $\Omega$ に添加して得られるものであり、
\[
  \Omega\bigl(f_1(x),\ldots,f_k(x)\bigr)
  \quad\text{または}\quad
  \Omega(f_1,\ldots,f_k)
\]
と表す。\footnote{これらの「特殊な体」がすでにすべての体を尽くすことは \S\ 4 で分かる。}
これら特殊な体の中には、$x_1,\ldots,x_n$ を $\Omega$ に添加して得られる体
\[
  \Omega(x_1,\ldots,x_n)
\]
も含めておく。これは $\Omega$ 係数の $x_1,\ldots,x_n$ の有理関数全体と同一である。

さらに、E. Steinitz に従って、中間体の概念を導入する。
\begin{quote}
「体 $\Mfield$ が $\Omega_1$ と $\Omega_2$ の間にあるとは、$\Mfield$ が $\Omega_1$ を含み、かつ $\Omega_2$ に含まれることをいう。」
\end{quote}
すると定義 I は次のようにも述べられる。

$n$ 個の不定元の有理関数体とは、$\Omega$ と $\Omega(x_1,\ldots,x_n)$ の間の中間体であり、少なくとも一つの関数において $n$ 個の不定元を実際に含むものである。

不定元の個数のほかに、有理関数の系にはさらに一つの特徴数、すなわち代数的に独立な関数の個数、あるいは代数的階数が、次の定義によって現れる。

\emph{定義 II: 有理関数の系が代数的階数 $\rho$ をもつとは、その系の中に代数的に独立な $\rho$ 個の関数を指定できるが、その系の任意の $(\rho+1)$ 個の関数は代数的に従属であることをいう。}\footnote{$\rho$ 個の関数 $f_1(x),\ldots,f_\rho(x)$ が代数的に独立であるとは、任意の関係
\[
F\bigl(f_1(x),\ldots,f_\rho(x)\bigr)=0\quad\text{$x_1,\ldots,x_n$ について恒等的に}
\]
から
\[
F(\lambda_1,\ldots,\lambda_\rho)=0\quad\text{$\lambda_1,\ldots,\lambda_\rho$ について恒等的に}
\]
が従うことをいう。反対の場合、それらは代数的に従属であると呼ばれる。}

$n$ 個の不定元の有理関数の（有限または無限の）系で、代数的階数 $\rho$ をもつものを、二重添字によって
\[
  \Ssys_{n\rho}
\]
と表す。特に
\[
  \Kfield_{n\rho}
\]
は、$n$ 個の不定元と代数的階数 $\rho$ をもつ体を意味する。次の不等式が成り立つ。
\[
  1\leq \rho\leq n.
\]

体 $\Kfield_{nn}$ および $\Kfield_{n\rho}$ の例をさらにいくつか挙げる。

1. 体 $\Kfield_{nn}$ として、Lagrange の \emph{Gattungsbereiche} がある。これは
\begin{quote}
「$x_1,\ldots,x_n$ の置換群 $G$ の置換を許す、$x_1,\ldots,x_n$ の有理関数（および整有理関数）全体」
\end{quote}
として定義できる。これらは常に対称関数の体を含み、したがって $n$ 個の代数的に独立な基本対称関数も含む。

2. 体 $\Kfield_{n\rho}$ として、射影不変式体がある。これは一つまたは複数の基本形式の有理不変式（および整有理不変式）全体から作られる。\footnote{行列式 $1$ の変換だけを考えるのが便利である。そのとき任意個の不変式の任意の有理結合も再び変換を許し、実際に体が得られる。斉次・等重の不変式だけを取り出したものは体をなさないが、系 $\Ssys_{n\rho}$ をなす。} ここでは $\rho<n$ である。不変式の代数的に独立な個数は、常に係数の個数より小さいからである。

射影群の部分群に関する不変式体についても、対応する主張が成り立つ。

\section*{\S\ 2. 代数的に従属な関数に関する補助定理.}

この節の補助定理により、\S\ 3 で、系の代数的階数がその本来的な特徴数であることが示される。すなわち補助定理は、$\rho$ 個の代数的に独立な関数が代数的に独立なままでいるように、いくつかの不定元を数値によって逐次特殊化しても、それら $\rho$ 個の関数に代数的に従属な任意の関数は、恒等的に消えることも不定になることもない、ということを示す。

したがって、系
\begin{equation}
  g_1(x),\ g_2(x),\ldots,g_\rho(x)
\tag{1}
\end{equation}
が代数的階数 $\rho$ をもち、ここで $\rho<n$ とし、$h(x)$ を関数 (1) に代数的に従属な有理関数とする。既知の定理により、(1) の関数行列の階数は $\rho$ に等しい。したがって不定元は、例えば次のように名づけられる。
\begin{equation}
 \frac{\partial(g_1,g_2,\ldots,g_\rho)}{\partial(x_1,x_2,\ldots,x_\rho)}
 =\frac{\Gamma(x)}{G(x)^{\rho+1}}\ne0,
\tag{2}
\end{equation}
ここで $G(x)$ は通常どおり、関数 (1) の共通分母を意味する。

\emph{いま数 $a$ を、積 $G(x)\cdot\Gamma(x)$ が $(x_i-a)$ で割り切れないように選ぶ}。ただし $i$ は $\rho+1,\rho+2,\ldots,n$ のいずれかとする。したがって $x_i=a$ としても、関数 (1) の共通分母は消えず、$\rho$ 個の関数
\begin{equation}
  [g_1(x)]_{x_i=a},\ldots,[g_\rho(x)]_{x_i=a}
\tag{3}
\end{equation}
もまた代数的に独立である。言い換えれば、系 (1) の代数的階数は、置換 $x_i=a$ のもとで保存される。

関数 (1) に関して $h(x)$ が満たす既約方程式を
\begin{equation}
\begin{aligned}
&A_0(g_1,\ldots,g_\rho)\,h(x)^\tau
 +A_1(g_1,\ldots,g_\rho)\,h(x)^{\tau-1}
 +\cdots \\
&\hspace{6em}+A_\tau(g_1,\ldots,g_\rho)=0
\end{aligned}
\tag{4}
\end{equation}
とする。ただし $A_0(g_1,\ldots,g_\rho)$ および $A_\tau(g_1,\ldots,g_\rho)$ は恒等的には消えない。これを多項式の間の恒等式に移すために、
\begin{equation}
  g_i(x)=G_i(x):G(x);\qquad h(x)=H_1(x):H_2(x)
\tag{5}
\end{equation}
と置く。そして、$h(x)$ が既約表示にあるという仮定により、$H_1(x)$ と $H_2(x)$ は互いに素であることに注意する。(5) によって (4) は
\begin{equation}
\begin{aligned}
&B_0(G,G_1,\ldots,G_\rho)\,G(x)^{\sigma_0}H_1(x)^\tau \\
&\quad+B_1(G,G_1,\ldots,G_\rho)\,G(x)^{\sigma_1}H_1(x)^{\tau-1}H_2(x)+\cdots \\
&\quad+B_\tau(G,G_1,\ldots,G_\rho)\,G(x)^{\sigma_\tau}H_2(x)^\tau=0
\end{aligned}
\tag{6}
\end{equation}
となる。ここで $B_k$ は $G,G_i$ の同次形式であり、非負指数 $\sigma_k$ の少なくとも一つは零でなければならない。

いま $H_1(x)$ が $(x_i-a)$ で割り切れると仮定する。すると (6) により
\[
  B_\tau(G,G_1,\ldots,G_\rho)\,G(x)^{\sigma_\tau}H_2(x)^\tau
\]
も $(x_i-a)$ で割り切れる。これは $G(x)$ についての仮定により排除されている。したがって $B_\tau$ の割り切れ性から
\[
  A_\tau=B_\tau:G^\lambda
\]
の割り切れ性も従い、関数 (3) の間に関係 $A_\tau=0$ が存在することになり、仮定された代数的独立性に反する。最後に、$H_2(x)$ は $H_1(x)$ と互いに素であるから、これも $(x_i-a)$ で割り切れることはできない。\footnote{既約表示に基づくこの最後の結論は、残りの仮定が保たれるとしても、置換 $x_i=a$, $x_k=b$ に対してはもはや可能ではない。その場合 $H_1(x)$ と $H_2(x)$ は同時に消えることがあり、商はその置換のもとで $\frac{0}{0}$ に等しい不定形となりうる。例えば
\[
 h(x)=\frac{x_3(\alpha x_1+\beta x_2)+x_4(\gamma x_1+\delta x_2)}
 {x_3(\alpha' x_1+\beta' x_2)+x_4(\gamma' x_1+\delta' x_2)}
\]
を $x_3=0$, $x_4=0$ と置換する場合である。} よって $H_1(x)$ が $(x_i-a)$ で割り切れるという仮定は矛盾に至る。同様に、$H_2(x)$ も $(x_i-a)$ で割り切れないことが示される。したがって関数 $[h(x)]_{x_i=a}=h'(x)$ は、恒等的に消えることも不定になることもない。

以上の結論を、再び既約表示に置いた系 (3) と関数 $h'(x)$ に適用できる。この手続きを続けることで、最終的に次の補助定理に到達する。

\emph{補助定理: 二つの系}
\[
  g_1(x),\ g_2(x),\ldots,g_\rho(x),
\]
\[
  g_1(x),\ g_2(x),\ldots,g_\rho(x),\ h(x)
\]
\emph{はいずれも代数的階数 $\rho$ をもつものとし、ここで $\rho<n$ である。例えば}
\[
 \frac{\partial(g_1,\ldots,g_\rho)}{\partial(x_1,\ldots,x_\rho)}
 =\frac{\Gamma(x)}{G(x)^{\rho+1}}\ne0
\]
\emph{と仮定する。数}
\[
  a_n,\ a_{n-1},\ldots,a_{\rho+1}
\]
\emph{は、置換}
\[
  x_n=a_n,\quad x_{n-1}=a_{n-1},\ldots,\quad x_{\rho+1}=a_{\rho+1}
\]
\emph{のもとで積 $G(x)\cdot\Gamma(x)$ が消えないように、すなわち系 (1) の代数的階数が保存されるように選ばれるものとする。逐次置換}
\[
\begin{aligned}
 [h(x)]_{x_n=a_n}&=h^{(n-1)}(x),\qquad
 [h^{(n-1)}(x)]_{x_{n-1}=a_{n-1}}=h^{(n-2)}(x),\ldots,\\
 [h^{(\rho+1)}(x)]_{x_{\rho+1}=a_{\rho+1}}&=h^*(x_1,x_2,\ldots,x_\rho)
\end{aligned}
\]
\emph{において、関数 $h(x)$, $h^{(n-1)}(x)$, \ldots, $h^{(\rho+1)}(x)$ をそれぞれ既約表示に置く。これらの仮定のもとで、関数 $h^*(x_1,\ldots,x_\rho)$ では分子も分母も恒等的には消えず、また関数は $\frac{0}{0}$ になることもない。}\footnote{例えば前注の関数 $h(x)$ では、逐次置換により
\[
 [h(x)]_{x_4=0}=h^{(3)}(x)=\frac{\alpha x_1+\beta x_2}{\alpha' x_1+\beta' x_2},\qquad
 h^*(x_1,x_2)=\frac{\alpha x_1+\beta x_2}{\alpha' x_1+\beta' x_2}
\]
が得られる。}

\section*{\S\ 3. 系 $\Ssys_{n\rho}$ から系 $\Ssys^*_{\rho\rho}$ への一対一対応.}

補助定理から、いくつかの不定元を数で置き換えることにより、系 $\Ssys_{n\rho}$ を系 $\Ssys^*_{\rho\rho}$ へ対応させることが直ちに得られる。このとき代数的階数 $\rho$ は本来的な特徴数として現れる。

$\Ssys_{n\rho}$ は代数的階数 $\rho$ をもつので、$\Ssys_{n\rho}$ から $\rho$ 個の関数
\begin{equation}
  g_1(x),\ldots,g_\rho(x)
\tag{1}
\end{equation}
が存在し、それらは代数的に独立である。さらに、$f(x)$ が $\Ssys_{n\rho}$ からの任意の関数を表すならば、系
\begin{equation}
  g_1(x),\ldots,g_\rho(x),\ f(x)
\tag{2}
\end{equation}
もまた代数的階数 $\rho$ をもち、系 (1) と (2) は補助定理の条件を満たす。\footnote{必要ならば、不定元を番号づけ直す。}

したがって、数
\[
  a_n,\ a_{n-1},\ldots,a_{\rho+1}
\]
を補助定理に従って、系 (1) の代数的階数 $\rho$ が置換によって保存されるように定める。このことは任意に多くの仕方で可能である。そのとき
\[
\begin{aligned}
 [f(x)]_{x_n=a_n}&=f^{(n-1)}(x),\qquad
 [f^{(n-1)}(x)]_{x_{n-1}=a_{n-1}}=f^{(n-2)}(x);\ldots,\\
 [f^{(\rho+1)}(x)]_{x_{\rho+1}=a_{\rho+1}}&=f^*(x_1,\ldots,x_\rho)
\end{aligned}
\]
で定義される関数、すなわち
\begin{equation}
  f^*(x_1,\ldots,x_\rho)
\tag{3}
\end{equation}
では、分子も分母も恒等的に消えることができない。

いま $f(x)$ に $\Ssys_{n\rho}$ のすべての（有限または無限個の）関数を走らせる。すべての関数 (3) 全体からなる系を考えると、それは次の性質をもつ。

1. 代数的階数 $\rho$ をもつ。実際、それは (1) から生じる関数
\[
  g_1^*(x_1,\ldots,x_\rho),\ldots,g_\rho^*(x_1,\ldots,x_\rho)
\]
を含み、これらは $a_n,a_{n-1},\ldots,a_{\rho+1}$ の選び方により代数的に独立である。したがってこの系を
\[
  \Ssys^*_{\rho\rho}
\]
と表す。\footnote{対応して、$f^*(x)$, $g^*(x)$, \ldots は、以下いつでも対応関数、すなわち $x_1,\ldots,x_\rho$ の関数を意味するものとする。}

2. 関数 $f(x)$ と $f^*(x)$ の対応は例外なく一対一である。

実際、$f_1(x)$ と $f_2(x)$ が $\Ssys_{n\rho}$ に属すならば、その差
\[
  d(x)=f_1(x)-f_2(x)
\]
も系 (1) に代数的に従属であり、さらに
\[
  d^*(x)=f_1^*(x)-f_2^*(x)
\]
である。$d(x)$ は系 (1) とともに補助定理の条件を満たすので、
\[
  d(x)\ne0
\]
から必ず
\[
  d^*(x)\ne0
\]
が従う。すなわち、$\Ssys_{n\rho}$ の異なる関数には、$\Ssys^*_{\rho\rho}$ の異なる関数が対応する。

3. $\Ssys^*_{\rho\rho}$ の関数の間の任意の関係
\[
  F\bigl(f_1^*(x),f_2^*(x),\ldots,f_k^*(x)\bigr)=0
  \quad\text{$x_1,\ldots,x_\rho$ について恒等的に}
\]
から、一意に対応する $\Ssys_{n\rho}$ の関数について
\[
  F\bigl(f_1(x),f_2(x),\ldots,f_k(x)\bigr)=0
  \quad\text{$x_1,\ldots,x_n$ について恒等的に}
\]
が従う。なぜなら、$f_1(x),f_2(x),\ldots,f_k(x)$ とともに、これらの任意の有理結合も系 (1) に代数的に従属だからである。

さらにこれは
\begin{equation}
\begin{aligned}
 h(x)&=F\bigl(f_1(x),\ldots,f_k(x)\bigr):\\
 h^*(x)&=F\bigl(f_1^*(x),\ldots,f_k^*(x)\bigr)
\end{aligned}
\tag{4}
\end{equation}
からも従う。補助定理により、したがって
\[
  h(x)\ne0\quad\text{ならば}\quad h^*(x)\ne0
\]
である。証明終わり。

要約すれば次を得る。

\emph{定理 I: 有理関数の（有限または無限の）任意の系 $\Ssys_{n\rho}$ に対して、任意に多くの仕方で、一対一の系 $\Ssys^*_{\rho\rho}$ を割り当てることができる。その際、$\Ssys^*_{\rho\rho}$ の関数間に存在するすべての関係は $\Ssys_{n\rho}$ においても保存され、また逆も成り立つ。特に (4) により、体 $\Kfield_{n\rho}$ には再び体 $\Kfield^*_{\rho\rho}$ が対応する。}

定理 I から、以後のすべての証明を簡単にする帰結を引き出す。\emph{系の関数間の有限または無限個の関係によって特徴づけられる系 $\Ssys_{n\rho}$ の性質は、対応する系 $\Ssys^*_{\rho\rho}$ について成り立てば、もとの系 $\Ssys_{n\rho}$ についても成り立つ。}

この帰結を略して「転移原理」と呼ぶ。

\providecommand{\Lsys}{\mathfrak{L}}
\providecommand{\Jdom}{\mathfrak{J}}

\clearpage
この対応の最も簡単な例は、代数的不変式の理論によって与えられる。実際、線形変換によって、基本形式のいくつかの係数を数値的に固定し、特殊な形式に対して成り立つすべての関係が一般の場合にも成り立つようにすることが常にできる。

\section*{\S\ 4. 体 $\Kfield_{n\rho}$ の有理基底.}

以下の研究では、基底概念のまわりにまとめられた問題に対して、\S\ 3 の「転移原理」を一貫して用いる。まず体について基底の存在を証明するが、定義は一般に述べる。

\emph{定義 III: $\Ssys_{n\rho}$ の有限個の関数が有理基底と呼ばれるのは、$\Ssys_{n\rho}$ の任意の関数が、$\Omega$ からの係数を用いて、それら有限個の関数の有理結合として表されるときである。}

有理基底は $\rho$ 個の代数的に独立な関数を含まなければならない。なぜなら、有理基底から導かれる系の代数的階数は、有理基底の代数的階数に等しいからである。

まず、有理基底の存在証明が転移原理を許すことを示す。すなわち、定義 III は次のように定式化できる。

$\Ssys_{n\rho}$ の関数
\[
  g_1(x),\ g_2(x),\ldots, g_k(x)
\]
が有理基底をなすことと、無限個の関係
\begin{equation}
  f(x)=\gamma\bigl(g_1(x),\ldots,g_k(x)\bigr)
\tag{1}
\end{equation}
が満たされることは同値である。ここで $f(x)$ は $\Ssys_{n\rho}$ のすべての関数を走り、$\gamma$ は $f$ に応じて変わる $k$ 個の引数の有理関数を表す。

したがって、対応系 $\Ssys^*_{\rho\rho}$ において
\[
  g_1^*(x),\ g_2^*(x),\ldots,g_k^*(x)
\]
が有理基底を与え、それにより関係
\[
  f^*(x)=\gamma\bigl(g_1^*(x),\ldots,g_k^*(x)\bigr)
\]
が生じるならば、定理 I によって、関係 (1) は $\Ssys_{n\rho}$ に対して成り立つ。したがって、\emph{有理基底に対する転移原理}は次のようにいう。

\emph{対応系 $\Ssys^*_{\rho\rho}$ に有理基底が存在することから、もとの系 $\Ssys_{n\rho}$ の有理基底が従う。}\footnote{転移原理は、ここで扱うすべての基底問題について同様に定式化される。}

すべての体 $\Kfield_{n\rho}$ の有理基底の存在を証明するためには、したがって体 $\Kfield^*_{\rho\rho}$ に限ればよい。ここでは E. Steinitz が用いた議論の直接の一般化が目的に至る。\footnote{E. Steinitz, \emph{Algebraische Theorie der Körper}, \S\ 24, Crelle's Journal 137 (1909).}

\begin{equation}
  g_1^*(x_1,\ldots,x_\rho),\ldots,
  g_\rho^*(x_1,\ldots,x_\rho)
\tag{2}
\end{equation}
を、$\Kfield^*_{\rho\rho}$ からの $\rho$ 個の代数的に独立な関数の系とする。$x_1,\ldots,x_\rho$ を消去すると、$\rho$ 個の関係
\begin{equation}
\begin{aligned}
  F_1\bigl(g_1^*(x),\ldots,g_\rho^*(x),x_1\bigr)&=0,\ldots,\\
  F_\rho\bigl(g_1^*(x),\ldots,g_\rho^*(x),x_\rho\bigr)&=0
\end{aligned}
\tag{3}
\end{equation}
が得られる。各関係 $F_i=0$ には不定元 $x_i$ が実際に現れる。そうでなければ、仮定に反して関数 (2) の間に代数的従属性が存在することになる。関係 (3) は
\[
  x_1,x_2,\ldots,x_\rho
\]
が
\[
  \Omega\bigl(g_1^*(x),\ldots,g_\rho^*(x)\bigr)
\]
に関して代数的元であることをいう。したがって、$x_1,\ldots,x_\rho$ を添加して得られる体
\[
  \Omega\bigl(g_1^*,\ldots,g_\rho^*,x_1,
  \ldots,x_\rho\bigr)=\Omega(x_1,
  \ldots,x_\rho)
\]
は、$\Omega(g_1^*,\ldots,g_\rho^*)$ の有限代数拡大である。既知の定理\footnote{例えば Weber, \emph{Lehrbuch d. Algebra} 1, \S\ 144（第1版）または \S\ 55（小版）を比較せよ。}により、$\Kfield^*_{\rho\rho}$ は、$\Omega(g_1^*,\ldots,g_\rho^*)$ と $\Omega(x_1,\ldots,x_\rho)$ の間の中間体として、それ自身 $\Omega(g_1^*,\ldots,g_\rho^*)$ の代数拡大であり、また $\Omega(x_1,\ldots,x_\rho)$ は $\Kfield^*_{\rho\rho}$ の代数拡大である。したがって $\Kfield^*_{\rho\rho}$ は、$\Omega(g_1^*,\ldots,g_\rho^*)$ に $\Kfield^*_{\rho\rho}$ からの有限個の元を添加することによって、または適当に選ばれた一つの関数を添加することによって得られる。すなわち有理基底が存在する。転移原理によって次を得る。

\emph{定理 II: 任意の体 $\Kfield_{n\rho}$ は有理基底をもつ。}

\section*{\S\ 5. 体 $\Kfield_{n\rho}$ のインボリューション形式とインボリューション基底.}

前の議論から、体そのものによって定まる有理基底をさらに得る。まず体 $\Kfield^*_{\rho\rho}$ についてそれを示す。

$\Kfield^*_{\rho\rho}$ を $\Omega$ と $\Omega(x_1,\ldots,x_\rho)$ の間の体とする。\S\ 4 と同様、$\Kfield^*_{\rho\rho}$ の中から $\rho$ 個の代数的に独立な関数
\[
  g_1^*(x),\ldots,g_\rho^*(x)
\]
を選び、$\Omega(g_1^*,\ldots,g_\rho^*)$ の有限代数拡大 $\Omega(x_1,\ldots,x_\rho)$ を作る。$\Omega(x_1,\ldots,x_\rho)$ の $\Kfield^*_{\rho\rho}$ に関する次数を $i$ とする。

不定係数 $u_1,\ldots,u_\rho$ をもつ線形形式
\[
  u(x)=u_1x_1+\cdots+u_\rho x_\rho
\]
を導入し、$u(x)$ の $\Kfield^*_{\rho\rho}$ に関する既約方程式を
\begin{equation}
 \Phi^*(z,u)=z^i+\sum_{\alpha+\alpha_1+\cdots+\alpha_\rho=i}
 g^*_{\alpha\alpha_1\ldots\alpha_\rho}(x)
 z^\alpha u_1^{\alpha_1}\cdots u_\rho^{\alpha_\rho}
\tag{1}
\end{equation}
とする。\footnote{Noether は総和を総和 $i$ をもつすべての添字系について書いている。ここでは主項 $z^i$ を分けている。} 係数
\[
  g^*_{\alpha\alpha_1\ldots\alpha_\rho}(x)
\]
は $\Kfield^*_{\rho\rho}$ に属し、体によって定まる。$\Phi^*(z,u)$ を体のインボリューション形式と呼ぶ。

これらの係数が有理基底をなすことを主張する。実際、(1) を $z$ の方程式として見れば、それは $\Kfield^*_{\rho\rho}$ に関して既約であり、したがって部分体
\[
  \Omega\bigl(g^*_{\alpha\alpha_1\ldots\alpha_\rho}(x)\bigr)
\]
に関してはなおさら既約である。ゆえに $\Phi^*(z,u)$ は、この部分体に関する $u(x)$ の既約方程式を与える。したがって $\Omega(x_1,\ldots,x_\rho)$ は、$\Omega(g^*_{\alpha\alpha_1\ldots\alpha_\rho}(x))$ の代数拡大、しかも次数 $i$ の拡大である。$\Kfield^*_{\rho\rho}$ はこの部分体と $\Omega(x_1,\ldots,x_\rho)$ の間にあるから、次数の一致は体の一致を意味することが知られている。転移原理により、$\Phi(z,u)$ の係数は対応して $\Kfield_{n\rho}$ の有理基底を与える。すなわち次を得る。

\emph{定理 III: インボリューション形式の係数は、体によって定まる有理基底、すなわちインボリューション基底をなす。体 $\Kfield^*_{\rho\rho}$ については、この基底は一意に定まる。}\footnote{体 $\Kfield_{n\rho}$ については、インボリューション形式は対応の選択に依存しうる。不定係数をもつ対応を用いれば一意性を得ることができる。}

この事実に非有理的な解釈を付け加える。これはインボリューションの幾何学的理論とのつながりを与える。

$\Phi^*(z,u)$ の一次因子への分解は、拡大体において
\begin{equation}
\begin{aligned}
 \Phi^*(z,u)&=(z-u_1x_1-\cdots-u_\rho x_\rho)
 (z-u_1x_1^{(1)}-\cdots-u_\rho x_\rho^{(1)})\cdots\\
 &\qquad\cdot
 (z-u_1x_1^{(i-1)}-\cdots-u_\rho x_\rho^{(i-1)})
\end{aligned}
\tag{4}
\end{equation}
によって与えられる。(1) と比較すれば、関数
\[
  g^*_{\alpha\alpha_1\ldots\alpha_\rho}(x_1,
  \ldots,x_\rho)\quad(\alpha+\alpha_1+\cdots+\alpha_\rho=i)
\]
は、量の行
\begin{equation}
\begin{array}{cccc}
  x_1 & x_2 & \cdots & x_\rho\\
  x_1^{(1)} & x_2^{(1)} & \cdots & x_\rho^{(1)}\\
  \vdots & \vdots & & \vdots\\
  x_1^{(i-1)} & x_2^{(i-1)} & \cdots & x_\rho^{(i-1)}
\end{array}
\tag{5}
\end{equation}
の基本対称関数と同一であることが分かる。したがって $\Kfield^*_{\rho\rho}$ の任意の関数は、行 (5) に関して対称なこれらの基本関数の有理結合であり、逆に行 (5) の任意の対称関数は $\Kfield^*_{\rho\rho}$ に属する。

したがって
\[
  f^*(x)=\frac{F^*(x)}{H^*(x)}
\]
が $\Kfield^*_{\rho\rho}$ のすべての関数を走るとき、無限個の関係の系
\begin{equation}
  \frac{F^*(x)}{H^*(x)}=
  \frac{F^*(x^{(1)})}{H^*(x^{(1)})}=\cdots=
  \frac{F^*(x^{(i-1)})}{H^*(x^{(i-1)})}
\tag{6}
\end{equation}
が存在する。要約すれば、
\[
  F^*(x^{(k)})H^*(x^{(l)})-
  F^*(x^{(l)})H^*(x^{(k)})=0
  \quad(k,l=0,1,\ldots,i-1),
\]
である。ここで $F^*(x^{(k)})$ も $H^*(x^{(k)})$ も恒等的には消えない。これらは $F^*(x)$ および $H^*(x)$ の共役であり、形式的に対称だからである。

逆に、無限個の関係
\begin{equation}
  F^*(x)H^*(\xi)-H^*(x)F^*(\xi)=0,
  \qquad H^*(\xi)\ne0
\tag{7}
\end{equation}
を満たす任意の量の行 $\xi$ について、量の行 (5) への所属が得られる。

実際、$G^*(x)$ を $\Phi^*(z,u)$ の係数の共通分母とすると、仮定 (7) により関数
\[
  G^*(\xi)\cdot\Phi^*(z,u;\xi)
\]
は、$\xi$ に関しても整であり、恒等的には消えない。すなわち $\Phi^*(z,u;\xi)$ の係数は不定にならず、かつ零点
\[
  z=u_1\xi_1+u_2\xi_2+\cdots+u_\rho\xi_\rho
\]
をもつ。(7) により、したがって $\Phi^*(z,u;x)$ も零点 $z=u(\xi)$ をもつ。すなわち $\xi$ は量の行 (5) に属する。$\xi$ が共役な量の行に関する方程式系
\[
  F^*(x^{(k)})H^*(\xi)-H^*(x^{(k)})F^*(\xi)=0;
  \qquad H^*(\xi)\ne0
\]
を満たす場合にも、(6) により対応する主張が成り立つ。

したがって、量の行 (5) は、方程式系
\[
  F^*(x)H^*(\xi)-H^*(x)F^*(\xi)=0;
  \qquad H^*(\xi)\ne0
\]
の解 $\xi$ として定義される。ここで $F^*(x):H^*(x)$ は $\Kfield^*_{\rho\rho}$ のすべての関数を走る。また行 $x$ は任意の共役な行で置き換えてもよい。

幾何学的には、各行 $x$ を点と解釈すると、これは次のことを意味する。(7) の解は、$\rho$ 次元線形空間において、各々 $i$ 個の点からなる $\infty^\rho$ 個の群を表し、群の任意の点がその個々の群を一意に定める。このような点群の系を「\emph{線形空間におけるインボリューション}」と呼ぶ。\footnote{(7) の除外された解、すなわち $F^*(\xi)=0$, $H^*(\xi)=0$ となるもの、およびインボリューション基底に対応する方程式系の除外された解は、インボリューションの基本点と呼ばれる。これには同次化によって生じるもの、すなわち「無限遠にある」点も含まれる。}

対応して、体 $\Kfield_{n\rho}$ は、曲線、曲面などからなる線形空間内のインボリューションを特徴づける。

以上により、$\Omega(x_1,\ldots,x_\rho)$ の $\Kfield^*_{\rho\rho}$ に関する次数 $i$ は、側条件 $H^*(\xi)\ne0$ を満たす (7) の解系の個数に等しい。これを $\Kfield^*_{\rho\rho}$ の解系と呼ぶことができる。したがって、インボリューション基底の存在を証明するために用いた議論は、次の非有理的な定式化も許す。

「$\mathfrak{L}^*$ が $\Kfield^*_{\rho\rho}$ の約体であり、体 $\mathfrak{L}^*$ のすべての解系が $\Kfield^*_{\rho\rho}$ の解系でもあるならば、$\mathfrak{L}^*$ と $\Kfield^*_{\rho\rho}$ は同一である。」

対応する命題は、転移原理により、$\Kfield_{n\rho}$ の約体 $\mathfrak{L}$ についても成り立つ。

\section*{\S\ 6. 体 $\Kfield_{n\rho}$ の最小基底.}

この節では既知の定理を概観する。ただし、転移原理と有理基底の存在によって可能となった拡張された定式化で述べる。

\emph{定義 IV: $\Kfield_{n\rho}$ の有理基底が最小基底と呼ばれるのは、それが最小個数 $\rho$ の関数だけを含むときであり、その場合それらは代数的に独立でなければならない。}

Lüroth, Castelnuovo, Enriques の定理\footnote{J. Lüroth: Beweis eines Satzes über rationale Kurven, \emph{Math. Ann.} 9 (1875). --- G. Castelnuovo: Sulla razionalità delle involuzioni piane, \emph{Math. Ann.} 44 (1893). --- F. Enriques: Sopra una involuzione non razionale dello spazio, \emph{Rend. Acc. Linc.}, vol. 21, 1912年1月21日.}から次の事実が得られる。

I. 任意の体 $\Kfield_{n1}$ は最小基底をもつ。それはその体の Lüroth 関数であり、一次分数変換を除いて一意に定まる。

II. 体 $\Kfield_{n2}$ については、係数体 $\Omega$ の代数拡大を許すならば、そして一般にはその場合に限って、常に最小基底が存在する。

III. 体 $\Kfield_{n\rho}$（$\rho\ge3$）については、一般には最小基底は存在しない。最小基底をもつ特殊な体 $\Kfield_{n\rho}$ の特徴づけはまだ試みられていない。

E. Steinitz が与えた体 $\Kfield^*_{11}$ に対する Lüroth の定理の証明を簡単に示す。\footnote{前掲書, \S\ 24.}

$\Phi^*(z,u)$ を $\Kfield^*_{11}$ のインボリューション形式とし、$\psi^*(x)$ を、実際に不定元 $x$ を含む $\Phi^*(z,u)$ の任意の係数とする。すると、$\Omega(x)$ の $\Omega(\psi^*(x))$ に関する次数は、$\Omega(x)$ の $\Kfield^*_{11}$ に関する次数に等しいことが分かる。そこから体の同一性が従う。さらに、$\Omega(\chi^*(x))$ が $\Omega(\psi^*(x))$ に等しいならば、$\chi^*(x)$ と $\psi^*(x)$ の相互の有理従属性は、二つの関数の一次分数従属性を意味する。

したがって、転移原理により、Lüroth の定理は次の形でも述べられる。

\emph{体 $\Kfield_{n1}$ の最小基底は、インボリューション基底の任意の一つの関数からなる。また $\Kfield_{n1}$ のインボリューション基底の任意の二つの関数は、一次分数変換によって互いに依存する。}

G. Castelnuovo は、幾何学的なインボリューション理論を用いて、有理基底が三つの関数から成る体 $\Kfield^*_{22}$ について事実 II を証明している。転移原理は係数体 $\Omega$ の有限代数拡大のもとでも有効であるから、有理基底の存在によって、すべての体 $\Kfield_{n2}$ について証明が完結する。

III については、Enriques が三次元空間内の非有理的インボリューション、すなわち最小基底をもたない体 $\Kfield_{33}$ を構成していることを注意しておく。

% ---- Continuation: Paper 06 §§7--15 ----
\providecommand{\Kfield}{\mathfrak{K}}
\providecommand{\Ssys}{\mathfrak{S}}
\providecommand{\Jdom}{\mathfrak{J}}
\providecommand{\Gdom}{\mathfrak{G}}
\providecommand{\Lsys}{\mathfrak{L}}
\providecommand{\Omegaint}{[\Omega]}
\providecommand{\eps}{\varepsilon}

\section*{\S\ 7. 有理関数の任意の系 $\Ssys_{n\rho}$、線形族 $\Lsys_{n\rho}$、および整性領域 $\Jdom_{n\rho}$. 有理基底の存在.}

体 $\Kfield_{n\rho}$ の有理基底の存在から、いま、有理関数および整有理関数の任意の系に対する有理基底を導く。これらの系を、加法、乗法、除法という基本演算によって順次互いに生じる順に並べる。

I. いつものように、$\Ssys_{n\rho}$ は、$n$ 個の不定元の有理関数（または整有理関数）からなり、代数的階数 $\rho$ をもつ任意の系を表す。$\Omega$ の量も系に属するものとして数える。

II. 系 $\Lsys_{n\rho}$ が次の条件を満たすとき、これを\emph{線形族}と呼ぶ。

1. $f(x)$ とともに、$\Omega$ の任意の量 $c$ に対して $c\cdot f(x)$ も常に $\Lsys_{n\rho}$ に含まれる。

2. $f(x)$ と $g(x)$ とともに、$f(x)+g(x)$ も常に $\Lsys_{n\rho}$ に含まれる。

III. 線形族 $\Jdom_{n\rho}$ がさらに次の条件を満たすとき、これを\emph{整性領域}と呼ぶ。

3. $f(x)$ と $g(x)$ とともに、$f(x)\cdot g(x)$ も常に $\Jdom_{n\rho}$ に含まれる。

IV. 整性領域 $\Kfield_{n\rho}$ がさらに次の条件を満たすとき、これを\emph{体}と呼ぶ。

4. $f(x)$ と $g(x)$ とともに、ただし $g(x)\ne0$ のとき、商 $f(x):g(x)$ も常に $\Kfield_{n\rho}$ に含まれる。

条件 1 から 4 は、\S\ 1 で体について述べた条件である。\footnote{$f(x)$ と $g(x)$ は零次、すなわち $\Omega$ の量であってもよい。}

これらの領域を、以下で順次互いから導く。

a) 任意の系 $\Ssys_{n\rho}$ は、次の要請によって線形族 $\Lsys_{n\rho}$ に拡張される。

$\Lsys_{n\rho}$ は、$\Ssys_{n\rho}$ に加えて、$\Ssys_{n\rho}$ の有限個の関数によって $\Omega$ の係数で線形に表される関数 $h(x)$ のすべて、かつそのようなものだけを含む。
\[
  h(x)=c_1f_1(x)+c_2f_2(x)+\cdots+c_kf_k(x),
\]
ここで $f_1(x),\ldots,f_k(x)$ は $\Ssys_{n\rho}$ のすべての関数を走り、$c_1,\ldots,c_k$ は $\Omega$ のすべての量を走る。

実際、系の関数の有理的結合を付け加えても、系の代数的階数は保たれる。$\Lsys_{n\rho}$ は条件 1 および 2 を満たすので線形族である。さらに、$\Ssys_{n\rho}$ を含む任意の線形族は、1 と 2 によって $\Lsys_{n\rho}$ も含む。したがって $\Lsys_{n\rho}$ は $\Ssys_{n\rho}$ を含む最小の線形族である。

b) 同様に、線形族 $\Lsys_{n\rho}$ から、それを含む最小の整性領域 $\Jdom_{n\rho}$ が次の要請によって生じる。

$\Jdom_{n\rho}$ は、$\Lsys_{n\rho}$ に加えて、$\Lsys_{n\rho}$ の有限個の関数 $f_1(x),\ldots,f_k(x)$ の、$\Omega$ 係数の\emph{整有理結合}として表される関数 $h(x)$ のすべて、かつそのようなものだけを含む。ここで $f_1(x),\ldots,f_k(x)$ は $\Lsys_{n\rho}$ のすべての関数を走る。

c) 最後に、整性領域 $\Jdom_{n\rho}$ から、それを含む最小の体 $\Kfield_{n\rho}$ が次の要請によって生じる。すなわち、$\Kfield_{n\rho}$ は $\Jdom_{n\rho}$ に加えて、$\Jdom_{n\rho}$ の二つの関数の商として表される関数 $h(x)$ のすべて、かつそのようなものだけを含む。

三つの拡張を一つの段階にまとめると、次を得る。

\emph{任意の系 $\Ssys_{n\rho}$ は、次の要請によって、それを含む最小の体 $\Kfield_{n\rho}$ へ拡張できる。}

\emph{$\Kfield_{n\rho}$ は、$\Ssys_{n\rho}$ に加えて、$\Ssys_{n\rho}$ の有限個の関数 $f_1(x),\ldots,f_k(x)$ の、$\Omega$ 係数の有理的結合として表される関数 $h(x)$ のすべて、かつそのようなものだけを含む。ここで $f_1(x),\ldots,f_k(x)$ は $\Ssys_{n\rho}$ のすべての関数を走る。}

この最後の事実から、任意の領域に対する有理基底の存在が直ちに従う。

$\Kfield_{n\rho}$ を $\Ssys_{n\rho}$ を含む最小の体とし、$\Kfield_{n\rho}$ の有理基底が
\begin{equation}
  h_1(x),\ h_2(x),\ldots,h_\sigma(x)
\tag{1}
\end{equation}
で与えられているとする。$\Kfield_{n\rho}$ の定義により、関係式
\begin{equation}
\begin{aligned}
 h_1(x)&=r_1\bigl(g_1(x),\ldots,g_\lambda(x)\bigr);\ldots;\\
 h_\sigma(x)&=r_\sigma\bigl(g_\mu(x),\ldots,g_\nu(x)\bigr),
\end{aligned}
\tag{2}
\end{equation}
が存在する。ただし $r_1,\ldots,r_\sigma$ はその引数の有理関数であり、
\begin{equation}
  g_1(x),\ldots,g_\lambda(x),\ldots,
  g_\mu(x),\ldots,g_\nu(x)
\tag{3}
\end{equation}
はすべて $\Ssys_{n\rho}$ に属する。関係式 (2) により、(3) も (1) と同じく $\Kfield_{n\rho}$ の有理基底を表す。また (3) の関数は $\Ssys_{n\rho}$ に属するから、これは同時に $\Ssys_{n\rho}$ の有理基底である。したがって次を得る。

\emph{定理 IV: 有理関数（または整有理関数）の任意の系 $\Ssys_{n\rho}$ は、その系の有限個の関数から成る、代数的階数 $\rho$ の有理基底をもつ。}

つぎに、特に $\Lsys_{n\rho}$ を線形族とし、
\begin{equation}
  g_1(x),\ldots,g_\rho(x),\ g_{\rho+1}(x),\ldots,g_\nu(x)
\tag{4}
\end{equation}
を $\Lsys_{n\rho}$ の有理基底とする。例えば
\[
  g_1(x),\ldots,g_\rho(x)
\]
が代数的に独立であるとする。このとき
\[
  g(x)=c_1g_{\rho+1}(x)+c_2g_{\rho+2}(x)+\cdots+c_{\nu-\rho}g_\nu(x)
\]
と置けば、$c_i$ を $\Omega$ の数として適当に選ぶことにより
\begin{equation}
  g_1(x),\ldots,g_\rho(x),\ g(x)
\tag{5}
\end{equation}
が $\Lsys_{n\rho}$ を含む最小の体 $\Kfield_{n\rho}$ の有理基底となるようにできる。ところが $\Lsys_{n\rho}$ は線形族であるから、$g(x)$ は $\Lsys_{n\rho}$ に属し、(5) は $\Lsys_{n\rho}$ の有理基底を表す。すなわち、

\emph{定理 V: 有理関数（または整有理関数）の任意の線形族 $\Lsys_{n\rho}$ の有理基底、従って特に任意の整性領域 $\Jdom_{n\rho}$ および任意の体 $\Kfield_{n\rho}$ の有理基底は、高々 $\rho+1$ 個の関数（かつ少なくとも $\rho$ 個）を含むように選ぶことができる。}

したがって、\S\ 4 で体について得られた有理基底の関数個数の制限は、体が線形族であるという性質に基づく。一方、最小基底が存在する場合に $\rho$ 個の関数まで減らせることは、体のすべての仮定を用いる。

なお、\S\ 3 (4) によって、系およびそれを含む最小系のすべての特徴的性質は写像によって保存される。

\section*{\S\ 8. 多項式からなる整性領域 $\Jdom_{n\rho}$ のインボリューション基底.}

以下の節では、元が整有理関数（多項式）である系 $\Ssys_{n\rho}$ を扱う。まず、多項式からなる整性領域 $\Jdom_{n\rho}$ を考える。これらの領域 $\Jdom_{n\rho}$ については、最初に割り切りの概念を定めなければならない。

$F(x),G(x)$ を $\Jdom_{n\rho}$ の二つの多項式とし、$L(x)$ をそれらの共通因子、すなわち
\[
  F(x)=L(x)\cdot F_1(x),\qquad G(x)=L(x)\cdot G_1(x)
\]
が成り立つものとする。このとき $L(x)$ が\emph{$\Jdom_{n\rho}$ における共通因子}であるとは、三つの多項式 $L(x),F_1(x),G_1(x)$ がすべて $\Jdom_{n\rho}$ に属すること、かつそのときに限る。

$\Jdom_{n\rho}$ の二つの多項式は、$\Jdom_{n\rho}$ における共通因子をもたないとき、\emph{$\Jdom_{n\rho}$ において互いに素}であるという。

$\Kfield_{n\rho}$ を $\Jdom_{n\rho}$ を含む最小の体とする。すると定義により、$\Kfield_{n\rho}$ の任意の関数は
\[
  f(x)=F_1(x):F_2(x)
\]
の形に表される。ここで $F_1(x),F_2(x)$ は $\Jdom_{n\rho}$ に属し、$\Jdom_{n\rho}$ において互いに素である。同様に、$\Kfield_{n\rho}$ の複数の関数は $\Jdom_{n\rho}$ 内の共通分母にそろえることができる。
\[
  f_1(x)=\frac{F_1(x)}{F(x)},\ldots,
  f_k(x)=\frac{F_k(x)}{F(x)}.
\]
$F(x)$ が\emph{$\Jdom_{n\rho}$ における最小共通分母}であるとは、$\Jdom_{n\rho}$ の多項式
\[
  F(x),F_1(x),\ldots,F_k(x)
\]
が $\Jdom_{n\rho}$ において互いに素であること、かつそのときに限る。

このような表示は複数の仕方で可能でありうる。なぜなら $\Jdom_{n\rho}$ では、一般に既約関数への一意分解の法則が成り立たないからである。

つぎに、$\Kfield_{n\rho}$ のインボリューション基底が $\Jdom_{n\rho}$ に対してもつ意味を調べる。そのため、インボリューション形式とインボリューション基底の概念を整性領域へ移す。$\Kfield_{n\rho}$ のインボリューション形式を
\[
  \Phi(z,u)=z^i+\sum{}' g_{\alpha\alpha_1\ldots\alpha_\rho}(x)
  z^\alpha u_1^{\alpha_1}\cdots u_\rho^{\alpha_\rho},
  \qquad
  (\alpha+\alpha_1+\cdots+\alpha_\rho=i)
\]
とし、$G(x)$ を $\Phi(z,u)$ の係数の $\Jdom_{n\rho}$ における最小共通分母とする。$G(x)$ を掛けると、$\Phi(z,u)$ は、係数が $\Jdom_{n\rho}$ の多項式であり、かつ $\Jdom_{n\rho}$ において互いに素である形式 $\Psi(z,u)$ へ移る。
\[
  G(x)\cdot\Phi(z,u)=\Psi(z,u)
  =G(x)z^i+
  \sum{}' G_{\alpha\alpha_1\ldots\alpha_\rho}(x)
  z^\alpha u_1^{\alpha_1}\cdots u_\rho^{\alpha_\rho},
  \qquad
  (\alpha+\alpha_1+\cdots+\alpha_\rho=i).
\]

$\Kfield_{n\rho}$ のインボリューション形式 $\Phi(z,u)$ に $\Jdom_{n\rho}$ の多項式を掛けることによって得られ、係数が $\Jdom_{n\rho}$ の多項式で、かつ $\Jdom_{n\rho}$ において互いに素であるような形式 $\Psi(z,u)$ を、$\Jdom_{n\rho}$ の\emph{インボリューション形式}と呼ぶ。また $\Psi(z,u)$ の係数を対応する\emph{インボリューション基底}と呼ぶ。

$\Kfield_{n\rho}$ のインボリューション基底の意味から、$\Jdom_{n\rho}$ の任意のインボリューション基底が同時に有理基底であることはまず明らかである。さらに調べるため、対応する最小の体 $\Kfield^*_{\rho\rho}$ で
\[
  z=u_1x_1+\cdots+u_\rho x_\rho
\]
を根にもつ既約方程式が $\Phi^*(z,u)=0$ で与えられるような写像領域 $\Jdom^*_{\rho\rho}$ を考える。

\S\ 5 によれば、
\[
  \Phi^*(z,u)=z^i+\sum{}' g^*_{\alpha\alpha_1\ldots\alpha_\rho}(x)
  z^\alpha u_1^{\alpha_1}\cdots u_\rho^{\alpha_\rho},
  \qquad
  (\alpha+\alpha_1+\cdots+\alpha_\rho=i)
\]
の係数は、$\Kfield^*_{\rho\rho}$ に関して共役な量の行
\[
  x,\quad x^{(1)},\ldots,x^{(i-1)}
\]
の基本対称関数を表す。$F^*(x)$ を $\Kfield^*_{\rho\rho}$ の任意の多項式とする。すると
\[
  F^*(x)=\frac1i\Bigl(F^*(x)+F^*(x^{(1)})+
  \cdots+F^*(x^{(i-1)})\Bigr)
\]
であるから、$F^*(x)$ はこれらの量の行の整対称関数である。従って、それは基本対称関数、すなわち関数
\[
  g^*_{\alpha\alpha_1\ldots\alpha_\rho}(x)
\]
の整有理結合として表される。

いま
\[
  G^*(x)\cdot\Phi^*(z,u)=\Psi^*(z,u)
  =G^*(x)z^i+
  \sum{}'G^*_{\alpha\alpha_1\ldots\alpha_\rho}(x)
  z^\alpha u_1^{\alpha_1}\cdots u_\rho^{\alpha_\rho},
  \qquad
  (\alpha+\alpha_1+\cdots+\alpha_\rho=i)
\]
が $\Jdom^*_{\rho\rho}$ のインボリューション形式であるならば、$\Kfield^*_{\rho\rho}$ の任意の多項式、とくに $\Jdom^*_{\rho\rho}$ の任意の多項式は、商
\[
  G^*_{\alpha\alpha_1\ldots\alpha_\rho}(x):G^*(x)
\]
の整有理結合である。

転移原理を考慮すると、次を得る。

\emph{定理 VI: 多項式からなる任意の整性領域 $\Jdom_{n\rho}$ に対して、少なくとも一つの特別な有理基底、すなわちインボリューション基底が存在し、$\Jdom_{n\rho}$ のすべての多項式をこの基底の関数で有理表示するとき、分母にはただ一つの固定された関数（対応するインボリューション形式の最高係数）の冪だけが現れる。}

\section*{\S\ 9. 多項式からなる相対的に整な領域 $\Gdom_{n\rho}$.}

以下では、多項式からなる特殊な整性領域を考える。これは整性領域と体との中間段階をなす。

多項式の整性領域 $\Gdom_{n\rho}$ が次の条件を満たすとき、これを「相対的に整な領域」と呼ぶ。

\emph{4a) $\Gdom_{n\rho}$ は、$F(x)$ と $G(x)$ とともに、ただし $G(x)\ne0$ のとき、商 $F(x):G(x)$ を、その商が不定元について整、すなわち多項式である場合に、かつその場合に限り含む。}

まず、相対的に整な領域が、有理関数体 $\Kfield_{n\sigma}$（$\sigma\ge\rho$）の多項式全体と一致することを証明する。

$\Kfield_{n\rho}$ を $\Gdom_{n\rho}$ を含む最小の体とする。$\Kfield_{n\rho}$ の任意の関数 $f(x)$ には、$\Gdom_{n\rho}$ の二つの多項式の商としての表示
\[
  f(x)=F_1(x):F_2(x)
\]
がある。ここでは零次の多項式も許す。$f(x)$ が多項式になるならば、条件 4a) によって $\Gdom_{n\rho}$ に属する。また $\Gdom_{n\rho}$ はそれ以上の多項式を含むことはできないので、\emph{任意の相対的に整な領域 $\Gdom_{n\rho}$ は、それを含む最小の体の多項式全体と同一である。}

逆に、$\Kfield_{n\sigma}$ を任意の有理関数体（したがって、ある与えられた系を含む最小の体とは限らない）とし、$\Jdom$ を $\Kfield_{n\sigma}$ に含まれる多項式全体とする。$\Jdom$ は整性領域の条件 1, 2, 3 および条件 4a) を満たすので、相対的に整な領域である。すなわち、\emph{体 $\Kfield_{n\sigma}$ に含まれる多項式全体は相対的に整な領域 $\Gdom_{n\rho}$（$\rho\le\sigma$）をなす。}\footnote{$\rho=0$ であってもよい。すなわち、$\Kfield_{n\sigma}$ の多項式全体が係数領域 $\Omega$ の元だけからなる場合である。}

さらに、相対的に整な領域が Hilbert の「相対的に整な関数の整性領域」と一致することを示す。\footnote{D. Hilbert: Mathematische Probleme. パリ講演 1900. 問題14 (Göttinger Nachrichten 1900).}

$\Gdom_{n\rho}$ の有理基底を
\begin{equation}
  G_1(x),\ldots,G_k(x)
\tag{1}
\end{equation}
とする。条件 1, 2, 3, 4a により、(1) の多項式の任意の有理結合で、不定元について整、すなわち多項式となるものは $\Gdom_{n\rho}$ に属する。逆に、$\Gdom_{n\rho}$ の任意の多項式は、有理基底の概念により、(1) の多項式で有理的に表される。従って $\Gdom_{n\rho}$ は、\emph{(1) の多項式の有理結合のうち、不定元について整、すなわち多項式であるもの全体と同一である。} こうして、特殊な有理基底から出発する Hilbert の定義が得られる。一方、われわれの定義は領域の抽象的性質だけを用いる。

相対的に整な領域については、\S\ 8 で一般の整性領域に対して与えた共通因子の定義が単純になる。

$F(x),G(x)$ を $\Gdom_{n\rho}$ の二つの多項式、$L(x)$ をそれらの共通因子とする。このとき $L(x)$ が\emph{$\Gdom_{n\rho}$ における共通因子}であるとは、$L(x)$ が $\Gdom_{n\rho}$ に属すること、かつそのときに限る。

というのは、商 $F(x):L(x)$ および $G(x):L(x)$ の属することは条件 4a) から従うからである。

相対的に整な領域のもう一つの本質的性質は、「$\Gdom_{n\rho}$ に関して代数的に整である」という概念を、代数的数体または体 $\Omega(x_1,\ldots,x_n)$ の場合とまったく同様に、任意の方程式によって定義できることである。

\emph{定義 V: 量 $\xi$ が $\Gdom_{n\rho}$ に関して代数的に整である、または相対的に代数的に整であるとは、最高係数が単位で、残りの係数が $\Gdom_{n\rho}$ の多項式であるような方程式を少なくとも一つ満たすことをいう。}

実際、相対的に代数的に整な量 $\xi$ が、例えば方程式
\[
  L(z)=z^\lambda+G_1(x)z^{\lambda-1}+\cdots+G_\lambda(x)=0
\]
を満たすとする。$L(z)$ は、最小の体 $\Kfield_{n\rho}$ に関して既約な整関数
\[
  M(z)=z^\mu+H_1(x)z^{\mu-1}+\cdots+H_\mu(x)
\]
で割り切れる。この割り切りから、Gauss の定理の既知の拡張\footnote{例えば J. König: \emph{Einleitung in die allgemeine Theorie der algebraischen Größen} (Leipzig, B.~G. Teubner, 1903), Kap. IX, \S\ 2, または Weber（小版）\S\ 20 を参照。}により、有理関数 $H_1(x),\ldots,H_\mu(x)$ は $\Kfield_{n\rho}$ の多項式であり、従って $\Gdom_{n\rho}$ に属する。したがって、相対的に代数的に整な量 $\xi$ の定義は既約方程式から得られる。とくに、$\Gdom_{n\rho}$ の任意の多項式 $F(x)$ は、方程式 $z-F(x)=0$ を満たすから、$\Gdom_{n\rho}$ に関して代数的に整である。

さらに、多項式の整性領域 $\Jdom_{n\rho}$ に対して、最小の体の場合と同様に、それを含む最小の相対的に整な領域 $\Gdom_{n\rho}$ も次の要請で定義されることを述べておく。

$\Gdom_{n\rho}$ は、$\Jdom_{n\rho}$ に加えて、$\Jdom_{n\rho}$ の二つの多項式の商として表される多項式 $H(x)$ のすべて、かつそのようなものだけを含む。

この定義からさらに次が従う。$\Jdom_{n\rho}$ の任意のインボリューション形式およびインボリューション基底から、基底のすべての多項式の $\Gdom_{n\rho}$ における共通因子を、あっても一つ取り除くことにより、$\Gdom_{n\rho}$ の対応するものが得られる。

なぜなら $\Jdom_{n\rho}$ と $\Gdom_{n\rho}$ は同じ最小包含体 $\Kfield_{n\rho}$ をもつからである。また $\Jdom_{n\rho}$ の任意のインボリューション形式は、$\Kfield_{n\rho}$ のインボリューション形式に $\Jdom_{n\rho}$ の多項式を掛けることによって生じるので、$\Gdom_{n\rho}$ の係数をもつ。ただし、それらは $\Gdom_{n\rho}$ における共通因子をもつことがある。\footnote{例えば、$x_1^2$ と $x_1x_2$ のすべての整有理結合からなる整性領域 $\Jdom_{22}$ に対して、インボリューション形式として
\[
\Psi(z,u)=x_1^2z^2-x_1^4u_1^2-2x_1^3x_2u_1u_2-x_1^2x_2^2u_2^2
\]
が得られる。$x_1^2$ は $\Jdom_{22}$ に属し、従ってなおさら最小の相対的に整な包含領域 $\Gdom_{22}$ に属し、従って $\Gdom_{22}$ における共通因子であるから、$\Gdom_{22}$ のインボリューション形式として
\[
\Psi(z,u)=z^2-x_1^2u_1^2-2x_1x_2u_1u_2-x_2^2u_2^2
\]
が得られる。}

最後に、相対的に整な領域 $\Gdom_{n\rho}$ の\emph{写像領域}は、\S\ 7 により再び多項式の整性領域 $\Jdom^*_{\rho\rho}$ であるが、それ自体は必ずしも相対的に整な領域ではないことを注意する。実際、$F(x)$ と $G(x)$ が $\Gdom_{n\rho}$ の二つの多項式で、その商が多項式でなく従って $\Gdom_{n\rho}$ に属さないならば、商 $F^*(x):G^*(x)$ も $\Jdom^*_{\rho\rho}$ に属さない。しかし、この商は、いくつかの不定元を特殊化して生じるため、多項式でありうる。その場合、$\Jdom^*_{\rho\rho}$ では条件 4a) が満たされない。\footnote{例えば $\Gdom_{22}$ を、$F=x_1^2x_2$ と $G=x_1^2(1+x_3)$ の有理結合であるすべての多項式からなるものとする。許容される特殊化 $x_3=0$ は $F^*:G^*=x_2$ を与えるが、$F:G$ は多項式でない。従って $\Jdom^*_{22}$ は相対的に整な領域ではない。}

\section*{\S\ 10. 第一種の相対的に整な領域とその整性基底.}

\S\ 9（定義 V）で導入した「相対的に代数的に整」という概念により、有限な整性領域である相対的に整な領域の一類を与えることができ、これによって D. Hilbert が提起した問題（数学の問題、問題 14）に、少なくともこの類について肯定的に答えることができる。次の定義を置く。

\emph{定義 VI: $\Gdom_{n\rho}$ の有限個の多項式であって、$\Gdom_{n\rho}$ の任意の多項式が、$\Omega$ の係数を用いてそれら有限個の整有理結合として表されるものを整性基底という。整性基底をもつ多項式の整性領域を有限整性領域という。}

ここで考える相対的に整な領域の類、すなわちインボリューション基底が整性基底であることが示される類を、次の定義に従って第一種の相対的に整な領域と呼ぶ。

\emph{定義 VII: 相対的に整な領域 $\Gdom_{n\rho}$ が第一種であるとは、その写像領域として相対的に整な領域 $\Gdom^*_{\rho\rho}$ をもち、$x_1,\ldots,x_\rho$ の任意の多項式（係数は $\Omega$）が $\Gdom^*_{\rho\rho}$ 上代数的に整であることをいう。}

この定義により、不定元 $x_1,\ldots,x_\rho$、従って線形形式 $u(x)=u_1x_1+\cdots+u_\rho x_\rho$ も $\Gdom^*_{\rho\rho}$ 上代数的に整である。根 $z=u(x)$ をもつ既約整関数、すなわち最小の体 $\Kfield^*_{\rho\rho}$ のインボリューション形式は、従って最高係数として単位をもち、残りの係数として $\Gdom^*_{\rho\rho}$ の多項式をもつ。これは同時に $\Gdom^*_{\rho\rho}$ のインボリューション形式であり、$\Gdom^*_{\rho\rho}$ にはそれ以外のインボリューション形式は存在しない。従って転移原理により、$\Gdom_{n\rho}$ も最高係数が単位であるインボリューション形式をもつ。さらに定理 VI によれば、$\Gdom_{n\rho}$ のすべての多項式を対応するインボリューション基底で有理表示するとき、分母に現れるのはこの最高係数だけであるから、インボリューション基底は整性基底になる。従って、

\emph{定理 VII. 第一種の任意の相対的に整な領域は有限整性領域である。その整性基底は、それを特徴づける写像領域から得られるインボリューション基底で与えられる。特に、第一種の相対的に整な領域 $\Gdom_{\rho\rho}$ の整性基底は、一意に定まる $\Gdom_{\rho\rho}$ のインボリューション基底によって与えられる。}

第一種の相対的に整な領域について判定条件を付け加える。$\Gdom_{n\rho}$ を第一種とし、$\Gdom_{n\rho}$ の整性基底が
\begin{equation}
  F_1(x),F_2(x),\ldots,F_k(x)
\tag{1}
\end{equation}
で与えられているとする。すると Hilbert の補助定理\footnote{D. Hilbert: Über die vollen Invariantensysteme, Math. Ann. 42 (1893), \S\ 1. そこで同次多項式についてだけ述べられた補助定理は、非同次多項式についても同様に成り立つ。}により、$\Gdom_{n\rho}$ の $\rho$ 個の代数的に独立な多項式
\begin{equation}
  G_1(x),G_2(x),\ldots,G_\rho(x),
\tag{2}
\end{equation}
が存在し、(1) の多項式、従って $\Gdom_{n\rho}$ の任意の多項式は、(2) の多項式上代数的に整である。同じことは、写像領域 $\Gdom^*_{\rho\rho}$ の任意の多項式について、対応する多項式
\begin{equation}
  G_1^*(x),G_2^*(x),\ldots,G_\rho^*(x)
\tag{3}
\end{equation}
に関して成り立つ。

従って定義 VII により、$x_1,\ldots,x_\rho$ の任意の多項式も (3) の多項式上代数的に整である。

逆に、任意の相対的に整な領域 $\Gdom_{n\rho}$ の写像領域 $\Jdom^*_{\rho\rho}$ が、$x_1,\ldots,x_\rho$ の任意の多項式を代数的に整にするような $\rho$ 個の多項式 (3) を含むと仮定する。このとき $\Jdom^*_{\rho\rho}$ は相対的に整な領域 $\Gdom^*_{\rho\rho}$ になり、さらに $\Gdom^*_{\rho\rho}$、従って $\Gdom_{n\rho}$ が第一種であることを示す。

実際、$F^*(x)$ を $\Jdom^*_{\rho\rho}$ を含む最小の体 $\Kfield^*_{\rho\rho}$ の多項式とする。仮定により、それは (3) の多項式上代数的に整である。対応する $\Kfield_{n\rho}$ の関数は $\Gdom_{n\rho}$ の $\rho$ 個の多項式上代数的に整であるから、多項式 $F(x)$ であり、従って $\Gdom_{n\rho}$ に属する。したがって $F^*(x)$ は $\Jdom^*_{\rho\rho}$ に属する。よって写像領域 $\Jdom^*_{\rho\rho}$ は $\Kfield^*_{\rho\rho}$ のすべての多項式を含み、相対的に整な領域 $\Gdom^*_{\rho\rho}$ である。$\Gdom^*_{\rho\rho}$、従って $\Gdom_{n\rho}$ が第一種であることは、定義 V と VII および仮定から直ちに従う。すなわち、

\emph{第一種の相対的に整な領域 $\Gdom_{n\rho}$ は、写像領域 $\Jdom^*_{\rho\rho}$ の中に $\rho$ 個の多項式が存在し、それらの上で $x_1,\ldots,x_\rho$ の任意の多項式が代数的に整である、という事実によって特徴づけられる。この仮定から、写像領域そのものが相対的に整な領域であることが従う。}\footnote{例えば、\S\ 9 の末尾の注意で述べた領域 $\Gdom_{22}$ では、許容される特殊化 $x_1=1$ によって二つの多項式 $F^*=x_2$, $G^*=1+x_3$ が得られ、これらの上で $x_2,x_3$ の任意の多項式が代数的に整である。従って $\Gdom_{22}$ は第一種であり、実際 $F(x)$ と $G(x)$ を整性基底にもつ。}

\section*{\S\ 11. 代数的整従属性に関する補助定理.}

いま、$x_1,\ldots,x_\rho$ の任意の多項式が、$x_1,\ldots,x_\rho$ の $\rho$ 個の多項式
\begin{equation}
  G_1^*(x),\ldots,G_\rho^*(x)
\tag{1}
\end{equation}
上代数的に整であるための判定条件を与える。

特に (1) の多項式が同次形式であるならば、(1) を消す任意の特殊な値系において、代数的に整従属なすべての形式も同時に消える。従って特に $x_1,\ldots,x_\rho$ も消える。

従って、形式 (1) は
\[
  x_1=0,
  \ldots,
  x_\rho=0
\]
以外の値系で同時に消えることはできない。それらの終結式は零であってはならない。

逆に、この条件も十分であり、次の形で非同次多項式へ拡張できる。

\emph{補助定理: $x_1,\ldots,x_\rho$ の任意の多項式が、$x_1,\ldots,x_\rho$ の $\rho$ 個の多項式}
\[
  G_1^*(x),\ldots,G_\rho^*(x)
\]
\emph{上代数的に整であるための十分条件は、これらの多項式 (1) の最高次元項}
\[
  H_1^*(x),\ldots,H_\rho^*(x)
\]
\emph{の、$x_1,\ldots,x_\rho$ に関する終結式が消えないことである。}
\begin{equation}
  R_0=\left[
  \begin{array}{ccc}
    H_1^*&\cdots&H_\rho^*\\
    x_1&\cdots&x_\rho
  \end{array}
  \right]\ne0.
\tag{2}
\end{equation}
\emph{同次形式 (1) に対しては、条件 (2) は必要である。}\footnote{終結式論については、F. Mertens, Zur Theorie der Elimination (I. and II. Teil), Sitzungsber. d. k. Ak. d. Wiss. Wien, Math.-naturw. Kl. Abt. IIa, Bd. 108 (1899) を参照。J. König, Theorie d. algebr. Größen, Kap. VI に部分的に再掲されている。}

証明のために、多項式
\begin{equation}
\begin{gathered}
  \bar G_1^*(x)=G_1^*(x)-\lambda_1,
  \ldots,
  \bar G_\rho^*(x)=G_\rho^*(x)-\lambda_\rho,\\
  \bar\Phi^*(x)=\Phi^*(x)-\mu
\end{gathered}
\tag{3}
\end{equation}
を作る。ここで $\lambda_1,\ldots,\lambda_\rho,\mu$ は不定元、$\Phi^*(x)$ は $x_1,\ldots,x_\rho$ の任意の多項式である。多項式 (3) の $x_1,\ldots,x_\rho,1$ に関する終結式\footnote{すなわち、多項式 (3) を追加の不定元 $x_0$ によって、$x$ に関する次数が保たれるように同次化したときの、$x_1,\ldots,x_\rho,x_0$ に関する同次終結式である。}を
\begin{equation}
  R(\lambda,\mu)=\left[
  \begin{array}{cccc}
   \bar G_1^*&\cdots&\bar G_\rho^*&\bar\Phi^*\\
   x_1&\cdots&x_\rho&1
  \end{array}
  \right]
\tag{4}
\end{equation}
と書く。これは $\lambda_1,\ldots,\lambda_\rho,\mu$ の整有理関数である。F. Mertens によれば、(4) を $\mu$ の冪で展開したとき、最高係数を明示的に指定できる。\footnote{同書 \S\ 3（定理 I の証明）および特に \S\ 6（終結式 $R$ の項のうち、与えられた冪積 $a_{\lambda\lambda}^{p_\lambda}a_{\mu\mu}^{p_\mu}\cdots a_{\sigma\sigma}^{p_\sigma}$ を因子として含むもの全体の決定）。本場合には $p_\lambda=0$, $p_\mu=0,\ldots,p_\rho=\tau$ とし、$a_{\sigma\sigma}=(-\mu)$ と置く。König, Kap. VI, \S\ 9 に再掲。}すなわち
\begin{equation}
  (-1)^\tau R(\lambda,\mu)
  =R_0^\sigma\mu^\tau+A_1(\lambda)\mu^{\tau-1}+\cdots+A_\tau(\lambda),
\tag{5}
\end{equation}
を得る。ここで $A_1(\lambda),\ldots,A_\tau(\lambda)$ は、$\lambda_1,\ldots,\lambda_\rho$ および多項式 (3) の係数の、整数係数の整関数を表す。この展開はまず、仮定 (2)、すなわち $R_0\ne0$ のもとでは、$R(\lambda,\mu)$ も恒等的に消えないことを示す。$x_1,\ldots,x_\rho,\lambda_1,\ldots,\lambda_\rho,\mu$ に関する恒等式
\[
  R(\lambda,
\mu)\equiv 0\pmod{\bar G_1^*(x),\ldots,
\bar G_\rho^*(x),\bar\Phi^*(x)}
\]
は、代入
\[
  \lambda_i=G_i^*(x),\qquad \mu=\Phi^*(x)
\]
によって
\begin{equation}
\begin{aligned}
  R(G_i^*(x),\Phi^*(x))
  &=R_0^\sigma\Phi^{*\tau}(x)
    +A_1(G_i^*(x))\Phi^{*\tau-1}(x)+\cdots \\
  &\qquad +A_\tau(G_i(x))=0
\end{aligned}
\tag{6}
\end{equation}
となる。仮定により $R_0$ は $\Omega$ の数であるから、これで補助定理が証明される。

\section*{\S\ 12. 多項式からなる正則系 $\Ssys_{n\rho}$ の整性基底.}

\S\ 11 の補助定理を \S\ 10 の結論と合わせると、直ちに次の事実が得られる。

\emph{$\Gdom_{n\rho}$ の写像領域 $\Jdom^*_{\rho\rho}$ に、最高次元項の同次終結式が恒等的に消えない、すなわち $R_0\ne0$ となる $\rho$ 個の多項式が存在するならば、$\Gdom_{n\rho}$ は第一種であり、従ってインボリューション基底を整性基底とする有限整性領域である。}

条件 $R_0\ne0$ は、補助定理によって、整性基底の存在の第二の証明を可能にする。この証明は本質的には D. Hilbert によるものである。\footnote{Über die vollen Invariantensysteme, \S\ 2. Hilbert の場合には、不変式体の整性基底の存在が別の方法で得られていることを仮定して、その基底を Kronecker の体の代数的に整な量の基本系として解釈することだけが問題となる。} ただしこの証明では、インボリューション基底をそれとして認識することはできない。しかし、示した性質をもつすべての系 $\Ssys_{n\rho}$ に一様に適用される。

実際、$R_0\ne0$ となり、従って代数的に独立である $\Jdom^*_{\rho\rho}$ の $\rho$ 個の多項式を
\begin{equation}
  G_1^*(x),\ldots,G_\rho^*(x)
\tag{1}
\end{equation}
とする。補助定理により、$x_1,\ldots,x_\rho$ の任意の多項式、従って特に $\Jdom^*_{\rho\rho}$ を含む最小の体 $\Kfield^*_{\rho\rho}$ の任意の多項式は、(1) の多項式上代数的に整である。$\Ssys_{n\rho}$ における対応する多項式を
\begin{equation}
  G_1(x),\ldots,G_\rho(x)
\tag{2}
\end{equation}
とすれば、転移原理により、$\Ssys_{n\rho}$ を含む最小の体 $\Kfield_{n\rho}$ の任意の多項式、従って特に $\Ssys_{n\rho}$ の任意の多項式は、(2) の多項式上代数的に整である。逆に、(2) の多項式上代数的に整である $\Kfield_{n\rho}$ の任意の関数は多項式である。したがって、$\Kfield_{n\rho}$ を（\S\ 4 に従って）$\Omega(G_1,\ldots,G_\rho)$ 上の代数体と見ると、$\Kfield_{n\rho}$ の多項式全体、すなわち相対的に整な領域 $\Gdom_{n\rho}$ は、この体の代数的に整な量全体と同一である。$G(x)$ を (2) を有理基底へ補う $\Kfield_{n\rho}$ の多項式とし、$D(x)$ を $G(x)$ が $\Omega(G_1,\ldots,G_\rho)$ に関して満たす次数 $k$ の既約方程式の判別式とする。すると $\Kfield_{n\rho}$ の任意の多項式、特に $\Ssys_{n\rho}$ の任意の多項式 $F(x)$ は、代数的に整な量として表示
\begin{equation}
  F(x)=\frac{\Gamma_1(G_i)G(x)^{k-1}+\Gamma_2(G_i)G(x)^{k-2}+\cdots+\Gamma_k(G_i)}{D(x)}
\tag{3}
\end{equation}
をもつ。

いま $F(x)$ に $\Ssys_{n\rho}$ のすべての多項式を走らせる。すると、対応する関数
\[
  v_1\Gamma_1(G_i)+v_2\Gamma_2(G_i)+\cdots+v_k\Gamma_k(G_i)
\]
からなる系に Hilbert の加群基底に関する定理 I（Math. Ann. 36）を適用することで、$\Ssys_{n\rho}$ の有限個の多項式
\begin{equation}
  F_1(x),F_2(x),\ldots,F_\lambda(x)
\tag{4}
\end{equation}
が存在し、$\Ssys_{n\rho}$ の任意の多項式が
\begin{equation}
  F(x)=A_1(G_i)F_1(x)+A_2(G_i)F_2(x)+\cdots+A_\lambda(G_i)F_\lambda(x)
\tag{5}
\end{equation}
の形に表されることが分かる。ここで $A_1,\ldots,A_\lambda$ は $G_i(x)$ の多項式である。従って、多項式 (2) と (4) は $\Ssys_{n\rho}$ の整性基底をなす。すなわち、

\emph{定理 VIII: 多項式系 $\Ssys_{n\rho}$ の写像系 $\Jdom^*_{\rho\rho}$ に、最高次元項の終結式が恒等的に消えない、すなわち $R_0\ne0$ となる $\rho$ 個の多項式が存在するならば、$\Ssys_{n\rho}$ は整性基底をもつ。}

この定理は少し一般化できる。$R_0\ne0$ となる $\rho$ 個の多項式 (1) が、$\Ssys^*_{\rho\rho}$ を含む最小の整性領域 $\Jdom^*_{\rho\rho}$ に属していれば十分である。実際、定理 VIII に至る議論は、その場合にも表示 (5) が成り立つことを示す。ところで定義により、最小の包含整性領域 $\Jdom_{n\rho}$ の任意の多項式 $L(x)$ は、$L(x)$ に依存する有限個の $\Ssys_{n\rho}$ の多項式によって整有理的に表される。従って、$\Ssys_{n\rho}$ の $\sigma$ 個の多項式
\begin{equation}
  K_1(x),K_2(x),\ldots,K_\sigma(x)
\tag{6}
\end{equation}
が存在し、多項式 $G_i(x)$ は (6) の整有理結合となる。よって (6) と (4) が整性基底を与える。次の用語を導入する。

\emph{定義 VIII: 多項式の系 $\Ssys_{n\rho}$ が正則であるとは、最小包含整性領域の写像領域 $\Jdom^*_{\rho\rho}$ に、その最高次元項の終結式が恒等的に消えない、すなわち $R_0\ne0$ となる $\rho$ 個の多項式が存在することをいう。}

従って次を得る。

\emph{定理 IX: 多項式からなる任意の正則系 $\Ssys_{n\rho}$ は整性基底をもつ。}\footnote{整性基底をもたない非正則系が存在することは、Hilbert が例によって示した（Gött. Nachr. 1891, p. 233）。次の例はさらに簡単である。
\[
  \Ssys_{22}=x_1,
  x_1x_2,
  x_1x_2^2,
  \ldots,
  x_1x_2^\nu,
  \ldots\quad\text{in inf.}
\]
ここで整性基底が存在することは、正整数 $\nu$ が存在して恒等式
\[
  x_1x_2^{\nu+\sigma}=\sum C\cdot x_1^{i_0}(x_1x_2)^{i_1}(x_1x_2^2)^{i_2}\cdots(x_1x_2^\nu)^{i_\nu}
  \quad(\sigma=1,2,\ldots)
\]
が成り立つことと同値である。すでに $\sigma=1$ について $x_1,x_2$ の次元を比較すると、指数に関する二つの条件式
\[
  1=i_0+i_1+\cdots+i_\nu,
  \qquad
  \nu+1=i_1+2i_2+\cdots+\nu i_\nu
\]
が得られるが、これは明らかに非負整数解をもたない。従って $\Ssys_{22}$ は整性基底をもたない。}

さらに、\S\ 5 のインボリューションの場合と同様に、正則系の幾何学的解釈を与える。そのため、そこでの (7) と同じく、方程式系
\begin{equation}
  L^*(x)-L^*(\xi)=0
\tag{7}
\end{equation}
を考える。ここで $L^*(x)$ は $\Jdom^*_{\rho\rho}$ のすべての多項式を走る。$x_0,\xi_0$ を用いて同次化すると、(7) は同次形式の間の方程式系
\begin{equation}
  \xi_0^{m}M^*(x)-x_0^{m}M^*(\xi)=0
\tag{8}
\end{equation}
へ移る。ここで $H^*(x)$ が $L^*(x)$ の最高次元項を表すならば、関係
\begin{equation}
  H^*(x)=M^*(x)\pmod{x_0}
\tag{9}
\end{equation}
が成り立つ。

$\Jdom^*_{\rho\rho}$ の\emph{基本点}（\S\ 5 の注を参照）とは、(8) の解で $x$ に依存しないもの、すなわち $x$ に共役な行とは異なるものをいう。したがって基本点とは、$\xi_1=0,\ldots,\xi_\rho=0$ とは異なる値系で、同時に
\[
  \xi_0^m=0,
  \qquad M^*(\xi)=0
\]
を満たすもの、また (9) によれば
\[
  \xi_0=0,
  \qquad H^*(\xi)=0
\]
を満たすものである。

いま $R_0\ne0$ となる、従って共通解系をもたない $\rho$ 個の形式 $H^*(\xi)$ が存在するならば、$\Jdom^*_{\rho\rho}$ は基本点をもてない。特に、$\Ssys^*_{\rho\rho}$ が同次形式から成る場合には、$R_0\ne0$ ならば $\Ssys^*_{\rho\rho}$ も基本点をもてない。そのような点は同時に $\Jdom^*_{\rho\rho}$ の基本点だからである。しかし逆も成り立つ。そのため、形式 $H^*(x)$ 全体からなる系 $\Jdom(H^*)$ を考える。これは再び整性領域である。

Hilbert の加群基底に関する定理 I と Hilbert の補助定理（\S\ 10 の引用を参照）から、$\Jdom(H^*)$ の $\rho$ 個の形式が存在し、それらの消滅が $\Jdom(H^*)$ のすべての形式の消滅を帰結することが従う。もし $\Jdom^*_{\rho\rho}$ が基本点をもたなければ、これら $\rho$ 個の形式は同時に消えることができず、その終結式は $R_0\ne0$ となる。

従って、

\emph{多項式の正則系 $\Ssys_{n\rho}$ は、最小包含整性領域の写像領域 $\Jdom^*_{\rho\rho}$ が基本点をもたないことによって特徴づけられる。同次形式の正則系 $\Ssys^*_{\rho\rho}$ については、系そのものがすでに基本点をもたない。}\footnote{前の注で与えた非正則系 $\Ssys_{22}$ は基本点
\[
  \xi_1:\xi_2:\xi_0=0:1:0
\]
をもつ。}

\section*{\S\ 13. 第一種の相対的に整な領域および正則系の例.}

1. 第一種の相対的に整な領域 $\Gdom_{nn}$ の第一の例として、\emph{$x_1,\ldots,x_n$ の多項式のうち、与えられた置換群 $G$ を許すもの全体}、すなわち Lagrange の種領域の多項式全体を考える。恒等式
\begin{equation}
  x_k^n-\sigma_1(x)x_k^{n-1}+\sigma_2(x)x_k^{n-2}+\cdots\pm\sigma_n(x)=0
  \qquad(k=1,2,
\ldots,n)
\tag{1}
\end{equation}
により、$x_1,\ldots,x_n$ は $\Gdom_{nn}$ に属する基本対称関数上代数的に整であり、従って定義 V により $\Gdom_{nn}$ 上代数的に整である。したがって $\Gdom_{nn}$ は第一種である。さらに、$\Gdom_{nn}$ は第一種の領域であり、その元が同次形式で、かつ $n=\rho$ であるから、\S\ 10 の結論により正則系をなす。実際、(1) によって、基本対称関数の終結式 $R_0$ は恒等的に消えない。終結式論の計算法則により、$R_0$ は $\pm1$ と計算される。$\Gdom_{nn}$ の\emph{インボリューション形式}は、$\Gdom_{nn}$ が第一種で $n=\rho$ であるため、対応する Lagrange の種領域のインボリューション形式として一意に定まる。この領域に関して $x_1,\ldots,x_n$ に共役な量は群 $G$ の置換によって生じるから、このインボリューション形式は
\begin{equation}
\begin{aligned}
  \Psi(z,u)&=\prod_i
  (z-u_1x_{i_1}-u_2x_{i_2}-\cdots-u_nx_{i_n})\\
  &=z^i+\sum{}'G_{\alpha\alpha_1\ldots\alpha_n}(x)
  z^\alpha u_1^{\alpha_1}\cdots u_n^{\alpha_n}
\end{aligned}
\tag{2}
\end{equation}
で与えられる。積は $G$ のすべての置換にわたって取る。式 (2) は「群 $G$ の Galois 形式」を表す。従って次の事実が成り立つ。

\emph{置換群 $G$ を許す $x_1,\ldots,x_n$ のすべての多項式は、群の Galois 形式の係数 $G_{\alpha\alpha_1\ldots\alpha_n}(x)$ の整有理結合である。}

2. \emph{正則系}の例として、多項式の系 $\Ssys_{n1}$、すなわち代数的に独立な多項式を一つだけ含む系を考える。

実際、$\Ssys^*_{11}$ を $\Ssys_{n1}$ の写像系とし、
\[
  F^*(x)=a_0x^k+a_1x^{k-1}+\cdots+a_k\qquad(a_0\ne0)
\]
を $\Ssys^*_{11}$ の多項式とする。このとき
\begin{equation}
  R_0=a_0\ne0.
\tag{3}
\end{equation}
従って $\Ssys_{n1}$ は定義 VIII により正則であり、したがって、

\emph{多項式の任意の系 $\Ssys_{n1}$、特に一つの不定元の無限個の多項式の任意の系は、整性基底をもつ。}\footnote{従って、整性基底をもたない最も簡単な非正則系は $\Ssys_{22}$ である。そのような系が実際に存在することは、Hilbert の例および \S\ 12 の注で与えた例が示している。}

3. いま系 $\Ssys_{n1}$ を相対的に整な領域 $\Gdom_{n1}$ に特殊化する。これは \S\ 12 により正則系であるから、第一種の相対的に整な領域である。従ってその整性基底はインボリューション基底で与えられ、この基底は同時に最小包含体 $\Kfield_{n1}$ のインボリューション基底でもある。しかし \S\ 6 により、$\Kfield_{n1}$ のインボリューション基底の任意の関数は同時に最小基底である。これを体の Lüroth 関数と呼んだ。そしてインボリューション基底の任意の関数は、この Lüroth 関数に一次分数変換によって依存する。現在の場合、Lüroth 関数は $\Gdom_{n1}$ に属するので多項式である。従って、インボリューション基底の他のすべての多項式は、この Lüroth 関数 $L(x)$ に一次分数的に依存し、かつ (3) によって代数的に整であるので、結局整かつ線形に依存する。したがって $L(x)$ が整性基底となる。$\Gdom^*_{11}$ のインボリューション形式は、$(u_1=1)$ と置けば
\[
  \Psi^*(z)=L^*(z)-L^*(x)
\]
であり、ここからも $L(x)$ が整性基底として現れる。よって、

\emph{相対的に整な領域 $\Gdom_{n1}$ は第一種である。その整性基底は一つの多項式、すなわち最小包含体の Lüroth 関数からなる。}

4. 相対的に整な領域 $\Gdom_{n1}$ の例として、$s$ 変数の二次形式の整有理な射影不変式を特に挙げておく。前述と一致して、それらがその形式の行列式 $|a_{ik}|$ の整有理結合として表されることはよく知られている。そしてこの行列式は、対応する不変式体の Lüroth 関数でもある。

\section*{\S\ 14. 整数多項式からなる系.}

得られた基底定理を、整数上の整性に関して精密化する。

そのため以下では、係数領域 $\Omega$ を（有限）代数的数体と仮定し、$[\Omega]$ は $\Omega$ の代数的整数全体を表すものとする。これを簡単に整数と呼ぶ。整数多項式とは、$[\Omega]$ からの係数をもつ多項式をいう。以下では $F(x),G(x),\ldots$ は整数多項式だけを表す。

\S\ 7 と同様に、整数多項式から成るさまざまな系を考える。

整数多項式の任意の系を整数系 $\Ssys_{n\rho}$ と呼ぶ。

整数線形族 $\Lsys_{n\rho}$ とは、$F_1(x)$ と $F_2(x)$ とともに、$[\Omega]$ の任意の整数 $c_1,c_2$ に対して $c_1F_1(x)+c_2F_2(x)$ も常に含む整数系である。

整数整性領域 $\Jdom_{n\rho}$ とは、$F(x)$ と $G(x)$ とともに $F(x)\cdot G(x)$ も常に含む整数線形族である。

整数相対整領域 $\Gdom_{n\rho}$ とは、$F(x)$ と $G(x)$ とともに、ただし $G(x)\ne0$ のとき、商 $F(x):G(x)$ を、それが整数多項式である場合に、かつその場合に限って常に含む整性領域である。

\S\ 7 と同様に、整数系 $\Ssys_{n\rho}$ を含む最小の整数整性領域 $\Jdom_{n\rho}$ は、$\Ssys_{n\rho}$ の有限個の多項式の整数係数の整有理結合であるすべての多項式 $H(x)$ からなる。

最小の整数相対整領域 $\Gdom_{n\rho}$ は、$\Ssys_{n\rho}$ の有限個の多項式の有理結合である整数多項式 $H(x)$ 全体からなる。

いま基底定理の転移へ進む。有理基底については、体、最小包含体、有理基底という「有理性」だけに基づく概念は何ら変化しないことに注意する。さらに、整数系 $\Ssys_{n\rho}$ の写像系 $\Ssys^*_{\rho\rho}$ は、任意に多くの仕方で再び整数系として選ぶことができる。整数線形族についても、生成元数の制限に導く議論は有効であるから、次を得る。

\emph{定理 V$'$: 任意の整数系 $\Ssys_{n\rho}$ は有理基底をもつ。整数線形族 $\Lsys_{n\rho}$ の有理基底は、特に高々 $\rho+1$ 個の多項式を含むように選ぶことができる。}

つぎに、多項式の整性領域のインボリューション基底（\S\ 8）の類似を調べる。そのため、\emph{量の行の対称関数}に関する定理を、\emph{整数上の整性}に関して次の補助定理によって拡張しなければならない。

\emph{補助定理. $n$ 行の量 $(yz\cdots t)$}
\begin{equation}
\begin{array}{cccc}
  y_1&z_1&\cdots&t_1\\
  y_2&z_2&\cdots&t_2\\
  \vdots&\vdots&&\vdots\\
  y_n&z_n&\cdots&t_n
\end{array}
\tag{1}
\end{equation}
\emph{の整数係数の整対称関数は、それらのうち有限個、すなわち基本対称関数}
\begin{equation}
\begin{gathered}
  \sigma_1(y),\sigma_2(y),\ldots,\sigma_n(y);\quad
  \sigma_1(z),\sigma_2(z),\ldots,\sigma_n(z);\quad\ldots;\\
  \sigma_1(t),\ldots,\sigma_n(t);
\end{gathered}
\tag{2}
\end{equation}
\emph{および関数}
\begin{equation}
  \chi_1(y,z,\ldots,t),\ldots,\chi_k(y,z,\ldots,t)
\tag{3}
\end{equation}
\emph{の整数係数の整関数である。ただし (3) の関数は (2) の関数上で代数的に整かつ整数的である。}

実際、$n$ 個の量 $y_1,\ldots,y_n$ は $\sigma_1(y),\ldots,\sigma_n(y)$ 上で代数的に整かつ整数上整であり、同じ主張が $z_1,\ldots,z_n$ と $\sigma_1(z),\ldots,\sigma_n(z)$ について、また以後同様に成り立つ。従って、(1) の量の行の対称関数の体は、(2) の関数によって定まる体上の代数体をなし、この体の代数的に整かつ整数的な量は、(1) の量の行の整数係数の整対称関数と一致する。従って (3) を Kronecker の基本系として取れば、補助定理が証明される。

いま $\Jdom_{n\rho}$ を整数整性領域、$\Jdom^*_{\rho\rho}$ を整数写像領域とし、
\[
  \Psi^*(z,u)=G^*(x)z^i+
  \sum{}'G^*_{\alpha\alpha_1\ldots\alpha_\rho}(x)
  z^\alpha u_1^{\alpha_1}\cdots u_\rho^{\alpha_\rho}
\]
を $\Jdom^*_{\rho\rho}$ のインボリューション形式とする。ここで $G^*(x)$ は零次元、すなわち $[\Omega]$ の数であってもよい。関数 (2) に対応する基本対称関数は、このとき商
\begin{equation}
  G^*_{\alpha\alpha_1\ldots\alpha_\rho}(x):G^*(x)
\tag{4}
\end{equation}
のうちのあるものによって与えられる。

補助定理により、関数 (3) は (2) 上で代数的に整かつ整数上整に依存するから、(3) に対応する $\Kfield^*_{\rho\rho}$ の関数は、代数的整数係数をもつ多項式に対する Gauss の定理の拡張により、
\begin{equation}
  K_\nu^*(x):(G^*(x))^\sigma
\tag{5}
\end{equation}
で与えられる。ここで $K_\nu^*(x)$ は $\Jdom^*_{\rho\rho}$ の整数多項式である。したがって、相対整領域の場合には $K_\nu^*(x)$ は領域に属し、一般には $\Jdom^*_{\rho\rho}$ の多項式の商である。

いま $F^*(x)$ を $\Jdom^*_{\rho\rho}$ の任意の整数多項式とすると、
\[
  F^*(x)=\frac1i\Bigl(F^*(x)+F^*(x^{(1)})+\cdots+F^*(x^{(i-1)})\Bigr)
\]
であるから、$i\cdot F^*(x)$ は量の行 $x,x^{(1)},\ldots,x^{(i-1)}$ の整数係数の整対称関数となる。

従って、補助定理と転移原理により、次を得る。

\emph{定理 VI$'$: 整数整性領域 $\Jdom_{n\rho}$ については、各インボリューション形式が特別な有理基底を定め、$\Jdom_{n\rho}$ の任意の多項式 $F(x)$ は表示}
\[
  F(x)=\frac{\Gamma(H(x),H_1(x),\ldots,H_s(x))}{i\cdot H(x)^i}
\]
\emph{をもつ。ここで $\Gamma$ は整数係数の整関数を表す。整数相対整領域 $\Gdom_{n\rho}$ で、写像領域が整数相対整領域 $\Gdom^*_{\rho\rho}$ である場合には、分母 $H(x)$ は対応するインボリューション形式の最高係数によって与えられる。}

\section*{\S\ 15. 第一種の整数相対整領域と整数正則系.}

整数系については、整性基底の代わりに整数整性基底が用いられる。すなわち、系の有限個の整数多項式であって、系の任意の整数多項式がそれら有限個の整数係数の整結合として表されるものをいう。

整数整性基底に関する定理は、本質的に先の定理と平行であり、証明だけがわずかに変更を要する。

まず、定義 V（\S\ 9）は直接に移される。

\emph{定義 V$'$: 量 $\xi$ が整数相対整領域 $\Gdom_{n\rho}$ に関して代数的に整かつ整数的であるとは、最高係数が $[\Omega]$ の単元で、残りの多項式が $\Gdom_{n\rho}$ に属するような方程式を少なくとも一つ満たすことをいう。}

代数的整数係数をもつ多項式に対する Gauss の定理の拡張により、\S\ 9 と同様、この定義は既約方程式から得られる。

定義 V$'$ から、さらに定義 VII の転移が得られる。

\emph{定義 VII$'$: 整数相対整領域 $\Gdom_{n\rho}$ が第一種であるとは、その写像領域として相対整領域 $\Gdom^*_{\rho\rho}$ をもち、$x_1,\ldots,x_\rho$ の任意の整数多項式が $\Gdom^*_{\rho\rho}$ 上で代数的に整かつ整数的であることをいう。}

これら第一種の領域の有限性を示すため、まず \S\ 10 と同じく、定義からインボリューション形式の最高係数が $[\Omega]$ の単元であることに注意する。

定理 VI$'$（\S\ 14）により、$\Gdom_{n\rho}$ の任意の多項式 $F(x)$ はしたがって表示
\[
  F(x)=\frac{\Gamma(H_1(x),\ldots,H_\sigma(x))}{i}
\]
をもつ。ここで、整数多項式 $\Gamma(H_1,\ldots,H_\sigma)$ からなる系に、Hilbert の整数加群基底に関する定理 II（Math. Ann. 36）を、もともとは整数有理数係数について述べられていたものを、代数的整数多項式への拡張として適用する。\footnote{$w_1,\ldots,w_k$ が $\Omega$ の代数的整数の基本系を表すとし、
\[
  \Gamma(H_1,\ldots,H_\sigma)=\Gamma_1(H)w_1+\cdots+\Gamma_k(H)w_k
\]
と置くならば、$\Gamma_1(H),\ldots,\Gamma_k(H)$ は整数有理数係数をもつ。形式 $v_1\Gamma_1(H)+\cdots+v_k\Gamma_k(H)$ からなる系に定理 II を適用すれば、この拡張の妥当性が直ちに示される。} これにより整数整性基底の存在が従う。すなわち、

\emph{定理 VII$'$: 第一種の任意の整数相対整領域は整数整性基底をもつ。}

さらに正則系の定義を移す。

\emph{定義 VIII$'$: 整数系 $\Ssys_{n\rho}$ が整数正則であるとは、最小包含整数整性領域の写像領域 $\Jdom^*_{\rho\rho}$ に、最高次元項の終結式が $[\Omega]$ の単元である、すなわち $R_0=\eps$ となる $\rho$ 個の多項式が存在することをいう。}

整性基底の存在証明を移すため、\S\ 11 の補助定理が次のより鋭い形を許すことに注意する。この形は、\S\ 11 の式 (5), (6) と、$A_1(\lambda),\ldots,A_\tau(\lambda)$ が $\lambda_1,\ldots,\lambda_\rho$ の整数係数の整関数であることから従う。

\emph{$x_1,\ldots,x_\rho$ の任意の整数多項式が、$\rho$ 個の整数多項式上で代数的に整かつ整数的に依存するための十分条件は、それらの多項式の最高次元項の終結式が $[\Omega]$ の単元、すなわち $R_0=\eps$ であることである。}

この補助定理から、\S\ 12 とまったく同じように、整数正則系の任意の多項式 $F(x)$ について \S\ 12 の表示 (3) が従う。ただし今は $\Gamma_1(G_i),\ldots,\Gamma_k(G_i)$ は $G_i$ の整数多項式である。さらに \S\ 12 と同様に、Hilbert の定理 II（Math. Ann. 36）を代数的整数多項式への拡張として適用し、最小包含整数整性領域の定義を用いれば、整数整性基底の存在が得られる。すなわち、

\emph{定理 IX$'$: 任意の整数正則系は整数整性基底をもつ。}

第一種の整数相対整領域の例として、\S\ 13 の例 1 と同様に、Lagrange の種領域の整数多項式全体 $\Gdom_{nn}$ を挙げる。$\Gdom_{nn}$ が整数の第一種であることは恒等式
\[
  x_k^n-\sigma_1(x)x_k^{n-1}+\cdots\pm\sigma_n(x)=0
  \qquad(k=1,2,
\ldots,n);
\]
から従う。また基本対称関数の終結式は値 $\pm1$ をもつので、$\Gdom_{nn}$ は整数正則系でもある。

さらにもう一つの例は、\S\ 14 の補助定理によって、$n$ 行の量の整数係数の整対称関数全体によって与えられる。

Erlangen, 1914年5月.

\section*{7. 有限群の不変式に対する有限性定理}
\begin{center}
Math. Ann. 77 (1916), pp. 89--92
\end{center}

以下では、有限群の不変式について、完全に初等的な有限性の証明を与える。この証明は対称関数の理論だけに基づくものであり、同時に完全系を実際に指定する。他方、Hilbert の加群基底定理（Ann. 36）に基づく通常の証明は、存在証明にすぎない。\footnote{例えば Weber, \emph{Lehrbuch der Algebra} (2nd ed.), vol. 2, \S\ 57 を参照。}

有限群 $\Hgrp$ は、行列式が 0 でない $h$ 個の線形変換 $A_1,\ldots,A_h$ から成るものとする。ここで $A_k$ は線形変換
\[
  x_1^{(k)}=\sum_{\nu=1}^{n}a_{1\nu}^{(k)}x_\nu,\ldots,
  x_n^{(k)}=\sum_{\nu=1}^{n}a_{n\nu}^{(k)}x_\nu,
\]
すなわち略記して $(x^{(k)})=A_k(x)$ を表す。したがって群 $\Hgrp$ は、成分 $x_1,
\ldots,x_n$ をもつ行 $(x)$ を、成分 $x_1^{(k)},
\ldots,x_n^{(k)}$ をもつ行 $(x^{(k)})$ へ移す。恒等変換は $A_1,
\ldots,A_h$ の中に含まれていなければならないから、行 $(x)$ 自身も行 $(x^{(k)})$ の中に含まれる。群の整有理（絶対）不変式とは、$x_1,
\ldots,x_n$ の整有理関数であって、$A_1,
\ldots,A_h$ の適用のもとで恒等的に変化しないものをいう。このような不変式 $f(x)$ に対しては、したがって
\begin{equation}
  f(x)=f(x^{(1)})=\cdots=f(x^{(h)})=
  \frac1h\sum_{k=1}^{h}f(x^{(k)}).
\tag{1}
\end{equation}
が成り立つ。

\subsection*{1.}
式 (1) は、$f(x)$ が数量行 $(x^{(1)}),
\ldots,(x^{(h)})$ の整有理対称関数であることをいう。しかも各項 $f(x^{(k)})$ はただ一つの行 $(x^{(k)})$ だけを含むから、これは通常\emph{一形式}の場合と呼ばれる最も単純な場合である。数量行の対称関数に関する既知の定理\footnote{2. への注を参照。}によって、$f(x)$ は、これらの行の基本対称関数、すなわち「Galois resolvent」
\[
\begin{aligned}
  \Phi(z,u)
  &=\prod_{k=1}^{h}\bigl(z+u_1x_1^{(k)}+\cdots+u_nx_n^{(k)}\bigr) \\
  &=z^h+
    \sum G_{\alpha\alpha_1\ldots\alpha_n}(x)
    z^{\alpha}u_1^{\alpha_1}\cdots u_n^{\alpha_n},
\end{aligned}
\qquad
\left(\begin{array}{c}
\alpha+\alpha_1+
\cdots+
\alpha_n=h\\
\alpha\ne h
\end{array}\right),
\]
の係数 $G_{\alpha\alpha_1\ldots\alpha_n}(x)$ によって整有理的に表される。ここで $G_{\alpha\alpha_1\ldots\alpha_n}(x)$ は、$x$ について次数 $\alpha_1+\cdots+\alpha_n$ の不変式である。したがって次が証明された。

\emph{Galois resolvent の係数 $G_{\alpha\alpha_1\ldots\alpha_n}(x)$ は群の不変式の完全系をなし、すべての不変式はこの有限個の不変式によって整有理的に表される。}

\subsection*{2.}
(1) から出発して、次の、なおいっそう初等的な考察を行うこともできる。\footnote{これは、上で述べた一形式の場合における数量行の対称関数の定理の証明に相当する。}同時にこれは第二の完全系へ導く。置く：
\[
  f(x)=a+b x_1^{\mu_1}\cdots x_n^{\mu_n}
       +\cdots+c x_1^{\nu_1}\cdots x_n^{\nu_n},
\]
ここで $a,b,\ldots,c$ は定数を表す。すると (1) から
\[
\begin{aligned}
 h\cdot f(x)
 &=h\cdot a+b\sum_{k=1}^{h}(x_1^{(k)})^{\mu_1}\cdots(x_n^{(k)})^{\mu_n}
     +\cdots
     +c\sum_{k=1}^{h}(x_1^{(k)})^{\nu_1}\cdots(x_n^{(k)})^{\nu_n} \\
 &=h\cdot a+b\cdot J_{\mu_1\ldots\mu_n}+\cdots+c\cdot J_{\nu_1\ldots\nu_n}
\end{aligned}
\]
を得る。したがって、すべての不変式は特殊な不変式
\[
  J_{\mu_1\ldots\mu_n}
  =\sum_{k=1}^{h}(x_1^{(k)})^{\mu_1}\cdots(x_n^{(k)})^{\mu_n}
\]
から整かつ線形に構成できるので、これら特殊な不変式について主張を証明すれば足りる。ところが、数値因子を別にすれば、$J_{\mu_1\ldots\mu_n}$ は式
\[
  S_\mu=\sum_{k=1}^{h}
  \bigl(u_1x_1^{(k)}+\cdots+u_nx_n^{(k)}\bigr)^{\mu}
\]
における $u_1^{\mu_1}
\cdots u_n^{\mu_n}$ の係数である。ただし $\mu=\mu_1+\cdots+\mu_n$ であり、この式は $h$ 個の線形形式
\[
  \xi_1=u_1x_1^{(1)}+\cdots+u_nx_n^{(1)},\ldots,
  \xi_h=u_1x_1^{(h)}+\cdots+u_nx_n^{(h)}
\]
の $\mu$ 番目の冪和である。無限に多くの冪和 $S_\mu$ は
\[
  S_1,S_2,\ldots,S_h
\]
の整有理結合であり、その係数は不変式
\[
  J_{\mu_1\ldots\mu_n}\qquad(\mu_1+\cdots+\mu_n\le h)
\]
によって与えられることが知られている。こうして第二の完全系が得られる。

\emph{群の不変式の完全系は、$\mu_1+\cdots+\mu_n\le h$ を満たすすべての不変式 $J_{\mu_1\ldots\mu_n}$ によって与えられる。ただし $h$ は群の位数を表す。}

1. との関係は、冪和 $S_1,
\ldots,S_h$ を $\xi_1,
\ldots,
\xi_h$ の基本対称関数で置き換えられることに注意すれば得られる。これにより、そこで与えた $G_{\alpha\alpha_1\ldots\alpha_n}(x)$ の完全系に至る。両方の結果から、すべての不変式は、$x$ に関する次数が群の位数 $h$ を超えないものによって整有理的に表されることが分かる。

\subsection*{3.}
いま証明したことから、有理表示に関する帰結が引き出される。任意の有理絶対不変式は、分子および分母に、群に関する分母の共役を掛ければただちに分かるように、必ずしも互いに素とは限らない二つの整有理絶対不変式の商として表せる。したがって、任意の有理絶対不変式は Galois resolvent $\Phi(z,u)$ の係数 $G_{\alpha\alpha_1\ldots\alpha_n}(x)$ によって有理的に表される。この定理は、有限性定理を経由しなくてもいくつかの方法で容易に証明できる。これはすでに Weber II, \S\ 58 に述べられている。\footnote{ただし、そこで与えられた証明は正しくない。式 (7) の関数 $\Psi(t)$ が不変式を係数にもつことは示されているが、それらの不変式が $\Phi(t)$ の係数によって有理的に表されることは示されていない。この欠陥は、$x_i^{(k)}$ を $\xi_k$ によって表す際に、Lagrange の補間公式ではなく既知の微分手続きを用いれば避けられる。これは、恒等式 $\Phi(-\xi_k,u)=0$ をすべての $u$ について微分して得られる関係である：
\[
 \left[\frac{\partial\Phi}{\partial u_i}
      -x_i^{(k)}\frac{\partial\Phi}{\partial z}\right]_{z=-\xi_k}=0.
\]
したがって、各有理関数 $\omega(x)$ に対しては、Weber の式 (7), (8) の代わりに次の式を置かなければならない：
\[
 \omega(x_1^{(k)},\ldots,x_n^{(k)})=
 \left.
 \omega\!\left(
   \frac{\partial\Phi/\partial u_1}{\partial\Phi/\partial z},\ldots,
   \frac{\partial\Phi/\partial u_n}{\partial\Phi/\partial z}
 \right)\right|_{z=-\xi_k},
\]
これは $\Phi(z,u)$ の係数だけを含み、不変式 $\omega(x)$ について全ての $k$ にわたり和を取れば、求める表示を与える。}

\subsection*{4.}
最後に、ここで得られた結果は、私の論文「有理関数の体と系」\footnote{Math. Ann. 76, p. 161 (1915).}に暗に含まれていることを述べておく。より正確には、1. で与えた完全系は定理 VI および VII に含まれており、この有限性定理に依存しない有理表示の証明は定理 III に含まれている。\footnote{相対不変式については、定理 VIII および IX に、通常のものとは異なる基礎にもとづく有限性の証明が含まれているが、これもまた存在証明にすぎない。その理由は、相対不変式が体をなさないからである。}しかし現在の特殊な場合には、一般理論に比べて証明の道筋も結果もかなり簡単になる。有限群の不変式へ一般的研究を適用することへ私を導いたのは E. Fischer との会話であり、実質的には、2. で与えた証明（一般の場合にはそこに現れる完全系は有理的に定義できない）および 3. への注も彼に由来する。

Erlangen, 1915 年 5 月。

\clearpage
\section*{8. 任意に多くの基本形式からなる系の不変式の整有理表示について}
\begin{center}
Math. Ann. 77 (1916), pp. 93--102
\end{center}

Schwarz 記念論文集への寄稿「任意に多くの基本形式からなる系の不変式について」の末尾で、D. Hilbert は次の予想を述べている。これらの不変式は、有限個の不変式からなる系 $(J,PJ)$ によって整有理的に表される、というものである。ここで $J$ は線形独立な基本形式の不変式の完全系を表し、不変式 $PJ$ は $J$ から極化過程によって生じる。

以下 I で、この予想が実際に正しいことを示す。それは、各々 $n$ 個の変数からなる任意に多くの行の形式は、固定された $n$ 行だけを含み、しかももとの形式から極化過程によって導かれる形式の極化の和として表される、という還元定理の単純な帰結である。この還元定理は、F. Mertens が初めて定式化した特殊な場合であり、A. Capelli、F. Mertens、J. Deruyts が与えた、Clebsch--Gordan の級数展開を各々 $n$ 変数からなる任意に多くの行の形式へ一般化したものの特殊な場合である。この一般化は、微分過程の形式的変換によって証明される。\footnote{F. Mertens: ``Über eine Formel der Determinantentheorie.'' (Formula 5), Sitzb. d. Ak. d. Wiss. Wien, vol. 91, II. Abt. (1885), さらに詳しく（応用を含めて）: ``Über ganze Funktionen von $m$ Systemen von je $n$ Unbestimmten''. Monatsh. f. Math. u. Phys. 4th year (1893). Capelli のいくつかの証明（1882 年以降）については、例えば Math. Ann. 37、または自筆版 ``Lezioni sulla teoria delle forme algebriche'' (1902) を参照。J. Deruyts: \emph{Essai d'une théorie générale des formes algébriques}, Mém. d. 1. Soc. Roy. des Sciences de Liège (2) 17 (1892).}

II では、還元定理を、形式の線形族の同値性に関する問題として扱うことにより、より概念的な証明も与える。この見方、また用いられる考察の本質的な一部は、「代数的加群系について」\footnote{E. Fischer: Über algebraische Modulsysteme und lineare homogene partielle Differentialgleichungen mit konstanten Koeffizienten, \S\ 1--4, J. f. M. 140 (1911).}という論文、およびそれに続く E. Fischer の未公刊の定理から取ったものであり、彼はその使用を好意的に許してくれた。

最後に III では、対応する考察が最も一般的な級数展開にも導くことを示す。

\subsection*{I.}
$\theta$ を、$N>n$ 個の各行
\[
  A_1,A_2,\ldots,A_N
\]
に関して斉次な形式とする。一般に行 $A_k$ は $n$ 個の成分
\[
  A_k^{(1)},A_k^{(2)},\ldots,A_k^{(n)}
\]
から成る。$\theta$ の係数は不定元でもよいし、与えられた有理性領域の量でもよい。さらに $P$ は、極化過程
\[
  P_{hk}=\left(A_h\frac{\partial}{\partial A_k}\right)
  =A_h^{(1)}\frac{\partial}{\partial A_k^{(1)}}+
   A_h^{(2)}\frac{\partial}{\partial A_k^{(2)}}+
   \cdots+
   A_h^{(n)}\frac{\partial}{\partial A_k^{(n)}}
\]
の、整数係数をもつ整有理関数を表すものとする。積 $P_{hk}P_{ij}\theta$ は、$P_{ij}\theta$ に $P_{hk}$ を適用することとして理解される。

いま述べた還元定理は、恒等式
\begin{equation}
  \theta=\sum PZ,
\tag{1}
\end{equation}
によって与えられる。ここで形式 $Z$ は行
\[
  A_1,A_2,\ldots,A_n
\]
だけを含み、$\theta$ から極化過程 $P$ によって導かれる。

次に、同じ位数かつ同じ変数の数をもつ $N$ 個の形式
\[
  F_1,F_2,\ldots,F_N
\]
から成る基本形式の系 $(F')$ を考える。その係数をそれぞれ $N$ 個の行
\[
  A_1,A_2,\ldots,A_N
\]
として取る。$F_1,
\ldots,F_n$ の、有限個の不変式から成る完全系を $(J)$ と書き、$F_1,
\ldots,F_n$ の任意の不変式が $(J)$ の不変式によって整有理的に表されるものとする。証明すべき Hilbert の予想は、$F_1,
\ldots,F_N$ の同時不変式 $S$ に対し、完全系が有限個の不変式 $(J,PJ)$ によって与えられるというものである。ここで $PJ$ は、極化過程 $P$ によって $J$ から導かれるすべての不変式を表す。

実際、$(F')$ の有理整数係数をもつ任意の同時不変式 $S$ は、$N$ 個の各行 $A_1,
\ldots,A_N$ に関して斉次な形式であり、係数は有理整数であるから、恒等式 (1) をこれに適用できる。極化過程 $P$ を $S$ に適用しても再び不変式が得られることは知られているので、形式 $Z$ も不変式である。すなわち、それらは行 $A_1,
\ldots,A_n$ だけを含むから、$F_1,
\ldots,F_n$ の同時不変式である。したがって仮定により、それらは不変式 $(J)$ の整有理関数である。よって恒等式 (1) は
\[
  S=\sum PY
\]
となる。ここで $Y$ は不変式 $(J)$ の冪積を表す。ところが $P$ の定義により、$Y$ に対する過程 $P$ の実際の実行は、単純な極化操作 $P_{hk}$ を有限回反復して逐次実行することにほかならず、積の微分法則によって、不変式 $(J)$ と $(PJ)$ の積の和へ導く。同じことは、$(J)$ と $(PJ)$ からなる積に $P_{hk}$ を適用する場合にも成り立つ。したがって $PY$ も $(J,PJ)$ の整有理結合だけを与える。\emph{これによって、すべての不変式 $S$ が $(J,PJ)$ によって整有理的に表されることが証明された。}

\subsection*{II.}
還元定理を証明する前に、形式の線形族についていくつかの事実を思い出す。$K$ 上の\emph{形式の線形族}とは、同じ次元をもつ形式の系であり、$\theta_1$ と $\theta_2$ を含めば常に $c_1\theta_1+c_2\theta_2$ も含むという意味で完全なものをいう。ここで $c_1,c_2$ は所定の有理性領域 $K$ の量であり、$K$ は $\theta$ たちの係数領域を含むものとする。これらの形式の中には、常に有限個、たとえば $\rho$ 個、$K$ 上線形独立なものがあり、任意の $\rho+1$ 個の形式は $K$ からの係数をもつ線形関係を満たす。この族は $K$ に関して\emph{階数 $\rho$} をもつ。\footnote{より一般の形式の線形族、すなわち階数が有限である必要のないものは、同じ次元という制限を外すか、または $K$ が $\theta$ たちの係数領域を含むという制限を外すことによって得られる。}以下では、$\mathcal L$ と $\mathcal T$ が $K$ 上の同じ階数 $\rho$ の二つの線形族で、$\mathcal T$ が $\mathcal L$ の部分族であるなら、$\mathcal L$ と $\mathcal T$ は同一（同値）である、という容易に分かる事実を用いる。

これを準備として、還元定理へ進む。I の恒等式
\[
  \theta=\sum PZ
\]
は、同じ極化過程 $P$ を両辺に適用すると、$P\theta$ についての類似の恒等式へ移る。したがって還元定理は、同時に、$\theta$ のすべての極化を特殊な極化 $PZ$ の和で表示することを与える。逆に、これら特殊な極化はもちろんすべての極化の中に含まれるから、還元定理はこれら二つの線形族の同値性を主張している。とくに、各々の行において $\theta$ と同じ次元をもつ形式からなる部分族の同値性も主張している。この同値性を証明するために、形式 $Z$ によって定まる中間の族を挿入する。簡単のため、まず三つの二元行の場合 $(N=3,n=2)$ について証明を与え、その後、完全に平行な一般証明を短く示す。

そこで、$\theta(x,y,z)$ がそれぞれ行
\[
  x_1x_2;\quad y_1y_2;\quad z_1z_2;
\]
について次元 $\alpha,\beta,p$ をもつとする。まず $\theta$ の係数は不定元 $a$（または有理数）とする。次の三つの形式の線形族を考える。

1. 族 $\mathcal L$ は、$\theta$ から極化過程 $P$ によって生じ、個々の行 $x,y,z$ において $\theta$ と同じ次元をもつすべての形式 $f(x,y,z)$ から成る。とくに $\mathcal L$ は $\theta$ を含む。$f(x,y,z)$ を $x,y,z$ および不定元 $a$ の形式と見るなら、係数体は有理数体 $R$ である。$K$ としても $R$ を取らねばならない。なぜなら、有理整数の $\theta_1,\theta_2$ についてだけ $c_1P_1+c_2P_2$ が再び過程 $P$ になるからである。$\mathcal L$ の $R$ に関する階数を $\rho$ とする。

2. 族 $\mathcal S$ は、
\[
  g(\lambda;x,y)=f(x,y,\lambda_1x+\lambda_2y)
\]
によって定義されるすべての形式 $g(\lambda;x,y)$ から成る。あるいは極化過程で表せば、
\[
\begin{aligned}
  g(\lambda;x,y)
  &=\frac1{p!}\left[(\lambda_1x+\lambda_2y)\frac{\partial}{\partial z}\right]^p f(x,y,z) \\
  &=g_0(xy)\lambda_1^p+g_1(x,y)\lambda_1^{p-1}\lambda_2+\cdots+g_p(xy)\lambda_2^p,
\end{aligned}
\]
ただし
\[
  i!(p-i)!\,g_i(xy)=
  \left(y\frac{\partial}{\partial z}\right)^i
  \left(x\frac{\partial}{\partial z}\right)^{p-i}f(xyz).
\]
\footnote{I と同じく、
\[
 \left(t\frac{\partial}{\partial x}\right)=t_1\frac{\partial}{\partial x_1}+t_2\frac{\partial}{\partial x_2},
\]
したがってとくに
\[
 \left[(\lambda_1x+\lambda_2y)\frac{\partial}{\partial z}\right]
 = (\lambda_1x_1+\lambda_2y_1)\frac{\partial}{\partial z_1}
 + (\lambda_1x_2+\lambda_2y_2)\frac{\partial}{\partial z_2},
\]
\[
 \left[z\left(\frac{\partial}{\partial\lambda_1}\frac{\partial}{\partial x}
 +\frac{\partial}{\partial\lambda_2}\frac{\partial}{\partial y}\right)\right]
 =z_1\left(\frac{\partial}{\partial\lambda_1}\frac{\partial}{\partial x_1}
 +\frac{\partial}{\partial\lambda_2}\frac{\partial}{\partial y_1}\right)
 +z_2\left(\frac{\partial}{\partial\lambda_1}\frac{\partial}{\partial x_2}
 +\frac{\partial}{\partial\lambda_2}\frac{\partial}{\partial y_2}\right).
\]}
したがって $g_i(xy)$ は恒等式 (1) の形式 $Z$ を与える。$g$ は $x,y,\lambda$ および不定元 $a$ に関する有理整数形式である。$\mathcal S$ の $R$ に関する階数を $\sigma$ とする。

3. 族 $\mathcal T$ は、$\mathcal S$ が $\mathcal L$ から得られるのと同じ仕方で、逆の極化過程によって $\mathcal S$ から得られる。$\mathcal T$ の形式 $h$ は
\[
\begin{aligned}
 h(x,y,z)
 &=\frac1{p!}\left[z\left(\frac{\partial}{\partial\lambda_1}\frac{\partial}{\partial x}
 +\frac{\partial}{\partial\lambda_2}\frac{\partial}{\partial y}\right)\right]^p g(\lambda;xy)\\
 &=h_0(xyz)+h_1(xyz)+\cdots+h_p(xyz)
\end{aligned}
\]
により定義される。ここで 2. により
\[
  h_i(xyz)=
  \left(z\frac{\partial}{\partial x}\right)^{p-i}
  \left(z\frac{\partial}{\partial y}\right)^i g_i(xy).
\]
個々の $h_i$、したがって $h$ も、形式 $PZ$ およびその和である。$\mathcal T$ の $R$ に関する階数を $\tau$ とする。

$\mathcal T$ の形式 $h(x,y,z)$ の構成が示すように、これらの形式は $\theta$ から極化過程 $P$ によって生じ、行 $x,y,z$ において個々に $\theta$ と同じ次元をもつ。したがって族 $\mathcal T$ は $\mathcal L$ の部分族である。上で述べた事実により、残るのは階数の等しさを証明することだけである。この証明では、$f(xyz)$ の特殊な構造、すなわちそれらが同一の形式 $\theta$ の極化であることは、もはや用いない。

\emph{a) $\mathcal L$ と $\mathcal S$ の階数を比較するために、補題 a) を用いる：$f(x,y,\lambda_1x+\lambda_2y)=0$ からは、必然的に $f(x,y,z)=0$ が従う。}これは、三つの二元行のあいだの恒等関係
\[
  (xy)z+(yz)x+(zx)y=0,
  \qquad (xy)=(x_1y_2-x_2y_1),\ldots,
\]
からただちに従う。この関係は
\[
  f[x,y,(xy)x+(xz)y]=(xy)^p f(x,y,z)
\]
を含意する。したがって $f(x,y,\lambda_1x+\lambda_2y)=0$（$\lambda$ について恒等的）から、置換 $\lambda_1=(zy)$, $\lambda_2=(xz)$ により、かつ $(xy)\ne0$ であることから、必然的に $f(x,y,z)=0$ が従う。

この補題により、$\mathcal S$ の形式と $\mathcal L$ の形式との間には一対一対応がある。実際、
\[
  f_1(x,y,\lambda_1x+\lambda_2y)=g(\lambda;xy),\qquad
  f_2(x,y,\lambda_1x+\lambda_2y)=g(\lambda;xy)
\]
からは、必然的に
\[
  f_1(xyz)=f_2(xyz)
\]
が従う。さらに、$\mathcal S$ における任意の関係は、$\mathcal L$ における一意に対応する形式のあいだにも保たれる（また逆も成り立つ）。なぜなら、
\[
  c_1g^{(1)}(\lambda;xy)+c_2g^{(2)}(\lambda;xy)+\cdots+c_rg^{(r)}(\lambda;xy)=0
\]
から、補題により必然的に
\[
  c_1f^{(1)}(x,y,z)+c_2f^{(2)}(x,y,z)+\cdots+c_rf^{(r)}(x,y,z)=0
\]
が従うからである。\emph{これにより $\mathcal L$ と $\mathcal S$ の階数の等しさ、すなわち $\rho=\sigma$ が証明された。}

\emph{b) $\mathcal S$ と $\mathcal T$ の階数を比較するために、対応する補題 b) を用いる：$h(x,y,z)=0$ からは、必然的に $g(\lambda;xy)=0$ が従う。}証明のため、3. により $h(x,y,z)=0$ は、個々の冪積
\[
 \left(\frac{\partial}{\partial\lambda_1}\frac{\partial}{\partial x_1}
 +\frac{\partial}{\partial\lambda_2}\frac{\partial}{\partial y_1}\right)^{p_1}
 \left(\frac{\partial}{\partial\lambda_1}\frac{\partial}{\partial x_2}
 +\frac{\partial}{\partial\lambda_2}\frac{\partial}{\partial y_2}\right)^{p_2}g(\lambda;x,y),
 \qquad p_1+p_2=p
\]
の消滅と同値であることに注意する。さらに微分すれば、冪積
\[
\begin{aligned}
&\left(\frac{\partial}{\partial x_1}\right)^{\alpha_1}
 \left(\frac{\partial}{\partial x_2}\right)^{\alpha_2}
 \left(\frac{\partial}{\partial y_1}\right)^{\beta_1}
 \left(\frac{\partial}{\partial y_2}\right)^{\beta_2} \\
&\quad\cdot
 \left(\frac{\partial}{\partial\lambda_1}\frac{\partial}{\partial x_1}
 +\frac{\partial}{\partial\lambda_2}\frac{\partial}{\partial y_1}\right)^{p_1}
 \left(\frac{\partial}{\partial\lambda_1}\frac{\partial}{\partial x_2}
 +\frac{\partial}{\partial\lambda_2}\frac{\partial}{\partial y_2}\right)^{p_2}
 g(\lambda;xy)
\end{aligned}
\]
もまた消滅する。したがって、これらの冪積の任意の線形結合が消滅し、それゆえ任意の形式
\[
  f\left(\frac{\partial}{\partial x},\frac{\partial}{\partial y},
  \frac{\partial}{\partial\lambda_1}\frac{\partial}{\partial x}
  +\frac{\partial}{\partial\lambda_2}\frac{\partial}{\partial y}\right)g(\lambda,xy)
\]
も消滅する。とくに
\[
  f(x,y,\lambda_1x+\lambda_2y)=g(\lambda;xy)
\]
となるように $f$ を選べば、さらに
\[
  g\left(\frac{\partial}{\partial x};
    \frac{\partial}{\partial x},
    \frac{\partial}{\partial y}\right)g(\lambda,x,y)=0
\]
を得る。あるいは、もし
\[
  g(\lambda,xy)=\sum G_i\lambda_1^{\nu_1}\lambda_2^{\nu_2}
  x_1^{\alpha_1}x_2^{\alpha_2}y_1^{\beta_1}y_2^{\beta_2}
\]
で、$G_i$ が不定元 $a$ に関する有理整数線形形式であるなら、
\[
  \sum \nu_1!\nu_2!\alpha_1!\alpha_2!\beta_1!\beta_2!\,G_i^2=0,
\]
したがって $g(\lambda;xy)=0$ となる。

補題 b) から、a) とまったく同様に、\emph{$\mathcal S$ と $\mathcal T$ の階数は等しい；$\sigma=\tau$}ことが従う。

したがって a) と b) により $\rho=\tau$ であり、従って、$\mathcal T$ は $\mathcal L$ の部分族であるから、これら二つの線形族の同値性が証明された。よって $\mathcal L$ の任意の形式、とくに $\theta$ は、$\rho$ 個の線形独立な形式 $h(x,y,z)$ の有理整数線形結合として表せる：
\[
  \theta=\sum c_i h^{(i)}(x,y,z)=\sum PZ.
\]
$\theta$ の不定係数 $a$ に対して導かれたこの恒等式は、その不定元 $a$ を任意の有理性領域の量で置き換えても有効である。したがって、\emph{三つの二元行の場合に対して、還元定理は一般に証明された。}

各々 $n$ 変数からなる $N$ 行の場合の一般証明は、完全に平行である。$\theta(A)$ が行 $A_i$ に関して次元 $\alpha_i$ をもつとする。再び次の三つの形式の線形族を考える。

1. 族 $\mathcal L$ は、$\theta(A)$ から極化過程 $P$ によって生じ、各行 $A_i$ において $\theta(A)$ と同じ次元をもつすべての形式 $f(A)$ から成る。

2. 族 $\mathcal S$ は、置換
\[
\begin{aligned}
 A_{n+1}&=\lambda_{n+1}^{(1)}A_1+\lambda_{n+1}^{(2)}A_2+
        \cdots+\lambda_{n+1}^{(n)}A_n,\\
 &\vdotswithin{=}\\
 A_N&=\lambda_N^{(1)}A_1+\lambda_N^{(2)}A_2+
        \cdots+\lambda_N^{(n)}A_n;
\end{aligned}
\]
によって $f(A)$ から生じるすべての形式 $g(\lambda;A_1\cdots A_n)$ から成る。この過程で、$g(\lambda;A_1\cdots A_n)$ における $\lambda$ の個々の冪積の係数は、恒等式 (1) の形式 $Z$ を与える。

3. 族 $\mathcal T$ は、次で定義されるすべての形式 $h(A)$ から成る：
\[
\begin{aligned}
 h(A)=&\left[A_{n+1}\left(
 \frac{\partial}{\partial\lambda_{n+1}^{(1)}}\frac{\partial}{\partial A_1}
 +\cdots+
 \frac{\partial}{\partial\lambda_{n+1}^{(n)}}\frac{\partial}{\partial A_n}\right)\right]^{\alpha_{n+1}} \\
&\cdots
\left[A_N\left(
 \frac{\partial}{\partial\lambda_N^{(1)}}\frac{\partial}{\partial A_1}
 +\cdots+
 \frac{\partial}{\partial\lambda_N^{(n)}}\frac{\partial}{\partial A_n}\right)\right]^{\alpha_N}
 g(\lambda;A_1\cdots A_n).
\end{aligned}
\]
ここで $h(A)$ は形式 $PZ$ の和である。

任意の $n+1$ 行のあいだの恒等関係により、補題 a) は有効なままである。補題 b) も同様に有効である。なぜなら、その証明には変数の数が入らなかったからである。したがって $\mathcal L$ と $\mathcal T$ は同じ階数をもち、かつ $\mathcal T$ は $\mathcal L$ の部分族であるから、両者は同一である。従って恒等式
\[
  \theta=\sum PZ
\]
は不定係数について成り立ち、ゆえに任意の有理性領域の量である係数についても成り立つ。\emph{これにより還元定理は一般に証明された。}

\subsection*{III.}
同様の考察が、\emph{一般級数展開}（$N\le n$ の場合）も与えることを簡単に示す。$Z$ を、各々 $n$ 個の量からなる $n$ 行 $A_1,
\ldots,A_n$ において別々に斉次な形式とし、$\Delta$ をこれらの行の行列式とする。まず (1) に対応する恒等式を証明する：
\begin{equation}
  Z=\sum PH \pmod{\Delta},
\tag{2}
\end{equation}
ここで形式 $H$ は行 $A_1,
\ldots,A_{n-1}$ だけを含み、$Z$ から過程 $P$ によって導かれる。

証明は、II で得られた関係を $\Delta$ を法とする合同式へ移すことから成る。簡単のため、三つの三元行 $x,y,z$ を取る。係数は不定元であるとする。

三つの線形族 $\mathcal L,\mathcal S,\mathcal T$ は、二元行の場合の II と同じように定義する。その階数を $\rho,\sigma,\tau$ とし、$\rho^*$ および $\tau^*$ は $\Delta$ を法とする $\mathcal L$ および $\mathcal T$ の階数を表すものとする（すなわち、$\Delta$ を法として線形独立な $\mathcal L$ の形式は $\rho^*$ 個存在するが、$\rho^*+1$ 個は存在しない）。すると II と同様に、次が成り立つ。

\emph{補題 a)：$f(x,y,\lambda_1x+\lambda_2y)=0$ からは、必然的に $f(xyz)\equiv0\pmod{\Delta}$ が従う（また逆も成り立つ）。}実際、四つの三元行のあいだの関係
\[
  \Delta_1x-\Delta_2y+\Delta_3z=\Delta t,
  \qquad \Delta_1=(yzt),\ldots,
\]
により、三つの三元行のあいだには常に $\Delta$ を法とする線形関係
\[
  \Delta_1x-\Delta_2y+\Delta_3z\equiv0\pmod{\Delta}
\]
がある。したがって II と同様に
\[
  f(x,y,-\Delta_1x+\Delta_2y)=\Delta_3^p f(x,y,z)\pmod{\Delta}
\]
を得る。$f(x,y,\lambda_1x+\lambda_2y)=0$（$\lambda$ について恒等的）から、$\Delta$ は既約であり $\Delta_3$ は $\Delta$ で割り切れないので、必然的に $f(x,y,z)\equiv0\pmod{\Delta}$ が従う。逆は $\Delta(x,y,\lambda_1x+\lambda_2y)=0$ から明らかである。したがって II と同様に $\rho^*=\sigma$ である。

補題 b) はその証明とともに有効であり、したがって $\sigma=\tau$ である。

$\tau=\tau^*$ を証明するために次を用いる。

\emph{補題 c)：$h(x,y,z)\equiv0\pmod{\Delta}$ からは、必然的に $h(x,y,z)=0$ が従う。}すなわち、極化として $h(x,y,z)$ は既知の $\Omega$ 過程の適用で消滅する。これは微分方程式
\[
  \Delta\!\left(\frac{\partial}{\partial x},\frac{\partial}{\partial y},\frac{\partial}{\partial z}\right)h(x,y,z)=0
\]
で表される。いま $h(x,y,z)=\Delta(x,y,z)\cdot k(x,y,z)$ とすれば、さらに微分することにより、
\[
  k\!\left(\frac{\partial}{\partial x},\frac{\partial}{\partial y},\frac{\partial}{\partial z}\right)
  \Delta\!\left(\frac{\partial}{\partial x},\frac{\partial}{\partial y},\frac{\partial}{\partial z}\right)
  =
  h\!\left(\frac{\partial}{\partial x},\frac{\partial}{\partial y},\frac{\partial}{\partial z}\right)h(x,y,z)
\]
の消滅も得られ、したがって $h(x,y,z)$ の消滅も得られる。よって $\rho^*=\sigma=\tau=\tau^*$ であり、$\mathcal T$ は $\mathcal L$ の部分族であるから、恒等式 (2) が証明された。同じ議論により $n$ 変数の場合の証明も得られる。\footnote{3. の後には、$\Omega(h)$ について、行列式因子のために消滅する対応する関係がある。作用素記法では、それは
\[
 \left(\lambda_1\frac{\partial}{\partial \xi}+\lambda_2\frac{\partial}{\partial \eta}\right),
 \qquad
 \left[z\left(\frac{\partial}{\partial\lambda_1}\frac{\partial}{\partial \xi}
 +\frac{\partial}{\partial\lambda_2}\frac{\partial}{\partial \eta}\right)\right]
\]
の積から生じる方程式であり、行列式因子が消えるために消滅する。}

恒等式 (2) を $\Delta$ の乗数に適用すると、展開
\[
  Z=\varphi_0+\Delta\varphi_1+\Delta^2\varphi_2+\cdots+\Delta^\mu\varphi_\mu
\]
が得られる。ここで各 $\varphi_i$ は $\Omega$ 過程の適用で消滅する。この過程を $i$ 回繰り返すと、一般に $\Omega(\Delta\cdot\psi)=c_1\psi+c_2\Delta\cdot\Omega(\psi)$ であるから、
\[
  c\cdot \Omega^i(Z)\equiv\varphi_i\pmod{\Delta};
\]
他方、(2) により
\[
  c\cdot\Omega^i(Z)\equiv\sum PH_i\pmod{\Delta}
\]
である。ここで $H_i$ は、$H$ が $Z$ から導かれるのと同じ仕方で、形式 $\Omega^i(Z)$ から導かれる。したがって補題 c)（$n$ 変数へ拡張したもの）により $\varphi_i=\sum PH_i$ を得る。従って
\begin{equation}
  Z=\sum PH+\Delta\cdot\sum PH_1+\Delta^2\cdot\sum PH_2+
  \cdots+\Delta^\mu\cdot\sum PH_\mu
\tag{3}
\end{equation}
\emph{これは、各々 $n$ 変数からなる $n$ 行の形式の、ただ $n-1$ 行の形式の極化による最も一般的な級数展開である。}

$N$ 行の $n$ 変数の形式 $(N<n)$ に対する級数展開は、よく知られているように、単純な置換によって得られる。$N=3$, $n>3$ に対し
\[
  u=\xi_1x+\xi_2y+\xi_3z,
  \quad v=\eta_1x+\eta_2y+\eta_3z,
  \quad w=\zeta_1x+\zeta_2y+\zeta_3z
\]
と置き、$H(u,v,w)=Z(\xi,\eta,\zeta)$ を (3) に従って展開する。すると $\xi,\eta,\zeta$ は組合せ $u,v,w$ においてだけ現れるので、
\[
  \left(\eta\frac{\partial}{\partial \xi}\right)=
  \left(v\frac{\partial}{\partial u}\right)\cdots,
  \qquad
  \Omega=\sum_{ikl}\left(\frac{\partial}{\partial u}
  \frac{\partial}{\partial v}
  \frac{\partial}{\partial w}\right)_{ikl}(xyz)_{ikl}
  =\nabla_{uvw}(xyz).
\]
特殊化 $\xi_1=\eta_2=\zeta_3=1$; $\xi_2=\xi_3=\cdots=\zeta_2=0$ のもとで、$H(u,v,w)$ は $H(x,y,z)$ へ移り、$\nabla_{uvw}(xyz)$ は、数値因子を除いて $\nabla_{xyz}(xyz)=\nabla_{xyz}$ へ移る。なぜなら、$(y\partial/\partial x)$ などを行列式 $(xyz)_{ikl}$ に適用するとそれらは消滅するからである。対応する主張は $N$ 行についても成り立つので、(3) は展開
\begin{equation}
  H=\sum P\Phi+\sum P\Phi_1+\cdots+\sum P\Phi_\mu
\tag{4}
\end{equation}
を与える。ここで $\Phi_i$ は、$\nabla^i_{A_1\ldots A_N}H$ から過程 $P$ によって生じ、行 $A_1,
\ldots,A_{N-1}$ だけを含む。ただし $\nabla^i$ に現れる行列式の組合せは別であり、それらはすでに述べたように、極化の際には考慮に入らない。

次に、$\nabla^i$ の中のこれらの行列式を不定元 $p$ で置き換えれば、得られる形式は再び (4) に従って展開できる。このようにして最終的に展開
\begin{equation}
  H=\left[\sum P\psi(p)\right]_{p=\text{$A_1\cdots A_N$ の行列式}},
\tag{5}
\end{equation}
が得られる。ここで $\psi$ は、もとの形式から過程 $P$ および $\nabla_{A_1\ldots A_N}(p)$, $\nabla_{A_1\ldots A_{N-1}}(p)$ などによって生じる。これらの展開およびその組合せ、例えば (1) に現れる形式 $Z$ に対して (3)、ついで (4) と (5) を適用することにより、各々 $n$ 個の共変な変数からなる任意に多くの行の形式について既知のすべての展開が尽くされる。

Erlangen, 1915 年 1 月 5 日。

\clearpage
\section*{9. 整超越数から成る最も一般の領域}
\begin{center}
Math. Ann. 77 (1916), pp. 103--128
\end{center}

E. Zermelo は整列定理を用いて「整超越数」の領域、すなわち次の抽象的性質をもつ数の領域 $\mathfrak G$ を構成した。
\begin{enumerate}
\item[I.] $\mathfrak G$ の二つの数の和、差、積は再び $\mathfrak G$ の数である。
\item[II.] すべての実数または複素数は、$\mathfrak G$ の二つの数の商である。
\item[III.] すべての有理整数（それぞれ代数的整数）は $\mathfrak G$ の元である。
\item[IV.] 整でない有理数（それぞれ整でない代数的数）は $\mathfrak G$ に属さない。
\footnote{E. Zermelo, Über ganze transzendente Zahlen, Math. Ann. 75, p. 434 (1914).}
\end{enumerate}

この構成は、整列定理が全ての数の\emph{代数的基底}の存在を含意するという事実に基づいている。すなわち、相互の間に代数関係が存在せず、しかもそれらによって他の全ての数を代数的に表せるような\emph{数} $\eta$ \emph{の系} $H$ が存在する。この基底数 $\eta$ は代数的に独立であるため、不定元とみなすことができる。したがって、求める領域は、不定元 $\eta$ の有理関数および代数関数から成る整域となり、\footnote{この理由から、本論文の考察は一部、著者の以前の論文 Körper und Systeme rationaler Funktionen, Math. Ann. 76, p. 161 (1915) の考察と平行している。}その係数は全ての代数的数から成る体 $K$ に属する。Zermelo によって構成された特殊な領域 $\mathfrak G_\eta$ は、次の基底性質によって特徴づけられる。

$\mathfrak G_\eta$ は次を含む。
\begin{enumerate}
\item[1)] すべての基底数 $\eta$。
\item[2)] $\eta$ のすべての多項式で、係数が整であるもの。\footnote{本論文を通じて「整」とは、係数が全ての代数的整数から成る整域 $[K]$ に属することを意味する。$\mathfrak G_\eta$ の定義については、1)--5) で有理整数係数だけを考えても十分である。しかし、より一般の領域については、その場合、基底性質がより単純でなくなる。}
\item[3)] $\eta$ のすべての代数的整関数で、係数が整であるもの。
\end{enumerate}
$\mathfrak G_\eta$ は次を排除する。
\begin{enumerate}
\item[4)] $\eta$ のすべての有理的に分数的な関数で、係数が整であるもの。
\item[5)] $\eta$ のすべての代数的に分数的な関数で、係数が整であるもの。\footnote{Zermelo, loc. cit., § 3.}
\end{enumerate}

いま、§§ 2, 3 の例と比較すれば、Zermelo 領域 $\mathfrak G_\eta$ と本質的に異なる領域 $\mathfrak G$ が存在することは容易に分かる。すなわち、どのような整列のもとでも基底性質 1)--5) がすべて成り立つことはなく、したがってどのような整列のもとでも Zermelo の構成によって生成することのできない領域 $\mathfrak G$ が存在する。言い換えると、$\mathfrak G_\eta$ の上へ一対一かつ同型的に写すことのできない領域 $\mathfrak G$ が存在する。そこで、整超越数から成る最も一般の領域という問題が生じる。この問題は以下で完全に答えられる。

そのためにはまず、基底性質のうちどれが抽象的な定義性質の帰結であり、どれが特殊な構成に由来するのかを決定しなければならない。任意に与えられた領域 $\mathfrak G$ に対して、基底性質 1), 2) が満たされるような整列が存在することを示す（§ 1）。したがって、すべての領域 $\mathfrak G$ は全ての数の代数的基底によって構成できる。これ以外の基底性質はどれも満たされる必要はない（§§ 2, 3）。特に、Zermelo 領域のさらに一つの抽象的性質、すなわち代数的整閉性はすべての領域 $\mathfrak G$ に属するわけではないことが従う。これは次のように定式化できる。
\begin{enumerate}
\item[V.] $\mathfrak G$ の有限個の数に代数的に整かつ整に依存するすべての数は $\mathfrak G$ に属する。
\end{enumerate}

I--IV に加えて V が成り立つ特殊な領域を、代数的に整な超越数から成る領域 $\mathfrak H$ と呼ぶ。最も一般の領域 $\mathfrak H$ については、次の単純な結果を得る（§ 4）。
\begin{enumerate}
\item[a)] 同じ基底 $H$ によって構成できるすべての領域 $\mathfrak H$ の交わりは、基底のすべての代数的整関数、すなわち Zermelo 領域 $\mathfrak G_\eta$ によって与えられる。この $\mathfrak G_\eta$ 自身も領域 $\mathfrak H$ である。（これは Zermelo 領域の抽象的特徴づけも与える。）
\item[b)] すべての領域 $\mathfrak H$ は、任意の系 $\mathfrak S(\eta)$ を $\mathfrak G_\eta$ に代数的に整に添加することによって生じる。\footnote{この表現の説明は § 4 にある。} ただし $\mathfrak S(\eta)$ は、添加によって代数的に分数的な数が領域に入らないというただ一つの条件を満たせばよい。
\end{enumerate}

代数的基底を基礎にとるとき、最も一般の領域 $\mathfrak G$ に対する結果はこれほど単純ではない。この場合、すべての領域 $\mathfrak G$ の交わりはそれ自身が領域 $\mathfrak G$ ではなく、したがって最も一般の領域の構成は見通しにくい（§ 5）。

領域 $\mathfrak H$ の場合と類似の結果を得るために、代数的基底を\emph{有理基底} $\Theta$ で置き換える。これは、すべての数を有理的に表すことを可能にし、かつ少なくとも一つの整列のもとでどの基底数も有限個の先行するものによって有理的に表せないような\emph{数} $\vartheta$ \emph{の系} $\Theta$ である。この有理基底 $\Theta$ には、さらに、$\vartheta$ のすべての整多項式から成る整域 $[\Theta]$ が代数的に分数的な数を含まないという付帯条件を課す必要がある。（この付帯条件をもつ基底の構成は § 7 で示す。）すると、最も一般の領域 $\mathfrak G$ ---これは有理的に整な超越数から成る領域とも呼べる--- について、再び次の単純な結果を得る（§ 6）。
\begin{enumerate}
\item[a)] 同じ有理基底 $\Theta$ によって構成できるすべての領域 $\mathfrak G$ の交わりは、基底のすべての整有理関数、すなわち領域 $[\Theta]$ によって与えられ、この $[\Theta]$ 自身が領域 $\mathfrak G$ である。
\item[b)] すべての領域 $\mathfrak G$ は、任意の系 $\mathfrak S(\vartheta)$ を $[\Theta]$ に有理的に整に添加することによって生じる。ここで $\mathfrak S(\vartheta)$ は、添加によって代数的に分数的な数が領域に入らないというただ一つの条件を満たせばよい。
\end{enumerate}

領域 $\mathfrak H$ と完全な類似を得るには、領域 $\mathfrak G$ の定義性質 III と IV から代数的数を省かなければならない。III と IV をこのように修正すると、有理基底による整元の構成は任意の体へ移すことができる（§ 10）。これに対し Zermelo の構成は体の代数的閉性を前提としている。

$H$ がすべての代数的基底を走り、また $\Theta$ が付帯条件を満たすすべての有理基底を走るとき、結果 a), b) はすべての領域 $\mathfrak H$ および $\mathfrak G$ の全体を与える。すべての同型な領域を一つの類にまとめるならば、この全体から、すべての類の全体を代表する部分集合を選ぶことができる。§ 9 により、少なくとも可算無限個の類がある（§ 8）。

\subsection*{§ 1. すべての領域 $\mathfrak G$ に対する基底性質 1), 2) の証明。}

$\mathfrak G$ を、抽象的性質 I--IV をもつ任意に与えられた整超越数の領域とする。E. Zermelo が全複素数体に対して用いた手続きに従い、$\mathfrak G$ の中から、\emph{$\mathfrak G$ のすべての数の代数的基底} $H$ を選ぶ。II により、すべての実数または複素数は $\mathfrak G$ の二つの数の商として表されるから、$H$ は同時に\emph{すべての}数の代数的基底である。したがって、任意に与えられた領域 $\mathfrak G$ に対し、全ての数の代数的基底を $\mathfrak G$ に属するように選ぶことができる。ゆえに基底性質 1) は満たされる。さらに III により、$\mathfrak G$ は全ての代数的整数から成る整域 $[K]$ を含み、したがって I により、基底元 $\eta$ の $[K]$ 係数多項式、すなわち $\eta$ のすべての整多項式を含む。これらは整域 $[H]$ をなす。ゆえに基底性質 2) も常に満たされる。言い換えると、すべての領域 $\mathfrak G$ は、全ての数の代数的基底 $H$ によって、$\mathfrak G$ が整域 $[H]$ を因子として含むように構成できる。

\textbf{注。} それでも、基底 $H$ から出発して、$H$ 自身に関しては基底性質 1), 2) が成り立たないが、別の基底 $H'$ に関しては成り立つ領域 $\mathfrak G$ を構成することはもちろん可能である。例えば、Zermelo 領域 $\mathfrak G_\eta$ の全構成において、ある基底元 $\eta$ を多項式 $f(\eta)$ で置き換えて領域 $\mathfrak G'$ を得るとする。このとき $\mathfrak G'$ は $\eta$ も $[H]$ も含まない。しかし基底 $H$ において基底元 $\eta$ を $\xi=f(\eta)$ で置き換えれば、$H$ は再び代数的基底 $H'$ となり、$\mathfrak G'$ は $H'$ に関して基底性質 1), 2) を満たす。

\subsection*{§ 2. 基底性質 4), 5) の排除。}

ここで、領域 $\mathfrak G$ は残りの基底性質を何も満たす必要がないことを示す。まず 4), 5) は、Zermelo の構成において整域 $[H]$ を整有理「関数式」の領域で置き換えることにより排除される。

Weber によれば、\footnote{H. Weber, Lehrbuch der Algebra. Kleine Ausgabe §§ 96, 97.} \emph{有理関数式}とは、次の一意に定まる表示における、有理数係数の任意の有理関数をいう。
\[
 A(\eta_1\cdots \eta_t)
   = a\cdot\frac{E_1(\eta_1\cdots \eta_t)}{E_2(\eta_1\cdots \eta_t)},
\]
ここで $E_1,E_2$ は、有理整数係数をもつ互いに素な原始多項式であり、$a$ は正の有理数（$A$ の絶対値）である。特に $a$ が\emph{有理整数}であるとき、$A$ は\emph{整有理関数式}と呼ばれる。整有理関数式の特殊な場合として有理整数および有理整数係数多項式が含まれ、一方、有理的に分数的な数および有理的に分数的な係数をもつ多項式は、分数的有理関数式に数えられる。

いま領域 $\mathfrak G$ を、一度に有限個の基底元 $\eta$ のすべての整有理関数式の全体 $\mathfrak F$ と、$\mathfrak F$ に代数的に整に依存するすべての関数、すなわち最高係数が 1 で残りの係数が整有理関数式であるようなある方程式を満たす関数、から成るものとする。

$\mathfrak G$ に対する性質 I--IV の証明は Zermelo の対応する証明とまったく平行であり、ここでは簡単に示すにとどめる。まず II と III は明らかである。なぜなら $\mathfrak G$ は Zermelo 領域 $\mathfrak G_\eta$ を因子として含むからである。有理整数係数多項式についての Gauss の定理により、整有理関数式の和、差、積は再び整有理関数式である。したがって $\mathfrak F$ は整域をなし、よく知られた議論\footnote{例えば Weber, § 97, 8 を参照。}により $\mathfrak G$ も整域となる。ゆえに I が成り立つ。さらに、拡張された Gauss の定理\footnote{Weber, § 20, 5.}から二つの補助命題が従う。すなわち「$z$ が\emph{ある}方程式によって $\mathfrak F$ に代数的に整に依存するならば、$z$ がすべての有理関数式の体に関して満たす\emph{既約}方程式によってもそうである」\footnote{Weber, § 97, 4 を参照。}、および「すべての有理数の体上既約で、ある代数的数が満たす方程式は、すべての有理関数式の体上でも既約である」である。したがって $\mathfrak G$ は分数的代数的数を含むことができない。よって IV、したがって抽象的性質 I--IV のすべてが検証された。

一方、どの整列のもとでも、$\mathfrak G$ が基底元のすべての有理的および代数的に分数的な関数を排除するように、$\mathfrak G$ から基底を選ぶことはできない。基底元として選ばれるべき領域の任意の関数を $z(\eta)$ とすれば、それは方程式
\[
 z^k+A_1(\eta)z^{k-1}+\cdots+A_k(\eta)=0
\]
によって定義される。ここで
\[
 A_i(\eta)=a_i\cdot\frac{E_1(\eta)}{E_2(\eta)}
\]
は整有理関数式である。すると関数 $\frac{a_k}{z}$ は $\mathfrak G$ に属し、同様に $\frac{a_k}{\sqrt z}$, $\frac{a_k}{\sqrt[3]z}$ なども属する。\footnote{まったく同様に、$\eta$ の任意の代数関数は、有理整数を掛けることによって関数式領域 $\mathfrak G$ の元へ変換できる。したがってこの領域は、代数的数の場合と同様に、任意の複素数を $\mathfrak G$ の元と有理整数との商として表せるという性質をもつ。} したがって基底性質 4), 5) は領域 $\mathfrak G$ の全体に対して排除される。

\textbf{注。} § 3 の例は、$\mathfrak G$ が、どの整列のもとでも有理的に分数的な関数式を含み得る一方、有理的または代数的に分数的な数を含まないこともあり得ることを示す。

\subsection*{§ 3. 基底性質 3) の排除。}

基底性質 3) を排除するため、一般化された\emph{加群領域}を用いる。この領域では、多項式の引数はもはや不定元ではなく、不定元のすべての代数的整関数であり、\footnote{「代数的整関数」とは常に\emph{整}係数をもつ関数を意味する。すなわち、最高係数が 1 で、残りの係数が代数的整数係数の $\eta$ の多項式、つまり $[H]$ の多項式であるような方程式を満たす関数である。特に $[H]$ のすべての多項式は代数的整関数に属する。}不定元自身が加群基底の役割を担う。

$\mathfrak G$ を\emph{すべての元}
\begin{equation}
 z=\alpha+g_1(\eta)\eta_1+g_2(\eta)\eta_2+\cdots+g_t(\eta)\eta_t
\tag{1}
\end{equation}
\emph{から成るものとする。ここで $\alpha$ はすべての代数的整数を走り、$\eta_1\cdots\eta_t$ はすべての基底元を走り、固定された各組合せ $\eta_1\cdots\eta_t$ に対して関数 $g_1(\eta)\cdots g_t(\eta)$ はそれぞれ独立に $\eta$ のすべての代数的整関数を走る。}\footnote{加群性は元 $z-\alpha$ に属する。}

この加群領域に対して性質 I--IV が成り立つことは容易に確かめられる。まず I が成り立つ。なぜなら二つの元 $z$ の和、差、積は再び (1) の形をとるからである。領域 $\mathfrak G$ はすべての基底元 $\eta$ も含み、さらに $\eta\cdot g(\eta)$ という形のすべての元も含む。ここで $g(\eta)$ は $\eta$ のすべての代数的整関数を走る。したがって $\eta$ のすべての代数的整関数、従って $\eta$ のすべての代数関数は、$\mathfrak G$ の二つの元の商として表せる。これにより II が証明される。表示 (1) には、$g_1=0\cdots g_t=0$ とすれば、すべての代数的整数も特に含まれている。しかし $z$ は分数的代数的数に等しくなり得ない。実際、$z$ が代数的数 $\beta$ を表すならば、
\[
 \beta=\alpha+g_1(\eta)\eta_1+\cdots+g_t(\eta)\eta_t
\]
が $\eta$ に関して恒等的に成り立つので、必ず $\beta=\alpha$ である。これは特殊化
\[
 \eta_1=0\cdots \eta_t=0
\]
によって示される。この特殊化のもとで、乗数 $g_1(\eta)\cdots g_t(\eta)$ は代数的整関数として有限に保たれる。なぜなら、それらは代数的整関数または代数的整数へ移るからである。\footnote{IV のこのより詳細な証明は、段落末の拡張例のために与える。ここでは、構成により $z$ が代数的整関数であるという注だけでも十分である。} したがって $\mathfrak G$ に対して条件 I--IV が満たされる。

一方、どの整列のもとでも、基底元のすべての代数的整関数が $\mathfrak G$ に属するように $\mathfrak G$ から基底を選ぶことはできない。領域の任意の関数 $z(\eta)$ を、基底元として選ぶべきものとする。したがってこれは実際に不定元 $\eta$ を含むものでなければならない。$z$ に依存するある有限の値以後の $\nu$ について、関数
\[
 \xi=\sqrt[\nu]{z-\alpha}
\]
はもはや $\mathfrak G$ に属せないことを示す。

実際、$\xi$ が $\mathfrak G$ に属すると仮定する。すると (1) により、$\xi$ は
\[
 \xi=\gamma+h_1(\eta)\eta_{i_1}+h_2(\eta)\eta_{i_2}
      +\cdots+h_\tau(\eta)\eta_{i_\tau}
\]
の形である。ここで再び $h_1(\eta)\cdots h_\tau(\eta)$ は $\eta$ の代数的整関数であり、$\eta_{i_1}\cdots\eta_{i_\tau}$ は任意の基底元で、特に $\eta_1\cdots\eta_t$ と一致してもよい。まず代数的整数 $\gamma$ が 0 であることを示さなければならない。$h_1(\eta)\cdots h_\tau(\eta)$ は代数的整関数であるから、残りの $\eta$ を任意の値にして $\eta_{i_1}=0\cdots\eta_{i_\tau}=0$ とおくと、$\xi$ は $\gamma$ に等しくなる。特に、
\[
 \eta_{i_1}=0\cdots \eta_{i_\tau}=0; \qquad
 \eta_1=0\cdots \eta_t=0
\]
とすれば、$\xi$ は $\gamma$ に等しい。しかしこの特殊化のもとで、(1) により関数 $(z-\alpha)$ は消え、したがって $\xi$ も消える。ゆえに $\gamma=0$ であり、
\[
 \xi=h_1(\eta)\eta_{i_1}+h_2(\eta)\eta_{i_2}
      +\cdots+h_\tau(\eta)\eta_{i_\tau}
\]
が成り立つ。これを $\nu$ 乗すると
\[
 \xi^\nu=z-\alpha=F_\nu(\eta)
\]
を得る。ここで $F_\nu(\eta)$ は、$\eta_{i_1}\cdots\eta_{i_\tau}$ に関する $\nu$ 次元の、恒等的には消えない同次形式で、その係数は $\eta$ の代数的整関数である。いま $z-\alpha$ は (1) により $\eta$ の代数的整関数であり、$[H]$ に関して既約な方程式
\[
 (z-\alpha)^\chi+B(\eta)(z-\alpha)^{\chi-1}+\cdots+C(\eta)=0
\]
を満たす。ここで最後の係数 $C(\eta)$、すなわち $(z-\alpha)$ とその共役の積、は多項式である。その次数を $\lambda$ とする。一方、$z-\alpha=F_\nu(\eta)$ から、$F_\nu(\eta)$ に対応する共役を掛けることにより
\[
 C(\eta)=G_{\nu\chi}(\eta)
\]
を得る。ここで $G_{\nu\chi}(\eta)$ は $\eta_{i_1}\cdots\eta_{i_\tau}$ に関する $\nu\chi$ 次元の同次形式で、係数は $\eta$ の多項式である。したがって $G_{\nu\chi}$ は $\eta$ に関して少なくとも次数 $\nu\chi$ をもつので、$\lambda\geq\nu\chi$ である。すなわち $\nu>\lambda/\chi$ となる関数 $\sqrt[\nu]{z-\alpha}$ は $\mathfrak G$ に属さない。したがって基底性質 3) は領域 $\mathfrak G$ の全体に対して排除される。

\textbf{注。} いま構成した領域をわずかに一般化すると、どの整列のもとでも分数的関数式をも含む領域 $\mathfrak G$ が得られる。$\mathfrak G$ をすべての元
\[
 z=\alpha+\gamma_1(\eta)\eta_1+\gamma_2(\eta)\eta_2
      +\cdots+\gamma_t(\eta)\eta_t
\]
から成るものとし、ここで $\gamma(\eta)$ は $\eta$ の代数的に整な、ただし必ずしも整ではない、すべての関数を走るものとする。性質 I--IV の証明は上と同じである。しかし $\mathfrak G$ は任意の関数 $z$ とともに $a\cdot(z-\alpha)$ も含む。ここで $a$ は任意の有理的または代数的に分数的な数である。したがってどの整列のもとでも分数的関数式が現れる。

\subsection*{§ 4. 代数的に整な超越数から成る最も一般の領域。}

以下を便利に述べるために、「元 $z$ が $\mathfrak T$ に代数的に整に依存する」という表現を導入する。これは、$\mathfrak T$ が任意に与えられた領域を表すとき、$z$ が、最高係数が 1 で、残りの係数が $\mathfrak T$ の元の整多項式であるようなある方程式を満たすことを意味する。\footnote{$\mathfrak T$ が整域の場合には、最高係数は 1 であり、残りの係数は $\mathfrak T$ の元である。}

領域 $\mathfrak T$ は、次の条件（序論を参照）を満たすとき、\emph{代数的整閉}であるという。
\begin{enumerate}
\item[V.] $\mathfrak T$ に代数的に整に依存するすべての元は $\mathfrak T$ に属する。
\end{enumerate}

まず、抽象的性質 V が Zermelo 領域 $\mathfrak G_\eta$ に属することを示す。元 $z$ が方程式 $f(z)=0$ によって $\mathfrak G_\eta$ に代数的に整に依存するとする。すると、$[H]$ に関する $f(z)$ の共役を掛け合わせることにより、$z$ は $[H]$ にも代数的に整に依存することがわかる。したがって $z$ は $\mathfrak G_\eta$ に属する。§ 2 で考えた関数式領域の代数的整閉性も同じように示される。

性質 V がすべての領域 $\mathfrak G$ に属するわけではないことは、§ 3 で考えた加群領域が示している。実際、それは基底性質 3) をもたない。より正確には、§ 3 は、加群領域が領域のどの超越元に関しても代数的整閉でないことを示した。一方、III と IV により、領域の任意の代数的元に関してはもちろんそうである。

これにより、まず領域 $\mathfrak G$ の全体から、抽象的性質 V ももつものを選び出す。これらを「代数的に整な超越数から成る領域 $\mathfrak H$」と呼ぶ。

$\mathfrak H$ をそのような領域とし、$H$ を $\mathfrak H$ のすべての数の代数的基底とする。II により、$H$ は同時にすべての数の代数的基底である。III, I, V により、$\mathfrak H$ は $H$ のほかに、$\eta$ のすべての代数的整関数、すなわち Zermelo 領域 $\mathfrak G_\eta$ を含む。したがって § 1 と同様に、代数的に整な超越数から成るすべての領域 $\mathfrak H$ は、すべての数の代数的基底によって、$\mathfrak H$ が Zermelo 領域 $\mathfrak G_\eta$ を因子として含むように構成できる。

いま $H$ をすべての数の任意の代数的基底とし、$\mathfrak M(H)$ を $H$ を含むすべての領域 $\mathfrak H$ の全体とする。$\mathfrak M(H)$ の任意の領域 $\mathfrak H$ は、いま証明したように、Zermelo 領域 $\mathfrak G_\eta$ を因子として含む。そして $\mathfrak G_\eta$ 自身が領域 $\mathfrak H$ であるから、$\mathfrak G_\eta$ は $\mathfrak M(H)$ のすべての領域 $\mathfrak H$ の最大公因子、すなわち交わりである。こうして $\mathfrak G_\eta$ は抽象的に定義された。また、$\mathfrak M(H)$ が交わりに関して閉じていることも直接に従う。すなわち、$\mathfrak M(H)$ の任意の部分系のすべての領域 $\mathfrak H$ の交わりは再び $\mathfrak M(H)$ に属する。実際、交わりとして I と V が成り立つ。II と III は $\mathfrak G_\eta$ が因子であることから従い、IV はその交わりが $\mathfrak H$ の因子であることから従う。

§ 2 の関数式領域の例が示すように、一般に領域 $\mathfrak H$ は $\mathfrak G_\eta$ 以外の関数を含む。そのような $\eta$ の有理関数または代数関数を $\psi(\eta)$ とする。すると I により、$\mathfrak H$ は $\psi$ と $\mathfrak G_\eta$ の有限個の関数とのすべての整有理結合も含む。このようにして生じる整域、すなわち $\mathfrak G_\eta$ 係数の $\psi$ のすべての多項式から成る整域を $[\mathfrak G_\eta,\psi]$ と書き、それへ至る過程を「$\psi$ の $\mathfrak G_\eta$ への有理的整添加」と呼ぶ。さらに V により、$\mathfrak H$ は $[\mathfrak G_\eta,\psi]$ に代数的に整に依存するすべての関数を含む。このようにして生じる領域を $\{\mathfrak G_\eta,\psi\}$ と書き、それへ至る過程を「$\psi$ の $\mathfrak G_\eta$ への代数的整添加」と呼ぶ。領域 $\{\mathfrak G_\eta,\psi\}$ は $\mathfrak G_\eta$ 係数の $\psi$ のすべての代数的整関数から成るので、よく知られた議論により整域である。\footnote{例えば Weber, § 97, 6, 7 を参照。}

$\{\mathfrak G_\eta,\psi\}$ も代数的整閉、すなわち条件 V を満たすことを示す。実際、$z$ が方程式
\[
 f(z;a\cdots c)=z^\nu+a z^{\nu-1}+\cdots+c=0
\]
によって $\{\mathfrak G_\eta,\psi\}$ に代数的に整に依存するとする。ここで $a\cdots c$ は $\{\mathfrak G_\eta,\psi\}$ の元として、$[\mathfrak G_\eta,\psi]$ の係数をもつ方程式
\[
\begin{gathered}
 a^\lambda+A_1a^{\lambda-1}+\cdots+A_\lambda=0,\\
 \vdots\\
 c^\mu+C_1c^{\mu-1}+\cdots+C_\mu=0
\end{gathered}
\]
を満たす。それらの根を
$a,a'\cdots a^{(\lambda-1)}\cdots c,c'\cdots c^{(\mu-1)}$ と書く。いま根のすべての組合せにわたって積
\[
 g(z)=f(z;a\cdots c)f(z;a'\cdots c')\cdots
      f(z;a^{(\lambda-1)}\cdots c^{(\mu-1)})
\]
を作れば、$g(z)$ は最高係数が 1 で、残りの係数が $[\mathfrak G_\eta,\psi]$ の元である。したがって $z$ は $\{\mathfrak G_\eta,\psi\}$ に属する。\footnote{このように、代数的整性を\emph{ある}方程式によって定義すれば、I と V の証明において基礎体に関する既約性の概念は不要になる。§ 2 も参照。そこでは IV の証明のためだけに、任意の方程式による代数的整性の定義から既約方程式へ移る補助命題が必要であった。この補助命題は現在の場合には一般に成り立たないので、IV は関数 $\psi$ に特別条件として課さなければならない。}

したがって領域 $\{\mathfrak G_\eta,\psi\}$ は性質 I と V をもつ。また $\mathfrak G_\eta$ がその領域の因子であるから、II と III ももつ。IV が成り立つかどうかは $\psi$ の選び方に依存する。例えば $\psi=\frac{1}{n\cdot\eta}$ のときは $\frac1n$ が領域に属するため、IV は成り立たない。しかし § 2 により、$\psi=\frac1\eta$ のとき、また $\psi=\frac{\eta}{n}$ のときは成り立つ。後者の場合、$\{\mathfrak G_\eta,\psi\}$ の関数が代数的数を表すなら、その値は $\eta=0$ としても変わらない。したがってそれは $\mathfrak G_\eta$ の関数へ移り、従って代数的整数へ移る。

ゆえに、$\psi$ に対し、代数的整添加によって代数的に分数的な数が領域に入らないという条件を課すと、$\{\mathfrak G_\eta,\psi\}$ は領域 $\mathfrak H$ となる。

まったく同様に、$\eta$ の有理関数または代数関数から成る任意の系 $\mathfrak S$ の $\mathfrak G_\eta$ への有理的整添加、また代数的整添加を定義する。有理的整添加は、一度に有限個の $\mathfrak G_\eta$ と $\mathfrak S$ の元の整多項式全体から成る整域 $[\mathfrak G_\eta,\mathfrak S]$ に導く。代数的整添加は、$[\mathfrak G_\eta,\mathfrak S]$ に代数的に整に依存するすべての元から成る領域 $\{\mathfrak G_\eta,\mathfrak S\}$ を与える。上とまったく同様に、この整域 $\{\mathfrak G_\eta,\mathfrak S\}$ は代数的整閉であり、また II と III も満たすことが示される。したがって、上と同様に次を得る。

$\mathfrak S$ に、代数的整添加によって代数的に分数的な数が領域に入らないという条件を課すと、すなわち $\mathfrak S$ が可容系となると、$\{\mathfrak G_\eta,\mathfrak S\}$ は領域 $\mathfrak H$ となる。\footnote{より一般に、同様に次が示される。$\mathfrak T$ を任意の領域とし、$\{\mathfrak T\}$ を $\mathfrak T$ に代数的に整に依存する元全体とすれば、$\{\mathfrak T\}$ は代数的整閉である。}

いま重要なのは逆も成り立つことである。すなわち、$\mathfrak S$ がすべての可容系を走るとき、領域 $\{\mathfrak G_\eta,\mathfrak S\}$ は実際に $\mathfrak M(H)$ のすべての領域 $\mathfrak H$ を尽くす。実際、$\mathfrak H$ を $\mathfrak M(H)$ の任意の領域とし、$\mathfrak L=\mathfrak H-\mathfrak G_\eta$ を $\mathfrak H$ のうち $\mathfrak G_\eta$ に属さないすべての関数から成る系とする。V により、$\mathfrak H$ は領域 $\{\mathfrak G_\eta,\mathfrak L\}$ を因子として含む。しかし $\{\mathfrak G_\eta,\mathfrak L\}$ も $\mathfrak H$ によって割り切られる。なぜなら、それは $\mathfrak G_\eta$ と $\mathfrak L$ を含み、これらは共通元をもたないので、$\mathfrak H$ も含むからである。したがって $\mathfrak H=\{\mathfrak G_\eta,\mathfrak L\}$ である。また IV が $\mathfrak H$ に対して成り立つので、$\mathfrak L$ はもちろん可容系である。\footnote{$\mathfrak L$ の部分系の添加ですでに十分なこともある。例えば § 2 の関数式領域は、多項式を除くすべての整有理関数式を $\mathfrak G_\eta$ に代数的整に添加することによって生じる。}

$H$ が可能なすべての代数的基底を走るとき、$\mathfrak M(H)$ は冒頭で証明したようにすべての領域 $\mathfrak H$ の全体を走る。したがって、代数的に整な超越数から成る最も一般の領域 $\mathfrak H$ の問題は、この段落の結果によって完全に答えられる。
\begin{enumerate}
\item[a)] $\mathfrak M(H)$ のすべての領域 $\mathfrak H$ の交わりは Zermelo 領域 $\mathfrak G_\eta$ によって与えられ、これはそれ自身が領域 $\mathfrak H$ である。$\mathfrak M(H)$ は交わりの形成に関して閉じている。
\item[b)] $\mathfrak M(H)$ のすべての領域 $\mathfrak H$ は、可容系 $\mathfrak S$ を $\mathfrak G_\eta$ に代数的整添加することによって生じる。$H$ が可能なすべての代数的基底を走るとき、$\mathfrak M(H)$ はすべての領域 $\mathfrak H$ の全体を走る。\footnote{可容系 $\mathfrak S$ は次のようにして見出せる。$\eta$ のすべての代数的に分数的な関数をある整列 $\Omega$ のもとに置き、$\mathfrak S$ には、$\mathfrak G_\eta$ およびすでに認められた関数に代数的整添加すると代数的に分数的な数を生成してしまうものを除き、すべての関数を入れる。この手続きは任意の場所で打ち切ることができる。$\Omega$ がすべての整列を走り、固定された各 $\Omega$ について任意の場所で手続きを打ち切るとき、すべての可容系 $\mathfrak S$ が尽くされる。}
\end{enumerate}

\section*{§ 5. 代数的基底に基づく、整超越数から作られる最も一般の領域。}

$H$ をすべての数の任意の代数的基底とし、$\mathfrak M(\mathfrak G)$ を $H$ を含み、したがって $[H]$ を含むすべての領域 $\mathfrak G$ の全体とする。$H$ が可能なすべての基底を走るとき、$\mathfrak M(\mathfrak G)$ はすべての領域 $\mathfrak G$ の全体を走る。したがって § 4 と同様に、$\mathfrak M(\mathfrak G)$ の領域 $\mathfrak G$ だけを考えればよい。

$\mathfrak M(\mathfrak G)$ の領域 $\mathfrak G$ はすべて $[H]$ を因子として含む。しかし $[H]$ 自身は領域 $\mathfrak G$ ではないので、§ 4 のように $[H]$ が最大公因子であるとは結論できない。それにもかかわらずそうであることは、§ 3 の加群領域をわずかに一般化して得られる $\mathfrak M(\mathfrak G)$ の可算個の領域 $\mathfrak G$ によって示される。それらは実際、$[H]$ 以外に共通元をもたない。

\emph{各固定された $\nu$ について、領域 $\mathfrak G_\nu$ $(\nu=1,2,\ldots\text{ in inf.})$ を、すべての量}
\begin{equation}
z=f(\eta)+g_1(\eta)\eta_1^\nu+g_2(\eta)\eta_2^\nu+\cdots+g_t(\eta)\eta_t^\nu
\tag{1}
\end{equation}
\emph{から成るものとする。ここで $f(\eta)$ は $\eta$ のすべての整多項式を走り、$\eta_1,\ldots,\eta_t$ は基底量を走り、固定された各組合せ $\eta_1^\nu,\ldots,\eta_t^\nu$ に対して、関数 $g_1(\eta),\ldots,g_t(\eta)$ は互いに独立に $\eta$ のすべての代数的整関数を走る。}

§ 3 と同様に、領域 $\mathfrak G_\nu$ は抽象的性質 I--IV をもつ。$\mathfrak G_\nu$ は代数的整関数だけを含み、$\nu=1$ のとき § 3 の領域と一致する。いま、領域 $\mathfrak G_\nu$ が $[H]$ の多項式 $f(\eta)$ 以外の共通元をもたないことを示す。実際、$z$ を $\eta$ の代数的整数であるが有理ではない関数とし、$[H]$ に関して既約な方程式
\begin{equation}
z^\chi+a_1(\eta)z^{\chi-1}+a_2(\eta)z^{\chi-2}+\cdots+a_\chi(\eta)=0
\tag{2}
\end{equation}
を満たすものとする。ここで $\chi>1$ であり、$\lambda$ を多項式 $a_1(\eta),\ldots,a_\chi(\eta)$ の次数の最大値とする。もし $z$ が $\mathfrak G_\nu$ に属するなら、量 $\xi=z-f(\eta)$ は、同じく $[H]$ に関して既約な方程式
\begin{equation}
\xi^\chi+b_1(\eta)\xi^{\chi-1}+b_2(\eta)\xi^{\chi-2}+\cdots+b_\chi(\eta)=0
\tag{3}
\end{equation}
を満たす。ここで (1) により、$b_1(\eta),b_2(\eta),\ldots,b_\chi(\eta)$ は $\xi$ の対称関数として、最低次元の項として少なくとも $\nu,2\nu,\ldots,\chi\nu$ 次元の項を含む。(2) と (3) を比較すると
\[
 b_1(\eta)=a_1(\eta)+\chi f(\eta)
\]
を得る。したがって、量 $\xi=z+\frac{a_1(\eta)}{\chi}$ が
\begin{equation}
\xi^\chi+c_2(\eta)\xi^{\chi-2}+\cdots+c_\chi(\eta)=0
\tag{4}
\end{equation}
を満たすならば、
\begin{equation}
 \xi-\zeta=\frac{b_1(\eta)}{\chi};
 \qquad
 c_i(\eta)\equiv b_i(\eta)\pmod{\frac{b_1(\eta)}{\chi}},
\tag{5}
\end{equation}
である。ここで $b_1(\eta)$ は少なくとも $\nu$ 次元の項から始まる。いま $c_i(\eta)$ は、(2) と (4) の比較が示すように、(2) の係数 $a_i(\eta)$ について次数 $\chi$ であり、したがって $\eta$ について高々次数 $\chi\lambda$ である。一方、すべての $b_i(\eta)$ は少なくとも $\nu$ 次元の項から始まる。ゆえに $\nu>\chi\lambda$ と選ぶと、(5) は
\[
 c_i(\eta)=0;\qquad
 \xi^\chi=0\quad\text{または}\quad
 \left(z+\frac{a_1(\eta)}{\chi}\right)^\chi=0
\]
を与える。したがって $z$ は仮定に反して $\eta$ の有理関数となる。ゆえに次が成り立つ。

\emph{$z$ が $\eta$ の代数的整数で非有理な関数であれば、$z$ は $\nu>\chi\lambda$ であるいかなる領域 $\mathfrak G_\nu$ にも属することができない。すなわち、}

\emph{$\mathfrak M(\mathfrak G)$ のすべての領域 $\mathfrak G$ の交わりは整域 $[H]$ によって与えられ、$[H]$ 自身は領域 $\mathfrak G$ ではない。したがって $\mathfrak M(\mathfrak G)$ は交わりの形成に関して閉じていない。}

したがって、代数的基底による最も一般の領域 $\mathfrak G$ の構成も、領域 $\mathfrak H$ の構成より見通しにくい。領域 $[H]$ は性質 I, III, IV をもつ。任意の系 $\mathfrak S$ を有理的に整に添加すると、得られる整域 $[H,\mathfrak S]$ については I と III が保たれ、IV は $\mathfrak S$ に課される特別条件となる。領域 $\mathfrak G$ が生じるためには、II も $\mathfrak S$ に対する条件として課されなければならない。別の言い方をすれば、$(H)$ 上有限な任意の体 $\mathfrak K$\footnote{$(H)$ は、代数的数係数の $\eta$ のすべての有理関数から成る体を意味する。}に対し、$\mathfrak K$ 上有限な $\mathfrak K$ 上の体 $\mathfrak L$ が存在し、$\mathfrak L$ のすべての関数が $[H,\mathfrak S]$ の関数の商として表されなければならない。この条件は § 4 の可容系に対する唯一の条件よりもはるかに扱いにくい。そこで、いま有理基底へ移る。

\section*{§ 6. 有理基底に基づく、整超越数から作られる最も一般の領域。}

$\mathfrak G$ を整超越数から成る任意の領域とする。$\mathfrak G$ の元を整列し、超限帰納法によって、有限個の先行する元によって有理的に表すことのできない元を $\mathfrak G$ から選ぶ。こうして得られる系を $\Theta$ と書く。$\mathfrak G$ のすべての元は $\Theta$ によって有理的に表される。実際、$z_0$ がそうでない最初の元であったなら、$z_0$ は $\mathfrak G$ の有限個の先行する元によって有理的に表されるはずであり、それらを帰納法の仮定により $\vartheta$ の有理関数で置き換えれば、$z_0$ の $\vartheta$ による表示が得られて矛盾する。また構成により、$\Theta$ のどの元も有限個の先行する元によって有理的に表すことはできない。したがって $\Theta$ は $\mathfrak G$ の有理基底である。

II により、すべての数は $\mathfrak G$ の二つの元の商であるから、すべての数は $\Theta$ によって有理的に表される。すなわち $\Theta$ は同時にすべての数の有理基底である。III と I により、$\mathfrak G$ は $\Theta$ のほかに、$\vartheta$ のすべての整多項式から成る整域 $[\Theta]$ を因子として含む。一方、IV により $[\Theta]$ は代数的に分数的な数を含まない。したがって、すべての数の有理基底 $\Theta$ は、$\Theta$ から得られる整域 $[\Theta]$ が代数的に分数的な数を含まないという付帯条件を満たす。\footnote{これが実際に条件であることは、例えば、二つの量 $\eta,\frac1{2\sqrt{\eta}}$ は基底元として許されないが、$\eta,\frac1{\sqrt{\eta}}$ は許されることから分かる。} $\Theta$ はすべての数の\emph{付帯条件付き有理基底}である。したがって § 2 と同様に、\emph{すべての領域 $\mathfrak G$ は、付帯条件付き有理基底によって、$\mathfrak G$ が整域 $[\Theta]$ を因子として含むように構成できる。}

$\Theta$ をすべての数の任意の付帯条件付き有理基底とする。その存在は領域 $\mathfrak G$ の存在により保証されている。$\mathfrak N(\mathfrak G)$ を、$\Theta$ を含み、したがって $[\Theta]$ を含むすべての領域 $\mathfrak G$ の全体とする。$\Theta$ が付帯条件付きの可能なすべての基底を走るとき、上により $\mathfrak N(\mathfrak G)$ はすべての領域 $\mathfrak G$ の全体を走る。したがって再び $\mathfrak N(\mathfrak G)$ の領域 $\mathfrak G$ だけを考えればよい。

いま $[\Theta]$ 自身は領域 $\mathfrak G$ である。実際、すべての数は $\Theta$ によって有理的に表されるので、$[\Theta]$ の二つの多項式の商として表される。またその構成により、$[\Theta]$ は残りの条件 I, III, IV も満たす。\emph{したがって $[\Theta]$ は $\mathfrak N(\mathfrak G)$ のすべての領域 $\mathfrak G$ の最大公因子、すなわち交わりである。さらに $\mathfrak N(\mathfrak G)$ は交わりに関して閉じている。すなわち、$\mathfrak N(\mathfrak G)$ の任意の部分系のすべての領域 $\mathfrak G$ の交わりも $\mathfrak N(\mathfrak G)$ に属する。} I は交わりに対して成り立つ。II と III は $[\Theta]$ が因子であることから従い、IV は交わりの領域がある領域 $\mathfrak G$ に因子として含まれることから従う。

$\vartheta$ の有理関数から成る任意の系 $\mathfrak S$ を $[\Theta]$ に有理的に整に添加すると、整域 $[\Theta,\mathfrak S]$ が得られる。ここで $\vartheta$ の代数関数は、有理基底の性質により、ある他の $\vartheta$ によって有理的に表すことができる。この整域は I, II, III をもち、したがって $\mathfrak S$ が「可容系」、すなわち有理的整添加によって代数的に分数的な数が領域に入らないという条件を満たすなら、領域 $\mathfrak G$ となる。逆に、$\mathfrak N(\mathfrak G)$ の任意の領域 $\mathfrak G$ は、$\mathfrak S$ を残余系 $\mathfrak G-[\Theta]$ とすれば、表示 $[\Theta,\mathfrak S]$ を許す。したがって、最も一般の領域 $\mathfrak G$ について、§ 4 と完全に類似して次が成り立つ。
\begin{enumerate}
\item[a)] \emph{$\mathfrak N(\mathfrak G)$ のすべての領域 $\mathfrak G$ の交わりは整域 $[\Theta]$ によって与えられ、$[\Theta]$ 自身が領域 $\mathfrak G$ である。$\mathfrak N(\mathfrak G)$ は交わりの形成に関して閉じている。}
\item[b)] \emph{$\mathfrak N(\mathfrak G)$ のすべての領域 $\mathfrak G$ は、可容系 $\mathfrak S$ を $[\Theta]$ に有理的整添加することによって生じる。$\Theta$ が付帯条件付きの可能なすべての有理基底を走るとき、$\mathfrak N(\mathfrak G)$ はすべての領域 $\mathfrak G$ の全体を走る。}\footnote{最も一般の可容系 $\mathfrak S$ については、§ 4 で述べたことが「代数的」を「有理的」に置き換えて成り立つ。}
\end{enumerate}

\emph{注。} $\Theta$ から $\Theta$ のすべての量の代数的基底 $H$ を抜き出すと、$H$ は同時にすべての数の代数的基底である。逆に、任意の代数的基底は、整列において $H$ を最初に置くことにより有理基底へ拡張できる。したがって、§ 5 の代数的基底による領域 $\mathfrak G$ の構成は次のように解釈できる。$[H]$ に添加すべき系 $\mathfrak S$ を二つの部分系 $\mathfrak S_1$ と $\mathfrak S_2$ に分解する。$\mathfrak S_1$ の添加は、$H$ を付帯条件付き有理基底へ拡張することに相当し、その基底に可容系 $\mathfrak S_2$ をさらに有理的に整に添加する。

\section*{§ 7. 付帯条件付きの最も一般の有理基底の構成。}

前節の結果にはなお補足が必要である。そこでは、すべての数に対する付帯条件付き有理基底の\emph{存在}は示されたが、与えられた整列からそれを構成する方法は示されていなかった。付帯条件のため、この構成は、領域 $\mathfrak G$ ではなく全複素数体が与えられると、より複雑な形をとる。\footnote{§ 6 b) では、$\Theta$ を、代数的基底とその基底量の代数的整関数から成る特殊な有理基底だけに走らせても十分である。その場合、付帯条件は自動的に満たされる。これを見るには、§ 6 により $\mathfrak G$ からの代数的基底を有理基底へ拡張し、得られた基底において、$\eta$ の各代数的に分数的な関数に、適当に選んだ $\eta$ の多項式を掛ければよい。すると基底性質は保たれる。この構成は全複素数体へそのまま移る。しかし、§ 6 から生じる付帯条件付きの最も一般の有理基底に関する問題は、それ自体としても興味をもつ。}

全複素数体を整列 $\Omega$ にかける。すると三種類の量が現れうる。
\begin{enumerate}
\item[1.] 量 $y$。これは、先行する基底量（3. で再帰的に定義される）への有理的整添加により、得られる整域に代数的に分数的な数が入ることを強制するものである。基底量が一つも先行しない場合、この領域は $y$ のすべての整多項式から成る。特に、すべての代数的に分数的な数は量 $y$ である。なぜなら得られる整域は量 $y=\beta$ を含むからである。
\item[2.] 量 $z$。これは性質 1 をもたないが、$y$ とは異なる有限個の先行する数によって有理的に表せるものである。したがって $z=\varphi(\xi_1,\ldots,\xi_t)$ という関係を満たす。ここで $\varphi$ は $K$ 係数の有理関数であり、$\xi$ の中に $y$ はない。特に、すべての代数的整数 $\alpha$ はここに属する。なぜなら性質 1 はそれらに対して成り立たず、関係 $z=\alpha$ を満たすからである。
\item[3.] 量 $\vartheta$。これらは 1 も 2 も成り立たない量で、基底量と呼ばれる。特に、整列における最初の超越数 $\vartheta_0$ はそのような基底量である。なぜなら $\vartheta_0$ の整多項式は代数的に分数的な数を表せず、また $\vartheta_0$ は先行する代数的数に有理的に依存することもできないからである。
\end{enumerate}

\emph{いま、すべての基底量 $\vartheta$ の系 $\Theta$ が、すべての数に対する付帯条件付き有理基底であることを示す。}

量 $y$ は $\vartheta$ の中に現れないので、整域 $[\Theta]$ は代数的に分数的な数を含むことができない。したがって付帯条件は満たされる。さらに、量 $z$ も $\vartheta$ の中に現れないので、各 $\vartheta$ は先行する基底量に有理的に独立である。帰納法の議論は § 6 とまったく同様に、関係 $z=\varphi(\xi)$（ただし $\xi$ の中に $y$ はない）が常に関係 $z=\psi(\vartheta)$ を伴うことを示す。したがってすべての量 $z$ は $\vartheta$ によって有理的に表せる。

残るのは、より困難な主張、すなわち $y$ も $\vartheta$ によって有理的に表せることを証明することだけである。まず $y$ が $\vartheta$ に代数的に依存することを示す。仮定により、各 $y$ について、少なくとも一つの代数的に分数的な数 $\beta$ が存在し、表示
\begin{equation}
\beta=f_0(\vartheta)+f_1(\vartheta)y+\cdots+f_\chi(\vartheta)y^\chi
\tag{1}
\end{equation}
をもつ。ここで $f_0(\vartheta),\ldots,f_\chi(\vartheta)$ は $[\Theta]$ の多項式であり、したがって代数的に分数的な数を表せない。もし $y$ が $\vartheta$ に代数的に独立であったなら、(1) は $y$ に関して恒等的に成り立ち、$\beta=f_0(\vartheta)$ となり、$[\Theta]$ の性質に反する。

代数的依存から有理的依存を導くため、$\Theta$ から $\Theta$ のすべての量の代数的基底 $H$ を抜き出す。すると各 $y$ は $\vartheta$ の代数関数として $\eta$ にも代数的に依存する。したがって $y=y(\eta)$ によって、$(H)$ 上の代数体 $(H,y(\eta))$ が定まる。すべての $\eta$ は同時に量 $\vartheta$ でもあるから、有理的依存の証明は、このような各体に対し、$\vartheta$ の有理関数である原始元が存在することを示すことに帰着する。より正確には、$y(\eta)$ に $\eta$ の多項式を掛けると性質 1 を失い、したがってそのような元になることを示す。

$(H)$ と $(H,y)$ の間には有限個の中間体しかないことが知られている。そのうち、$\vartheta$ によって有理的に表せる原始元をもつものを $\Omega_0=(H),\Omega_1,\ldots,\Omega_\sigma$ とする。つまり体の各元が $\vartheta$ の有理関数となるものを列挙する。この条件は少なくとも常に $\Omega_0=(H)$ について満たされる。$\Omega_i$ に属する原始関数を $\frac{g_i(\vartheta)}{h_i(\vartheta)}$ とし、$g_i,h_i$ は $[\Theta]$ の多項式とする。すると、$\alpha_i$ を適当に選ばれた整数として、$[\Theta]$ の多項式
\[
 \xi_i=g_i(\vartheta)+\alpha_i h_i(\vartheta)
\]
は $\Omega_i$ の原始関数、または $\Omega_i$ 上の代数体の原始関数であり、\footnote{§ 5 の終わりを参照。}$\Omega_i$ の任意の関数 $\psi(\eta)$ は表示
\begin{equation}
\psi(\eta)=
\frac{a_1(\eta)\xi_i^{\kappa-1}+a_2(\eta)\xi_i^{\kappa-2}+\cdots+a_\kappa(\eta)}
     {a(\eta)}
=\frac{A(\eta,\xi_i)}{a(\eta)}
\tag{2}
\end{equation}
をもつ。ここで $a(\eta),a_1(\eta),\ldots,a_\kappa(\eta)$ は $\eta$ の整多項式であり、したがって $A(\eta,\xi_i)$ と $a(\eta)$ は $[\Theta]$ の量である。

$y$ が $\Omega_i$ に関して満たす既約方程式は、(2) により
\[
 f^{(i)}(\eta)y^{\lambda_i}+f_1^{(i)}(\eta,\xi_i)y^{\lambda_i-1}
+\cdots+f_{\lambda_i}^{(i)}(\eta,\xi_i)=0
\]
の形であり、この方程式のすべての係数は $[\Theta]$ に属する。すると関数 $y_i=f^{(i)}(\eta)\cdot y$ は、$\Omega_i$ に関して同じく既約な方程式
\[
 y_i^{\lambda_i}+f_1^{(i)}(\eta,\xi_i)y_i^{\lambda_i-1}
+\cdots+\bigl(f^{(i)}(\eta)\bigr)^{\lambda_i-1}
 f_{\lambda_i}^{(i)}(\eta,\xi_i)=0
\]
を満たし、それによって $[H,\xi_i]$ に代数的に整に依存する。同じことは $y_i$ に $\eta$ の任意の多項式を掛けたものについても成り立つ。特に、関数
\begin{equation}
Y=f^{(0)}(\eta)f^{(1)}(\eta)\cdots f^{(\sigma)}(\eta)\cdot y
\tag{3}
\end{equation}
は、すべての領域 $[H,\xi_i]$ に代数的に整に依存し、それぞれ $\Omega_i$ に関して既約な方程式
\begin{equation}
Y^{\lambda_i}+F_1^{(i)}(\eta,\xi_i)Y^{\lambda_i-1}
+\cdots+F_{\lambda_i}^{(i)}(\eta,\xi_i)=0
\qquad (i=0,1,\ldots,\sigma)
\tag{4}
\end{equation}
によってそうである。いま $Y$ が、求める、$\vartheta$ によって有理的に表せる原始元であることを示す。例えば、代数的数 $\gamma$ が $Y$ と $[\Theta]$ によって
\begin{equation}
\gamma=k_0(\vartheta)+k_1(\vartheta)Y+\cdots+k_\nu(\vartheta)Y^\nu
\tag{5}
\end{equation}
と表されると仮定する。ここで $k_0(\vartheta),\ldots,k_\nu(\vartheta)$ は $[\Theta]$ の多項式である。いま、(5) の係数領域から生じる体 $\mathfrak K=(H;k_0(\vartheta),\ldots,k_\nu(\vartheta))$ に関して $Y$ が満たす既約方程式を決定する。それは、$(H)$ と $(H,Y)$ の間にある $\mathfrak K$ と $(H,Y)$ の交体に関する $Y$ の既約方程式によって与えられることが知られている。ところが (3) により $(H,Y)$ は $(H,y)$ と同一である。また $\mathfrak K$ のすべての元、したがって交体のすべての元は有限個の $\vartheta$ によって有理的に表せるので、この交体は必ず $\Omega_0,\Omega_1,\ldots,\Omega_\sigma$ の一つ、例えば $\Omega_\tau$ である。したがって求める既約方程式は (4) により次数 $\lambda_\tau$ であり、
\begin{equation}
Y^{\lambda_\tau}+F_1^{(\tau)}(\eta,\xi_\tau)Y^{\lambda_\tau-1}
+\cdots+F_{\lambda_\tau}^{(\tau)}(\eta,\xi_\tau)=0
\tag{6}
\end{equation}
である。その係数は $[\Theta]$ の多項式である。(6) により、(5) は
\begin{equation}
\gamma=K_0(\vartheta)+K_1(\vartheta)Y+\cdots+
K_{\lambda_\tau-1}(\vartheta)Y^{\lambda_\tau-1}
\tag{7}
\end{equation}
の形に書き直せる。ここで $K_0(\vartheta),K_1(\vartheta),\ldots$ は $\mathfrak K$ に属し、他方で $[\Theta]$ の多項式である。しかし $\gamma$ は代数的数であるから、関係 (7) は $Y$ において次数 $\lambda_\tau$ に達しておらず、したがって $Y$ に関して\emph{恒等的に}満たされる。ゆえに
\[
 \gamma=K_0(\vartheta)
\]
であり、すなわち $\gamma$ は $[\Theta]$ に属し、したがって代数的整数である。

$\gamma$ が $Y$ と $[\Theta]$ によって表されるすべての代数的数を走るとき、$(H,Y)$ と対応する体 $\mathfrak K$ との交体は常に $\Omega_0,\Omega_1,\ldots,\Omega_\sigma$ だけを走り、上の議論はすべて有効である。すなわち、\emph{$Y$ を $[\Theta]$ に有理的整添加しても、代数的に分数的な数は表されない。} したがって $Y$ はもはや性質 1 をもたない。2) または 3) により、$Y$ は量 $z$ または $\vartheta$ であり、したがって有限個の $\vartheta$ によって有理的に表される。(3) により同じことは $y$ にも成り立つ。\emph{$\Theta$ がすべての数の有理基底であることが証明された。} 逆に、付帯条件付きの任意の有理基底は、整列において $\Theta$ を先に置くことにより 1), 2), 3) に従って構成できる。したがって次が証明された。\emph{1), 2), 3) によって与えられた構成は、すべての数に対する付帯条件付きの最も一般の有理基底に導く。}

\section*{§ 8. 領域 $\mathfrak G$ および $\mathfrak H$ の類への分割。}

最も一般の領域 $\mathfrak G$ および $\mathfrak H$ の問題と並んで、さらに\emph{本質的に異なる}領域の問題が生じる。これは、以下でより正確にする意味で、同じ構成原理によって生成できない領域のことである。このことは、これらの領域を類へ分割することへ導く。

\emph{二つの領域の元を一対一に対応させることができ、しかも一方の領域の任意の二元の和および積が、対応する他方の元の和および積に常に対応するならば、その二つの領域は同じ類に属する、または同型であるという。} この定義から、任意の同型関係のもとで零元と単位元はそれ自身に対応し、有限個の元の間のすべての代数関係は対応する元について保たれ、逆も成り立つことが従う。特に単位元のほか、すべての有理数はそれ自身に対応する。一方、代数的数全体、同様に代数的整数全体はそれ自身へ移るが、個々の代数的数は共役に対応してもよい。

まず、固定された代数的基底 $H$ を含むすべての領域 $\mathfrak G$ の全体 $\mathfrak M(\mathfrak G)$ からの任意の領域は、別の代数的基底 $H^*$ に属する全体 $\mathfrak M^*(\mathfrak G)$ からのある領域に同型であることに注意する。実際、全複素数体は任意の代数的基底から代数拡大によって生じ、したがって基底と同じ濃度をもつ。\footnote{Steinitz, §§ 23, 24.} したがって $H$ と $H^*$ は同じ濃度をもち、求める同型写像は基底量 $\eta_i$ を $\eta_i^*$ に対応させることで得られる。\footnote{もちろんこの結論は、二つの有理基底 $\Theta$ と $\Theta^*$ に属する集合 $\mathfrak N(\mathfrak G)$ と $\mathfrak N^*(\mathfrak G)$ には成り立たない。また、$\Theta$ と $\Theta^*$ の相互の有理的表現可能性から $\mathfrak N(\mathfrak G)$ と $\mathfrak N^*(\mathfrak G)$ の同型を推論することもできない。しかし、体 $(\Theta)$ と $(\Theta^*)$ の間の同型関係は得られる。} 他方、$\mathfrak M(\mathfrak G)$ は、§§ 2, 3, 5 の例から従うように、同型でない、すなわち本質的に異なる領域 $\mathfrak G$ を含む。実際、すべての代数関係は同型写像のもとで保たれるから、代数的基底は再び代数的基底に対応しなければならない。しかし §§ 2, 3 は、関数式領域および加群領域が、Zermelo 領域のすべての基底性質を満たすような代数的基底をもたないことを示した。§ 2 の関数式領域は代数的整閉であるから、部分集合 $\mathfrak M(\mathfrak H)$ も本質的に異なる領域 $\mathfrak H$ を含む。

いま、基底構成と同様に、$\mathfrak M(\mathfrak G)$ から部分集合 $\mathfrak L(\mathfrak G)$ を取り出す方法を示す。この部分集合の中のすべての領域 $\mathfrak G$ は本質的に異なり、一方で $\mathfrak M(\mathfrak G)$ の任意の領域、したがって任意の領域 $\mathfrak G$ は $\mathfrak L(\mathfrak G)$ のある領域に同型である。$\mathfrak M(\mathfrak G)$ を整列 $\Omega$ にかける。すると $\mathfrak M(\mathfrak G)$ の領域 $\mathfrak G$ は、先行するもののいずれかに同型である場合も、そうでない場合もある。先行するどの領域にも同型でない領域 $\mathfrak G$ は集合 $\mathfrak L(\mathfrak G)$ をなし、構成により、これは本質的に異なる領域 $\mathfrak G$ だけを含む。帰納法の議論はまた、$\mathfrak M(\mathfrak G)$ の任意の領域が $\mathfrak L(\mathfrak G)$ のある領域に同型であることを示す。もしこれが失敗する最初の領域が $\mathfrak G_1$ であったなら、$\mathfrak G_1$ は $\mathfrak L(\mathfrak G)$ に属するか、または先行する領域 $\mathfrak G_0$ に同型である。後者は仮定により $\mathfrak L(\mathfrak G)$ のある領域に同型である。したがって $\mathfrak G_1$ についても同じことが成り立つはずである。任意の領域 $\mathfrak G$ は何らかの集合 $\mathfrak M^*(\mathfrak G)$ に属し、したがって $\mathfrak M(\mathfrak G)$ のある領域に同型であるから、\emph{領域 $\mathfrak G$ の類全体は $\mathfrak L(\mathfrak G)$ によって代表される。}

$\mathfrak L(\mathfrak G)$ の領域 $\mathfrak G$ において、基底量 $\eta$ を別の基底の $\eta^*$ で置き換えると、上で示したように $\mathfrak G$ に同型な領域が生じる。逆に、$\mathfrak G^*$ を任意の領域とし、それが $\mathfrak L(\mathfrak G)$ の領域 $\mathfrak G$ に同型であると仮定する。基底量 $\eta$ が $\mathfrak G^*$ の量 $\eta^*$ に対応するなら、これらは再び代数的基底をなし、$\mathfrak G^*$ は、$\mathfrak G$ が $\eta$ の関数として含むのと同じ代数関数、またはその共役を、$\eta^*$ の関数として含む。\emph{したがって $\mathfrak G$ に同型な任意の領域は、代数的基底 $H$ を $\mathfrak G$ の中、および $(H)$ に関するその共役の中で可能なすべての代数的基底に走らせることによって生じる}\footnote{$\mathfrak G$ に共役なこれらの領域の元は、II により確かに $\mathfrak G$ の元によって有理的に表せるが、必ずしも整かつ有理的に表せるわけではない。したがって共役は $\mathfrak G$ と異なり得る。}\emph{---それらが $\mathfrak G$ と異なる場合には。} $\mathfrak L(\mathfrak G)$ の固定された各領域 $\mathfrak G_i$ から、このようにして類 $\mathfrak R_i$ が得られる。これは $\mathfrak G_i$ に同型なすべて、かつそれだけの領域 $\mathfrak G$ を含む。ただし、各領域が何度も、あるいは無限にしばしば現れる可能性がある。$\mathfrak R_i$ からは、上の手続きにより、$\mathfrak G_i$ に同型な各領域を一度だけ含む部分集合 $\mathfrak R_i^*$ を再び取り出すことができる。$\mathfrak R_i^*$ がすべての類を走ると、すべての領域 $\mathfrak G$ の全体が、各領域一度かぎりで得られる。類 $\mathfrak R_i$ を特徴づけるには、もちろん $\mathfrak G_i$ の代わりに、その類の任意の他の領域を用いることができ、その領域から上のように $\mathfrak R_i$ のすべての領域を導ける。要約すると次の通りである。

\emph{$\mathfrak M(\mathfrak G)$ から部分集合 $\mathfrak L(\mathfrak G)$ を取り出すことができ、その中の領域 $\mathfrak G$ は類全体と一対一に対応する（重複なし）。$\mathfrak G_i$ から、対応する類 $\mathfrak R_i$ は、$H$ を $\mathfrak G_i$ の中および $(H)$ に関するその共役の中で可能なすべての代数的基底に走らせることによって生じる。$\mathfrak R_i^*$ が $\mathfrak R_i$ の互いに異なるすべての領域を表し、$\mathfrak R_i^*$ がすべての類を走ると、すべての領域 $\mathfrak G$ の全体が各領域一度かぎりで生じる。}\footnote{ここで $\mathfrak M(\mathfrak G)$ は §§ 6, 7 に従って構成されるべきである。§ 7 により代数的基底 $H$ を、それから生じるすべての付帯条件付き有理基底 $\Theta$ へ拡張し、各 $\Theta$ に対して § 6 b) により、$\Theta$ を含むすべての領域 $\mathfrak G$ の全体 $\mathfrak N(\mathfrak G)$ を形成する。}

代数的に整な超越数から成る領域 $\mathfrak H$ についても類似の主張が成り立つ。

\section*{§ 9. 領域 $\mathfrak G$ の可算無限個の類の例。}

§ 4 b) と § 6 b) で与えた構成の濃度は、それぞれの注から、すべての整列の濃度、したがって連続体の濃度を $c$ とすれば $2^c$ に等しいと読み取れる。しかしそこから、一度だけ数えたすべての領域 $\mathfrak G$ の全体の濃度について、ましてすべての類の濃度について、何らかの結論を引き出すことはできない。類の数が有限でないことは、§ 5 のものに類似する可算無限個の領域 $\mathfrak G$ の例によって示す。これらはすべて本質的に異なることが分かり、したがって可算無限個の類を与える。\footnote{§ 6 の領域 $\mathfrak G$ も同様に可算無限個の類を与えるが、証明はやや複雑である。}

§ 5, (1) と類似に、領域 $\mathfrak G_\sigma$ を量
\begin{equation}
 z=a_0+a_1(\eta)+\cdots+a_{\sigma-1}(\eta)
 \pmod{\mathfrak M_\sigma(\eta)}
\tag{1}
\end{equation}
の全体から成るものとする。\emph{ここで $a_i(\eta)$ は $\eta$ に関する $i$ 次元の同次形式を表す。また加群 $\mathfrak M_\sigma$ は、$\eta$ に関する $\sigma$ 次元のすべての冪積を基底として含み\footnote{もちろん、個々の $z$ の中では、加群において有限個の $\eta$ の冪積だけが現れる。}、$\eta$ のすべての代数的整関数を乗数として含む。} 二つの領域 $\mathfrak G_\sigma$ と $\mathfrak G_\tau$ $(\sigma\ne\tau)$ の間の同型関係のもとで、加群 $\mathfrak M_\sigma$ と $\mathfrak M_\tau$ は必ず同型に対応しなければならないことを示す。すると不可能性は容易に従う。

$\mathfrak G_\tau$ の不定元を $\xi$ と書く。同型関係のもとで、これら $\xi$ に $\mathfrak G_\sigma$ の量 $f(\eta)$ が対応するとする。同型関係は商を取ることのもとで保たれるので、$\xi$ のすべての代数的整関数から成る領域 $\mathfrak G_\xi$ は、$f(\eta)$ 上代数的に整なすべての関数から成る代数的整閉領域 $\mathfrak T$ に対応し、したがって $\eta$ においても代数的に整である。

(1) が $\mathfrak M_\tau$ の量に対応すると仮定する。すると $\mathfrak M_\tau$ の加群性により、(1) に $\mathfrak T$ の任意の量 $\psi(\eta)$ を掛けたものも $\mathfrak G_\sigma$ に属さなければならない。すなわち式で書けば
\begin{equation}
(a_0+a_1(\eta)+\cdots+a_{\sigma-1}(\eta))\psi(\eta)
\equiv b_0+b_1(\eta)+\cdots+b_{\sigma-1}(\eta)
\pmod{\mathfrak M_\sigma}
\tag{2}
\end{equation}
である。すべての $\eta$ を 0 に等しく置くと $a_0\psi(0)=b_0$ となり、(2) は
\begin{equation}
(a_0+a_1(\eta)+\cdots+a_{\sigma-1}(\eta))(\psi(\eta)-\psi(0))
\equiv c_1(\eta)+\cdots+c_{\sigma-1}(\eta)
\pmod{\mathfrak M_\sigma}
\tag{3}
\end{equation}
となる。$\psi(\eta)$ のほかに、関数
\[
 \chi_\nu(\eta)=\sqrt[\nu]{\psi(\eta)-\psi(0)}
\]
もすべての $\nu$ について $\mathfrak T$ に属する。また $\chi_\nu(0)=0$ であるから、(3) と類似の式は
\begin{equation}
(a_0+a_1(\eta)+\cdots+a_{\sigma-1}(\eta))\chi_\nu(\eta)
\equiv c_1^{(\nu)}(\eta)+\cdots+c_{\sigma-1}^{(\nu)}(\eta)
\pmod{\mathfrak M_\sigma}
\quad(\nu=1,2,\ldots)
\tag{4}
\end{equation}
である。$\chi_1=\psi(\eta)-\psi(0)$ とその $\kappa-1$ 個の共役との積を $A(\eta)$ とし、その次数を $\lambda$ とする。$\chi_1(0)=0$ なので、$A(\eta)$ は少なくとも一次元の項から始まらなければならない。(4) から
\begin{equation}
a_0\chi_\nu(\eta)\equiv0\pmod{\mathfrak M_1};\qquad
 a_0^\nu\chi_1(\eta)\equiv0\pmod{\mathfrak M_\nu};\qquad
 a_0^{\nu\kappa}A(\eta)\equiv0\pmod{\mathfrak M_{\nu\kappa}}
\tag{5}
\end{equation}
を得る。したがって $a_0\ne0$ なら、§ 3 と同様に $\lambda>\nu\kappa$ である。しかし $\nu<\frac{\lambda}{\kappa}$ は $\mathfrak T$ の代数的整閉性に矛盾するので、必ず $a_0=0$ である。いま (4) から
\begin{equation}
a_1(\eta)\chi_\nu(\eta)\equiv c_1^{(\nu)}(\eta)\pmod{\mathfrak M_2};\qquad
[a_1(\eta)]^{\nu\kappa}A(\eta)
\equiv[c_1^{(\nu)}(\eta)]^{\nu\kappa}\pmod{\mathfrak M_{\nu\kappa+1}}
\tag{6}
\end{equation}
を得る。ここから、$A(\eta)$ が少なくとも一次元の項から始まるので、$c_1^{(\nu)}(\eta)=0$ である。したがって (5) と完全に同様に、すべての $\nu$ について
\[
 a_1(\eta)\chi_\nu(\eta)\equiv0\pmod{\mathfrak M_2};\qquad
 [a_1(\eta)]^{\nu\kappa}A(\eta)\equiv0\pmod{\mathfrak M_{2\nu\kappa}}
\]
となり、従って $a_1(\eta)=0$ である。このように続けると、(5) と (6) を交互に適用して
\[
 a_0=0,\qquad a_1(\eta)=0,\quad \ldots,\quad a_{\sigma-1}(\eta)=0
\]
を得る。$\mathfrak G_\tau$ に関しても同様の考察が成り立つから、\emph{$\mathfrak G_\sigma$ と $\mathfrak G_\tau$ $(\sigma\ne\tau)$ の任意の同型関係のもとで、加群 $\mathfrak M_\sigma$ と $\mathfrak M_\tau$ は同型に対応する。}

例えば $\sigma>\tau$ と仮定する。不定元 $\xi$ には $\mathfrak G_\sigma$ の量 $f(\eta)$ が対応していた。そして $\mathfrak G_\tau$ のすべての量は $\xi$ の多項式 mod. $\mathfrak M_\tau$ によって表されるから、$\mathfrak M_\tau$ と $\mathfrak M_\sigma$ の対応は、$\mathfrak G_\sigma$ のすべての量が $f(\eta)$ の多項式 mod. $\mathfrak M_\sigma$ として得られることを意味する。$\sigma>\tau\ge1$ であるから、不定元 $\eta$ は確かに $\mathfrak M_\sigma$ に属さない。したがって、$\eta$ が $f(\eta)$ によって mod. $\mathfrak M_\sigma$ で整かつ有理的に表せる以上、線形項が消えない $f(\eta)$ が存在しなければならない。そこで
\[
 \xi:f(\eta)\equiv a_0+a_1(\eta)\pmod{\mathfrak M_2}\quad[a_1(\eta)\ne0],
\]
\[
 \xi^\tau:[f(\eta)]^\tau\equiv a_0^\tau\pmod{\mathfrak M_1}\quad\text{$a_0\ne0$ のとき},
\]
\[
 \hphantom{\xi^\tau:[f(\eta)]^\tau}\equiv [a_1(\eta)]^\tau
 \pmod{\mathfrak M_{\tau+1}}\quad\text{$a_0=0$ のとき}
\]
とする。いま $\xi^\tau$ は $\mathfrak M_\tau$ に属するが、$[f(\eta)]^\tau$ は 0 次または $\tau$ 次の消えない項から始まり、$\tau<\sigma$ であるため $\mathfrak M_\sigma$ に属し得ない。これは上で導いた $\mathfrak M_\sigma$ と $\mathfrak M_\tau$ の対応に反する。$\tau>\sigma$ の場合は $\eta$ と $\xi$ を入れ替えれば同時に扱われる。したがって、領域 $\mathfrak G_\sigma$ と $\mathfrak G_\tau$ は\emph{同型でない、すなわち本質的に異なる}ことが示された。ゆえに\emph{領域 $\mathfrak G_i$ $(i=1,2,\ldots\text{ in inf.})$ は、領域 $\mathfrak G$ の可算無限個の類の例を与える。}

\emph{注。} 任意の二つの領域 $\mathfrak G_i$ は非同型であるが、二つの領域について、それぞれが他方の一部に同型であるということは十分起こり得る。\footnote{この振る舞いは体にも起こり得る。Steinitz, loc. cit., 結論を参照。} 例えば $\tau=1$ とし、$\sigma$ を任意とする。
\[
\mathfrak G_1:a_0+\mathfrak M_1(\xi);\qquad
\mathfrak G_\sigma:a_0+a_1(\eta)+\cdots+a_{\sigma-1}(\eta)+\mathfrak M_\sigma(\eta).
\]
対応 $\xi=\eta$ は $\mathfrak G_\sigma$ を $\mathfrak G_1$ の真の因子にする。逆に、$\mathfrak G_1$ と $\mathfrak G_\sigma$ の間の一対一対応 $\xi=\eta^\sigma$, $\eta=\sqrt[\sigma]{\xi}$ は、$\mathfrak G_1$ を $\mathfrak G_\sigma$ の真の因子にする。したがって各 $\mathfrak G_i$ はそれ自身の真の因子にも同型である。交わりの概念とは対照的に、交わり類、すなわち、その各領域が他のすべての類のすべての領域の真の因子に同型であるような類、という概念は存在しない。少なくとも一意に定まる概念としては存在しない。

\section*{§ 10. 任意の体の整な量。}

前節までの展開は、領域 $\mathfrak G$ については、領域 $\mathfrak G$ の定義性質 III と IV に代数的数を含めないなら、全複素数の全体が体をなすという事実だけに依拠していた。\footnote{§ 7 でのみ、有理数の「素体」が無限個の元を含むという仮定をさらに用いた。しかし有限個の元だけをもつ素体の場合、この段落は、これから見るように、単に消える。} 領域 $\mathfrak H$ については、さらにこの体が代数的閉であるという事実に依拠していた。したがって得られた結果は、Weber および E. Steinitz\footnote{loc. cit., 序論。} に従って「二つの演算、加法と乗法をもつ元の系であり、それらが結合律と交換律に従い、分配律で結びつき、逆演算を無制限かつ一意に許すもの」\footnote{零による除法だけは除かなければならない。}として抽象的に定義される\emph{最も一般の体}へ直接一般化できる。

そのような体はすべて単位元を含み、したがって単位元から加法、乗法およびその逆演算によって導かれる「素体」を含む。この素体は、有理数の型（標数 $0$）であるか、または素数 $p$ の剰余類系の型（標数 $p$）である。\footnote{Steinitz, § 4.} そのとき\emph{素体の整な量}は、単位元から加法および減法によって生じる量としてよく定義される。標数 $p$ の素体では、整な量はすでに素体を尽くし、素体の分数的量は存在しない。任意の代数的閉体は「絶対代数体」を含む。これは素体に、素体上代数的なすべての元を添加して得られる。したがって標数 $0$ の体では、これは全代数的数の体の型である。すると\emph{絶対代数体の代数的に整な量}は、素体の整な量（または同じことだが単位元）上代数的に整なすべての量としてよく定義される。標数 $p$ の体では、したがって絶対代数体の代数的に分数的な量も存在しない。これらの準備的注意の後、「整」量および「代数的に整」量の定義を任意の体、あるいは代数的閉体に移すことができる。\emph{体 $\mathfrak R$ の整な量とは、$\mathfrak R$ の量から成る整域 $\mathfrak G$ で、その商体\footnote{すなわち、$\mathfrak G$ の量の対の商すべてから成る体。}が $\mathfrak R$ を尽くし、単位元を含むが素体の分数的量を含まないものとして定義される。}

\emph{代数的閉体 $\mathfrak R$ の代数的に整な量とは、$\mathfrak R$ の量から成る代数的整閉な整域 $\mathfrak H$ で、その商体が $\mathfrak R$ を尽くし、単位元を含むが絶対代数体の代数的に分数的な量を含まないものとして定義される。}

標数零の体については、前節までで得られた結果がそのまま適用される。ただし「代数的数」を、素体、あるいは絶対代数体の量で置き換えるだけでよい。標数 $p$ の体については、素体が分数的量を含まないため、結果はさらに単純化する。すなわち有理基底および添加される系に対する付帯条件が単に消える。\footnote{この段落の第一の注を参照。} したがって結果は次の通りである。\emph{抽象的定義は、標数 $p$ の体の整な量として、任意の有理基底から導かれるすべての整域、体自身、およびこれらの領域と体との間にある任意の整域を許す。対応して、素標数の代数的閉体の代数的に整な量としては、代数的基底から導かれる領域と体自身との間にある任意の代数的整閉な整域、二つの境界領域を含むものが得られる。} したがって素標数の最も一般の体では、対応する素体の場合と同様に、整と分数との区別は消える。

\medskip
\noindent Erlangen, 1915 年 3 月 30 日。

\clearpage
\section*{10. 同型写像の関数方程式}
\begin{center}
Math. Ann. 77 (1916), pp. 536--545
\end{center}

Dedekind によれば、\footnote{例えば Dirichlet-Dedekind, \emph{Vorlesungen über Zahlentheorie}（第4版）, Supplement XI, \S\ 161 を参照。群論にならった ``isomorphic mapping'' または ``isomorphism'' という名称は後の文献にはじめて現れる。Dedekind は「体の置換」と呼んでいる。}\emph{体} $\mathfrak A$ \emph{からある系} $\mathfrak B$ \emph{への同型写像}とは、次のような一価の対応をいう。すなわち、$\mathfrak A$ の各元に対して $\mathfrak B$ の元がただ一つ対応し、しかも $\mathfrak A$ の任意の二元の和、差、積、商には、それぞれ一意に対応する元の和、差、積、商が再び対応する、という対応である。このとき $\mathfrak B$ は後から同じく体であることが分かる。

そのような写像が、上の意味で一価な関数 $f(z)$ によって媒介されるものとする。$z$ が体 $\mathfrak A$ のすべての元 $x,y,\ldots$ を走るとき、$f(z)$ は次の関数方程式によって特徴づけられる。
\[
\begin{array}{ll}
(1)\quad f(x+y)=f(x)+f(y); & (2)\quad f(x-y)=f(x)-f(y);\\[0.45em]
(3)\quad f(x\cdot y)=f(x)\cdot f(y); & (4)\quad f\!\left(\dfrac{x}{y}\right)=\dfrac{f(x)}{f(y)}.
\end{array}
\]

\emph{以下では、} $x,y,\ldots$ \emph{がすべての実数および複素数を走るときの、これらの関数方程式の最も一般の一価解} $f(z)$ \emph{を与える。すなわち、全複素数体の同型写像を与える最も一般の関数} $f(z)$ \emph{を与える。}\footnote{この問題は Landau 氏から折にふれて私に提示された。後に気づいたことだが、解の一部はすでに H. Lebesgue, \emph{Sur les transformations ponctuelles, transformant les plans en plans}, Atti Acad. Torino, 1906/07 に現れており、また A. Ostrowski, \emph{Über einige Fragen der allgemeinen Körpertheorie}, Crelle's Journal 143 (1913), \S\ 2 にも暗黙に含まれている。} 任意の体に対する解もまったく同様に得られる。

G. Hamel\footnote{G. Hamel, \emph{Eine Basis aller Zahlen und die unstetigen Lösungen der Funktionalgleichung:} $f(x+y)=f(x)+f(y)$, Math. Ann. 60, p. 459 (1905).} は、すべての数の\emph{線形}基底によって関数方程式 (1) の最も一般の解を構成した。ここで (1) から (4) の解、したがって有理演算および代数演算が不変となる写像関数に対して与える構成は、すべての数の\emph{有理}基底および\emph{代数的}基底を用いる。1 で同型写像をより詳しく論じた後、2 で $f(z)$ に対する必要条件を与え、3 でそれらが十分条件でもあることを認め、それによって実際の構成に至る。\footnote{私の論文 \emph{Die allgemeinsten Bereiche aus ganzen transzendenten Zahlen}, Math. Ann. 77, p. 103 (1915), とくに \S\S\ 8, 10 の平行した考察を参照。} 既知の例外 $f(z)=z$ または $\overline z$ を除いて、(1) から (4) の解は不連続である。4 ではさらに、それらが「極端に不連続」であることを示す。これは G. Hamel が (1) の実解について証明した事実であるが、複素領域では必ずしも満たされない。\footnote{Hamel の用語「全く不連続」は、他の文献では別の意味で用いられている。}

\subsection*{1.}

まず、関数方程式 (1) から (4) からいくつかの帰結を引き出す。これは一般の体 $\mathfrak A$ に対して直接行う。(1) から、一価性により $f(0)=0$ が従う。$f(y)$ が消えなければ、元の $y$ も消えることはできない。したがって関数方程式 (1) から (4) は、\emph{すべての値} $f(z)$ \emph{から成る像の系} $\mathfrak B$ \emph{が再び体をなす}ことを示す。逆に、初期条件 $f(0)=0$ を課せば、(2) から関数 $f(z)$ の一価性が従う。したがって、\emph{一価性の要求と初期条件} $f(0)=0$ \emph{とは同値である}。(4) から $f(z)$ は恒等的には消えないことが従い、さらに (3) から $f(z)$ は $z=0$ に対してだけ消える。実際、$f(y)=0$ かつ $y\ne0$ ならば、$x/y$ も $\mathfrak A$ に属するので
\[
f(x)=f\!\left(\frac{x}{y}\cdot y\right)=f\!\left(\frac{x}{y}\right)\cdot f(y)=0
\]
となり、$f(z)$ が恒等的に消えることになるが、これは (4) に反する。したがって $f(x)=f(y)$ から $x=y$ が従う。すなわち $f(z)$ の各値にも $z$ のただ一つの値が対応する。\emph{$f(z)$ は} $z$ \emph{の一対一対応関数である}。関数方程式が示すように、逆関数によって媒介される写像もまた同型である。すべての有理演算は有限個の加法、乗法およびそれらの逆演算から構成されるから、次のようにも言える。

\emph{同型写像によって、$\mathfrak A$ の元と $\mathfrak B$ の元との間に一対一の関係が確立され、その関係において、$\mathfrak A$ の有限個の元の間に成り立つすべての有理関係
\[
F(x,y,\ldots,t)=0
\]
は $\mathfrak B$ において保存され、逆もまた成り立つ。}

さらに注意しておくべきことは、体に対しては、$f(z)$ はすでに関数方程式 (1) と (3)、$f(z)$ が恒等的には消えないという要求、および初期条件 $f(0)=0$ によって特徴づけられる、ということである。実際、(1) において $x$ を $\mathfrak A$ に属する元 $(x-y)$ で置き換えれば (2) が従う。そして (2) と $f(0)=0$ から $f(z)$ の一価性が得られる。上で示したように、さらに (3) において $x$ を $\mathfrak A$ に属する元 $x/y$ で置き換えれば、$y\ne0$ のとき $f(y)$ は消えないことが従う。したがって (4) の妥当性と同時に一対一性が得られる。体の代わりに整域 $\mathfrak J$ を基礎に置くならば、$x/y$ は $\mathfrak J$ に属するとは限らず、(3) から一対一性を結論することはできない。しかしここでは、次の最もよく知られた定式化で十分である。\emph{同型写像とは、$\mathfrak J$ と $\mathfrak J'$ の元の間の一対一対応であり、和と積には常に対応する元の和と積が対応するような対応である。}\footnote{第1項については Dedekind, loc. cit., \S\ 161 を参照。}

\subsection*{2.}

以下では簡単のため、$\mathfrak A$ の代わりに、すべての実数および複素数から成る体 $\mathfrak C$ を基礎に置く。まず問題となるのは、$f(z)$ がとる値の系の考察である。(3) または (4) から $f(1)=1$ が従い、さらに関数方程式によって、すべての\emph{有理数} $\alpha$ について $f(\alpha)=\alpha$ が従う。したがって、写像のもとで各有理数は自分自身に対応する。\emph{代数的数} $\beta$ は有理数係数の既約方程式 $F(\beta,\alpha)=0$ を満たすから、1 と先の注意により、$f(\beta)$ についても $F(f(\beta),\alpha)=0$ が成り立つ。したがって、\emph{各代数的数は自分自身またはその共役の一つへ移る}。また一対一性により、代数的数全体には再びすべての代数的数から成る体が対応する。

\emph{超越数}の挙動を認識し、同時に共役値を分離するために、\emph{すべての数の有理基底}\footnote{「整超越数」に関する私の引用論文 \S\ 6 を参照。}を基礎に置く。すなわち、すべての数を有理的に、つまり有理数係数で\footnote{「有理関数」とは、常に係数も有理数であるものを意味する。}表すことを可能にする数 $\vartheta$ の系 $\Theta$ であって、少なくとも一つの整列のもとで、どの基底数も有限個の先行する基底数によって有理的に表せないような系である。この基底の存在は、線形基底および代数的基底の場合とまったく同様に整列定理から従う。

$\Theta$ に整列 $\Omega$ を与える。すると各数は、有限個の先行する数によって有理的に表せるか、または先行するものから有理的に独立である。後者の数が基底 $\Theta$ をなし、帰納法により、実際にすべての数が有限個の $\vartheta$ によって有理的に表せることが分かる。

各基底数 $\vartheta$ に対して「切片体 $\mathfrak K(\vartheta)$」を定義する。これは整列において $\vartheta$ に先行する基底数のすべての有理結合から成る。最初の基底数の切片体としては、すべての有理数から成る体 $\mathfrak R$ を考える。$\vartheta$ は $\mathfrak K(\vartheta)$ に関して二種類の挙動を示し得る。すなわち、$\mathfrak K(\vartheta)$ に代数的に従属であるか、代数的に独立であるかである。1 により、$\mathfrak C$ で成り立つすべての関係は像の体でも成り立ち、逆もまた成り立つから、$\vartheta$ に対応する値 $f(\vartheta)$ の全体は像領域の基底をなし、$\Theta$ 自身の整列によって整列される。したがって、特に切片体 $\mathfrak K(\vartheta)$ には切片体 $\mathfrak K(f(\vartheta))$ が対応する。そして $\vartheta$ が $\mathfrak K(\vartheta)$ に代数的に従属または独立であるに従い、$f(\vartheta)$ も $\mathfrak K(f(\vartheta))$ に関して同じ性質をもつ。従属の場合には、$\vartheta$ と $f(\vartheta)$ の既約方程式も対応する。

それぞれ $\mathfrak K(\vartheta)$ から代数的に独立である値 $\vartheta$ の全体は、再び帰納法が示すように、\emph{すべての数の代数的基底}\footnote{E. Zermelo, \emph{Über ganze transzendente Zahlen}, \S\ 1, Math. Ann. 75 (1914) を参照。}をなす。すなわち、すべての数を代数的に表すことを可能にし、しかもそれら自身の間には代数関係が存在しないような数 $\eta$ の系 $H$ である。この代数的基底 $H$ には、像領域において再び代数的基底 $Z$ が対応する。一対一性により、この $Z$ は $H$ と同じ濃度をもつ。すなわち連続体の濃度である。代数拡大によって濃度は増加しないことが知られているからである。\footnote{E. Steinitz, \emph{Algebraische Theorie der Körper}, Crelle's Journal 137 (1910), \S\ 24 を参照。}

しかし $f(\vartheta)$ の値が分かれば、すべての値の系に対して $f(z)$ はただちに得られる。なぜなら、$z=\Psi(\vartheta)$ から常に $f(z)=\Psi(f(\vartheta))$ が従うからである。

\subsection*{3.}

いま、上で得られた必要条件が実際に $f(z)$ の構成も与えることを示す。代数的基底 $H$ に一対一に対応する基底 $Z$ を、次のように構成して $H$ に対応させる。すべての数の代数的基底 $Z'$ を与える第二の整列 $\Omega'$ を置く。$Z$ を $Z'$ の部分集合とし、その元を $\zeta$ とする。$Z$ は $Z'$ と同じ濃度、したがって $H$ と同じ濃度をもつものとする。$H$ の各元 $\eta_i$ に、同じ添字をもつ一つの $\zeta_i$ を一対一に割り当てる。その規則は、$\zeta_i$ が整列 $\Omega'$ において、$\eta_i$ に先行する $H$ のどの元にもまだ割り当てられていない $Z$ の元のうち最初のものである、というものである。

ここで $f(z)$ を次のように定義する。
\begin{enumerate}
\item[a)] すべての有理数 $\alpha$ に対して、$f(0)=0$, $f(\alpha)=\alpha$ とする。
\item[b)] 代数的基底 $H$ のすべての量に対して、$f(\eta_i)=\zeta_i$ とする。
\item[c)] 有理基底 $\Theta$ の残りのすべての量 $\vartheta$、すなわち切片体 $\mathfrak K(\vartheta)$ に代数的に従属し、したがって $\mathfrak K(\vartheta)$ で既約な方程式 $F(\vartheta;\vartheta_1,\ldots,\vartheta_r)=0$ を満たすものについては、$F(f(\vartheta);f(\vartheta_1),\ldots,f(\vartheta_r))=0$ の根として定める。
\item[d)] 残りのすべての値 $z$、すなわち $\Theta$ の定義により $z=\psi(\vartheta_{k_1},\ldots,\vartheta_{k_r})$ と表せるものについては、$f(z)=\psi(f(\vartheta_{k_1}),\ldots,f(\vartheta_{k_r}))$ とする。\footnote{(d) には (a) が特殊な場合として含まれる。}
\end{enumerate}

このように定義された $f(z)$ が実際に関数方程式 (1) から (4) を満たすことを示すには、第1項により、(1) と (3) について示せば十分である。なぜなら $\mathfrak C$ は体であり、$f(z)$ は (a) と (b) により恒等的には消えず、さらに初期条件 $f(0)=0$ を満たすからである。

有理基底 $\Theta$ によって $x$, $y$, $x+y$ を
\[
x=\psi_1(\vartheta_1,\ldots,\vartheta_t),\quad
 y=\psi_2(\vartheta_1,\ldots,\vartheta_t),\quad
 x+y=\psi_3(\vartheta_1,\ldots,\vartheta_t)
\]
と表す。このとき、$f(z)$ が関数方程式 (1)
\[
f(x)+f(y)-f(x+y)=0
\]
を満たすことの証明は、(d) により、
\[
\psi_1(\vartheta)+\psi_2(\vartheta)-\psi_3(\vartheta)=\frac{H_1(\vartheta)}{H_2(\vartheta)}=0
\]
から常に
\[
\psi_1(f(\vartheta))+\psi_2(f(\vartheta))-\psi_3(f(\vartheta))
=\frac{H_1(f(\vartheta))}{H_2(f(\vartheta))}=0
\]
が従うことと同一である。ここで $H_1,H_2$ はその引数の整有理関数を意味する。関数方程式 (3) もまったく同じ形の条件に導く。したがって、すべての関数方程式が満たされることは、次の条件が成り立つことと同一である。この条件は第1項および第2項により必要でもある。すなわち、すべての整有理関数 $H(\vartheta)$ について、$H(\vartheta)=0$ から常に $H(f(\vartheta))=0$ が従い、$H(\vartheta)\ne0$ から $H(f(\vartheta))\ne0$ が従う、という条件である。後者は分母が現れるために必要である。

これが実際に成り立つことは帰納法で示される。例えば $H(\vartheta_1,\ldots,\vartheta_n)$ において、基底量 $\vartheta_n$ が $\vartheta_1,\ldots,\vartheta_n$ のうち整列で最後のものであるとする。\footnote{$\vartheta$ について0次の式 $H(\vartheta)$ は写像によって自分自身に移るので、さらに調べる必要はない。} すると $H(\vartheta_1,\ldots,\vartheta_n)=0$ は、$\vartheta_n$ が切片体 $\mathfrak K(\vartheta_n)$ に関して満たす関係とみなせる。同様に $H(f(\vartheta_1),\ldots,f(\vartheta_n))=0$ は、$f(\vartheta_n)$ が $\mathfrak K(f(\vartheta_n))$ に関して満たす関係とみなせる。

もしすべての関係 $H(\vartheta)=0$ が $f(\vartheta)$ についても成り立ち、かつ逆も成り立つのでなければ、ある最初の基底量、例えば $\vartheta_n$ が存在して、$\mathfrak K(\vartheta_n)$ に関して成り立つ少なくとも一つの関係が像の体では保存されないか、またはその逆が起こるはずである。ただし、それ以前の切片体 $\mathfrak K(\vartheta_n)$ と $\mathfrak K(f(\vartheta_n))$ は同型である。

特に (a) により、$\Theta$ の最初の基底量の切片体、すなわちすべての有理数から成る体 $\mathfrak R$ は、その像の体と同型である。いま $H(\vartheta_n)$ が次の形であるとする。
\[
H(\vartheta_n)=\vartheta_n^{\mu}+a_1(\vartheta)\vartheta_n^{\mu-1}+\cdots+a_\mu(\vartheta),
\]
ここで $a_i(\vartheta)$ は $\mathfrak K(\vartheta_n)$ の量である。このとき二つの場合がある。

1) $\vartheta_n$ が $\mathfrak K(\vartheta_n)$ から代数的に\emph{独立}である場合。このとき $H(\vartheta_n)=0$ は $\vartheta_n$ について恒等的に満たされる。\footnote{$x,y,\ldots$ が $\vartheta$ によって表せるため、このような形の関係は十分起こり得る。} したがってすべての係数について $a_i(\vartheta)=0$ であり、$\mathfrak K(\vartheta_n)$ と $\mathfrak K(f(\vartheta_n))$ の同型性により、$a_i(f(\vartheta))=0$ も成り立つ。ゆえに $H(\vartheta_n)=0$ から常に $H(f(\vartheta_n))=0$ が従う。逆に $H(f(\vartheta_n))=0$ とする。$\vartheta_n$ は $\mathfrak K(\vartheta_n)$ から代数的に独立であるので代数的基底 $H$ に属し、したがって $f(\vartheta_n)$ は (b) により代数的基底 $Z$ に属し、$\mathfrak K(f(\vartheta_n))$ から代数的に独立である。したがってすべての $a_i(f(\vartheta))$ が消え、同型性によりすべての $a_i(\vartheta)$ が消える。ゆえに $H(f(\vartheta_n))=0$ から $H(\vartheta_n)=0$ が従い、また $H(\vartheta_n)\ne0$ から $H(f(\vartheta_n))\ne0$ が従う。

2) $\vartheta_n$ が $\mathfrak K(\vartheta_n)$ に代数的に\emph{従属}する場合。このとき $H(\vartheta_n)$ は、$\mathfrak K(\vartheta_n)$ に関して既約な方程式 $F(\vartheta_n)$ によって割り切れる。すなわち、$t$ について恒等的に
\[
H(t;a_i)=F(t;\vartheta_n)\cdot Q(t;a_i)
\]
が成り立つ。ところが現れるすべての係数は $\mathfrak K(\vartheta_n)$ に属するので、同型な体 $\mathfrak K(f(\vartheta_n))$ において対応する関係
\[
H(t;a_i(f))=F(t;f(\vartheta_n))\cdot Q(t;a_i(f))
\]
が成り立つ。しかも (c) により、$f(\vartheta_n)$ は $F(t;f(\vartheta_n))$ の根となるように選ばれていた。したがって $H(\vartheta_n)=0$ から $H(f(\vartheta_n))=0$ が従う。

逆に $H(f(\vartheta_n))=0$ であるとする。このとき、$H(t;a_i(f(\vartheta)))$ が $\mathfrak K(f(\vartheta_n))$ に関して既約な関数 $F(t;f(\vartheta_n))$ によって割り切れることから、同型な体 $\mathfrak K(\vartheta_n)$ に対して、$H(t;a_i(\vartheta))$ が $F(t;\vartheta_n)$ によって割り切れることが同様に従う。そして $F(t;f(\vartheta_n))$ は $\vartheta_n$ の既約方程式から同型関係によって生じたものであり、この関係は一対一であるから、$F(t;\vartheta_n)$ は必然的に $\vartheta_n$ を根にもつ既約関数を表す。したがって $H(f(\vartheta_n))=0$ から $H(\vartheta_n)=0$ が従い、また $H(\vartheta_n)\ne0$ から $H(f(\vartheta_n))\ne0$ が従う。

したがって、同型な切片体 $\mathfrak K(\vartheta_n)$ と $\mathfrak K(f(\vartheta_n))$ から、$\vartheta_n$ および $f(\vartheta_n)$ をそれぞれ添加することによって、やはり同型な体が生じる。これがある $\vartheta_n$ について成り立たないという仮定は矛盾に導く。

したがって、(a) から (d) によって定義された関数 $f(z)$ は、実際に関数方程式 (1) から (4) の解を与え、しかも 2 により最も一般の解を与えることが証明された。

$f(z)$ の構成に用いられたすべての考察は、すべての数の全体 $\mathfrak C$ が体をなすという事実だけを用いている。したがって、それらは任意の抽象的に定義された体 $\mathfrak A$ に対して成り立つ。ただし、$\mathfrak R$ という有理数体に対応するものは $\mathfrak A$ の「素体」である。これは、$\mathfrak A$ の単位元から加法と乗法およびそれらの逆演算によって導かれる体である。したがって (a) から (d) において有理数を素体の元で置き換えれば、そのようにして得られる関数 $f(z)$ は、\emph{任意の抽象的に定義された体の最も一般の同型写像}を与える。

\subsection*{4.}

最後に、すべての実数および複素数から成る体 $\mathfrak C$ に対して定義された関数 $f(z)$、すなわち $f(z)=z$ または $\overline z$ を除いて不連続であることが知られている関数が、\emph{極端に不連続}であることを示す。これは、任意の実数または複素数 $z_0$ の任意の近傍に、$f(z)$ が任意に指定された値 $Z_0$ にいくらでも近づくような値 $z$ が存在する、という意味である。

この極端な不連続性は、G. Hamel が前掲箇所で証明したように、すべての実数について定義された関数方程式 (1) の実解 $\varphi(x)$ に属する。しかし、下に挙げる例が示すように、複素数上で定義された解には必ずしも属さない。Hamel は次のように論じる。$\varphi(x)$ が $C\cdot x$ と異なるならば、線形基底の値\footnote{すべての（実）数の線形基底は、すべての（実）数を有理係数で線形に表せるような（実）数の系であって、有限個の基底数の間には有理係数の線形関係が存在しないような系として与えられる。}のうち少なくとも二つ、例えば $x_1,x_2$ が存在し、$\varphi(x_1):x_1\ne\varphi(x_2):x_2$ となる。このとき二つの方程式
\[
\alpha x_1+\beta x_2=a,
\qquad
\alpha\varphi(x_1)+\beta\varphi(x_2)=b
\]
は実数 $\alpha,\beta$ に解をもつ。したがって $a$ と $b$ は有理数 $\alpha,\beta$ によって任意に近似できる。一方、$\alpha,\beta$ が有理数であるため、$\alpha\varphi(x_1)+\beta\varphi(x_2)$ は実際に $\varphi(\alpha x_1+\beta x_2)$ を表す。この議論は複素値系では失敗する。実際、複素数上でも、(1) の解であって不連続ではあるが極端には不連続でないものを与えることができる。例えば
\[
\varphi(z)=\varphi(x+iy)=\varphi(x)+i\bigl(y+\varphi(x)\bigr),
\]
ここで $\varphi(x)$ は実数上で極端に不連続な実解を意味する。あるいは、さらに簡単に $\varphi(z)=\varphi(x)$ としてもよい。どちらの場合も $\varphi(z)$ の実部は極端に不連続であるが、虚部は実部によってともに決定される。

この挙動、そして同時に関数 $f(z)$ の挙動は、\emph{階数}の概念を導入すると完全に明らかになる。$z=x+iy$, $\varphi(z)=X+iY$ とおき、$\varphi(z)$ の階数を四、三、二、一と呼ぶのは、それぞれ、線形独立な関係
\[
c_1x+c_2y+c_3X+c_4Y=0
\]
が存在しない、または一つ、二つ、三つ存在する場合である。もちろん $c$ は実数である。

1) $\varphi(z)$ が階数四であるとする。このとき線形基底の値が少なくとも四つ、例えば $z_1,z_2,z_3,z_4$ 存在して、行列式
\[
\begin{vmatrix}
x_1&x_2&x_3&x_4\\
y_1&y_2&y_3&y_4\\
X_1&X_2&X_3&X_4\\
Y_1&Y_2&Y_3&Y_4
\end{vmatrix}
\]
が消えない。したがって方程式
\[
\begin{aligned}
\alpha x_1+\beta x_2+\gamma x_3+\delta x_4&=a,\\
\alpha y_1+\cdots+\delta y_4&=b,\\
\alpha X_1+\cdots+\delta X_4&=A,\\
\alpha Y_1+\cdots+\delta Y_4&=B
\end{aligned}
\]
は実数 $\alpha,\ldots,\delta$ に解をもち、$a,b,A,B$ は有理数 $\alpha,\beta,\gamma,\delta$ によって任意に近似できる。さらに
\[
\alpha(X_1+iY_1)+\beta(X_2+iY_2)+\gamma(X_3+iY_3)+\delta(X_4+iY_4)
=\varphi(\alpha z_1+\beta z_2+\gamma z_3+\delta z_4)
\]
である。「一般には」、したがって解 $\varphi(z)$ は複素数上でも極端に不連続である。$z_0=a+ib$ の近傍における $\varphi(z)$ の\emph{不連続値}は二次元線形多様体をなす。

2) $\varphi(z)$ が階数三であるとする。このとき線形基底の値系が少なくとも三つ、例えば $z_1,z_2,z_3$ 存在して、上の行列式の階数は三である。関係
\[
c_1a+c_2b+c_3A+c_4B=0
\]
を満たす値 $a,b,A,B$、そしてそれらだけが、任意に近似され得る。\emph{不連続値}は $z_0=a+ib$ によって定まる\emph{一次元}線形多様体をなす。挙げた例が示すように、この場合は実際に起こり得る。

3) $\varphi(z)$ が階数二であるとする。常に二つの複素数 $z_1,z_2$ を
\[
\begin{vmatrix}x_1&x_2\\ y_1&y_2\end{vmatrix}\ne0
\]
となるように選べるから、関係は
\[
X=\lambda_1x+\lambda_2y,
\qquad
Y=\mu_1x+\mu_2y
\]
の形をもつ。ここから生じる式
\[
\varphi(z)=(\lambda_1x+\lambda_2y)+i(\mu_1x+\mu_2y)
\]
は、複素数の領域における関数方程式 (1) の最も一般の\emph{連続}解を表す。

4) $\varphi(z)$ が階数一であるとする。この場合は複素数の領域では起こり得ない。実数の領域では\emph{連続}解 $\varphi(x)=C\cdot x$ を意味する。

いま、すべての実数および複素数に対して定義された関数方程式 (1) から (4) の\emph{解} $f(z)$ \emph{は、階数二または階数四しかもち得ない}ことを示す。階数二は、$f(1)=1$ および $f(i)=\pm i$ の特殊化を 3) と合わせると分かるように、二つの\emph{連続}関数 $f(z)=z$ または $\overline z$ だけを与える。階数一はすでに上で排除されたので、階数三を排除すればよい。そこで、少なくとも一つの関係
\[
c_1x+c_2y+c_3X+c_4Y=0
\]
が存在し、したがって $f(z)$ が階数三、または場合によってはより低い階数であると仮定し、この仮定のもとでは常に後者の場合が起こることを示す。再び $f(1)=1$, $f(i)=\pm i$ の特殊化から、関係は
\[
c_3(X-x)+c_4(Y\mp y)=0
\]
へ移る。ここで $c_3$ と $c_4$ は同時には消えない。まず例えば $c_4=0$ とする。このとき関係 $(X-x)=0$ は、純虚数 $z$ にも純虚数 $f(z)$ が対応することを意味する。したがって任意の実数 $y$ に対して $f(iy)=i\cdot\mathfrak Y$ と書ける。ここで $\mathfrak Y$ は再び実数である。ところが $f(iy)=\pm i f(y)$ であるから、$f(y)=\pm\mathfrak Y$ も従う。よって各実数 $z$ にも実数 $f(z)$ が対応し、$(X-x)$ によりすべての実値は自分自身へ移る。したがって $f(z)=z$ または $\overline z$ のみであり、実際には階数二である。$c_3=0$, $c_4\ne0$ の場合もまったく同様に結論される。

次に $c_3$ と $c_4$ がともに零でないとする。このとき関係は
\[
(Y\mp y)=c\cdot(X-x)
\]
の形に書ける。ここで $c$ は実数で、零でなく有限である。これは、$y=\pm cx$ に対して $Y=c\cdot X$ も成り立つことを意味する。したがって関数方程式を考慮すると、任意の実数 $r$ について
\[
f\{r\cdot(1\pm ic)\}=X(1+ic)=f(r)\cdot(1+i f(c))
\]
が成り立つ。ここで $X$ は再び実数である。しかし第1項により、$f(r)$ の因子が恒等的に消えることはできないので、割り算によって
\[
f(r)=X(\sigma+i\tau)
\]
を得る。ここで $\sigma$ と $\tau$ は実数で、$r$ と $X$ には依存しない。

特に $r=1$ に対しても実数 $X_0$ が属するので、条件 $1=X_0(\sigma+i\tau)$ は必然的に $\tau=0$ を示す。したがって $f(r)=\sigma\cdot X$ である。ゆえに各実数 $z$ にも実数 $f(z)$ が対応する。すなわち $y=0$ から必然的に $Y=0$ が従う。これにより、実数 $z$ に対する線形関係は $X-x=0$ に移る。\footnote{ここでは $c\ne0$、したがって $c_3$ も $c_4$ も消えないという事実を用いている。これに対し上では $c_4\ne0$ だけを用いた。そのため、二つのうち一つが消える場合を先に扱っておく必要があった。} 各実値は自分自身に対応する。したがって再び $f(z)=z$ または $\overline z$ のみであり、実際には階数二である。こうして階数三はすべての場合に排除された。一方、最初の各項の研究が示すように、不連続解は実際に存在し、それらはしたがって必然的に階数四でなければならない。ゆえに 1) により次の定理が成り立つ。\emph{「関数方程式 (1) から (4) のすべての解} $f(z)$ \emph{は、連続解} $f(z)=z$ \emph{および} $\overline z$ \emph{を除いて、実数上および複素数上で極端に不連続である。」}\footnote{Lebesgue は前掲箇所で、$f(z)$ の実部および虚部がともに $x,y$ の「非可測」関数であることを示している。しかし第4項の初めの例 $\varphi(z)=\varphi(x)+i(y+\varphi(x))$ が示すように、そこから複素領域での極端な不連続性はまだ従わない。}

\medskip
\noindent Erlangen, 1915年10月30日。

\newpage

\begin{center}
{\Large\bfseries 11. 指定された群をもつ方程式}\par
\vspace{1.2em}
\emph{Math. Ann.} 78 (1918), pp. 221--229
\end{center}

\vspace{1.5em}

指定された群をもつ方程式を構成する問題は、二つの方向から攻めることができる。すなわち、根を特徴づける方向、簡単にいえば「非有理的」方向と、係数を特徴づける「有理的」方向である。

函数論的・算術的に働く「非有理的」方向には、有理数の領域におけるすべての Abel 体は円分体であるという Kronecker の定理、および二次数体に関する相対 Abel 体についての対応する定理が属する。これらの定理は、一方では、これらの特殊な有理性領域に関してすべての Abel 群が実現可能であることを与える。他方では、方程式の全体を与え、同時にそれらによって定義される数体の構造についてより深い洞察を与える。しかし各々の定理は特殊な有理性領域に限られており、新しい有理性領域ごとに全く新しい扱いを必要とする。

代数的に働く「有理的」方向は Hilbert の既約性定理に基づく。この定理によれば、係数の表示においてパラメータ表示に制限してよい。ここには、例えば任意の有理性領域に関して、交代群をもつ方程式が任意に多く存在することが含まれる。ここでは当然、上に述べたような、方程式で定義される体の構造についての洞察は欠ける。しかしこの「有理的」方向の利点は、群の実現可能性がすべての有理性領域に対して同時に証明される点にある。ただし、これまで知られているパラメータ表示は一般には方程式全体を与えない。\footnote{可解方程式については、係数のパラメータ表示を根のパラメータ表示で置き換えることができる。最近 F. Mertens は、次数 8 のある群について、最も一般的な根の表示を与えた: ``Gleichungen mit Quaternionengruppen'', Sitzb. d. Ak. d. Wiss., Wien, Abt. IIa, Bd. 125, S. 735。このノートによれば、それ以前に G. Bucht も、3 次および 4 次方程式、ならびに上記の群をもつ 8 次方程式について、常に作れる標準形に対する最も一般的な根の式を与えていた: ``Über einige algebraische Körper achten Grades'', Arkiv for Math., Astron. och Physik, Bd. 6, Nr. 30。[校正時の追記]}

以下はこの「有理的」方向への一つの寄与である。任意の有理性領域において、指定された群をもつ方程式の全体を与えるパラメータ表示が存在するための十分条件を与える。その条件とは、群に属する不変体、すなわち Lagrange の \emph{Gattungsbereich} が、その領域の独立な函数によって有理的に表示されること、つまり「最小基底」をもつことである。別の言い方をすれば、それが $n$ 個の不定元のすべての有理函数の体と同型であることである。

これは、\S 2 でなお示すように、例えば 3 次および 4 次方程式に現れるすべての群について成り立つ。これらの 3 次および 4 次方程式の実際の構成と詳しい議論は F. Seidelmann の学位論文\footnote{\emph{Die Gesamtheit der kubischen und biquadratischen Gleichungen mit Affekt bei beliebigem Rationalitätsbereich}. Erlangen, 1916.}にあり、その抜粋がこの報告の直後に続く。

Lagrange の \emph{Gattungsbereich} の最小基底の問題は、すでに論文 ``Körper und Systeme rationaler Funktionen'' (Math. Ann. 76) で述べたように、ここで説明した関連とは独立に、E. Fischer が私に対して提起したものであり、これが本研究の契機となった。\footnote{``Rationale Funktionenkörper'' に関する予備報告の第 2 部を参照。Jahresb. d. d. Math. Vereinig., Bd. 22, 1913.}

\subsection*{\S 1. パラメータ表示のための十分条件}

\textbf{1.} 指定された群をもつ方程式の構成をパラメータ表示に帰着できることは、Hilbert の既約性定理から従う。その内容は次の通りである。方程式
\begin{equation}\tag{1}
f(x)=x^n+F_1(\lambda_1,\ldots,\lambda_r)x^{n-1}+\cdots+F_n(\lambda_1,\ldots,\lambda_r)=0
\end{equation}
があり、ここで $F_1(\lambda_1,
\ldots,\lambda_r),\ldots,F_n(\lambda_1,
\ldots,\lambda_r)$ は、独立パラメータ $\lambda_1,
\ldots,\lambda_r$ の有理函数であり、係数は数体\footnote{数体の代わりに一般の有理性領域を置くこともできる。} $\Omega$ に属するとする。この方程式が、$\Omega$ 係数の $\lambda_1,
\ldots,\lambda_r$ のすべての有理函数から成る体 $\Omega(\lambda_1,
\ldots,\lambda_r)$ に関して群 $\Gamma$ をもつならば、パラメータ $\lambda$ を任意に多くの有理数で置き換えて、得られる数値方程式
\begin{equation}\tag{2}
g(x)=x^n+a_1x^{n-1}+\cdots+a_n=0
\end{equation}
が $\Omega$ に関してちょうど群 $\Gamma$ をもつようにできる。そのような $g(x)$ を、より正確には $g_\Gamma$ と記す。

そこで、次のようなパラメータ表示 (1) を与えられるかが問われる。すなわち、$\lambda_1,
\ldots,\lambda_r$ が $\Omega$ のすべての量を走るとき、ただし群の縮小を引き起こす $\Omega$ の量は除くとして、得られる数値方程式 (2) がすべての $g_\Gamma$ を走るような表示である。\footnote{ここには $g(x)=0$ の重根の出現も数え入れなければならない。} 言い換えると、非線型方程式系
\begin{equation}\tag{3}
F_1(\lambda_1,
\ldots,\lambda_r)=a_1;\quad \ldots;\quad F_n(\lambda_1,
\ldots,\lambda_r)=a_n
\end{equation}
が、$g_\Gamma$ に対応する各値系 $a_1,
\ldots,a_n$ に対して、しかも $\Omega$ に属する量 $\lambda_1,
\ldots,\lambda_r$ において解をもつべきである。方程式系 (3) は非線型であるから、$g_\Gamma$ の全体を一つのパラメータ表示で得るという要求には、初めから制限を加えなければならない。実際、この種の非線型系では、代数関係 $H(a_1,
\ldots,a_n)=0$ を満たす特異な値系 $a_1,
\ldots,a_n$ が常に現れ得て、そのような値系に対しては解が存在しない。\footnote{これは「非線型方程式系における交代」についての論文で、近く一般に実行される予定である。} その場合、パラメータ表示 (1) は失敗し、補足表示が必要になる。しかし本場合では、これらの特異値 $a_1,
\ldots,a_n$ は簡単に特徴づけられる。

\textbf{2.} このようなパラメータ表示を得るために、さらに、$\lambda_i$ が、根を不定元と見なしたとき、恒等的に群に属する自然な無理性であることを要求する。これは次の意味である。(1) においてパラメータ $\lambda_1,
\ldots,\lambda_r$ を、$\Omega$ 係数の、不定元 $x_1,
\ldots,x_n$ の適切に選んだ有理函数 $\varphi_1(x),\ldots,\varphi_r(x)$ で置き換えると、(1) は
\begin{equation}\tag{4}
h(x)=x^n-\sigma_1(x)x^{n-1}+\sigma_2(x)x^{n-2}\cdots \pm \sigma_n(x)
\end{equation}
に移るものとする。ここで $\sigma_1(x),\ldots,\sigma_n(x)$ は $x_1,
\ldots,x_n$ の基本対称函数を意味する。そして $h(x)$ は、$\Omega(\lambda_1,
\ldots,\lambda_r)$ からこのようにして生じる有理性領域 $\Omega(\varphi_1(x),\ldots,\varphi_r(x))$ に関して、群 $\Gamma$ をもつものとする。さらにパラメータの独立性により、$\varphi_1(x),\ldots,\varphi_r(x)$ も代数的に独立な函数でなければならない。

この要求を精密にするためには、$\Omega(\varphi_1(x),\ldots,\varphi_r(x))$ と群 $\Gamma$ との関係を明らかにする必要がある。そのため、まず $h(x)$ が $h_\Gamma$ となるような最も一般の有理性領域 $K$ を定める。$\Omega_\Gamma$ で、$\Gamma$ を許す $x_1,
\ldots,x_n$ の有理整係数のすべての有理函数から成る体を表す。これは $\Gamma$ に属する「Lagrange の \emph{Gattungsbereich}」、または $\Gamma$ の不変体とも呼ばれる。また、$\Omega(\Omega_\Gamma)$ で、$\Omega_\Gamma$ を $\Omega$ に添加して得られる対応する体を表す。

したがって $\Omega_\Gamma$ は、$x_1,
\ldots,x_n$ の有理整係数のすべての有理函数の体 $\mathfrak R(x_1,
\ldots,x_n)$ の一つの因子である。すると群の定義により、最も一般の領域 $K$ は、$\Omega_\Gamma$ を因子として含み、かつ $\mathfrak R(x_1,
\ldots,x_n)$ との交わりがちょうど $\Omega_\Gamma$ になるような任意の体によって与えられる。これに対応して、$\Omega$ を含む任意の領域 $K$ について、$\Omega(x_1,
\ldots,x_n)$ との交わりは $\Omega(\Omega_\Gamma)$ によって与えられる。とくに、われわれの要求により $\Omega(\varphi_1(x),\ldots,\varphi_r(x))$ は一つの領域 $K$ になるから、$\Omega(\varphi_1(x),\ldots,\varphi_r(x))$ と $\Omega(x_1,
\ldots,x_n)$ との交わりは $\Omega(\Omega_\Gamma)$ によって与えられる。他方、$\Omega(\varphi_1(x),\ldots,\varphi_r(x))$ は $\Omega(x_1,
\ldots,x_n)$ の因子であるから、ちょうど
\[
\Omega(\varphi_1(x),\ldots,\varphi_r(x))=\Omega(\Omega_\Gamma)
\]
となる。したがって要求は、Lagrange の \emph{Gattungsbereich} $\Omega(\Omega_\Gamma)$ が、代数的に独立な函数 $\varphi_1(x),\ldots,\varphi_r(x)$ を $\Omega$ に添加することで生じるべきである、ということを意味する。このとき $r=n$ でなければならない。なぜなら $\Omega_\Gamma$ は $n$ 個の代数的に独立な基本函数を含む一方、$n$ 個より多くの函数は代数的に独立ではあり得ないからである。$\varphi_1(x),\ldots,\varphi_n(x)$ は $\Omega(\Omega_\Gamma)$ の\emph{最小基底}を成す、と言うことにする。また、$\Omega_\Gamma$ は $\Omega$ 係数の最小基底をもつ、とも言う。

\textbf{3.} いま重要なのは、この考察全体を反転でき、実際に求めるパラメータ表示の十分条件\footnote{要求は 2 で 1 より拡張されたので、条件は必要条件である必要はない。}へと導けることである。そこで、$\varphi_1(x),\ldots,\varphi_n(x)$ が $\Omega(\Omega_\Gamma)$ の最小基底であると仮定する。さらに $\Omega$ を、$\varphi_1(x),\ldots,\varphi_n(x)$ の係数を含む最小の体とする。このとき $\Omega$ は必然的に代数的数体であり、とくにすべての有理数の体 $\mathfrak R$ であり得る。$\sigma_1(x),\ldots,\sigma_n(x)$ は $\Omega(\Omega_\Gamma)$ に属するから、$x_1,
\ldots,x_n$ に関する恒等式が存在する。ここで $F_i$ はその引数の有理函数を意味する。
\begin{equation}\tag{5}
\sigma_1(x)=F_1(\varphi_1(x),\ldots,\varphi_n(x));\quad\ldots;\quad
\sigma_n(x)=F_n(\varphi_1(x),\ldots,\varphi_n(x)).
\end{equation}
逆に、$\varphi_1(x),\ldots,\varphi_n(x)$、したがってまた
$\varphi(x)=\mu_1\varphi_1(x)+\cdots+\mu_n\varphi_n(x)$ は、$\Omega(\sigma_1(x),\ldots,\sigma_n(x))$ に関して既約な方程式
\begin{equation}\tag{6}
G(\varphi(x);\sigma_1(x),\ldots,\sigma_n(x))=0
\end{equation}
によって、$\sigma_1(x),\ldots,\sigma_n(x)$ に代数的に依存する。この式もまた $x$ における恒等式である。

(5) により、この恒等式 (6) は
\begin{equation}\tag{7}
G\bigl(\varphi(x);F_1(\varphi_1,
\ldots,\varphi_n),\ldots,F_n(\varphi_1,
\ldots,\varphi_n)\bigr)
=H(\varphi_1(x),\ldots,\varphi_n(x))=0
\end{equation}
に移る。そして $\varphi_1(x),\ldots,\varphi_n(x)$ の代数的独立性により、(7) は $x$ においてだけでなく $\varphi_i$ においても恒等的に満たされる。すなわち、独立パラメータ $\lambda$ において恒等的に
\begin{equation}\tag{8}
H(\lambda_1,
\ldots,\lambda_n)=G\bigl(\mu_1\lambda_1+\cdots+\mu_n\lambda_n;
F_1(\lambda_1,
\ldots,\lambda_n),\ldots,F_n(\lambda_1,
\ldots,\lambda_n)\bigr)=0
\end{equation}
が成り立つ。恒等式 (8) から、方程式
\begin{equation}\tag{9}
f(x)=x^n-F_1(\lambda_1,
\ldots,\lambda_n)x^{n-1}+F_2(\lambda_1,
\ldots,\lambda_n)x^{n-2}\cdots \pm F_n(\lambda_1,
\ldots,\lambda_n)=0
\end{equation}
は $\Omega(\lambda_1,
\ldots,\lambda_n)$ に関して群 $\Gamma$ をもつことが従う。実際、この恒等式は、既知の有理的な量 $\lambda=\mu_1\lambda_1+
\cdots+
\mu_n\lambda_n$ が、$\Gamma$ に属する根の函数であることを示している。しかし、$\lambda_i=\varphi_i(x)$ という代入が示すように、$\Gamma$ の真の部分群に属する函数がさらに有理的に知られることはあり得ない。また同じ代入により、$\lambda$ はそのすべての共役とは異なる。さらに $\Omega^*$ が $\Omega$ を含む任意の数体であれば、$f(x)$ は $\Omega^*(\lambda_1,
\ldots,\lambda_n)$ に関して群 $\Gamma$ をもつ。なぜなら 2 により、$\lambda_i=\varphi_i(x)$ と代入しても、$\Omega^*$ を添加した場合に群は縮小されないからである。

つぎに、1 で修正された意味において、$f(x)$ が実際にすべての $g_\Gamma$ の全体を走ることを示さなければならない。そのため、(5) により $F_i(\lambda)$ は引数 $\lambda$ の代数的に独立な函数であることに注意する。したがって方程式系
\begin{equation}\tag{10}
F_1(\lambda_1,
\ldots,\lambda_n)=-a_1;\quad F_2(\lambda_1,
\ldots,\lambda_n)=a_2;\quad\ldots;\quad F_n(\lambda_1,
\ldots,\lambda_n)=\pm a_n
\end{equation}
は、任意の値系 $a_1,
\ldots,a_n$ に対して、まだ述べる特異な場合を除き、解をもつ。$\Omega^*$ に関してある $g_\Gamma$ に対応する任意の値系 $a_1,
\ldots,a_n$ に対して、対応する $\lambda_i$ は、群に属する自然な無理性\footnote{さらに下をも参照。}として、$\Omega^*$ の量になる。したがって $\lambda_1,
\ldots,\lambda_n$ がすべての値系を走れば、$g(x)$ はすべての方程式の全体を、修正された意味で走る。その中には、$\Omega^*$ の値が対応する $g_\Gamma$ の全体が含まれている。

逆に、Hilbert の既約性定理により、$\Omega^*$ の値 $\lambda$ には「一般に」$g_\Gamma$ も対応するから、1 で修正された意味において、パラメータ表示 (9) により $\Omega^*$ に関するすべての $g_\Gamma$ の全体が実際に得られる。

最後に、パラメータ表示が失敗する特異値系を求める必要がある。最小基底の函数を、分子と分母に分けて
\[
\varphi_i(x)=\frac{\psi_i(x)}{\chi_i(x)}
\]
と書く。これを恒等式 (5) に代入すると、これらが（約分なしに）
\begin{equation}\tag{11}
\sigma_k(x)=\frac{G'_k(x)}{H_k(x)}\quad\text{または}\quad H_k(x)\cdot\sigma_k(x)=G_k(x)
\end{equation}
に移るものとする。したがってすべての $H_k(x)$ の積は、すべての分母 $\chi_i(x)$ の積で割り切れる。根系 $\alpha_1,
\ldots,\alpha_n$ について
\[
H_1(\alpha)H_2(\alpha)\cdots H_n(\alpha)\ne0
\]
が成り立ち、したがっていかなる分母 $\chi_i(\alpha)$ も消えないなら、(11) を $H_k(\alpha)$ で割ることができ、
\begin{equation}\tag{12}
\sigma_k(\alpha)=\frac{G_k(\alpha)}{H_k(\alpha)}=F_k(\varphi_1(\alpha),\ldots,\varphi_n(\alpha))
\end{equation}
を得る。ここで $\varphi_i(\alpha)$ は分母が消えないので有限値を取り、$g_\Gamma$ である方程式については $\Omega^*$ の量となる。$\alpha_1,
\ldots,\alpha_n$ を $a_k=(-1)^k\sigma_k(\alpha)$ となるように定めるなら、(12) は (10) の解系を表す。\footnote{方程式の解法はこうして、独立函数から成る系
$\varphi_1(\alpha_1,
\ldots,\alpha_n)=\lambda_1;\ldots;\varphi_n(\alpha_1,
\ldots,\alpha_n)=\lambda_n$
の解法に帰着される。}

したがって特異になり得るのは、
\[
H_1(\alpha)H_2(\alpha)\cdots H_n(\alpha)=0
\]
となるような係数系 $a_k=(-1)^k\sigma_k(\alpha)$ だけである。すなわち、(11) が表示 $0=0$ に移る場合である。その場合、このように定義される特異な $a_1,
\ldots,a_n$ の中に、方程式 $g_\Gamma$ に対応するものが含まれるか、またどの数体についてそうであるかを、別に調べる必要がある。\footnote{例えば F. Seidelmann が前掲書で示したように、3 次および 4 次方程式の交代群についてのみ、このような、特異値系に対応する方程式が現れる。しかもそれは、三乗根を含む数体 $\Omega^*$ についてのみである。この場合にはそれぞれ簡単な補足表示を与えることができる。} まとめると次の定理を得る。

\smallskip
\noindent\textbf{定理.} \emph{群 $\Gamma$ に属する Lagrange の \emph{Gattungsbereich} が、$\Omega$ 係数の最小基底\footnote{一つの最小基底の存在から、元のものから可逆有理変換で導かれる任意に多くの最小基底の存在がただちに従う。}をもつならば、$\Omega$ を含む任意の数体 $\Omega^*$ について、$\Omega^*$ に関して指定された群 $\Gamma$ をもつ方程式の全体は、パラメータ表示 (2) によって有理的に構成できる。ただし、パラメータ表示に関して特異なものは例外となることがあり、その場合は別に与えることができる。パラメータは、群の縮小を引き起こさない $\Omega^*$ のすべての値を走らなければならない。最小基底の係数領域 $\Omega$ がすべての有理数の体 $\mathfrak R$ に等しいならば、この構成は任意の数体に対して可能である。}

さらに Hilbert の既約性定理により、群の縮小を引き起こす $\Omega^*$ の値を除いても、パラメータ多様体の次元は低下しない。したがってなお次が成り立つ。\emph{最小基底の存在から、$\Omega^*$ に関して群 $\Gamma$ をもつ方程式の多様体は、Affekt をもたない方程式の多様体と等しいことが従う。多様体は Affekt によって低下しない。}

\subsection*{\S 2. 特殊な方程式への応用}

いま、3 次および 4 次方程式に現れるすべての群について、最小基底の存在を示す。そのためにまず一般に、$n$ 次方程式についてこの問題が $n-2$ 個の不定元の問題に還元されることを示す。第一の還元は、二番目に高い項を消去するための線型 Tschirnhaus 変換に対応する。第一の基底函数として $\varphi_1(x)=\sigma_1(x)$ を取る。さらに $\Omega_\Gamma$ は、$x_1,
\ldots,x_n$ の、$\Gamma$ を許すすべての同次形式から有理的に導かれることに注意する。$p(x_1,
\ldots,x_n)$ をそのような形式とする。このとき
\[
q(x)=p\left(x_1-\frac{\sigma_1(x)}{n},\ldots,x_n-\frac{\sigma_1(x)}{n}\right)=p\left(x_i-\frac{\sigma_1(x)}{n}\right)
\]
も $\Omega_\Gamma$ に属し、さらに
\[
p(x)\equiv p\left(x_i-\frac{\sigma_1(x)}{n}\right)\pmod{\sigma_1(x)},
\quad\text{または}\quad
p(x)=q(x)+\sigma_1(x)\cdot p_1(x)
\]
が成り立つ。ここで $p_1(x)$ は $\Omega_\Gamma$ に属し、しかも $p(x)$ より低い次元をもつ。この手続きを続けると、$\Omega_\Gamma$ のすべての同次形式、したがって一般に $\Omega_\Gamma$ のすべての有理函数は、$\varphi_1(x)=\sigma_1(x)$ と同次形式
\[
q(x)=p\left(x_i-\frac{\sigma_1(x)}{n}\right)=p(y_i)
\]
によって有理的に表されることが分かる。これらの結合 $y_i$ の間には線型関係 $\sum y_i=0$ がある。したがって $q(x)$ は $(n-1)$ 個の引数にのみ依存する。

第二の還元は、$q(x)$ がその引数の同次形式であることに基づく。第二の基底函数として
\[
\varphi_2(x)=\frac{\tau_3(x)}{\tau_2(x)}
\]
を取る。ただし
\[
\tau_3(x)=\sigma_3\left(x_i-\frac{\sigma_1(x)}{n}\right)=\sigma_3(y_i),
\qquad
\tau_2(x)=\sigma_2\left(x_i-\frac{\sigma_1(x)}{n}\right)=\sigma_2(y_i)
\]
である。したがって $\varphi_2(x)$ は、$(n-1)$ 個の引数、すなわち $\sum y_i=0$ を満たす $y_i$ の、一次の同次分数函数である。$\varphi_2(x)$ によって、すべての $q(x)$ は、同じく $\Omega_\Gamma$ に属する函数
\[
r(x)=\frac{q(x)}{\varphi_2(x)^\lambda}
=\frac{q(x)\tau_2(x)^\lambda}{\tau_3(x)^\lambda}
=\frac{p(y_i)\sigma_2(y_i)^\lambda}{\sigma_3(y_i)^\lambda}
\]
で有理的に表される。ここで $\lambda$ は $q(x)$ の次数を表す。したがって $r(x)$ は 0 次同次となり、$\sum y_i=0$ の下での比 $y_1:y_2:\cdots:y_n$ という $(n-2)$ 個の引数だけに依存する。こうして $\Omega_\Gamma$ は
\[
\varphi_1(x)=\sigma_1(x),
\qquad
\varphi_2(x)=\frac{\tau_3(x)}{\tau_2(x)}
\]
と、これら $(n-2)$ 個の引数に依存する $\Omega_\Gamma$ のすべての函数 $r(x)$ から成る体によって有理的に表される。\footnote{第二項をもたない $n$ 次方程式
\[
f(x)=x^n+a_2x^{n-2}+a_3x^{n-3}+\cdots+a_n=0
\]
について、この第二の還元に対応するのは、$a_2\ne0$, $a_3\ne0$ なら常に可能な代入
\[
y=\frac{x a_2}{a_3}
\]
である。これにより $f(x)$ は、因子を除いて、
\[
g(y)=y^n+\frac{a_2^3}{a_3^2}y^{n-2}+\frac{a_2^3}{a_3^2}y^{n-3}
+\frac{a_4a_2^4}{a_3^4}y^{n-4}+\cdots+a_n\frac{a_2^n}{a_3^n}=0
\]
に移る。すなわち、ただ $n-2$ 個のパラメータだけを含む方程式
\[
g(y)=y^n+c\,y^{n-2}+c\,y^{n-3}+c_4y^{n-4}+\cdots+c_n=0
\]
になる。} これで望む還元が達成される。

3 次および 4 次方程式については、したがって、一つまたは二つの引数の函数体への還元が得られる。これらには Lüroth と Castelnuovo の定理により常に最小基底が存在する。\footnote{J. Lüroth: ``Beweis eines Satzes über rationale Kurven'', Math. Ann. 9 (1875); G. Castelnuovo: ``Sulla razionalità delle involuzioni piane'', Math. Ann. 44 (1893)。二つを超える不定元の体では最小基底が存在する必要がないことを、F. Enriques が反例によって示した。Rend. Acc. Linc., Vol. 21, 1912 年 1 月 21 日。} 一つの不定元の場合には、最小基底は有理演算によって導かれ、したがってとくに $\Omega_\Gamma$ に対して有理数係数だけを含む。二つの不定元の場合、Castelnuovo によれば、なお代数的数体 $\Omega$ の係数をもつ基底である可能性がある。しかし実際の研究（F. Seidelmann 前掲書を参照）は、この場合にも各 $\Omega_\Gamma$ について有理数係数の基底が得られることを示す。これはほとんど計算なしに、すでに次のことから分かる。すなわち、四元群については有理数係数の基底
\[
\varphi_1(x)=\sigma_1(x);\quad
\varphi_2(x)=\frac{vw}{u};\quad
\varphi_3(x)=\frac{uw}{v};\quad
\varphi_4(x)=\frac{uv}{w},
\]
を選べる。ただし
\[
u=x_1+x_2-x_3-x_4;\quad v=x_1-x_2+x_3-x_4;\quad w=x_1-x_2-x_3+x_4
\]
である。さらに、交代群は三つの基底函数 $\varphi_2,\varphi_3,\varphi_4$ の巡回群へ移り、各八元群はこれらの函数のうち二つの対称群へ移る。これにより、3 次および 2 次方程式の群への還元がなされる。巡回群については、有理数係数の基底は Abel によってすでに知られていた。ただしそのようには定式化されていなかった。\footnote{Weber, \emph{Algebra}（小版）, \S 93, 式 (12) を参照。} したがって次が証明された。\emph{指定された群をもつ 3 次および 4 次方程式の全体は、パラメータ表示に関して特異なものを例外とする可能性を除き、任意の有理性領域についてパラメータ表示により有理的に構成できる。その多様体は Affekt をもたない方程式の多様体と等しい。}

実際の研究（F. Seidelmann 前掲書を参照）は、三乗根の体を含む有理性領域に関する交代群の場合にだけ補足表示が必要となり、それ以外のすべての場合にはパラメータ表示が端的に全体を与えることを示す。

なお、すべての Abel 群についても最小基底は知られている。\footnote{E. Fischer: ``Die Isomorphie der Invariantenkörper der endlichen Abelschen Gruppen linearer Transformationen'', Gött. Nachr. 1915, および ``Zur Theorie der endlichen Abelschen Gruppen'', Math. Ann. 77 (1915)。} しかしこの場合は、群の指標から導かれる円分体を係数にもつ基底である。したがって、指定された群をもつ方程式の構成について、この場合に新しい結果は得られない。

\medskip
\noindent Göttingen, 1916年7月。

\clearpage

\begin{center}
{\Large\bfseries 12. 任意の微分式の不変量}\par
\vspace{1.5em}
Nachr. v. d. Ges. d. Wiss. zu Göttingen 1918, pp. 37--44
\end{center}

\vspace{3em}
\begin{center}
1918年1月25日の会合で F. Klein により提出。
\end{center}

微分式とは、\(n\) 個の変数 \(x_1,\ldots,x_n\) とその微分
\(dx_1,\ldots,dx_n\) との函数
\[
f(x,dx)=f(x_1,\ldots,x_n;\,dx_1,\ldots,dx_n)
\]
であり、二つの \(n\) 個ずつの引数の系に関して解析的であるが、必ずしも実であるとは仮定しないものをいう。いま基礎に置くのは、変数のすべての解析的変換から成る群であり、これに微分のすべての線形変換から成る群が対応する：
\begin{equation}
x_i=x_i(y_1,\ldots,y_n);\qquad
dx_i=\sum_k\frac{\partial x_i}{\partial y_k}\,dy_k;\qquad
\delta x_i=\sum_k\frac{\partial x_i}{\partial y_k}\,\delta y_k,
\tag{1}
\end{equation}
この変換によって \(f(x,dx)\) は \(g(y,dy)\) に移るものとする。\(f(x,dx)\) の不変量とは、この群に関する不変量であり、さらにこの群が高次微分 \(d^2x,d\delta x,\ldots\) および \(f(x,dx)\) の導函数に誘導する群に関する不変量である。すなわち、任意の \(f\) について、(1) により、変数および現れるすべての微分において恒等的に
\[
\begin{aligned}
&J\!\left(f,\frac{\partial f}{\partial dx}\cdots
\frac{\partial^{\rho+\sigma}f}{\partial x^\rho\partial dx^\sigma}
\cdots dx,\delta x,d^2x,\ldots\right)\\
&\qquad =
J\!\left(g,\frac{\partial g}{\partial dy}\cdots
\frac{\partial^{\rho+\sigma}g}{\partial y^\rho\partial dy^\sigma}
\cdots dy,\delta y,d^2y,\ldots\right)\footnotemark
\end{aligned}
\]
を満たす（解析的）函数 \(J\) をいう。
\footnotetext{\(\displaystyle
\frac{\partial^{\rho+\sigma}f}{\partial x^\rho\partial dx^\sigma}\)
は、以下一貫して
\[
\frac{\partial^{\rho+\sigma}f}
{\partial x_{i_1}\cdots\partial x_{i_\rho}
 \partial dx_{k_1}\cdots\partial dx_{k_\sigma}}
\]
の略記として用いる。ここで添字は \(1\) から \(n\) までの任意の数を表す。}
同時微分式系の不変量もまったく同様に定義される。とくに、そのような不変量が第一微分 \(dx,\delta x\) だけを含み、かつそれらの微分に関する導函数だけを含んで、変数に関する導函数を含まないならば、函数に変数が現れること、および微分の線形変換は考慮に入らない。この特殊な不変量は、不定係数をもつ微分 \(dx,\delta x\) のすべての線形変換の群に関する不変量であり、同時系の\emph{射影不変量}と呼ぶことにする。

二次の斉次微分形式については、Christoffel（Crelle 70）および Ricci（Math. Annalen 54）が、不変な構成から成る一つの系を与えている。その系により、すべての不変量はこの同時系の射影不変量になる。したがって、不変量全体の問題および同値性の問題は、線形不変式論の問題へ還元される。これを私は「還元定理」と呼びたい。完全系は可算無限個の形式から成るが、任意の有限位数 \(\rho\) の不変量、すなわち変数に関して \(\rho\) 次を超える導函数を含まない不変量については、有限個の形式だけから成る。

以下では、この「還元定理」が任意の微分式に対して成り立つことを示す。ここでも、問題となるのは可算無限個の不変な構成から成る完全系であるが、各有限位数については有限個しか現れない。ただし Christoffel と Ricci は、不変量の生成および還元定理の証明を、見通しにくい消去に依拠させているのに対し、私は不変な微分過程と変分過程によって不変量を直接生成する。後者は、単に別のパラメータに関する微分過程とみなすことができる。この不変な微分過程のアルゴリズムは、Lagrange（\emph{Mécanique analytique}, 第2部, IV節）に続いて、Riemann（全集 XXII）および Lipschitz（Crelle 70, 72）が二次微分形式について発展させたが、還元定理までは到達しなかった。還元定理の証明に用いる補助手段は、Riemann が二次形式について導入した「正規座標」（全集 XIII）である。これは、対応する変分問題において一点から出る極値曲線を直線へ変換するが、斉次微分式
\(f(x,\varkappa\cdot dx)=\Phi(\varkappa)\cdot f(x,dx)\footnote{容易に、\(\Phi(\varkappa)=\varkappa^c\) となることがわかる。ここで \(c\) は任意の実または複素量を表す。}\) に対してのみ存在する。しかし非斉次の場合はこの場合へ還元できる。

\begin{center}
\textbf{I. 不変量の生成。}
\end{center}

\(f(x,dx)\) の射影不変量の、一般には終わらない列は、極化形によって与えられる。微分の変換が線形であるため、\(f(x,dx)=g(y,dy)\) から
\(f(x,dx+\lambda\delta x)=g(y,dy+\lambda\delta y)\) も従う。これを \(\lambda\) で展開すると、不変性を特徴づける関係が得られる：
\begin{equation}
\begin{gathered}
f(dx)=g(dy),\\
f_\delta=\sum\frac{\partial f(dx)}{\partial dx}\,\delta x
       =\sum\frac{\partial g(dy)}{\partial dy}\,\delta y,\\
\vdots\\
f_{\delta^\sigma}=
\sum\frac{\partial^\sigma f(dx)}{\partial dx^\sigma}\,\delta x^\sigma
=
\sum\frac{\partial^\sigma g(dy)}{\partial dy^\sigma}\,\delta y^\sigma,\\
\vdots
\end{gathered}
\tag{2}
\end{equation}

各極化形は、変分過程 \(\bdelta\) を適用することにより、一般不変量の終わらない列を与える：
\begin{equation}
\begin{gathered}
\bdelta^\rho f(x,dx)=\bdelta^\rho g(y,dy),\\
\vdots\\
\bdelta^\rho f_{\delta^\sigma}
=\bdelta^\rho\sum\frac{\partial^\sigma f(dx)}{\partial dx^\sigma}\,
\delta x^\sigma
=\bdelta^\rho\sum\frac{\partial^\sigma g(dy)}{\partial dy^\sigma}\,
\delta y^\sigma,\\
\vdots\\[0.2em]
(\rho=0,1,2,\ldots).
\end{gathered}
\tag{3}
\end{equation}
関係 (2) と (3) から、(1) によって \(f(x,dx)\) の導函数の変換法則を読み取ることができるが、以下ではこれは用いない。不変量 (3) から、線形結合によって「\(\rho\) 次変分の正規形」を得ることができる。これは \(\rho=1\) では Lagrange に、\(\rho=2\) では Riemann に見られるもので、\((\rho+1)\) 次の混合微分、すなわち \(d^\rho\delta,d^{\rho-1}\delta^2,\ldots\) が消去されていることによって特徴づけられる。その形は次である：
\begin{equation}
\begin{gathered}
\Omega_1=\delta f(dx)-d f_\delta(dx),\\
\Omega_2=\delta^2 f-\delta d f_\delta+\frac12 d^2 f_{\delta^2},\\
\vdots\\
\Omega_\rho=\delta^\rho f-\delta^{\rho-1}d f_\delta
+\frac12\delta^{\rho-2}d^2 f_{\delta^2}
+\cdots+\frac{(-1)^\rho}{\rho!}\,d^\rho f_{\delta^\rho},\\
\vdots
\end{gathered}
\tag{4}
\end{equation}

不変量 (4) から、Riemann の一つの Ansatz を一般化することにより、第一微分だけを含む不変量に到達できる。このような不変量を「基本関数」と呼ぶことにする。実際、\(\Omega_1\) において微分を微分商で置き換えると、\(\Omega_1=0\)（\(\delta x\) に関して恒等的）は、ちょうど
\[
f\!\left(\frac{dx}{dt}\right)
\]
に属する変分問題の Lagrange 方程式を与える。そこで私は、方向 \((d+\lambda\delta)\) がこれらの Lagrange 方程式を満たす、という不変条件を課す：
\begin{equation}
\Omega_1(d+\lambda\delta)=
\bdelta f(dx+\lambda\delta x)
-(d+\lambda\delta)f_{\bdelta}(dx+\lambda\delta x)=0
\quad(\text{\(\delta x\) に関して恒等的}),
\tag{5}
\end{equation}
ここから、置換 \(\bdelta=d+\lambda\delta\) によって、斉次な \(f(dx)\) についてさらに
\begin{equation}
(d+\lambda\delta)f(dx+\lambda\delta x)=0
\tag{6}
\end{equation}
が従う。ただし一次斉次の場合は例外である。この場合には、方程式系 (5) が \(n-1\) 個の独立条件しか表さないため、(6) を新たな条件として加えることにする。(5)、または (5) と (6) において、\(\lambda\) の零次・一次・二次の係数を零とおくと、三つの不変な方程式系
\(\varphi=0\), \(\psi=0\), \(\chi=0\)（\(\delta x\) に関して恒等的）が得られ、これにより \(d^2x,d\delta x,\delta^2x\) を第一微分によって表すことができる。さらに次の不変な方程式系
\begin{equation}
\begin{gathered}
d^\sigma\varphi=0,\qquad d^\sigma\psi=0,\qquad
d^\sigma\chi=0,\qquad d^{\sigma-1}\delta\chi=0,\\
d^{\sigma-2}\delta^2\chi=0,\ldots,\delta^\sigma\chi=0
\qquad(\sigma=0,1,2,\ldots)
\end{gathered}
\tag{7}
\end{equation}
は、すべての高次微分を第一微分によって表すのにちょうど十分である。\footnote{線形形式 \(\sum A_i\,dx_i\) だけは特別な処理を要する。ここでは方程式系 (5) が第二微分を含まないからである。} したがって、函数 (4) と不変な方程式系 (7) とを組み合わせることにより、第一微分だけを含む基本関数 \([\Omega_\rho]\) が得られる。ただし \([\Omega_1]\) は恒等的に消える。

同様に、基本関数 \(h\) の微分 \(\delta h(dx,\delta x)\) から、(7) によって再び基本関数を得ることができ、これを \(h\) の「共変導関数」\([h^{(1)}]\) と呼ぶことにする。一般に \([h^{(\nu)}]\) は \([h^{(\nu-1)}]\) の共変導関数を表す。この二つの過程により、基本関数の二重無限列が得られる：
\begin{equation}
\begin{array}{cccccc}
[\Omega_2], & [\Omega_2^{(1)}], & [\Omega_2^{(2)}], & \ldots,
& [\Omega_2^{(\nu)}], & \ldots\\
\vdots &&&& \vdots&\\
{}[\Omega_\rho], & [\Omega_\rho^{(1)}], & \ldots, & \ldots,
& [\Omega_\rho^{(\nu)}], & \ldots\\
\vdots & \vdots &&& \vdots&
\end{array}
\tag{8}
\end{equation}

さらに、高次微分を第一微分によって表す方程式の左辺は、第一微分と共変的である、すなわち同じ線形変換を受けることがわかる。したがって、高次微分の代わりに、これらの左辺 \(p,q,r,\ldots\) を不変量の引数として導入すると、その不変量は、導函数を別にすれば、互いに共変的な量だけに依存する。上で述べた函数 (8) の形成は、単にこれらの新しい量を導入することに帰着する。任意の不変量はそれによって不変量の和へ移り、その和から、ちょうど \(dx,\delta x\) を含み、\(p,q,r,\ldots\) を含まない項を取り出すのである。

以下では、斉次性の仮定
\[
f(\varkappa\cdot dx)=\Phi(\varkappa)\cdot f(dx)
\]
の下で、基本関数 (8) が求める完全系を尽くすことを示す。

\begin{center}
\textbf{II. 正規座標、同値性および還元定理。}
\end{center}

私は、\(f\!\left(\frac{dx}{dt}\right)\) に属する変分問題を積分することによって正規座標に到達する：
\begin{equation}
\begin{gathered}
\Omega_1=\delta f\!\left(\frac{dx}{dt}\right)
-\frac{d}{dt}f_\delta\!\left(\frac{dx}{dt}\right)=0
\quad(\text{\(\delta x\) に関して恒等的}),\\
\frac{d}{dt}f\!\left(\frac{dx}{dt}\right)=0
\quad(\text{帰結または追加条件として}).
\end{gathered}
\tag{9}
\end{equation}
\(f(dx)\) の斉次性の仮定により、この方程式系はパラメータ置換 \(t=\varkappa\cdot\tau\) のもとで自分自身へ移る。したがって、(9) によって \(\frac{d^\rho x_i}{dt^\rho}\) を第一導函数の函数
\(\varphi_\rho^{(i)}\!\left(\frac{dx}{dt}\right)\) として表すなら、\(\varphi_\rho^{(i)}\) は \(\rho\) 次斉次である：
\begin{equation}
\varphi_\rho^{(i)}\!\left(\varkappa\cdot\frac{dx}{dt}\right)
=\varkappa^\rho\varphi_\rho^{(i)}\!\left(\frac{dx}{dt}\right).
\tag{10}
\end{equation}
いま (9) を積分するにあたり、\(t=t_0\) で初期値
\[
x=(x)_0,\qquad \frac{dx}{dt}=\left(\frac{dx}{dt}\right)_0
\]
を与え、
\begin{equation}
u_i=(t-t_0)\left(\frac{dx_i}{dt}\right)_0
\tag{11}
\end{equation}
とおく。ここで \(f\!\left(\frac{dx}{dt}\right)=\mathrm{const.}\) により、パラメータ \((t-t_0)\) は「極値曲線の長さ」に比例する。すると、\((t-t_0)\) の冪による展開は、(10) により
\begin{equation}
x_i-(x_i)_0=u_i+\frac12\varphi_2^{(i)}(u)+\cdots+
\frac{1}{\rho!}\varphi_\rho^{(i)}(u)+\cdots .
\tag{12}
\end{equation}
を与える。この展開は変数変換 \(x_i=x_i(u)\) と見ることができ、このとき量 \(u\) がまさに Riemann の正規座標である。したがって、(11) と (12) により、これらの座標は \((x)_0\) を通る極値曲線を直線に変換する。さらに任意の変数変換 \(x=x(y)\) に対して、(11) により、対応する正規座標の線形変換 \(u=u(v)\) が対応し、その係数は任意の値系 \((x)_0\) だけに依存する。したがって微分 \(du,\delta u\) は \(u\) 自身と同じ変換を受ける。

(12) によって \(f(x,dx)\) は \(F(u,du)\) に移るものとする。すなわち、\(u\) の冪で展開して
\begin{equation}
F(u,du)=F_0(u,du)+F_1(u,du)+\cdots+F_\rho(u,du)+\cdots,
\tag{13}
\end{equation}
ここで \(F_\rho(u,du)\) は \(u\) における \(\rho\) 次元の斉次形式を表す。変換 \(u=u(v)\) の線形性のため、\(F(u,du)=G(v,dv)\) によって、個々の斉次成分も対応する成分へ移る：
\begin{equation}
F_\rho(u,du)=G_\rho(v,dv).
\tag{14}
\end{equation}
したがって、\((x)_0\) に任意の固定された定数値系を取ると、(13) と (14) は、\(f(dx)\) と \(g(dy)\) の同値性が、線形変換に関する対応する函数系 \(F_\rho(u,du)\) の同値性と同じ意味であることを示す。さらに \(F_\rho(u,du)\)、または対応する \(F_\rho(\delta u,du)\) は、点 \(x=(x)_0\) で取った \(f(dx)\) の不変量を表し、実際、この点 \(x=(x)_0\) に対する基本関数の完全系を成す。すなわち
\[
F_\rho(\delta u,du)=\frac1{\rho!}\sum
\frac{\partial F_\rho}{\partial \delta u^\rho}\,\delta u^\rho;
\qquad
\frac{\partial F_\rho}{\partial \delta u^\rho}
=\left(\frac{\partial F(du)}{\partial u^\rho}\right)_0,
\]
であり、この最後の導函数は、(12) によって、\(x=(x)_0\) で取った \(f(dx)\) の導函数を用いて、\(G(dv)\) から対応して形成される導函数が \(g(dy)\) の導函数を用いて表されるのと全く同じ仕方で表すことができる。さらに
\[
\left(\frac{\partial^{\rho+\sigma}F(du)}
{\partial u^\rho\partial du^\sigma}\right)_0
=
\frac{\partial^{\rho+\sigma}F_\rho(\delta u,du)}
{\partial\delta u^\rho\partial du^\sigma}
\]
であるから、点 \(x=(x)_0\) で取った \(f(dx)=F(du)\) の任意の不変量は、高次微分を \(dx,\delta x\) と共変的な \(p,q,r,\ldots\) で置き換えるだけで、\(F_\rho(\delta u,du)\) の射影不変量となる。ところが \((x)_0\) はパラメータと見ることができる。点 \(x=(x)_0\) における各不変量には、\((x)_0\) をふたたび \(x\) で置き換えれば、\(f(dx)\) の不変量が対応する。したがって、\(F_\rho(\delta u,du)\) が \(x=(x)_0\) で取った不変量 \(\Psi_\rho(\delta x,dx)\) になるならば、ここで \(x=(x)_0\) では \(du,\delta u\) が \(dx,\delta x\) へ移るので、\(\Psi_\rho\) 自身が完全系の基本関数を成す。残るのは、この基本関数系を \(f(dx)\) の明示的な不変量、すなわち函数 (8) によって表すことである。それが可能であることは、それらが微分過程によって生成されることから示される。最高位の項については、これらの過程は単純な極化過程へ移り、Clebsch--Gordan の級数展開に類似した変形を、無視した項に関する帰納法と組み合わせることで、主張が証明される。同時系の不変量を扱う場合には、\(f(dx)\) に属する正規座標をこれらの式に導入することから分かるように、基本関数 (8) にさらに残りの微分式の共変導関数を付け加えなければならない。

(14) に還元されていた \(f(dx)\) と \(g(dy)\) の同値性は、したがって、\(\Psi_\rho\) の同値性、または同じことだが函数 (8) の同値性、すなわち微分の線形変換に関する同値性とも同じ意味になる。このとき特別な可積分条件は不要であり、それは (14) への還元可能性から従う。これにより還元定理はすべての部分において証明された。とくに、函数列
\([\Omega_2],[\Omega_3],\ldots,[\Omega_\rho],\ldots\) が恒等的に消えることは、\(f(dx)\) が定数係数をもつ式へ変換できるための必要十分条件である。

最後に、すべての解析的変換から成る群の代わりに部分群を基礎に置くなら、対応する \(u\) の線形変換の群も射影群の部分群へ移る。\(f(dx)\) の不変量は、また函数 (8) の不変量となり、線形変換に関するものではあるが、今度はこの部分群に関するものとなる。したがって、斉次化によって、非斉次函数 \(f(dx)\) の場合はアフィン群へ還元できる。この場合、完全系は一つ多い変数について形成された函数 (8) から導き出される。

還元定理から、二次形式の場合、より一般に \(p\) 次元の形式の場合が、次のことにより得られる。ここでは函数系が \([\Omega_2]\)、より一般には \([\Omega_p]\) とその共変導関数で終わり、後続のすべての基本関数がこれらの射影不変量となるのである。これにより、二次形式における Riemann の曲率形式 \([\Omega_2]\) の特別な位置が説明される。二次形式、または \(p\) 次元の形式の同値性の問題は、まさに \([\Omega_2]\)、または \([\Omega_2]\ldots[\Omega_p]\)、およびそれらの共変導関数の、線形変換に関する同値性へ帰着する。そして \([\Omega_2]\)、または \([\Omega_2]\ldots[\Omega_p]\) が恒等的に消えることは、定数係数をもつ形式へ変換できるために必要十分である。これにより、二次形式について最もよく知られた結果の一つが述べられたことになる。

より詳しい叙述は、近く \emph{Mathematische Annalen} に掲載される予定である。

\clearpage

\begin{center}
{\Large\bfseries 13. 不変な変分問題}\par
\vspace{1.5em}
\emph{Nachrichten von der Gesellschaft der Wissenschaften zu Göttingen} 1918, pp. 235--257
\end{center}

\vspace{3em}
\begin{center}
1918年7月26日の会合で F. Klein により提出。\footnote{原稿の最終版が提出されたのは9月末になってからである。}
\end{center}

ここで扱うのは、（Lie の意味での）連続群を許す変分問題である。このことから対応する微分方程式について生じる帰結は、\S{}1 で定式化され、後続の節で証明される定理において最も一般的な表現を得る。変分問題から生じるこれらの微分方程式については、Lie の研究対象であった、群を許す任意の微分方程式についてよりも、はるかに精密な主張を行うことができる。したがって以下は、形式的変分計算の方法と Lie の群論の方法とを結合するものである。特殊な群と特殊な変分問題については、この方法の結合は新しいものではない。有限群の特殊な場合については Hamel と Herglotz を、無限群の特殊な場合については Lorentz とその弟子（例えば Fokker）、Weyl、Klein を挙げておく。\footnote{Hamel: \emph{Math. Ann.} vol. 59 および \emph{Zeitschrift f. Math. u. Phys.} vol. 50. Herglotz: \emph{Ann. d. Phys.} (4) vol. 36, とくに \S{}9, p. 511. Fokker, Verslag d. Amsterdamer Akad., 1917年1月27日。さらに文献については Klein の第二ノート、Göttinger Nachrichten, 1918年7月19日を参照。つい最近現れた Kneser の論文（\emph{Math. Zeitschrift} vol. 2）では、類似の方法による不変量の構成が問題になっている。} とくに、Klein の第二ノートと本稿の展開とは相互に影響し合っている。この点については Klein のノートの結語を参照してよい。

\section*{\S{}1. 予備的事項と定理の定式化}

以下に現れるすべての関数は、考察する領域で、解析的であるか、少なくとも連続で、有限回連続微分可能で、一価であると仮定する。

「変換群」とは、通常どおり、各変換に対してその逆変換が同じ系に含まれ、かつその系の任意の二つの変換の合成がふたたびその系に属するような変換の系をいう。群は、その変換が、\(\rho\) 個の本質的なパラメータ \(\eps\) に解析的に依存する最も一般の変換に含まれているとき、有限連続群 \(\G_\rho\) と呼ばれる（すなわち、これら \(\rho\) 個のパラメータは、より少数のパラメータの \(\rho\) 個の関数として表されてはならない）。これに対応して、無限連続群 \(\G_{\infty\rho}\) とは、その最も一般の変換が \(\rho\) 個の本質的な任意関数 \(p(x)\) とそれらの導関数に、解析的に、または少なくとも連続かつ有限回連続微分可能な仕方で依存する群をいう。両者の中間に、任意関数には依存しないが無限個のパラメータに依存する群がある。最後に、任意関数にもパラメータにも依存する群を混合群と呼ぶ。\footnote{Lie は \emph{Grundlagen für die Theorie der unendlichen kontinuierlichen Transformationsgruppen}（Berichte der Königl. Sächs. Gesellschaft der Wissenschaften, 1891）で、無限連続群を、変換が偏微分方程式系の最も一般の解で与えられ、かつそれらの解が有限個のパラメータだけに依存しないような変換群として定義している。これにより、上に挙げた型のうち、有限群とは異なる一つの型が得られる。しかし逆に、無限個のパラメータという極限の場合が、必ずしも微分方程式系を満たす必要はない。}

\(x_1,\ldots,x_n\) を独立変数、\(u_1(x),\ldots,u_\mu(x)\) をそれらに依存する関数とする。\(x\) と \(u\) をある群の変換に従わせるとき、変換が可逆であるという仮定により、変換された量の中には再びちょうど \(n\) 個の独立なもの \(y_1,\ldots,y_n\) が含まれていなければならない。これらに依存する残りの量を \(v_1(y),\ldots,v_\mu(y)\) と書く。変換には、\(x\) に関する \(u\) の導関数、すなわち \(\p u/\p x,\ \p^2u/\p x^2,\ldots\) も現れてよい。\footnote{添字は、和においても可能なかぎり省略する。したがって、例えば \(\p^2u/\p x^2\) は \(\p^2u_\nu/\p x_i\p x_k\) などを表す。} ある関数が群の不変量であるとは、関係式
\[
P\!\left(x,u,\frac{\p u}{\p x},\frac{\p^2u}{\p x^2},\ldots\right)
 =
P\!\left(y,v,\frac{\p v}{\p y},\frac{\p^2v}{\p y^2},\ldots\right)
\]
が成り立つことをいう。とくに、積分 \(I\) が群の不変量であるとは、関係式
\begin{equation}
\begin{aligned}
I&=\int\!\cdots\!\int f\!\left(x,u,\frac{\p u}{\p x},\frac{\p^2u}{\p x^2},\ldots\right)\dd x\\
 &=\int\!\cdots\!\int f\!\left(y,v,\frac{\p v}{\p y},\frac{\p^2v}{\p y^2},\ldots\right)\dd y
\end{aligned}
\tag{1}
\end{equation}
が成り立つことである。ここで積分は任意の実 \(x\)-領域と、対応する \(y\)-領域にわたる。\footnote{簡単のため \(\dd x,\dd y\) と書いて、それぞれ \(\dd x_1\cdots \dd x_n,\dd y_1\cdots \dd y_n\) を表す。変換に現れるすべての引数 \(x,u,\eps,p(x)\) は実とするが、係数は複素数であってよい。しかし最終結果は \(x,u\)、パラメータ、任意関数に関する恒等式であるから、現れる関数を解析的と仮定すれば複素値についても成立する。さらに、結果の大部分は積分を用いずに基礎づけることができるので、その証明においても実数への制限は必要ではない。これに対し、\S{}2 の末尾と \S{}5 の初めの考察は、積分なしには実行できないようである。}

他方、必ずしも不変とは限らない任意の積分 \(I\) に対して、私は第一変分 \(\delta I\) を作り、変分計算の規則に従って部分積分で変形する。よく知られているように、\(\delta u\) および境界に現れるそのすべての導関数を境界で零とし、その他では任意とすると、次を得る：
\begin{equation}
\delta I=\int\!\cdots\!\int \delta f\,\dd x
=\int\!\cdots\!\int\left(\sum_i\psi_i\!\left(x,u,\frac{\p u}{\p x},\ldots\right)\delta u_i\right)\dd x,
\tag{2}
\end{equation}
ここで \(\psi\) はラグランジュ式を表す。すなわち、対応する変分問題 \(\delta I=0\) のラグランジュ方程式の左辺である。この積分関係には、境界項も書き出すことによって生じる、\(\delta u\) とその導関数に関する積分を含まない恒等式が対応する。部分積分が示すように、これらの境界項は発散の積分、すなわち
\[
\Div A=\frac{\p A_1}{\p x_1}+\cdots+\frac{\p A_n}{\p x_n}
\]
という形の式の積分である。ただし \(A\) は \(\delta u\) とその導関数に関して線形である。したがって
\begin{equation}
\sum_i\psi_i\,\delta u_i=\delta f+\Div A.
\tag{3}
\end{equation}
を得る。とくに \(f\) が \(u\) の一階導関数だけを含むなら、単積分の場合には恒等式 (3) はいわゆる「ラグランジュの中央方程式」と同一である：
\begin{equation}
\sum_i \psi_i\,\delta u_i
=\delta f-\frac{\dd}{\dd x}\left(\sum_i\frac{\p f}{\p u_i'}\,\delta u_i\right),
\qquad \left(u_i'=\frac{\dd u_i}{\dd x}\right),
\tag{4}
\end{equation}
一方、\(n\)-重積分に対して (3) は
\begin{equation}
\sum_i\psi_i\,\delta u_i
=
\delta f
-\frac{\p}{\p x_1}\left(\sum_i\frac{\p f}{\p \left(\frac{\p u_i}{\p x_1}\right)}\,\delta u_i\right)
-\cdots
-\frac{\p}{\p x_n}\left(\sum_i\frac{\p f}{\p \left(\frac{\p u_i}{\p x_n}\right)}\,\delta u_i\right)
\tag{5}
\end{equation}
となる。単積分で、\(u\) の \(\kappa\) 階導関数まで含む場合、(3) は次で与えられる：
\begin{equation}
\begin{aligned}
\sum_i\psi_i\,\delta u_i
={}&\delta f
-\frac{\dd}{\dd x}\Bigg\{\sum_i\left[
\binom{1}{1}\frac{\p f}{\p u_i^{(1)}}\delta u_i
+\binom{2}{1}\frac{\p f}{\p u_i^{(2)}}\delta u_i^{(1)}
+\cdots+
\binom{\kappa}{1}\frac{\p f}{\p u_i^{(\kappa)}}\delta u_i^{(\kappa-1)}
\right]\Bigg\}\\
&+\frac{\dd^2}{\dd x^2}\Bigg\{\sum_i\left[
\binom{2}{2}\frac{\p f}{\p u_i^{(2)}}\delta u_i
+\binom{3}{2}\frac{\p f}{\p u_i^{(3)}}\delta u_i^{(1)}
+\cdots+
\binom{\kappa}{2}\frac{\p f}{\p u_i^{(\kappa)}}\delta u_i^{(\kappa-2)}
\right]\Bigg\}\\
&+\cdots
+(-1)^\kappa\frac{\dd^\kappa}{\dd x^\kappa}
\left\{\sum_i\binom{\kappa}{\kappa}
\frac{\p f}{\p u_i^{(\kappa)}}\delta u_i\right\}.
\end{aligned}
\tag{6}
\end{equation}
対応する恒等式は \(n\)-重積分に対しても成り立つ。とくに \(A\) は \((\kappa-1)\) 階導関数までの \(\delta u\) を含む。(4), (5), (6) が実際にラグランジュ式 \(\psi_i\) を定義することは、右辺の組合せにおいて \(\delta u\) の高階導関数がすべて消去される一方で、部分積分により一意的に導かれる関係 (2) が満たされることから従う。

以下で問題となるのは、次の二つの定理である。

I. 積分 \(I\) が \(\G_\rho\) に関して不変であるなら、ラグランジュ式の \(\rho\) 個の線形独立な組合せは発散になる。逆に、このことから \(I\) の \(\G_\rho\) に関する不変性が従う。この定理は、無限個のパラメータという極限の場合にも成り立つ。

II. 積分 \(I\) が、任意関数が \(\sigma\) 階導関数まで現れる \(\G_{\infty\rho}\) に関して不変であるなら、ラグランジュ式とその \(\sigma\) 階までの導関数との間に \(\rho\) 個の恒等的な関係が存在する。ここでも逆が成り立つ。\footnote{ある自明な例外については \S{}2 の第二注を参照。}

混合群については、両定理の主張が適用される。したがって、互いに独立な従属関係と発散関係が同時に現れる。

これらの恒等式から対応する変分問題へ移り、したがって \(\psi=0\) とおくと、\footnote{少し一般化して \(\psi_i=T_i\) とおくこともできる。\S{}3 の第一注を参照。} 定理 I は一次元の場合、すなわち発散が全微分となる場合に、\(\rho\) 個の第一積分が存在することを述べる。ただしそれらの間に非線形の従属が生じることはありうる。\footnote{\S{}3 の末尾を参照。} 多次元の場合には、近年しばしば「保存則」と呼ばれている発散方程式が得られる。定理 II は、ラグランジュ方程式のうち \(\rho\) 個が残りの方程式の帰結であることを述べる。

定理 II の最も簡単な例（逆は除く）は Weierstrass のパラメータ表示である。ここでは一階の同次性の下で、独立変数 \(x\) を \(x\) の任意関数で置き換え、\(u\) を変えない \((y=p(x);\ v_i(y)=u_i(x))\) とき、積分が不変であることが知られている。したがって任意関数が一つ現れるが、導関数は現れない。これに対応して、ラグランジュ式そのものの間に有名な線形関係
\[
\sum_i \psi_i\,\frac{\dd u_i}{\dd x}=0
\]
がある。さらに別の例は、物理学者の「一般相対性理論」によって与えられる。ここではすべての \(x\) の変換 \(y_i=p_i(x)\) の群が問題となり、\(u\)（\(g_{\mu\nu}\) および \(q\) と記される）は、それによって二次および一次の微分形式の係数に誘導される変換を受ける。その変換には任意関数 \(p(x)\) の一階導関数が含まれる。これに対応して、ラグランジュ式とその一階導関数との間の既知の \(n\) 個の従属関係がある。\footnote{例えば Klein の叙述を参照。}

とくに、群を特殊化して、変換に \(u(x)\) の導関数が現れないようにし、さらに変換された独立量が \(u\) ではなく \(x\) だけに依存するようにすると、\S{}5 で示されるように、\(I\) の不変性から \(\sum_i\psi_i\delta u_i\) の相対不変性\footnote{すなわち、\(\sum_i\psi_i\delta u_i\) は変換の下である因子を獲得する。} が従う。また、パラメータを適切な変換に従わせるなら、定理 I に現れる発散についても同様である。さらに上に述べた第一積分も群を許すことが従う。定理 II についても同様に、任意関数によって結合された従属関係の左辺の相対不変性が得られ、その帰結として、発散が恒等的に消え、かつ群を許す関数も得られる。これは物理学者の相対性理論において、従属関係とエネルギー定理との連関を媒介する関数である。\footnote{Klein の第二ノートを参照。} 最後に、群論的な形で、定理 II は「一般相対性」において真性エネルギー定理が失敗することに関する Hilbert の関連する主張の証明を与える。これらの補足的な注を加えれば、定理 I は力学などで知られている第一積分に関するすべての定理を含み、定理 II は「一般相対性理論」の可能な限り最大の群論的一般化と述べることができる。

\section*{\S{}2. 発散関係と従属関係}

\(\G\) を、有限または無限の連続群とする。恒等変換が、パラメータ \(\eps\)、それぞれ任意関数 \(p(x)\) の値ゼロに対応するよう、常に整えることができる。\footnote{例えば Lie, \emph{Grundlagen}, p. 331 を参照。任意関数が関与する場合、パラメータの特殊値 \(a^\sigma\) は固定された関数 \(p^\sigma,\ \p p^\sigma/\p x,\ldots\) に置き換えられるべきであり、それに対応して値 \(a^\sigma+\eps\) は \(p^\sigma+p(x),\ \p p^\sigma/\p x+\p p/\p x\) などに置き換えられる。} したがって最も一般の変換は次の形をもつ：
\[
y_i=A_i\!\left(x,u,\frac{\p u}{\p x},\ldots\right)=x_i+\D x_i+\cdots,\qquad
v_i(y)=B_i\!\left(x,u,\frac{\p u}{\p x},\ldots\right)=u_i+\D u_i+\cdots,
\]
ここで \(\D x_i,\D u_i\) は、\(\eps\)、それぞれ \(p(x)\) とその導関数における最低次の項を表し、これらの量に関して線形であると仮定する。後で示すように、これは一般性を制限しない。

いま積分 \(I\) が \(\G\) に関して不変であり、したがって関係 (1) が満たされているとする。とくに \(I\) は、\(\G\) に含まれる infinitesimal 変換
\[
y_i=x_i+\D x_i,\qquad v_i(y)=u_i+\D u_i
\]
に関しても不変である。このとき関係 (1) は
\begin{equation}
\begin{aligned}
0=\D I
&=\int\!\cdots\!\int f\!\left(y,v(y),\frac{\p v}{\p y},\ldots\right)\dd y\\
&\qquad-\int\!\cdots\!\int f\!\left(x,u(x),\frac{\p u}{\p x},\ldots\right)\dd x
\end{aligned}
\tag{7}
\end{equation}
となる。第一の積分は、\(x\)-領域に対応する \(x+\D x\)-領域にわたって拡張される。しかしこの積分は、infinitesimal な \(\D x\) に対して有効な変換により、\(x\)-領域上の積分へ変換することもできる：
\begin{equation}
\begin{aligned}
\int\!\cdots\!\int f\!\left(y,v(y),\frac{\p v}{\p y},\ldots\right)\dd y
&=\int\!\cdots\!\int f\!\left(x,v(x),\frac{\p v}{\p x},\ldots\right)\dd x\\
&\quad+\int\!\cdots\!\int \Div(f\cdot \D x)\,\dd x.
\end{aligned}
\tag{8}
\end{equation}
したがって、infinitesimal 変換 \(\D u\) の代わりに変分
\begin{equation}
\bd u_i=v_i(x)-u_i(x)=\D u_i-\sum_\lambda\frac{\p u_i}{\p x_\lambda}\D x_\lambda
\tag{9}
\end{equation}
を導入すると、(7), (8) は
\begin{equation}
0=\int\!\cdots\!\int\{\bd f+\Div(f\cdot\D x)\}\dd x.
\tag{10}
\end{equation}
となる。関係 (10) は任意の領域上で積分したとき成立するので、被積分関数は恒等的に消えなければならない。したがって \(I\) の不変性に対する Lie の微分方程式は、関係
\begin{equation}
\bd f+\Div(f\cdot\D x)=0
\tag{11}
\end{equation}
へ移る。ここで (3) によって \(\bd f\) をラグランジュ式で表すと
\begin{equation}
\sum_i\psi_i\bd u_i=\Div B,\qquad (B=A-f\cdot\D x)
\tag{12}
\end{equation}
を得る。したがってこの関係は、任意の不変積分 \(I\) に対して、現れるすべての引数に関する恒等式を表す。これは \(I\) に対する Lie の微分方程式の求める形である。\footnote{\(\Div(f\cdot\D x)=0,\ \bd u=0\) となる自明な場合には、(12) は \(0=0\) になる。この場合は、\(\D x,\D u\) が \(u\) の導関数にも依存するときにだけ起こりうる。したがってこれらの infinitesimal 変換は常に群から切り離すべきであり、定理の定式化においては残りのパラメータまたは任意関数の数だけを数えるべきである。残りの infinitesimal 変換がなお常に群をなすかどうかは未決としておく必要がある。}

まず \(\G\) が有限連続群 \(\G_\rho\) であると仮定する。仮定により \(\D u\) と \(\D x\) はパラメータ \(\eps_1,\ldots,\eps_\rho\) に関して線形であるから、(9) により \(\bd u\) とその導関数にも同じことが成り立つ。したがって \(A\) と \(B\) は \(\eps\) に関して線形である。そこで
\[
B=B^{(1)}\eps_1+\cdots+B^{(\rho)}\eps_\rho,\qquad
\bd u=\bd u^{(1)}\eps_1+\cdots+\bd u^{(\rho)}\eps_\rho
\]
とおく。ここで \(\bd u^{(1)},\ldots\) は \(x,u,\p u/\p x,\ldots\) の関数である。(12) から求める発散関係
\begin{equation}
\sum_i\psi_i\bd u_i^{(1)}=\Div B^{(1)},\ldots,\qquad
\sum_i\psi_i\bd u_i^{(\rho)}=\Div B^{(\rho)}
\tag{13}
\end{equation}
が従う。したがって、ラグランジュ式の \(\rho\) 個の線形独立な組合せが発散になる。線形独立性は、(9) により \(\bd u=0,\D x=0\) が \(\D u=0,\D x=0\) を含意し、したがって infinitesimal 変換の間の従属関係を含意することから従う。しかし仮定により、そのような従属関係はいかなるパラメータ値に対しても存在しない。そうでなければ、infinitesimal 変換から積分により再び生じる \(\G_\rho\) は、\(\rho\) より少ない本質的パラメータに依存することになるからである。さらに \(\bd u=0,\ \Div(f\D x)=0\) という可能性は除外してある。これらの結論は、無限個のパラメータという極限の場合にも有効である。

次に \(\G\) が無限連続群 \(\G_{\infty\rho}\) であるとする。このときも \(\bd u\) とその導関数、したがって \(B\) も、任意関数 \(p(x)\) とその導関数に関して線形である。\footnote{\(p\) が \(u,\p u/\p x,\ldots\) に依存しないと仮定することが制限でないことは、逆定理によって示される。} (12) とはまだ独立に、\(\bd u\) の値を代入すると
\[
\sum_i\psi_i\bd u_i
=\sum_{\lambda,i}\psi_i\left\{
a_i^{(\lambda)}(x,u,\ldots)p^{(\lambda)}(x)
+b_i^{(\lambda)}(x,u,\ldots)\frac{\p p^{(\lambda)}}{\p x}
+\cdots+c_i^{(\lambda)}(x,u,\ldots)
\frac{\p^\sigma p^{(\lambda)}}{\p x^\sigma}\right\}
\]
となる。いま、部分積分の公式と同様に、恒等式
\[
\varphi(x,u,\ldots)\frac{\p^\tau p(x)}{\p x^\tau}
=(-1)^\tau\frac{\p^\tau\varphi}{\p x^\tau}\,p(x)
\quad \bmod \text{ 発散}
\]
によって、\(p\) の導関数は \(p\) 自身と、\(p\) およびその導関数に関して線形な発散とに置き換えられる。したがって
\begin{equation}
\begin{aligned}
\sum_i\psi_i\bd u_i
={}&\sum_\lambda\left\{
a_i^{(\lambda)}\psi_i-\frac{\p}{\p x}\bigl(b_i^{(\lambda)}\psi_i\bigr)
+\cdots+(-1)^\sigma
\frac{\p^\sigma}{\p x^\sigma}\bigl(c_i^{(\lambda)}\psi_i\bigr)
\right\}p^{(\lambda)}+\Div\Gamma
\end{aligned}
\tag{14}
\end{equation}
を得る。ただし括弧内では \(i\) について和をとる。これを (12) と組み合わせると
\begin{equation}
\sum_\lambda\left\{
a_i^{(\lambda)}\psi_i-\frac{\p}{\p x}\bigl(b_i^{(\lambda)}\psi_i\bigr)
+\cdots+(-1)^\sigma
\frac{\p^\sigma}{\p x^\sigma}\bigl(c_i^{(\lambda)}\psi_i\bigr)
\right\}p^{(\lambda)}
=\Div(B-\Gamma)
\tag{15}
\end{equation}
となる。いま (15) の \(n\)-重積分を任意の領域上で作り、\(p(x)\) を、\(B-\Gamma\) に現れるすべての導関数とともに境界で消えるように選ぶ。発散の積分は境界積分に帰着されるので、(15) の左辺の積分も、十分多くの導関数とともに境界で消える任意の \(p(x)\) について消える。通常の議論から、被積分関数はすべての \(p(x)\) について消えることが従い、\(\rho\) 個の関係
\begin{equation}
\sum_i\left\{
a_i^{(\lambda)}\psi_i-\frac{\p}{\p x}\bigl(b_i^{(\lambda)}\psi_i\bigr)
+\cdots+(-1)^\sigma
\frac{\p^\sigma}{\p x^\sigma}\bigl(c_i^{(\lambda)}\psi_i\bigr)
\right\}=0,\qquad(\lambda=1,2,\ldots,\rho)
\tag{16}
\end{equation}
が得られる。これらが、\(I\) が \(\G_{\infty\rho}\) に関して不変であるときの、ラグランジュ式とその導関数との間の求める従属関係である。線形独立性は前と同様に従う。というのも、逆定理は (12) へ戻すからであり、また infinitesimal 変換から有限変換へふたたび結論できるからである。このことは \S{}4 でより詳しく行う。したがって、\(\G_{\infty\rho}\) については、\(\rho\) 個の任意変換がすでに infinitesimal 変換に現れる。(15) と (16) からさらに \(\Div(B-\Gamma)=0\) が従う。

「混合群」に対応して \(\D x\) と \(\D u\) を \(\eps\) と \(p(x)\) に関して線形にとるなら、一度 \(p(x)\) を零に、もう一度 \(\eps\) を零におくことにより、発散関係 (13) と従属関係 (16) の双方が存在することがわかる。

\section*{\S{}3. 有限群の場合の逆}

逆を証明するには、まず本質的に前の考察を逆順にたどらなければならない。(13) が成り立つことから、\(\eps\) を掛けて加えると (12) が成り立つ。さらに恒等式 (3) により関係
\[
\bd f+\Div(A-B)=0
\]
を得る。したがって \(\D x=\frac1f(A-B)\) とおけば (11) に到達し、積分により最終的に (7), \(\D I=0\) を得る。すなわち \(I\) は \(\D x,\D u\) で決まる infinitesimal 変換に関して不変である。ここで \(\D u\) は (9) により \(\D x\) と \(\bd u\) から決まり、\(\D x,\D u\) はパラメータに関して線形になる。しかし \(\D I=0\) は、よく知られているように、同時系
\begin{equation}
\frac{\dd x_i}{\dd t}=\D x_i,\qquad
\frac{\dd u_i}{\dd t}=\D u_i,\qquad
\left\{\begin{array}{l}
x_i=y_i,\\
u_i=v_i
\end{array}\right\}
\quad\text{for }t=0.
\tag{17}
\end{equation}
の積分により生じる有限変換に関する \(I\) の不変性を含意する。これらの有限変換は、\(t\eps_1,\ldots,t\eps_\rho\) という組合せ、すなわち \(a_1,\ldots,a_\rho\) という \(\rho\) 個のパラメータを含む。\(\rho\) 個、しかも \(\rho\) 個だけの線形独立な発散関係 (13) があるという仮定から、さらに、有限変換は、それが導関数 \(\p u/\p x\) を含まないかぎり、常に群をなすことが従う。反対の場合、Lie の括弧過程によって生じる少なくとも一つの infinitesimal 変換が、残りの \(\rho\) 個の線形結合ではなくなる。そして \(I\) もこの変換を許すので、\(\rho\) 個を超える線形独立な発散関係が存在することになる。あるいは、この infinitesimal 変換は \(\bd u=0,\ \Div(f\D x)=0\) という特殊な形であろう。しかしその場合、\(\D x\) または \(\D u\) は導関数に依存し、仮定に反する。\(\D x\) または \(\D u\) に導関数が現れる場合にこの事態が起こりうるかどうかは未決としておかなければならない。その場合、群性をふたたび得るためには、上で決められた \(\D x\) に、\(\Div(f\D x)=0\) を満たすすべての関数 \(\D x\) をさらに加えなければならないが、こうして加えられるパラメータは約束により数えない。これで逆が証明された。

この逆からまた、\(\D x\) と \(\D u\) を実際にパラメータに関して線形と仮定してよいことも従う。もし \(\D u\) と \(\D x\) が \(\eps\) の高次形式であったなら、\(\eps\) の冪積の線形独立性により、(13) と完全に対応する関係が、ただしより多数に、従うであろう。逆定理により、それらは infinitesimal 変換がパラメータを線形に含むような群に関する \(I\) の不変性を含意する。この群がちょうど \(\rho\) 個のパラメータを含むべきなら、\(\eps\) の高次項から本来得られた発散関係の間に線形従属がなければならない。

さらに、\(\D x\) と \(\D u\) が \(u\) の導関数も含む場合、有限変換は \(u\) の無限に多くの導関数に依存しうることに注意すべきである。実際この場合、(17) を積分して
\[
\frac{\dd^2x_i}{\dd t^2},\quad \frac{\dd^2u_i}{\dd t^2}
\]
を決めるとき、
\[
\D\!\left(\frac{\p u}{\p x_\lambda}\right)
=\frac{\p\,\D u}{\p x_\lambda}
-\sum_\nu\frac{\p u}{\p x_\nu}
\frac{\p\,\D x_\nu}{\p x_\lambda}
\]
へ導かれるので、一般には各段階で \(u\) の導関数の数が増える。例は
\[
f=\frac12 u'^2,\qquad
\psi=-u'',\qquad
\psi\cdot x=\frac{\dd}{\dd x}\{(u-u'x)\},\qquad
\bd u=x\eps,
\]
\[
\D x=-\frac{2u}{u'}\,\eps,\qquad
\D u=\left(x-\frac{2u}{u'}\right)\eps.
\]
発散のラグランジュ式は恒等的に消えるから、逆は最後に次を示す。\(I\) が \(\G_\rho\) を許すなら、\(I\) と境界積分だけ、すなわち発散の積分だけ異なる任意の積分も、同じ \(\bd u\) をもつ \(\G_\rho\) を許す。ただしその infinitesimal 変換は一般に \(u\) の導関数を含む。したがって前の例に対応して
\[
f^*=\frac12\left\{u'^2-\frac{\dd}{\dd x}\left(\frac{u^2}{x}\right)\right\}
\]
は infinitesimal 変換 \(\D u=x\eps,\ \D x=0\) を許す。一方、\(f\) に対応する infinitesimal 変換には \(u\) の導関数が現れる。

変分問題へ移る、すなわち \(\psi_i=0\) とおくと、\footnote{\(\psi_i=0\)、またはやや一般に \(\psi_i=T_i\)（ここで \(T_i\) は新たに付加された関数）は、物理学では「場の方程式」と呼ばれる。\(\psi_i=T_i\) の場合、恒等式 (13) は \(\Div B^{(\lambda)}=\sum_i T_i\bd u_i^{(\lambda)}\) という方程式に移り、物理学ではこれもなお保存則と呼ばれる。} (13) は方程式
\[
\Div B^{(1)}=0,\ldots,\qquad \Div B^{(\rho)}=0
\]
になる。これはしばしば「保存則」と呼ばれる。一次元の場合にはこれは
\[
B^{(1)}=\mathrm{const.},\ldots,\qquad B^{(\rho)}=\mathrm{const.}
\]
を与える。ここで \(B\) は、\(\D u\) と \(\D x\) が \(f\) に現れる \(\kappa\) 階導関数より高い導関数を含まないかぎり、(6) により高々 \((2\kappa-1)\) 階導関数までの \(u\) を含む。\(\psi\) は一般に \(2\kappa\) 階導関数を含むので、\footnote{\(f\) が \(\kappa\) 階導関数に関して非線形である場合。} こうして \(\rho\) 個の第一積分の存在が得られる。それらの間に非線形従属がありうることは、再び上の \(f\) が示す。線形独立な \(\D u=\eps_1,\ \D x=\eps_2\) には、線形独立な関係
\[
u''=\frac{\dd}{\dd x}u',\qquad
u''u'=\frac12\frac{\dd}{\dd x}(u')^2
\]
が対応する。一方、第一積分 \(u'=\mathrm{const.},\ u'^2=\mathrm{const.}\) の間には非線形従属がある。これは \(\D u,\D x\) が \(u\) の導関数を含まない初等的な場合である。\footnote{そうでない場合、さらに各 \(\lambda\) について \(u''u'^{\lambda-1}=\mathrm{const.}\) があり、これは
\[
u''(u')^{\lambda-1}=\frac1\lambda\frac{\dd}{\dd x}(u')^\lambda
\]
に対応する。}

\section*{\S{}4. 無限群の場合の逆}

まず \(\D x\) と \(\D u\) の線形性の仮定が制限でないことを示す必要がある。ここではこれは、逆定理を使わなくても、\(\G_{\infty\rho}\) が形式的に \(\rho\) 個、しかも \(\rho\) 個だけの任意関数に依存するという事実から従う。すなわち非線形の場合には、変換を合成すると最低次の項が加わり、任意関数の数が増える。例えば
\[
\begin{aligned}
y&=A\!\left(x,u,\frac{\p u}{\p x},\ldots;p\right)\\
&=x+\sum a(x,u,\ldots)p^\nu
+b(x,u,\ldots)p^{\nu-1}\frac{\p p}{\p x}
+c\,p^{\nu-2}\left(\frac{\p p}{\p x}\right)^2+\cdots
+d\left(\frac{\p p}{\p x}\right)^\nu+\cdots,
\end{aligned}
\]
ただし
\[
p^\nu=(p^{(1)})^{\nu_1}\cdots(p^{(\rho)})^{\nu_\rho},
\]
であり、対応して
\[
v=B\!\left(x,u,\frac{\p u}{\p x},\ldots;p\right)
\]
であるとする。このとき
\[
z=A\!\left(y,v,\frac{\p v}{\p y},\ldots;q\right)
\]
との合成により、最低次の項は
\[
z=x+\sum a(p^\nu+q^\nu)+
b\left\{p^{\nu-1}\frac{\p p}{\p x}+q^{\nu-1}\frac{\p q}{\p x}\right\}
+c\left\{p^{\nu-2}\left(\frac{\p p}{\p x}\right)^2
+q^{\nu-2}\left(\frac{\p q}{\p x}\right)^2\right\}+\cdots
\]
を与える。ここで \(a\) と \(b\) 以外のある係数が零でなく、したがって
\[
p^{\nu-\sigma}\left(\frac{\p p}{\p x}\right)^\sigma
+q^{\nu-\sigma}\left(\frac{\p q}{\p x}\right)^\sigma
\]
という項が \(\sigma>1\) について実際に現れるなら、それは単一の関数の微分商としても、そのような商の冪の積としても書けない。したがって仮定に反して任意関数の数が増えている。\(a\) と \(b\) 以外の係数がすべて零なら、指数 \(\nu_1,\ldots,\nu_\rho\) の値に応じて、第二項は第一項の微分商である（例えば \(\G_{\infty1}\) では常にそうである）ので線形性が実際に現れるか、あるいはここでも任意関数の数が増える。したがって \(p(x)\) に関する線形性により、infinitesimal 変換は線形偏微分方程式系を満たす。そして群性が成り立つので、それらは Lie の意味での「infinitesimal 変換の無限群」をなす（\emph{Grundlagen}, \S{}10）。

逆は有限群の場合と同様に得られる。従属関係 (16) の存在から、\(p^{(\lambda)}(x)\) を掛けて加えると、恒等変換 (14) により
\[
\sum_i\psi_i\bd u_i=\Div\Gamma
\]
が導かれる。ここから \S{}3 と同様に \(\D x\) と \(\D u\) の決定、およびこれらの infinitesimal 変換に関する \(I\) の不変性が得られる。これらは実際に \(\rho\) 個の任意関数とその \(\sigma\) 階までの導関数に線形に依存する。これらの infinitesimal 変換が、\(\p u/\p x,\ldots\) を含まないとき、確かに群をなすことは \S{}3 と同じく従う。そうでなければ合成によりさらに多くの任意関数が現れるが、仮定では従属関係 (16) は \(\rho\) 個だけでなければならないからである。したがってそれらは「infinitesimal 変換の無限群」をなす。しかしそのような群は（\emph{Grundlagen}, 定理 VII, p. 391）、Lie の意味でそれによって定義されるある「有限変換の無限群 \(\G\)」の最も一般の infinitesimal 変換からなる。各有限変換は infinitesimal なものから生成される（\emph{Grundlagen}, \S{}7）。\footnote{ここから、とくに、\(\G_{\infty\rho}\) の infinitesimal 変換 \(\D x,\D u\) から生成される群 \(\G\) は、再び \(\G_{\infty\rho}\) に戻ることが従う。というのも \(\G_{\infty\rho}\) は、任意関数に依存し \(\D x,\D u\) と異なる infinitesimal 変換を含まず、またパラメータに依存する独立な変換も含みえないからである。もしそうであれば混合群になってしまう。しかし前の議論により、有限変換は infinitesimal 変換によって決定される。} したがってこれは同時系
\[
\frac{\dd x_i}{\dd t}=\D x_i,\qquad
\frac{\dd u_i}{\dd t}=\D u_i,\qquad
\left\{\begin{array}{l}
x_i=y_i,\\
u_i=v
\end{array}\right\}
\quad\text{for }t=0.
\]
の積分によって生じる。ただし、任意の \(p(x)\) をさらに \(t\) に依存させる必要があるかもしれない。したがって \(\G\) は実際に \(\rho\) 個の任意関数に依存する。とくに \(p(x)\) を \(t\) に依存しないものとしてよいなら、この依存は任意関数 \(q(x)=t\,p(x)\) に解析的である。\footnote{この後者の場合が常に起こるのではないかという問題は、Lie によって別の定式化で提起された（\emph{Grundlagen}, \S{}7 および \S{}13 末尾）。} \(\p u/\p x,\ldots\) が現れるなら、同じ結論を引き出す前に infinitesimal 変換 \(\bd u=0,\ \Div(f\D x)=0\) を付加する必要があるかもしれない。

Lie の例（\emph{Grundlagen}, \S{}7）に従って、ここで、明示公式にまで到達できるかなり一般的な場合を示す。その公式は、任意関数の導関数が \(\sigma\) 階まで現れることも示している。したがってこの場合、逆は完全である。これは、すべての \(x\) の変換の群と、それによって「誘導される」\(u\) の変換の群が対応する infinitesimal 変換群である。すなわち \(\D u\)、したがって \(u\) が、\(\D x\) に現れる任意関数だけに依存するような \(u\) の変換である。なお \(\D u\) には導関数 \(\p u/\p x,\ldots\) が現れないと仮定する。したがって
\[
\D x_i=p^{(i)}(x),\qquad
\D u_i=
\sum_{\lambda=1}^{n}\left\{
a_i^{(\lambda)}(x,u)p^{(\lambda)}
+b_i^{(\lambda)}\frac{\p p^{(\lambda)}}{\p x}
+\cdots+
c_i^{(\lambda)}\frac{\p^\sigma p^{(\lambda)}}{\p x^\sigma}
\right\}.
\]
infinitesimal 変換 \(\D x=p(x)\) は任意の \(g(y)\) をもつすべての変換 \(x=y+g(y)\) を生成するので、特に \(p(x)\) を \(t\) の関数として、一次元パラメータ群
\begin{equation}
x_i=y_i+t\,g_i(y)
\tag{18}
\end{equation}
が生成されるよう決定できる。これは \(t=0\) で恒等変換となり、\(t=1\) で求める \(x=y+g(y)\) になる。(18) を微分すると
\begin{equation}
\frac{\dd x_i}{\dd t}=g_i(y)=p^{(i)}(x,t)
\tag{19}
\end{equation}
を得る。ここで \(p(x,t)\) は (18) の反転により \(g(y)\) から決まる。逆に (18) は、\(t=0\) で \(x_i=y_i\) という副条件により、(19) から生じる。この副条件により積分は一意に定まる。(18) により、\(\D u\) における \(x\) は積分定数 \(y\) と \(t\) で置き換えられる。そのとき \(g(y)\) はちょうど \(\sigma\) 階導関数まで現れる。なぜなら
\[
\frac{\p p}{\p x}=\sum_k\frac{\p g}{\p y_k}\frac{\p y_k}{\p x}
\]
において \(\p y/\p x\) は \(\p x/\p y\) によって表され、一般に \(\p^\sigma p/\p x^\sigma\) は
\[
\frac{\p g}{\p y},\ldots,\frac{\p^\sigma g}{\p y^\sigma},\ldots,\frac{\p^\sigma x}{\p y^\sigma}
\]
における値で置き換えられるからである。\(u\) の決定については、したがって方程式系
\[
\frac{\dd u_i}{\dd t}
=
F_i\!\left(g(y),\frac{\p g}{\p y},\ldots,\frac{\p^\sigma g}{\p y^\sigma},u,t\right),
\qquad (u_i=v_i\ \text{for }t=0),
\]
を得る。この中では \(t\) と \(u\) だけが変数であり、\(g(y),\ldots\) は係数域に属する。したがって積分により
\[
u_i=v_i+
B_i\!\left(v,g(y),\frac{\p g}{\p y},\ldots,\frac{\p^\sigma g}{\p y^\sigma},t\right)_{t=1}.
\]
が与えられる。こうして、任意関数の導関数 \(\sigma\) 階にちょうど依存する変換が得られる。(18) により \(g(y)=0\) の場合に恒等変換が含まれる。また群性は、いま述べた手続きがすべての変換 \(x=y+g(y)\) を与え、それによって \(u\) の誘導変換が一意に決定されるので、群 \(\G\) が尽くされるという事実から従う。

逆定理から、任意関数が \(u,\p u/\p x,\ldots\) ではなく \(x\) のみに依存すると仮定しても制限でないことも付随的に従う。後者の場合、恒等変換 (14)、したがって (15) において、\(p^{(\lambda)}\) のほかに \(\p p^{(\lambda)}/\p u,\ \p p^{(\lambda)}/\p(\p u/\p x),\ldots\) という導関数が現れるであろう。ここで \(p^{(\lambda)}\) を順に、\(u,\p u/\p x,\ldots\) に関する 0 次、1 次、さらに高次のものとし、その係数を \(x\) の任意関数とするなら、再び従属関係 (16) が生じる。ただし数は増える。しかし前の逆定理により、それらは \(x\) のみに依存する任意関数と組み合わせることによって先の場合に戻すことができる。同様に、混合群には従属関係と独立な発散関係の同時出現が対応することが示される。\footnote{\S{}3 と同じく、ここでも逆定理から、\(I\) のほか、それと発散の積分だけ異なる任意の積分 \(I^*\) も、同じ \(\bd u\) をもつ無限群を許すことが従う。ただしここでは \(\D x\) と \(\D u\) は一般に \(u\) の導関数を含む。Einstein は、エネルギー定理のより簡単な定式化を得るため、このような積分 \(I^*\) を一般相対性理論に導入した。Klein の第二ノートの記法をそのまま用いて、この \(I^*\) が許す infinitesimal 変換を与える。積分 \(I=\int\!\cdots\!\int K\,\dd\omega=\int\!\cdots\!\int \mathfrak{R}\,\dd S\) は、すべての \(w\) の変換の群と、それによって \(g_{\mu\nu}\) に誘導される群を許す。これに対応する従属関係（Klein の (30)）は
\[
\sum \mathfrak{R}_{\mu\nu}g_{\mu\nu}^{\alpha\beta}
+2\sum\frac{\p g_{\mu\nu}^{\alpha\beta}\mathfrak{R}_{\mu\nu}}{\p w^\sigma}=0.
\]
いま \(I^*=\int\!\cdots\!\int \mathfrak{R}^*\,\dd S\)、ただし \(\mathfrak{R}^*=\mathfrak{R}+\Div\) であり、したがって \(\mathfrak{R}_{\mu\nu}^*=\mathfrak{R}_{\mu\nu}\) である。ここで \(\mathfrak{R}_{\mu\nu}^*,\mathfrak{R}_{\mu\nu}\) はそれぞれラグランジュ式を表す。よって与えられた従属関係は \(\mathfrak{R}_{\mu\nu}^*\) に対する従属関係でもある。さらに \(p^\tau\) を掛けて加え、積の微分の変形を逆にたどると
\[
\sum \mathfrak{R}_{\mu\nu}p^{\mu\nu}
+2\Div\!\left(\sum g^{\mu\sigma}\mathfrak{R}_{\mu\tau}p^\tau\right)=0,
\]
\[
\delta\mathfrak{R}^*+\Div\!\left(\sum 2g^{\mu\sigma}\mathfrak{R}_{\mu\tau}p^\tau
-\frac{\p\mathfrak{R}^*}{\p g_{\mu\nu}^{\sigma}}p^{\mu\nu}\right)=0.
\]
Lie の微分方程式
\(\delta\mathfrak{R}^*+\Div(\mathfrak{R}^*\D w)=0\) と比較すると、\(I^*\) が許す infinitesimal 変換として
\[
\D w^\sigma=\frac1{\mathfrak{R}^*}\left(\sum 2g^{\mu\sigma}\mathfrak{R}_{\mu\tau}p^\tau
-\frac{\p\mathfrak{R}^*}{\p g_{\mu\nu}^{\sigma}}p^{\mu\nu}\right),
\qquad
\D g^{\mu\nu}=p^{\mu\nu}+\sum g^{\mu\nu}_{\sigma}\D w^\sigma,
\]
を得る。したがってこれらの infinitesimal 変換は \(g^{\mu\nu}\) の一階および二階導関数に依存し、任意の \(p\) を一階導関数まで含む。}

\section*{\S{}5. 関係の各成分の不変性}

群 \(\G\) を最も単純で通常考察される場合、すなわち変換に \(u\) の導関数を含まず、変換された独立変数が \(u\) ではなく \(x\) だけに依存する場合に特殊化すると、公式中の個々の成分の不変性を推論できる。まずよく知られた議論により、
\[
\int\!\cdots\!\int\left(\sum_i\psi_i\delta u_i\right)\dd x
\]
の不変性が得られ、したがって \(\sum_i\psi_i\delta u_i\) の相対不変性が得られる。ここで \(\delta\) は任意の変分を表す。すなわち一方では
\[
\delta I
=\int\!\cdots\!\int \delta f\!\left(x,u,\frac{\p u}{\p x},\ldots\right)\dd x
=\int\!\cdots\!\int \delta f\!\left(y,v,\frac{\p v}{\p y},\ldots\right)\dd y.
\]
他方、\(\delta u,\delta(\p u/\p x),\ldots\) が境界で消えるとき、\(\delta u,\delta(\p u/\p x),\ldots\) の線形同次変換により、それに対応する \(\delta v,\delta(\p v/\p y),\ldots\) も境界で消えるので、
\[
\int\!\cdots\!\int \delta f\!\left(x,u,\frac{\p u}{\p x},\ldots\right)\dd x
=
\int\!\cdots\!\int\left(\sum_i\psi_i(u,\ldots)\delta u_i\right)\dd x,
\]
\[
\int\!\cdots\!\int \delta f\!\left(y,v,\frac{\p v}{\p y},\ldots\right)\dd y
=
\int\!\cdots\!\int\left(\sum_i\psi_i(v,\ldots)\delta v_i\right)\dd y.
\]
したがって境界で消える \(\delta u,\delta(\p u/\p x),\ldots\) について
\[
\int\!\cdots\!\int\left(\sum_i\psi_i(u,\ldots)\delta u_i\right)\dd x
=
\int\!\cdots\!\int\left(\sum_i\psi_i(v,\ldots)\delta v_i\right)\dd y
\]
\[
=
\int\!\cdots\!\int\left(\sum_i\psi_i(v,\ldots)\delta v_i\right)
\left|\frac{\p y_i}{\p x_k}\right|\dd x.
\]
第三の積分で \(y,v,\delta v\) を \(x,u,\delta u\) で表し、それを第一の積分に等置すれば、境界で消え、その他では任意の \(\delta u\) について関係
\[
\int\!\cdots\!\int\left(\sum_i\chi_i(u,\ldots)\delta u_i\right)\dd x=0
\]
を得る。ここから、よく知られているように、被積分関数は任意の \(\delta u\) に対して消える。したがって関係
\[
\sum_i\psi_i(u,\ldots)\delta u_i
=
\left|\frac{\p y_i}{\p x_k}\right|
\left(\sum_i\psi_i(v,\ldots)\delta v_i\right)
\]
は \(\delta u\) に関して恒等的に成り立つ。これは \(\sum_i\psi_i\delta u_i\) の相対不変性を表し、したがって
\[
\int\!\cdots\!\int\left(\sum_i\psi_i\delta u_i\right)\dd x
\]
の不変性を表す。\footnote{これらの議論は、\(y\) も \(u\) に依存する場合には失敗する。そのとき \(\delta f(y,v,\p v/\p y,\ldots)\) には \(\sum(\p f/\p y)\delta y\) という項も含まれ、発散変換はラグランジュ式へ導かないからである。\(u\) の導関数を許す場合にも同様に失敗する。このとき \(\delta v\) は \(\delta u,\delta(\p u/\p x),\ldots\) の線形結合となり、さらに発散変換を行った後に初めて \(\int\!\cdots\!\int(\sum\chi_i(u,\ldots)\delta u_i)\dd x=0\) という恒等式を導くので、右辺にはやはりラグランジュ式が現れない。\(\int\!\cdots\!\int(\sum\psi_i\delta u_i)\dd x\) の不変性からすでに発散関係の存在を結論できるかという問題は、逆定理により、同じ \(\D u,\D x\) に至る必要はないが同じ \(\bd u\) には至る群に関する \(I\) の不変性をそれが含意するかという問題と同値である。単積分で \(f\) に一階導関数だけが現れる特殊な場合には、有限群について、ラグランジュ式の不変性から第一積分の存在を推論できる（例えば Engel, Gött. Nachr. 1916, p. 270 を参照）。}

これを導かれた発散関係と従属関係に適用するには、まず \(\D u,\D x\) から導かれる \(\bd u\) が、実際に変分 \(\delta u\) の変換法則を満たすことを示さなければならない。ただし、\(\bd v\) におけるパラメータ、または任意関数は、\(y,v\) における同様の infinitesimal 変換群に対応するように決定されているものとする。\(T_q\) を \(x,u\) を \(y,v\) に移す変換とし、\(T_p\) を \(x,u\) における infinitesimal 変換とする。このとき \(y,v\) における同様のものは \(T_r=T_qT_pT_q^{-1}\) で与えられ、ここでパラメータまたは任意関数 \(r\) は \(p\) と \(q\) によって決まる。式で書けば次のようになる：
\[
T_p:\ \xi=x+\D x(x,p),\qquad u^*=u+\D u(x,u,p);
\]
\[
T_q:\ y=A(x,q),\qquad v=B(x,u,q);
\]
\[
T_qT_p:\ \eta=A(x+\D x(x,p),q),\quad
v^*=B(x+\D x(p),\,u+\D u(p),q).
\]
これから \(T_r=T_qT_pT_q^{-1}\) が生じるので、
\[
\eta=y+\D y(r),\qquad v^*=v+\D v(r),
\]
である。ただし \(T_q\) の逆により \(x\) を \(y\) の関数とみなし、infinitesimal な項だけを残す。したがって恒等式
\begin{equation}
\eta=y+\D y(r)=y+\sum\frac{\p A(x,q)}{\p x}\D x(p)
\tag{20a}
\end{equation}
および
\begin{equation}
v^*=v+\D v(r)=v+\sum\frac{\p B(x,u,q)}{\p x}\D x(p)
+\sum\frac{\p B(x,u,q)}{\p u}\D u(p)
\tag{20b}
\end{equation}
がある。ここで \(\xi=x+\D x\) を \(\xi-\D \xi\) に置き換え、\(\xi\) が \(x\) に戻り、したがって \(\D x\) が消えるようにすると、第一式 (20) により \(\eta\) も \(y=\eta-\D\eta\) に戻る。この置換の下で \(\D u(p)\) が \(\bd u(p)\) へ移るなら、\(\D v(r)\) も \(\bd v(r)\) へ移り、第二式 (20) は
\[
v+\bd v(y,v,\ldots,r)=v+\sum\frac{\p B(x,u,q)}{\p u}\,\bd u(p),
\]
\[
\bd v(y,v,\ldots,r)=\sum\frac{\p B}{\p u_\kappa}\,\bd u_\kappa(x,u,p)
\]
を与える。したがって \(\bd v\) がパラメータまたは任意関数 \(r\) に依存するものとして取られているかぎり、変分に関する変換公式は実際に満たされる。\footnote{結論が成り立つためには \(y\) が \(u\) などに依存しないと仮定しなければならないことが、ここでも見られる。例として、Klein が示した \(\delta g^{\mu\nu}\) と \(\delta q_\sigma\) を挙げられる。これらは \(p\) をベクトル変換に従わせると、変分に関する変換を満たす。}

とくに \(\sum_i\psi_i\bd u_i\) は相対不変であることが従う。したがって (12) により、発散関係は \(y,v\) でも満たされているので、\(\Div B\) の相対不変性も従う。さらに (14) と (13) により、\(\Div\Gamma\) と、\(p^{(\lambda)}\) と組み合わされた従属関係の左辺の相対不変性も従う。変換された公式では、任意の \(p(x)\)（またはパラメータ）は至る所で \(r\) に置き換えられる。ここから、\(\Div(B-\Gamma)\) の相対不変性も従う。すなわち、恒等的には消えないが、その発散は恒等的に消える関数系 \(B-\Gamma\) の発散の相対不変性である。

一次元の場合および有限群の場合には、\(\Div B\) の相対不変性から、第一積分の不変性についてさらに結論を引き出すことができる。infinitesimal 変換に対応するパラメータの変換は、(20) により線形同次となる。またすべての変換は可逆であるから、\(\eps\) も変換されたパラメータ \(\eps^*\) に関して線形同次である。\(\psi=0\) とおいてもこの可逆性は確かに保たれる。なぜなら (20) には \(u\) の導関数が現れないからである。
\[
\Div B(x,u,\ldots,\eps)=\frac{\dd y}{\dd x}\,\Div B(y,v,\ldots,\eps^*)
\]
で \(\eps^*\) の係数を比較すると、量 \(\frac{\dd}{\dd y}B^{(\lambda)}(y,v,\ldots)\) も \(\frac{\dd}{\dd x}B^{(\lambda)}(x,u,\ldots)\) の線形同次関数となる。したがって
\[
\frac{\dd}{\dd x}B^{(\lambda)}(x,u,\ldots)=0
\quad\text{または}\quad B^{(\lambda)}(x,u)=\mathrm{const.}
\]
から
\[
\frac{\dd}{\dd y}B^{(\lambda)}(y,v,\ldots)=0
\quad\text{または}\quad B^{(\lambda)}(y,v)=\mathrm{const.}
\]
が従う。したがって \(\G_\rho\) に対応する \(\rho\) 個の第一積分は、それぞれ群を許すので、さらに先の積分も単純化される。この最も単純な例は、\(f\) が \(x\) を含まない、またはある \(u\) を含まない場合であり、それぞれ infinitesimal 変換 \(\D x=\eps,\D u=0\)、\(\D x=0,\D u=\eps\) に対応する。そのとき \(\bd u=-\eps\,\dd u/\dd x\)、または \(\eps\) である。そして \(B\) は \(f\) と \(\bd u\) から微分と有理演算によって導かれるので、\(B\) も \(x\)、または \(u\) を含まず、対応する群を許す。\footnote{第一積分の存在がすでに \(\int(\sum\psi_i\delta u_i)\dd x\) の不変性から従う場合、それらは完全な群 \(\G_\rho\) を許さない。例えば \(\int(u''\delta u)\dd x\) は infinitesimal 変換 \(\D x=\eps_2,\D u=\eps_1+x\eps_3\) を許す。一方、\(\D x=0,\D u=x\eps_3\) に対応する第一積分 \(u-u'x=\mathrm{const.}\) は、\(u\) と \(x\) を明示的に含むので、他の二つの infinitesimal 変換を許さない。この第一積分に対応して、\(f\) に対する導関数を含む infinitesimal 変換がある。こうして、\(\int\!\cdots\!\int(\sum\psi_i\delta u_i)\dd x\) の不変性は、いずれにせよ \(I\) の不変性より少ないことしか達成しないことがわかる。これは前の注で提起された問題に関係する。}

\section*{\S{}6. Hilbert の主張}

以上から最後に、真性エネルギー定理の失敗と「一般相対性」との関連に関する Hilbert の主張（Klein の第一ノート、Göttinger Nachr. 1917、回答の第一段落）の証明が、一般化された群論的定式化で得られる。

積分 \(I\) が \(\G_{\infty\rho}\) を許すとし、\(\G_\sigma\) を、任意関数の特殊化によって生じる任意の有限群、したがって \(\G_{\infty\rho}\) の部分群とする。無限群 \(\G_{\infty\rho}\) には従属関係 (16) が対応し、有限群 \(\G_\sigma\) には発散関係 (13) が対応する。逆に、任意の発散関係の存在から、ある有限群に関する \(I\) の不変性が従う。この有限群が \(\G_\sigma\) と同一であるのは、\(\bd u\) が \(\G_\sigma\) から生じるものの線形結合である場合、かつその場合に限る。したがって \(\G_\sigma\) に関する不変性は、(13) と異なる発散関係へは導きえない。しかし (16) の存在から、任意の \(p(x)\) について、\(\G_{\infty\rho}\) の infinitesimal 変換 \(\D u,\D x\) に関する \(I\) の不変性が従うので、とくに特殊化によってこれらから生じる \(\G_\sigma\) の infinitesimal 変換に関して、したがって \(\G_\sigma\) に関して、\(I\) はすでに不変である。したがって発散関係
\[
\sum_i\psi_i\bd u_i^{(\lambda)}=\Div B^{(\lambda)}
\]
は、(16) の従属関係の帰結でなければならない。(16) はまた
\[
\sum_i\psi_i a_i^{(\lambda)}=\Div\chi^{(\lambda)}
\]
とも書ける。ここで \(\chi^{(\lambda)}\) はラグランジュ式とその導関数の線形結合である。\(\psi\) は (13) と (16) の双方に線形に現れるので、発散関係は特に従属関係 (16) の線形結合でなければならない。したがって
\[
\Div B^{(\lambda)}=\Div\left(\sum \alpha_\mu\chi^{(\mu)}\right),
\]
かつ \(B^{(\lambda)}\) 自身は、\(\chi\)、すなわちラグランジュ式とその導関数から、および発散が恒等的に消える関数から線形に構成されている。例として \S{}2 の末尾に現れる \(B-\Gamma\) があり、ここでは \(\Div(B-\Gamma)=0\) で、同時にその発散は不変量の性格をもつ。\(B^{(\lambda)}\) が示された仕方でラグランジュ式とその導関数から構成されるような発散関係を「非真性」と呼び、残りのすべてを「真性」と呼ぶことにする。

逆に、発散関係が従属関係 (16) の線形結合であり、したがって「非真性」であるなら、\(\G_\sigma\) に関する不変性は \(\G_{\infty\rho}\) に関する不変性から従う。\(\G_\sigma\) は \(\G_{\infty\rho}\) の部分群となる。したがって、有限群 \(\G_\sigma\) に対応する発散関係が非真性になるのは、\(\G_\sigma\) が、\(I\) が不変であるような無限群の部分群である場合、かつその場合に限る。

群を特殊化すると、Hilbert のもとの主張がこれから従う。「平行移動群」とは有限群
\[
y_i=x_i+\eps_i,\qquad v_i(y)=u_i(x)
\]
をいう。したがって
\[
\D x_i=\eps_i,\qquad \D u_i=0,\qquad
\bd u_i=-\sum_\lambda\frac{\p u_i}{\p x_\lambda}\eps_\lambda.
\]
平行移動群に関する不変性は、知られているように、
\[
I=\int\!\cdots\!\int f\!\left(x,u,\frac{\p u}{\p x},\ldots\right)\dd x
\]
において \(f\) に \(x\) が明示的に現れないことを意味する。対応する \(n\) 個の発散関係
\[
\sum_i\psi_i\frac{\p u_i}{\p x_\lambda}=\Div B^{(\lambda)}
\qquad(\lambda=1,2,\ldots,n)
\]
を「エネルギー関係」と呼ぶことにする。なぜなら、変分問題に対応する「保存則」\(\Div B^{(\lambda)}=0\) が「エネルギー定理」に、\(B^{(\lambda)}\) が「エネルギー成分」に対応するからである。したがって次を得る。\(I\) が平行移動群を許すなら、エネルギー関係が非真性となるのは、\(I\) が、平行移動群を部分群として含む無限群に関して不変である場合、かつその場合に限る。\footnote{古典力学のエネルギー定理、および古い「相対性理論」のエネルギー定理（そこでは \(\sum \dd x^2\) がそれ自身へ移る）は「真性」である。ここでは無限群が現れないからである。}

このような無限群の例は、すべての \(x\) の変換と、\(u(x)\) に誘導される変換からなる群で与えられる。この群では任意関数 \(p(x)\) の導関数だけが現れる。平行移動群は特殊化 \(p^{(i)}(x)=\eps_i\) によって生じる。ただし、この手段と、\(I\) を境界積分で変えることによって生じる群とにより、すでにこのような群の最も一般のものが与えられているかどうかは、未決にしておかなければならない。示された種類の誘導変換は、例えば \(u\) を「全微分形式」の係数変換に従わせることによって生じる。すなわち、\(\dd x\) のほか高次微分も含む形式
\[
\sum a\,\dd^2x_i+\sum b\,\dd x_i\dd x_k+\cdots
\]
である。より特殊な誘導変換、すなわち \(p(x)\) が一階導関数にだけ現れるものは、通常の微分形式 \(\sum c\,\dd x_{i_1}\cdots \dd x_{i_\lambda}\) の係数変換によって与えられる。これらが通常考えられてきたものである。

さらに示された種類の群として、対数項が現れるため係数変換ではありえないものが、例えば次で与えられる：
\[
y=x+p(x),\qquad
v_i=u_i+\lg(1+p'(x))=u_i+\lg\frac{\dd y}{\dd x};
\]
\[
\D x=p(x),\qquad
\D u_i=p'(x),\qquad
\bd u_i=p'(x)-u_i'p(x).
\]
ここで従属関係 (16) は
\[
\sum_i\left(\psi_i u_i'+\frac{\dd\psi_i}{\dd x}\right)=0
\]
となり、非真性エネルギー関係は
\[
\sum_i\left(\psi_i u_i'+
\frac{\dd(\psi_i+\mathrm{const.})}{\dd x}\right)=0
\]
となる。この群の最も簡単な不変積分は
\[
I=\int \frac{e^{-2u_1}}{u_1'-u_2'}\,\dd x.
\]
最も一般の \(I\) は Lie の微分方程式 (11)
\[
\bd f+\frac{\dd}{\dd x}(f\cdot \D x)=0
\]
を積分することによって決まる。\(\D x\) と \(\bd u\) の値を代入し、\(f\) が \(u\) の一階導関数だけに依存すると仮定すると、これは
\[
\frac{\p f}{\p x}p(x)
+\left\{\sum_i\frac{\p f}{\p u_i}
-\sum_i\frac{\p f}{\p u_i'}u_i'+f\right\}p'(x)
+\left\{\sum_i\frac{\p f}{\p u_i'}\right\}p''(x)=0
\]
（\(p(x),p'(x),p''(x)\) に関する恒等式）となる。すでに二つの関数 \(u(x)\) に対して、この方程式系は導関数を実際に含む解をもつ。すなわち
\[
f=(u_1'-u_2')\,
\Phi\!\left(u_1-u_2,\frac{e^{-u_1}}{u_1'-u_2'}\right),
\]
ここで \(\Phi\) は示された引数の任意関数である。

Hilbert は、真性エネルギー定理の失敗が「一般相対性理論」の特徴である、とその主張を述べている。この主張が文字どおり成り立つためには、「一般相対性」という名称を通常より広く理解し、上に述べた \(n\) 個の任意関数に依存する群にも拡張しなければならない。\footnote{これにより、物理学で普通に使われる「相対性」という語は「群に関する不変性」に置き換えられるべきだ、という Klein の注意がふたたび正しいことが確認される。（Über die geometrischen Grundlagen der Lorentzgruppe. \emph{Jhrber. d. d. Math. Vereinig.} vol. 19, p. 287, 1910; \emph{Physikalische Zeitschrift} に再録。）}

% --- Paper 14 local macro support ---
\providecommand{\p}{\partial}
\providecommand{\frS}{\mathfrak S}
\providecommand{\frp}{\mathfrak p}
\providecommand{\fra}{\mathfrak a}
\providecommand{\frb}{\mathfrak b}
\providecommand{\frc}{\mathfrak c}
\providecommand{\frf}{\mathfrak f}
\providecommand{\frr}{\mathfrak r}
\providecommand{\frm}{\mathfrak m}
\providecommand{\frD}{\mathfrak D}
\providecommand{\calA}{\mathcal A}
\providecommand{\calB}{\mathcal B}
\providecommand{\calN}{\mathcal N}
% --- end Paper 14 local macro support ---

\section*{14. 一変数の代数関数の算術的理論と、他の諸理論および数体論との関係}
\begin{center}
\emph{Jahresbericht der Deutschen Mathematiker-Vereinigung} 38 (1919), pp. 182--203
\end{center}

以下の報告は、代数関数の個々の理論のあいだに連結環を得たいという願いから生まれた。同時にそれは、算術的理論の議論が本質的に欠けている Brill--Noether の報告\footnote{古い時代および新しい時代における代数関数論の発展。\emph{Jahresbericht der Deutschen Mathematiker-Vereinigung} vol. 3 (1894)［Brill--Noether として引用する］。}への補足と見ることもできる。\footnote{Kronecker の判別式に関する研究の議論を例外とする。（\emph{J. f. M.} vol. 91.）}

個々の理論の説明としては、Brill--Noether のほか、そこに掲げられた文献を含めて、次の \emph{Enzyklopädie} の論説が考慮される。

Riemann の理論については Wirtinger (II B 2), nos. 1--19；

Weierstrass の理論については Wirtinger (II B 2), nos. 23, 32, 33, 34；

Brill--Noether の代数幾何学的理論については Wirtinger (II B 2), nos. 21--31；Berzolari (III C 4), nos. 23--38；

Dedekind--Weber および Hensel--Landsberg の算術的理論については Hensel (II C 5)［Hensel として引用する］；\footnote{Dedekind--Weber に関する部分は A. Ostrowski による。}

より一般に、代数量の算術的理論については Landsberg (I C 5)、多変数代数関数の算術的理論については Jung (II C 6)。

数体論については Hilbert の「Zahlbericht」\footnote{代数的数体の理論。\emph{Jahresbericht der Deutschen Mathematiker-Vereinigung} vol. 4 (1894/95).}および Dirichlet--Dedekind の数論講義［Dedekind--数論として引用する］\footnote{Braunschweig, Vieweg, 4th ed. 1894.}を参照されたい。

\subsection*{目次}
\begin{enumerate}
\item さまざまな理論の基礎。
\item 代数的数と代数関数とのあいだの平行性（共役元、基底定理、加群概念）。
\item 代数的数と代数関数のイデアル論における平行性と相違。
\item 級数展開の算術的基礎づけとイデアル論との関係。
\item 代数関数の算術的理論と幾何学的理論との比較。異なる不変性概念。
\item 比較の続き。特異点。
\item 剰余定理とイデアル論の諸定理との比較。
\item 代数関数と代数的数における超越的問題。
\end{enumerate}

\subsection*{1. さまざまな理論の基礎}
Riemann の理論は、存在定理（Dirichlet の原理、Schwarz--Neumann の方法）に依拠する、唯一の純粋に超越的な理論である。ここでは代数関数の積分が第一のものであり、代数関数そのものは第二のものである。Riemann 面上の関数は、その零点と極によって特徴づけられる。これは算術的理論における多角形（または除数）に対応する。一方 Weierstrass は、関数そのものを、その零点および極とともに語る。存在定理を別にすれば、主要問題は、与えられた極をもつすべての関数を設定することであり、これは第一種および第二種の積分によって遂行される。さらに、無限遠点を指定したときになお入ってくる任意定数の数が問題になる。この後者の問いには Riemann--Roch の定理が答える。算術的理論ではこれに多角形類（または除数類）の次元の問いが対応し、代数幾何学的理論では完全線形系の次元の問いが対応するが、ここでも答えは Riemann--Roch の定理によって与えられる。関数の実際の構成は、算術的理論では基底定理によって、代数幾何学的理論では剰余定理によって行われる。とくに、第一種の被積分関数の零点には、算術的理論では微分類が対応し、代数幾何学的理論では特殊群の理論、すなわち \((n-3)\) 次の随伴曲線（\(\varphi\)-曲線）によって切り出される点群の束が対応する。

Riemann の理論が歴史的には他の理論に先行したが、存在定理の厳密な証明は後に初めて得られた。それと同様に、後にもその超越的な補助手段によって、今日に至るまで純粋に代数的または算術的な処理ではなお到達できない代数的問題が処理された。ここではとりわけ Hurwitz の特異対応の理論を挙げなければならない。\footnote{A. Hurwitz, Über algebraische Korrespondenzen und das erweiterte Korrespondenzprinzip. \emph{Math. Ann.} vol. 28.（Brill--Noether 第10節を参照。）}

代数的数体の理論における超越的な問題設定との類似は、最後の節で論じる。

残りの理論では、代数関数が第一のものであり、積分は第二のものである。この点で Weierstrass の理論は、冪級数展開と剰余定理という超越的基礎の上に立つ。しかしその後は、考慮すべきすべての関数を明示的に与えることによって、代数的に進む。

Hensel--Landsberg の形の算術的理論も、級数展開という超越的基礎を用いる。しかしこの基礎の上で個々の点に除数を対応させることにより、その後は純粋に算術的に進むので、Weierstrass と Dedekind--Weber との融合と見ることができる。後者と対照的に、この理論は、問題となるすべての基底関数を実際に構成する方法を同時に与える。細部では Dedekind--Weber に比べていくつかの簡略化があり、その本質は「分数的」除数の導入にある。最近 Hensel\footnote{\emph{Mathematische Zeitschrift} vol. 4。この基礎づけは彼の報告にも基礎として用いられている。} は、級数展開の理論を算術的に基礎づけた（収束証明のための優関数法を除く）。これによって理論はほとんど純粋に算術的なものになった。

Brill--Noether の代数幾何学的理論は、枝の「相続く点」という概念だけを前提にする。この概念は通常点における級数展開に帰着され\footnote{「特異」な相続く点については Brill--Noether 第6節、とくに no. 10, 11 を参照。通常の相続く点とは、例えば接線と曲線との交点、またはより高次の接触点（分岐点）である。}、Weierstrass の元概念に対応する。その後は代数的・有理的な方法、本質的には三元消去論の方法によって進み、これもまた明示的表示にまで到達する。

Dedekind--Weber の本来の純粋に算術的な理論\footnote{\emph{J. f. M.} 92（D.--W. として引用する）。}は、その諸定理がふたたび存在証明の性格をもつという点で Riemann の理論に近い。この理論は方法の完全な純粋性を保つ。\footnote{この理由から、例えば Hensel--Landsberg の理論よりも、他の理論との比較に適している。後者は以下の比較では時折現れるだけである。--- Abel 積分とその反転の理論は、連続性を前提とする限り、D.--W. にも代数的理論のさらなる展開（M. Noether, \emph{Ann.} 37）にも扱われていない。} 変数が不定元の役割だけを果たすこの理論の算術的基礎づけは、2 で詳しく説明するように、代数的数の理論と平行して進む。\footnote{Hensel--Landsberg の理論と数論、とくに \(p\)-進数の理論との比較は Hensel に見られる。} この基礎づけは、点概念を算術的に導入する手段を与え（Hensel no. 10a を参照）、したがって連続性に関するいかなる考察も用いずに絶対 Riemann 面の概念を設定する手段を与える。理論はそれによって形式的な性格を得る。例えば微分の概念も純粋に形式的に導入される。点概念によって得られた基礎の上で、代数幾何学的理論の方法と概念との近い親族性を示すことができる。というのも、後者も元概念を除けば超越的補助手段から自由だからである。これは次の節で行う。

\subsection*{2. 代数的数と代数関数との平行性。（共役元、基底定理、加群概念）}
代数的数体は、有理数体 \(R\) 上の \(n\) 次の体として定義される。一変数の代数関数体は、任意の複素係数をもつ一つの不定元 \(z\) のすべての有理関数からなる体 \(K(z)\) 上の \(n\) 次の体として定義される。したがって両方の体について、基礎体の特殊な性質を完全に捨象するすべての定理が共通に成り立つ。すなわち、基底、ノルム、トレース、判別式に関する定理である。\footnote{D.--W., § 1 および 2；Hensel no. 4（注 5 付き）。} Dedekind--Weber はこれらすべての概念を、共役元を用いずに導入する。しかし、Steinitz が Kronecker に従って示したように、共役体の概念は完全に形式的な仕方で定義できることを注意しておくべきである。\footnote{Steinitz, Algebraische Theorie der Körper.（\emph{J. f. M.} 137, § 6 および 8.）} すなわち、\(K\) を任意の体、\(\varphi(x)\) を \(K\) 上既約な \(x\) の多項式とすると、\(\varphi(x)\) を法とする剰余類の系は体をなす。\(x\) の剰余類に元 \(\xi\) を対応させると、同型対応のもとで、多項式 \(f(x)\) によって表される剰余類に元 \(f(\xi)\) が対応する。\(\varphi(x)\) で割り切れるすべての多項式は零類を表すので、\(\varphi(\xi)\) は零類に対応する。したがって元 \(\xi\) は記号的に \(\varphi(x)=0\) の零点と見なすことができ、代数体 \(K(\xi)\) を生成する。この手続きを反復すると、\(n\) 個の共役体 \(K(\xi),K(\xi'),\ldots,K(\xi^{(n-1)})\) に至る。したがって、とくに代数関数の場合にも、純粋に算術的な理論を離れることなく共役元について語ることができる。

二つの基礎体 \(R\) と \(K(z)\) は、その中ですべての整な量の系が定義されているという性質を共有する。それはそれぞれ、有理整数および一つの不定元の整有理関数である。これらの整な量は、任意の二つ \(a\) と \(b\) が最大公約元 \(d\) をもち、それが \(ma+nb\) の形に表されるという特徴的性質をもつ。ここで \(d,m,n\) もまた \(R\) または \(K(z)\) の整な量である。しかも \(d\) は、この形に表される絶対値最小の数、または次数最小の多項式である。有理整数および一変数関数の素因数への一意分解を導くこの性質だけに、代数的体または上体の整な量に関する諸定理は基づいている。最大公約元の存在を示す考察は、代数体の整な量の基本系または基底の存在も与える。整な量からなる体の基底だけを考えると、その判別式は有理整数または一つの不定元の整有理関数になる。絶対値最小の数、または次数最小の関数に対応する各基底は、基本系をなす。

基底の問題について、より一般の体も含めてより精密な見通しを与えるのが加群概念の導入である。この概念は、文献で元の性質に応じて異なる形で現れる定義を含むように定式化する。基本整性領域 \(\frS\) に関する加群とは、二つの元の差がふたたびその系に属し、また一つの元と \(\frS\) の任意の量との積もその系に属するような元の系である。（\(\frS\) が有理整数からなる場合、第二の条件は第一の条件から従う。）\footnote{加群概念は Dedekind に由来する。彼は数論の第二版で初めて数加群（\(\frS=\) 有理整数）を導入した。整係数線形形式の加群もそこに暗黙に含まれている。Dedekind--Weber には「関数加群」が現れる。その元は代数関数であり、\(\frS\) は一つの不定元のすべての整有理関数からなる。--- \(n\) 個の不定元の斉次形式、より一般に多項式の加群は、Hilbert（\emph{Annalen} 36）に初めて現れる。\(\frS\) はすべての多項式からなり、元も多項式であるから、この加群はわれわれのイデアル定義に入る。Kronecker の多項式の「加群系」は、抽象的定義からではなく基底から出発する。}

有限個の元からなる基底をもち、その基底によって各元が \(\frS\) の係数を用いて線形に表される加群を有限（有限生成）という。とくに一項加群は \(c\cdot\alpha\) の形である。ここで \(\alpha\) は基底元を表し、\(c\) は \(\frS\) のすべての量を走る。\emph{元そのものが \(\frS\) に属する加群を \(\frS\) におけるイデアルという。} ここから有限イデアルの概念、とくに一項イデアル（主イデアル）の概念が従う。

整な量とイデアルについてのすべての基底定理は、次の定理に基づいている。

\emph{\(M\) を \(\frS\) に関する加群とし、その元が \(\frS\) の係数をもつ \(n\) 個の不定元の線形形式からなるとする。もし \(\frS\) のすべてのイデアルが有限ならば、\(M\) も有限加群である。とくに \(\frS\) のすべてのイデアルが主イデアルならば、\(M\) は線形独立な形式からなる基底をもつ。}\footnote{この定理は、異なる係数領域 \(\frS\) について文献に別々に現れる。有理整数については、例えば Hilbert, Zahlbericht § 3 および 4 を参照。複数の不定元の多項式については、Hilbert, Über die vollen Invariantensysteme, § 2 の末尾（\emph{Math. Ann.} 42）を参照。}

実際、\(M\) の元が \(l(x)=a_1x_1+\cdots+a_nx_n\) で与えられているなら、係数 \(a_1\) 全体は \(\frS\) の一つのイデアルを走る。仮定によりそれは \(b_1a^{(1)}+\cdots+b_\rho a^{(\rho)}\) の形である。したがって \(M\) に属する線形形式
\[
m(x)=l(x)-b_1l^{(1)}(x)-\cdots-b_\rho l^{(\rho)}(x)
\]
は \((n-1)\) 個の不定元だけに依存する。ここから有限回反復することで主張が従う。各場合に \(\rho=1\) ならば、基底はそれぞれ \(n,n-1,\ldots,1\) 個の不定元の \(n\) 個の形式からなり、これによって零でない形式の線形独立性が得られる。

基本系についての定理は、これからさまざまな場合に次のように従う。

代数体が、整な代数元 \(\delta\) を基礎体に添加することによって得られ、\(\frS\) が基礎体のすべての整な量を表すとする。このとき上体の各整な量は次の表示をもつ（例えば Zahlbericht § 3 を参照）：
\[
\omega=\frac{a_1\delta^{\,n-1}+a_2\delta^{\,n-2}+\cdots+a_n}{\Delta(\delta)},
\]
ここで \(a_i\) は \(\frS\) の量であり、\(\Delta(\delta)\) は \(\delta\) の判別式を表す。置換
\[
x_i=\frac{\delta^{\,n-i}}{\Delta(\delta)}
\]
によって、すべての整な量の加群は線形形式の加群に移される。したがって、\(\frS\) のすべてのイデアルが有限であるすべての場合に、基本系の存在が従う。同じ考察は、より一般に、上体の整な量の領域におけるすべてのイデアルが有限であることを示す。さて、とくに有理整数および一変数関数については、すべてのイデアルが主イデアルである。したがって、代数的数体および一変数代数関数体における基本系は、\(n\) が次数を表すとき、ちょうど \(n\) 個の量からなることが従う。また、これらの体では上の議論によりすべてのイデアルが有限であるから、同時に相対体の基本系も得られる。それは一般に相対次数が示す数よりも多くの量からなる。\footnote{線形形式の加群に関する Steinitz の研究によれば、\(\frS\) が（有限）代数的数体の整数からなるとき、\(\nu\) を相対次数とすれば \(\nu+1\) 個の基底元で足りる（\emph{Math. Ann.} 71, § 4）。} さらに、\(\frS\) が複数の不定元のすべての多項式の領域を表すならば、\(\frS\) の各イデアルも有限である。\footnote{Hilbert の加群基底定理（\emph{Math. Ann.} 36）による。ここで一貫性のためイデアルと呼んだ多項式の加群は、通常「形式加群」または単に「加群」と呼ばれる。こうした加群 \(N\) が有限であるという Hilbert の定理も、線形形式に関する定理に帰着できる。すなわち、\(x_1,\ldots,x_r,x_{r+1}\) の \(r+1\) 個の不定元において、\(N\) に含まれ \(x_{r+1}^{\,n}\) の項を含む次数 \(n\) の多項式 \(f\) をとる（これは線形変換によって常に達成できる）。すると \(N\) のすべての多項式は \(f\) を法として \(a_1(x)x_{r+1}^{\,n-1}+\cdots+a_n(x)\) の形に表されるものに合同である。したがって \(x_{r+1}^{\,n-i}=X_i\) とおけば、それらは線形形式の加群をなす。この加群については、\(r\) 個の不定元の領域のすべてのイデアルが有限であると仮定してよい。これは一つの不定元について成り立つからである。} したがって、多変数の代数関数体についても基本系の存在が得られ、それはまた次数が示す数よりも多くの量からなる。

\subsection*{3. 代数的数と代数関数のイデアル論における平行性と相違}
すべてのイデアルが素イデアルの冪の積として一意的に表されるというイデアル論の基本定理の証明は、代数的数と代数関数について共通に実行できる。それは、有理整数および一変数関数についてすべてのイデアルが主イデアルであるという性質だけに基づく。\footnote{Hurwitz: Über die Theorie der Ideale（\emph{Gött. Nachr.} 1894）の序論の注を参照。Weber の代数学教科書では、Kronecker に従って、汎関数の導入により両方の体について平行に証明が行われている。［Kronecker の理論とイデアル論との関係は「拡張された Gauss の定理」によって媒介される；Hurwitz, loc. cit.］} 証明の主要点は、イデアル \(\frc\) がイデアル \(\fra\) で割り切れること（すなわち \(\frc\) の各元が \(\fra\) に含まれること）から、積表示 \(\frc=\fra\frb\) が従うことを示す点にある。\footnote{Dedekind（数論）でとくに強調されている。補遺 XI, 定理 VII, § 177；p. 554 の注を参照。この定理は代数体を扱っているという事実を用いる。多項式のイデアル（加群）についてはもはや成り立たず、そこでは積の代わりに最小公倍元が現れる。} Hurwitz ではこの証明は「拡張された Gauss の定理」に基づく。これは、二つの多項式 \(\varphi\) と \(\psi\) の積の各係数を割る整な量 \(\omega\) が、\(\varphi\) と \(\psi\) のすべての係数の積も割るというものである。Dedekind では、この定理と同値な加群定理に基づく。\footnote{数論、補遺 XI；定理 VI, § 173。Dedekind は二つの定理の同値性を指摘している（Über die Begründung der Idealtheorie, \emph{Göttinger Nachr.} 1895）。}

さらに進むと、補完加群の概念（Hensel no. 12）、およびそこから分岐イデアル（基礎イデアル、異なるもの）\(\frD\) を共通に定義できる（そのノルムが体判別式 \(D\) となる）。そして、分岐イデアル \(\frD\) は \emph{すべての主イデアル \(F'(\delta)\) の最大公約元} であるという基本定理が成り立つ。ここで \(\delta\) は体のすべての整な量を走り、\(F(\delta)=0\) はその都度対応する既約方程式を表す。これから、\emph{体判別式は、素イデアルの平方で割り切れる有理素数、または \(z\) の一次因子をすべて、かつそれだけを含む} ことが従う。\(F'(\delta)\) 自身については次の表示を得る：
\[
F'(\delta)=\frf\cdot\frD,
\]
さらにノルムをとると
\[
\Delta(\delta)=R^2\cdot D;
\]
ここで \(\frf\) は \(\delta\) から導かれる環（整域、整性領域）の導手を表す。すなわち \(\frf\) は、任意の整な量を掛けた後に加群 \([1,\delta,\ldots,\delta^{n-1}]\) で割り切れる（整な）量全体からなる。\(R^2=\operatorname{Norm}\frf\) は、基底 \(1,\delta,\ldots,\delta^{n-1}\) を整な量の基底に移す置換行列式の平方を表す。\footnote{Dedekind, Über die Diskriminante endlicher Körper（\emph{Abhandl. der Gött. Ges. d. Wissensch.} vol. 29, 1882）。そこで代数的数について与えられた定義は、代数関数に直接移すことができる。証明もほとんどそのまま移すことができる。D.--W. における証明の順序は異なる。下を見よ。引用した定理については Zahlbericht, § 31 および 32 も参照。（そこで \(F'(\delta)\) は整数 \(\delta\) の「異なるもの」と呼ばれる。）}

さらなる展開における\emph{相違}は、基礎体の特殊な性質によって条件づけられる。例えば代数的数の場合、\emph{基礎体の元が同時に指数でもある}という事実により、分岐イデアルが、\(p\) の \(e\) 乗に現れる素イデアル \(\frp\) の \((e-1)\) 乗より高い冪で割り切れる可能性が生じる（すなわち \(e\) が \(p\) で割り切れる場合である）。これはさらに慣性体と分岐体の概念\footnote{Zahlbericht, § 39--45 を参照。}へ導くが、これらは代数関数の理論には類似物をもたない。しかし何よりも、基礎体に\emph{無限に多くの定数が存在する}ため、代数関数ではイデアル類数の有限性が消え、それとともに類数の決定に関わるすべての問題も消える。それにもかかわらず、ある意味で、すべての（超越的な）素形式の領域\footnote{Klein, Zur Theorie der Abelschen Funktionen（\emph{Ann.} 36, 1890）。}を類体の超越的類似物と見なすことができる。ここでは各点が一つの形式、すなわち主イデアルによって生成されるからである。

他方、基礎体 \(K(z)\) に無限に多くの定数が存在することは、代数関数について、代数的数には類似物をもたない定理と問題を導く。第一に、このことから証明の配置にある簡略化が生じる。例えば D.--W. による、イデアルの素イデアルへの一意分解の証明では、任意の一次形式 \((z-c)\) が無限に多くの剰余類、すなわちすべての定数をもつという事実を用いることによって、「拡張された Gauss の定理」またはそれと同値な定理を避けることができる（これに対し、素数を法とする剰余類の数は有限である）。\footnote{D.--W. § 9。pp. 212, 213 の注を参照。}

さらに、\(z\) の整有理関数を係数にもつ任意の多項式が \((z-\alpha)\) を法として一次因子に分解するという事実（これは整係数多項式を \(p\) を法として考える場合には成り立たない）が、本当に新しい結果を導く。この事実は、すべての定数の代わりにすべての代数的数だけをとってもなお成り立つ。\footnote{D.--W. は序論で、この場合にも全理論がそのまま保たれることを述べている。} すなわち、ここからすべての素イデアルが一次のイデアルであること\footnote{D.--W., § 9, no. 7。}、さらに、体判別式 \(D\) が、\(\delta\) がすべての整な量を走るときのすべての判別式 \(\Delta(\delta)\) の交わりになること\footnote{D.--W., p. 224。}が従う。この二点はいずれも代数的数の場合と対照的である。続いてここから、分岐イデアルがすべての主イデアル \(F'(\delta)\) の交わりであることもふたたび従う。

さらに別の相違は、代数的数では有理数体 \(R\) が絶対的に特別なもの、すなわちその体に含まれる素体（単位から導かれる体）であるのに対し、基礎体 \(K(z)\) は相対的なものであるという事実によって条件づけられる。任意の非定数関数 \(\xi\) を独立変数として選ぶことができ、そのとき体は \(K(\xi)\) 上の代数体となる。各素イデアル \(\frp\) が一次のイデアルであり、したがって各関数が \(\frp\) を法として定数に合同であるという定理を、任意の関数を独立変数にとれるという可能性と合わせると、代数的数を超えて進む最も重要な概念、すなわち純粋に算術的な形での点概念に導かれる（Hensel 10a を参照）。「ここでは各点が、関数体に同型な定数体によって表される。」\footnote{よく知られているように、二つの体が「同型」であるとは、和、差、積、商も対応するように一対一対応に置けることをいう。} こう定義された点全体が「絶対 Riemann 面」をなす。この点の定義は体だけに依存し、特殊な変数には依存しない。しかし存在証明は、ある関数を独立変数 \(z\) として区別することによって行われる。点 \(P\) は、各関数に \(\frp\) を法とする定数剰余を割り当てることによって、素イデアル \(\frp\) により生成される。したがって、とくに \(\frp\) で割り切れる関数は \(P\) で消え、逆も成り立つ。\(\frp\) が \(z\) におけるすべての素イデアルを走ると、これにより \(z\) が有限であるすべての点が得られる。残りの点は独立変数 \(\xi=1/z\) を導入することによって生じ、主イデアル \(\xi\) を割る素イデアルに対応する。もちろん、点概念に付随する概念、例えば多角形、多角形類、多角形類の次元なども、代数的数には類似物をもちえない。

独立変数を自由に選べることは、分解 \(F'(\vartheta)=\frf\cdot\frD\) にさらに別の解釈を導くことも注意しておくべきである。すなわち、\(\vartheta\) を独立変数と見なすと、対応して
\[
\left(\frac{\p F}{\p z}\right)=\frf\cdot\fra
\]
を得る。なぜなら \(\vartheta\) から導かれる環は \(z\) と \(\vartheta\) に関して対称的に定義されていたからである。そして \(\frf\) は二つの主イデアル \(\frac{\p F}{\p z}\) と \(\frac{\p F}{\p\vartheta}\) の最大公約元となる。導手 \(\frf\) は同時に、\(\vartheta,z\) に関する「二重点イデアル」となる。\footnote{D.--W. p. 263。}

\subsection*{4. 級数展開の算術的基礎づけとイデアル論との関係}
Hensel--Landsberg では、級数展開という超越的補助手段と関数論から借りた点概念が、イデアル論の位置を占める。Hensel が最近与えた代数的数と代数関数のための級数展開の算術的基礎づけ\footnote{Eine neue Theorie der algebraischen Zahlen（\emph{Math. Zeitschr.} vol. 2, 1918）および Neue Begründung der arithmetischen Theorie der algebraischen Funktionen einer Variablen（\emph{Math. Zeitschr.} vol. 4, 1919）。より詳しい説明については Hensel no. 44ff. を参照。}は、したがってイデアル論の等価物と見なすことができる。実際、この算術的基礎づけは、各素イデアルについて、この素イデアルが主イデアルになるような（超越的）拡大体に移ることにほかならない。そうすれば通常の分解定理が成り立つ。そして、この考察は一つの素数 \(p\)、または一つの一次関数 \(z-a\) に現れるすべての素イデアルについて行えば十分である。

まず、\(z-a\)、または \(p\)（\(p\)-進数）の整冪級数展開全体からなる超越拡大体を導入する。ただし、それぞれ負冪は高々有限個だけ許す。この体において、代数的数体または関数体を定義する既約方程式は \(\rho\) 個の既約因子の積に分解する。これは \(p\)、または \(z-a\) が \(\rho\) 個の「一次的」イデアル \(\mathfrak q_1,\ldots,\mathfrak q_\rho\) の積に分解することに対応する。各一次的イデアルを素イデアルの冪としてさらに表示すること、
\[
\mathfrak q_i=\frp_i^{\,\ell_i},
\]
は、Hensel では超越拡大体に対して代数的なさらに一つの体を導入することに対応する。そして代数関数の場合、この体は純方程式によって定義できる。代数的数の場合には、より一般に、代数的に解ける方程式によって定義される。代数関数の場合の簡略化は、ここでは慣性体と分岐体（no. 3 参照）が現れないことによる。

これらの研究に現れる素元 \(\pi\) に関する「有限で終わる級数」\footnote{Neue Begründung \(\ldots\), 公式 (11)。}（これはまだ超越的な収束証明を必要としない）
\[
v=c_0+c_1\pi+\cdots+c_{\sigma-1}\pi^{\sigma-1}-v_\sigma\pi^\sigma,
\]
は、\(c_0,\ldots,c_{\sigma-1}\) が定数で \(v_\sigma\) が体の整元を表すとき、D.--W. にまったく対応する形で見いだされる。\footnote{D.--W., p. 231 および § 15。} そこではこれはイデアル論から導かれ、一つの点における関数の零の位数（または無限大になる位数）を定義するための基礎を与える。

\subsection*{5. 代数関数の算術的理論と幾何学的理論との比較。異なる不変性概念}
「不変」とは、体だけによって特徴づけられる任意の概念を意味する。算術的理論では、この不変性は一般にすでに定義の中に与えられている。その定義は、ある基底や 区別された変数に関してではなく、体のすべての元に関して述べられるからである。上で与えた点概念がその例である。これに対して残りの理論は、ある特殊な変数または基底から出発し、その後で概念がその特殊な選択に依存しないことを示す。ここでは不変性は双有理変換のもとでの不変性となる。

すなわち、体が一方では代数元 \(s\) を \(K(z)\) に添加することによって、他方では \(\sigma\) を \(K(\xi)\) に添加することによって生じるとする。ここで \(f(z,s)=0\)、また \(F(\xi,\sigma)=0\) は、それぞれ \(s\)、また \(\sigma\) の既約方程式を表す。あるいは、さらに一般に、\(K(\eta)\) に有限個の元を添加し、対応する既約方程式
\[
g_1(\eta,\tau_1)=0,\qquad g_2(\eta,\tau_1,\tau_2)=0,\ldots .
\]
をもつとする。すると定義により、すべての元は \(z\) と \(s\) によって、\(\xi\) と \(\sigma\) によって、また \(\eta,\tau_1,\tau_2,\ldots\) によって、有理的に表される。とくに \(z\) と \(s\) は \(\xi\) と \(\sigma\) によって有理的に表され、その逆も成り立つ。同様に \(z\) と \(s\) は \(\eta,\tau_1,\tau_2,\ldots\) によって有理的に表され、その逆も成り立つ。したがって平面曲線 \(f(z,s)=0\) は「双有理変換」によって平面曲線 \(F(\xi,\sigma)=0\) に移り、あるいは \(g_1(\eta,\tau_1)=0; g_2(\eta,\tau_1,\tau_2)=0,\ldots\) によって定義される変数 \(\eta,\tau_1,\tau_2,\ldots\) の空間内の曲線に移る。したがって不変性は、一つの特殊曲線について述べられた定義が、そこから双有理変換によって得られるすべての曲線（任意の次元数の空間内の曲線）について、すなわち Riemann の意味での「類」のすべての曲線について有効なままである、という事実の中に現れる。ここで「点」としては、\(f(a,b)=0\) となるような整合的な値系 \(z=a,s=b\) が定義される。双有理変換のもとで、この点は一般にふたたび \(F(\xi,\sigma)=0\) の整合的な値系 \(\alpha,\beta\) に移る。しかし特殊な場合、\(a,b\) が「特異」点であったときには、\(F(\xi,\sigma)=0\) 上の複数の「点」に対応することがある。そして \(f(z,s)=0\) の「特異」点が、分離して、または相続いて位置する有限個の単純点に対応するような曲線を常に指定できる（特異点の解消）。\footnote{例えば Brill--Noether 第6節を参照。「特異」点についてのより詳しい情報は次節に現れる。} 同様に、点群もまた点群に移る。そして適当な変換で対応する単純点の数だけ特異点を数えれば、その点数は不変である。任意の曲線は「通常多重点」をもつ曲線に双有理的に変換できるので、\footnote{同上。} 幾何学的理論はその概念をそのような特殊曲線についてだけ展開し、その後で、最も一般の特異点に至るものも含め、任意の双有理変換のもとでの不変性を証明する。

とくに、超楕円的な場合を除いて、その「類」の 特別な曲線として、\((p-1)\) 次元空間内のいわゆる「\(\varphi\) の正規曲線」を導入できる。これは特異点をもたない。ここで \(\varphi\) は、斉次化された \(f(z,s)=0\) に対する \(p\) 個の「\((n-3)\) 次の随伴曲線」である。\(f(z,s)=0\) から \(F(\xi,\sigma)=0\) への双有理変換には、\(\varphi\) の線形変換が対応する。したがってこれは射影的意味での不変性である。体のすべての関数を \(\varphi\) の有理関数として表す表示は、幾何学的理論では「不変表示」と呼ばれる。\footnote{例えば Brill--Noether 第8節を参照。多変数代数関数について、この不変性を線形変換のもとでのそのような不変性に還元できるかという問題は扱われていない。}

したがって \(\varphi\) の線形系、またそれらによって切り出される「特殊群」は、各双有理変換のもとで自分自身に移り、したがって体だけによって特徴づけられる。それは算術的理論における微分類に対応する。そこでは不変性が、体の各関数を独立変数とする性質によって与えられる定義の中に、ふたたび直接に示されている。

\subsection*{6. 比較の続き。特異点}
点概念の異なる把握と結びついて、代数幾何学的理論で特別な困難を生む曲線の「特異点」は、算術的理論の基礎づけにはまったく現れない。またその結果、幾何学的理論の場合とは違い、「随伴関数」は算術的理論を組み立てるためには必要ではなく、理論の特殊でより高次の部分に属する。一方、幾何学的理論の基礎づけには分岐点が現れない。実際、Aronhold の標準形による積分の表示（Hensel no. 27, 公式 80）でも、分岐点は消えてしまう。

5 で与えた特異点の定義は、次のように述べられる。\(f(z,s)=0\)（または高次元空間内の曲線）が \(z=a\), \(s=b\)（または \(z=a\), \(s_i=b_i\)）に特異点をもつとは、この値系が絶対 Riemann 面の複数の点（算術的意味での点）によってとられるか、または一つもしくは複数の点でより高い位数でとられることをいう。後者の場合、それは \(z\) と \(s\) の双方の分岐多角形の点である。これは幾何学的理論の「相続く点」に対応し、その点は \(f(z,s)=0\)（または高次元空間内の曲線）の（高次の）尖点となる。

体の任意の関数 \(\eta\) は、仮定により \(z\) と \(s\) によって有理的に表される。したがって、同型性により、\(z=a\), \(s=b\) となる絶対 Riemann 面の各点は、\(z\) と \(s\) によって表された各関数 \(\eta\) に、その \(z=a\), \(s=b\) での値を対応させることによって生じる。その点が特異であるためには、少なくとも一つの関数 \(\eta\) が、\(z=a\), \(s=b\) で複数の値をとるか、一つの値系を複数回とらなければならない。そして \(z\) と \(s\) によって有理的に表されるということから、これはその表示において \(z=a\), \(s=b\) で分子と分母が同時に消える場合にだけ起こる。したがって上の定義は、\(f\), \(\frac{\partial f}{\partial z}\), \(\frac{\partial f}{\partial s}\) が \(z=a\), \(s=b\) で消えるという通常の定義と同一である。実際、この場合 \(\eta=\frac{z-a}{s-b}\) は指定された種類の関数であり、\(z=a\), \(s=b\) で多価である。反対の場合には、たとえ \(\frac{0}{0}\) になっても、各関数は一つの値だけをとる。

算術的点概念では特異点は起こりえない。一つの関数にでも異なる値系が対応すれば、二つの点は異なるものと見なされるからである。しかし、算術的点概念の基礎づけに用いられるイデアル論においても、特異点の概念は現れない。これは、体のすべての整な量の全体を常に同時に考えるからである。\(z\) に関して代数的に整である \(\delta\) に対し、\(F(z,\delta)=0\) の有限部分にあるすべての特異点は、（第二の定義と no. 3 の末尾に従えば）「二重点イデアル」\(\frf\) によって生成される。さて、基本定理により分岐イデアル \(\frD\) はすべての主イデアル \(F'(\delta)\) の最大公約元であるから、任意の素イデアル \(\frp\) に対して、対応する \(\frf\) が \(\frp\) で割り切れないような \(\delta\) を指定できる。したがって、\(\frf\) は \(F'(\delta)\) と \(F'(z)\) の最大公約元であるので、少なくとも二つの主イデアル \(F'(\delta)\), \(F'(z)\) の一方は \(\frp\) で割り切れない。したがって、\emph{すべての \(F'(\delta)\) と \(F'(z)\) が同時に消える値系は存在しない}。その結果、値系 \(z=a,\delta_i=b_i\) は絶対 Riemann 面のただ一つの点でだけとられる。体の\emph{整代数関数の全体}（これは無限次元空間内の曲線と解釈してよい）は\emph{有限部分に特異点をもたない}。イデアル論で考慮されるのは有限値だけである。

しかし、すべての整な量の基底関数 \(\omega_1,\ldots,\omega_n\) も有限部分に特異点をもたない。上のことにより、算術的意味での各点は、すべての整関数に同型に対応させられた定数系によってすでに一意的に決定される。したがって、任意の整関数 \(s_i\) によって定義される「空間曲線」が \(s_i=\alpha_i\) において有限部分に特異点をもつならば、\(s_i=\alpha_i\) で複数の値系をとる整関数 \(\eta\) が少なくとも一つ存在しなければならない。いま \(\omega_1,\ldots,\omega_n\) がすべての整な量の基底を表すなら、整数的有理係数 \(a_i\) による \(\delta=a_1\omega_1+\cdots+a_n\omega_n\) のため、\(\omega\) の各値系には \(\delta\) の一つの値系だけが対応する。したがって \(\omega_1,\ldots,\omega_n\) によって定義される空間曲線は有限部分に特異点をもちえない。さらに、\emph{\(z\) が有限である絶対 Riemann 面の各点は、\(\omega\) に同型に対応させられた定数系によってすでに一意的に定義される。基底関数とはまさに、\(s_i=\alpha_i\) で多価になる関数 \(\eta\) である。} したがって、例えば基底 \(1,\delta,\ldots,\delta^{n-1}\) の代わりに整な量の基底を導入することによって、特異点の困難は回避される。

幾何学的理論では、絶対 Riemann 面のすべての点を一意的に生成する課題、より一般には与えられた点で無限大になるすべての関数を形成する課題は、随伴形式とその商によって遂行される。形式 \(\psi\) が随伴であるとは、基礎曲線 \(f(z,s)=0\) の各 \(i\) 重点で \(\psi=0\) が少なくとも \((i-1)\) 重点をもつことをいう。高次の特異点は、まず通常多重点に解消されなければならない。最も一般の関数 \(\eta=\frac{\psi}{\chi}\) について、分子と分母が \(f(z,s)=0\) のすべての特異点で消えなければならないことは上で示した。それらが随伴でなければならないことは、いたるところ有限な微分を考察すれば分かる。逆に、随伴性の条件が最も一般の関数を形成するために十分であることは、後で詳しく論じる「剰余定理」が示す。したがって随伴関数は、算術的理論の基底関数の位置を占める。

算術的には、特異点をすべて有限部分にあるものと仮定すれば（これは変換によって常に達成できる）、随伴形式は「二重点イデアル」\(\frf\) で割り切れる体の関数の分子に対応する。ここから、基底量 \(\omega_1,\ldots,\omega_n\) を随伴形式の商として表すこと、およびさまざまな表示方法の関係が、形式的に次のように示される。

\(R(z)\) を、基底 \(1,\delta,\ldots,\delta^{n-1}\) から基底 \(\omega_1,\ldots,\omega_n\) への置換行列式とする。すると逆に、各 \(\omega\) は \(z\) と \(\delta\) の整関数を \(R(z)\) で割った商として表される：
\[
\omega_i=\frac{G_i(z,\delta)}{R(z)}.
\]
\(R(z)\omega_i\) は \(\delta\) から導かれる環に属するから、導手の定義により \(R(z)\) は \(\frf\) で割り切れる［\((R(z))^2=\operatorname{Norm}\frf\) からは、この可除性は \((R(z))^2\) についてだけ従う］。いま \(\omega_i\) は整関数としてすべての有限な \(z\) で有限にとどまるので、分子関数 \(G_i(z,\delta)\) も \(\frf\) で割り切れなければならない。したがって、すべての \(\omega_i\) の分子と分母は随伴である。しかしこの \(\omega\) の表示から、体のすべての関数が随伴形式の商によって表示されることが従う。

以上により、幾何学的理論は、すべての整な量を同時に考える算術的理論と対比して、\emph{一つの量 \(\delta\) から導かれる環の理論}（整域、整性領域）として特徴づけられる。したがって、その環の導手、すなわち特異点が決定的な役割を果たすことは明らかである。環（整域）の算術的理論と幾何学的理論とのあいだには、さらに別の類似も与えられる。例えば幾何学的理論は、随伴曲線との交わりにおいて特異点へ落ちる交点をどこにも数えず、それとは異なる点群だけを考える。これに対応して Dedekind\footnote{Über die Anzahl der Idealklassen in den verschiedenen Ordnungen eines endlichen Körpers.（Festschrift Braunschweig 1877）；§ 3--5。（これらの節で代数的数について述べられた定理は、一変数代数関数について文字通り成り立つ。）} は、導手イデアル \(\frf\) と互いに素な環イデアルだけを考え、このイデアル領域に対して、素イデアルへの一意分解を含むすべての可除性法則を立てている。

また、幾何学的理論の基礎にある「代数関数の基本定理」\footnote{M. Noether, Über einen Satz aus der Theorie der algebraischen Funktionen. \emph{Math. Ann.} vol. 6 (1873).}、あるいは少なくともその直接の帰結は、ある意味で環の算術的理論に類似物をもつ。この定理は、表示可能性
\[
\psi(z,s)=A(z,s)f(z,s)+B(z,s)\varphi(z,s)
\]
の条件を与える。ここで現れるすべての量は \(s,z\) に関して整有理である（より一般には \(x_1,x_2,x_3\) に関して整かつ斉次である）。これらの条件から、\footnote{Brill--Noether, no. 55.} \(f(z,s)=0\) によって、商 \(\frac{\psi}{\varphi}\) が \(f\) に随伴であるとき、常に整有理関数 \(B(z,s)\) に等しいことが従う。これに対応するのが、no. 3 の定理、すなわちこれらすべての比較の基礎にある定理である。一つの \(\delta\) から導かれる環の導手 \(\frf\) は、\(f(z,\delta)=0\) の二重点イデアルに等しい。なぜなら、導手の定義により、\(\frf\) で割り切れる各代数的整量はその環に属し、したがって \(z\) と \(\delta\) によって整かつ有理的に表示できるからである。一方、引用した定理は、そのような量が随伴関数であると言っている。\footnote{Dedekind に従って、ここでは一貫して \(\delta\) が代数的整数であると仮定している（これは変換によって常に達成できる）。幾何により正確に対応する場合、すなわち \(s\) が \(z\) に代数的に分数的に依存してよい \(f(z,s)=0\) は、Hensel--Landsberg が扱っている。Hensel no. 26 および 29 を参照。}

\subsection*{7. 剰余定理とイデアル論の諸定理との比較}
今述べた「代数関数の基本定理」の帰結は、幾何学的理論において直接に「剰余定理」へ導き、したがって幾何学的理論全体の基礎へ導く。しかし、ある点では、剰余定理は他の基礎よりも再び初等的な性格をもつ。というのも、それはイデアルの素イデアルへの一意分解、より正確にはこの分解定理の基礎にある事実からの直接の帰結として算術的に述べることができ、したがって高次の環の理論には依存しないからである。この算術的定式化へは次のように導かれる。

基礎曲線 \(f(z,\delta)=0\) が、変数の一次分数変換によって常に達成できるように、特異点および考慮している有限個の点群を有限部分にもつものと仮定する。さらに \(\delta\) は \(z\) に代数的に整に依存していると仮定する。これは変数に整な後続の線形変換によって常に達成できる。このとき、\(z=a,\delta=b\) となる絶対 Riemann 面の点 \(P\) を生成する素イデアル \(\frp\) に対して、\(f(z,\delta)=0\) の点 \(z=a,\delta=b\) が対応する。より一般に、イデアル \(\fra=\frp_1^{\,e_1}\cdots\frp_\rho^{\,e_\rho}\) は、対応する重複度をもつ点 \(z=a_i,\delta=b_i\) からなる点群を生成する。そして \(f(z,\delta)=0\) の有限部分にあるすべての点はこのように生成される。点 \(P\) では \(\frp\) で割り切れるすべての関数 \(g_1,g_2,\ldots\) が消え、逆にこの関数 \(g\) の全体がイデアル \(\frp\) を生成するから、\(f\) 上の対応する点は
\[
g_1(z,\delta)=0,\quad g_2(z,\delta)=0,\ldots
\]
と \(f=0\) との交わりである。イデアル \(\fra\) に対応する点群についても同じである。実際、各イデアルが二つの主イデアルの最大公約元であるという定理により、二つのそのような曲線で常に一つの点群を切り出すのに十分である。とくに主イデアルに対応する点群は曲線 \(g(z,\delta)=0\) によって切り出され、逆も成り立つ。主イデアルは完全交差に対応する。

二つの点群 \(A\) と \(B\) は、幾何学的理論において、同じ点群 \(R\) を加えることにより完全交差に完成できるとき、\emph{余剰対応} であると呼ばれる。これらの点群がそれぞれイデアル \(\fra,\frb,\frr\) によって生成されるなら、これらのイデアルのあいだには関係
\[
\fra\cdot\frr=(\alpha),\qquad \frb\cdot\frr=(\beta),
\]
があり、これはイデアル \(\fra\) と \(\frb\) が同値であることを示す。逆に、二つのイデアルの同値性 \(\fra=\eta\frb\) からは、各イデアルが別のイデアルを掛けることによって主イデアルに変換できるという定理に基づいて、常にこの関係が従う。\footnote{Dedekind, 数論 § 181 を参照。すなわち、\(\fra\frr\) が主イデアル \((\alpha)\) となるように \(\frr\) を選ぶと、\(\frb\frr=\eta^{-1}\cdot(\alpha)\) となる。そして \(\frb\frr\) は整な量だけを含むから、\(\eta^{-1}\cdot(\alpha)\) についても同じことが成り立ち、したがってこれも主イデアル \((\beta)\) である。} \emph{したがって、同値なイデアルと余剰対応する点群とは同一の概念である。}

剰余定理はいま次のように言う。\emph{\(A\) と \(B\) が余剰対応であるなら、これらは同じ剰余をもつ随伴曲線（特異点の外で）によって切り出される。} 関数の零点と極は常に余剰対応であるから、これは同時に、体の最も一般の関数が随伴形式の商として表されることを証明する。イデアル論の言葉では、剰余定理は単に、\(\fra\) と \(\frb\) のほかに \(\fra\frf\) と \(\frb\frf\) も同値であるということである。ここで \(\frf\) は二重点イデアルを表す。というのも、上によって、常にあるイデアル \(\frm\) が存在して
\[
\fra\frf\frm=(\alpha),\qquad \frb\frf\frm=(\beta)
\]
が主イデアルになるからである。\(\alpha=0\) と \(\beta=0\) は、点群 \(A\) と \(B\)、または特異点と異なるその部分群 \(A'\) と \(B'\) を切り出す随伴曲線である。\footnote{\(\fra=\frf\fra_1,\ \frb=\frf\frb_1,\ \frr=\frf\frr_1\) であり、\(\fra_1,\frb_1,\frr_1\) が互いに素であるなら、すでに \(\fra_1\frf\mathfrak n,\ \frb_1\frf\mathfrak n\) が主イデアルとなる。ここで \(\mathfrak n\) は \(\frf\) と互いに素である。任意のイデアルは、与えられたイデアルと互いに素なイデアルを掛けることによって主イデアルに変換できるからである（D.--W. p. 214；または数論 § 178, 定理 XI）。} したがって、剰余定理の証明は、算術的には、点を素イデアルに対応させられる可能性と、同値なイデアルを同一のイデアルを掛けることで主イデアルに変換できるという証明の中に含まれている。

余剰対応する二つの点群は、必ずしも同じ数の点からなる必要はないことに注意すべきである。これは同値なイデアルが同じ次数をもつ必要がないのと同様である。例えば、すべての主イデアルは同値であり、完全交差によって生成されるすべての点群は余剰対応である。同値な多角形（除数）\footnote{多角形 \(A\) は絶対 Riemann 面の有限個の点からなる。二つの多角形 \(A\) と \(B\) は、ある関数 \(\eta\) の零点と極であるとき同値と呼ばれる。このとき \(\eta\) は多角形商 \(\eta=\frac{A}{B}\) によって表される。（D.--W. § 17 および 18。）}との相違は、\(z\) が無限大になる点が \(z\) の素イデアルによって生成されないという事実によって説明される。したがって、\(\fra\) と \(\frb\) の同値性を表すイデアル商 \(\eta=\frac{\fra}{\frb}\) としての関数の表示には、\(z\) が無限大になる点での \(\eta\) の零点や極は現れない。例えば主イデアル \((\alpha)\) は、整関数 \(\alpha\) の零点に対応する素イデアルの積になる。\(\alpha\) の極は現れない。なぜなら \(\alpha\) が無限大になるのは、\(z\) が無限大になる点に限られるからである。この点で、イデアル論は幾何学的理論における形式論の正確な類似である。そこでも同様に零点だけがあり、極はない。

形式から関数への移行は、幾何学的には同じ次数の形式の商を作ることによって行われる。算術的には、これはイデアル類から、不変的に定義された多角形（除数）類への移行に対応する。というのも、ここではもはや特別な変数が区別されていないので、関数を多角形商で表示するとすべての零点と極が現れるからである。イデアル類（余剰対応点群の類）と多角形類との関係は次のとおりである。\(\calA\) と \(\calB\) をイデアル \(\fra\) と \(\frb\) によって生成される多角形とし、\(\calN\) を \(z\) が無限大になる点の多角形（\(z\) および \(z\) の整関数の分母多角形）とする。このとき、\(\fra\) と \(\frb\) の同値性から \(\calA\) と \(\calB\) の同値性が従うのは、\(\fra\) と \(\frb\) が同じ次数をもつ場合、かつその場合に限る。一般には、\(\calA\calN^\nu\) と \(\calB\calN^{\nu'}\)（\(\nu\ne\nu'\)）の同値性だけが従う。したがって、イデアル（または余剰対応点群）を類に分けることは、無限個の多角形類を一つの類にまとめることである。それはまた、多角形を \(\calN\) の冪を法として類に分けることとも見なせる。\footnote{Hensel--Landsberg, 第25講, p. 438。--- 二つの変数 \(z\) と \(s\) を多角形商で表すなら、これはこれら各変数の二元斉次化と理解できる。とくに \(z\) と \(s\) が同じ分母多角形をもつように選べば、幾何で通常用いられる三元斉次化が生じる。}

異なる数の点からなる余剰対応点群の場合の剰余定理は、純粋に幾何学的な問題でだけ現れる。代数関数への応用では、実際には多角形類に対応する場合、すなわち同じ数の点をもつ余剰対応点群の場合だけが現れる。一つの点群によって生成される「完全線形系」、すなわちその点群に余剰対応し同じ点数をもつ点群全体は、対応する多角形によって生成される多角形類に正確に対応する。「二つの完全系の和」と「二つの多角形類の積」という概念も一致する。そのような完全系が有限個の線形独立な点群だけを含むことは、剰余定理の直接の帰結である。剰余定理は、系の次元（線形独立な点群の数）を、任意だが固定された次数をもち、剰余群を通るすべての随伴曲線の次元に帰着する。この次元の実際の決定、すなわち Riemann--Roch の定理に到達するために、幾何学的理論はまず剰余定理から「還元定理」（固定点に関する定理）\footnote{Brill--Noether, no. 58。}を導き、そこから Riemann--Roch の定理を導く。

これとまったく対応して、算術的理論は、各多角形類が有限次元であることを示し、整な量の特殊な基底、すなわち補完加群によって定義される正規基底によってその次元を決定する。こうして、補完加群と分岐多角形との関係の帰結として Riemann--Roch の定理を得る。D.--W. の提示にも還元定理が現れる\footnote{D.--W., § 30, no. 2。}が、幾何学的理論におけるような基礎としてではない。no. 6 で述べた基底関数と随伴関数との対応の結果として、ここでも両理論の補助手段のあいだにある対応を立てることができる。要約すれば、二つの理論における概念形成は、本質的には一致していると言える。ただし、定式化を別にして、その配置は異なる。二つの理論は互いに独立に生じた。幾何学的理論のほうが十年早い。さらに注意すべきことは、Riemann--Roch の定理はもともと、Riemann と Roch においても Weierstrass においても、階数定理として現れるということである。問題となるのは行列 \((\varphi_i(x_j))\) の階数であり、ここで \(\varphi\) は \(n-3\) 次の \(p\) 個の随伴曲線を走り、\(x\) は剰余群の点を走る。

\subsection*{8. 代数関数と代数的数における超越的問題}
Hilbert は、とくに代数関数と代数的数における超越的問題の類似を指摘した。\footnote{Mathematische Probleme, problem 12.（\emph{Göttinger Nachrichten} 1900）} これは本質的には存在問題である。したがって、与えられた Riemann 面（したがって与えられた分岐点）に対して代数関数が存在するかという Riemann の問題は、指定された群と判別式をもつ代数的数体の存在問題に対応する。\footnote{これによって示唆された、群から出発する代数関数の扱いは、R. König によって、より一般的な問題の特殊な場合として遂行された。とくに Riemannsche Funktionen- und Differentialsysteme in der Ebene（\emph{J. f. M.} vol. 148）および \emph{Math. Ann.} 78, 79 を参照。代数体の代わりに、König には \(K(z)\) に関する特殊な有限加群がある。そこでは「類の性質」だけが存在する。} 代数関数の場合には存在証明を一般に実行できるのに対し、代数的数体の領域では、それは特殊な基礎体（有理数、二次数体）について、しかも Abel 群についてだけ達成されている。\footnote{三次および四次方程式に現れる群の単純な場合と、八次方程式のいくつかのきわめて特殊な群だけは、任意の基礎体について扱うことができる：G. Bucht, Über einige algebraische Körper 8. Grades（\emph{Arkiv f. Math., Astron. och Physik}, vol. 6, no. 30）。E. Noether, Gleichungen mit vorgeschriebener Gruppe およびそこで引用された Seidelmann の学位論文（\emph{Math. Ann.} 78）も参照。} これは、有理数上の各 Abel 数体が 1 の根の体から組み立てられるという Kronecker の定理、および相対 Abel 数体についての Fueter\footnote{\emph{Math. Ann.} 75。} と Hecke\footnote{\emph{Math. Ann.} 76。} の対応する定理に関わる。したがって、求める数体は実際に、指数関数、モジュラー関数という超越的手段によって与えられる。これは代数関数についての純粋な存在証明と対照的である。

代数関数の存在証明には、いたるところ有限な積分の存在が先立つ。これに類似するのは、素数 \(\ell\) に関して定義される代数的数体の「特異主数」の存在である。その \(\ell\) 乗根は相対的に不分岐であり、類体の最初の構成要素を与える。\footnote{Furtwängler, Reziprozitätsgesetze für Potenzreste mit Primzahlexponenten in algebraischen Zahlkörpern. \emph{Math. Ann.} 67 および 72。} この存在証明は、数体において指定された剰余指標をもつ素イデアルが常に存在するという定理の本質的な助けを借りて行われる。したがってこの定理は、いたるところ有限な積分の存在が依拠する境界値問題の類似物と見なすことができる。

いたるところ有限な積分は、二つの多角形が同じ類に属するかどうかを決定する超越的判定条件を与える（同じ数の点からなる二つの点群が余剰対応であるかどうかを決定する）。すなわち Abel の定理による判定である。この定理は、第一種積分を「余剰対応する道」に沿って延長したものが消えると言う。あるいは、対応する微分定理とその逆によっても判定できる。\footnote{例えば M. Noether（\emph{Ann.} 37）の定式化を参照。剰余定理は歴史的には Abel の定理から生じたので、幾何学的理論では積分和に従って点群と線形系の「和」を語る。これは算術的理論における類の「積」と対照的である。}

代数的数体に対する Abel の定理の類似物としては、「\(\ell\) 乗冪剰余に対する相互法則」を取らなければならない。これは、体そのもの、または少なくとも適当な上体について、二つのイデアルが同じ類に属するための判定条件を与える。その内容は、この上体において、特異主数が同じ類のすべてのイデアルに関して同じ剰余指標をもつ、と述べることもできる。この後者の事実は体そのものでも成り立つが、そこでは相互法則とは一致しない。\footnote{Furtwängler, loc. cit.}

最後に、Hensel の \(p\)-進数の理論は、代数関数の冪級数展開の類似物を与える。これについては Hensel の報告を参照されたい。


% Paper 15 current invariant-theory macros
\providecommand{\frH}{\mathfrak H}
\providecommand{\frK}{\mathfrak K}
\providecommand{\frM}{\mathfrak M}
\providecommand{\frN}{\mathfrak N}
\providecommand{\frG}{\mathfrak G}
\providecommand{\frS}{\mathfrak S}
\providecommand{\tmod}[1]{\;(\operatorname{mod} #1)}

\clearpage
\begin{center}
{\Large\bfseries 15. 二元形式の整係数不変式系の有限性}\par
\vspace{1em}
\emph{Nachrichten von der Gesellschaft der Wissenschaften zu Göttingen} 1919, pp. 138--156
\end{center}

\vspace{2em}
\begin{center}
1919年3月27日の会合で F. Klein により提出。
\end{center}

以下では、既知の有限性定理の拡張として、次の定理を証明する。
\emph{二元基本形式系のすべての整係数多項式不変式は、そのうち有限個を用いた、整数係数の多項式的有理関数である。}
ここで通常どおり、「不変式」とは常に、基本形式または基本形式系の多項式的有理不変式、すなわち、すべての線形変換の群によって誘導される係数変換に関する不変式を意味する。

証明は、Mertens--Hilbert の有限性証明を整性の面で強めたものに基づく。\footnote{F. Mertens, Wiener Sitzungsberichte Jan. 1889; D. Hilbert, \emph{Math. Ann.} vol. 33 (1889) [1888年3月の日付]. Mertens, Crelle vol. 100 に続いて互いに独立に生じた二つの証明は、内容上は一致している。}
この有限性証明は次のように特徴づけられる。基本形式をその線形因子に分解する、より正確には各基本形式に線形形式の積を対応させる。すると各不変式 \(I\) はこれらの線形形式の不変式、したがって同次化された根の関数に移る。この関数は三つの条件を満たさなければならない。1) 不変性の条件、2) 重み条件、3) 対称性条件である。条件 2) と 3) は、\(I\) が基本形式の係数の多項式的かつ有理的な関数となるために必要である。条件 1) は、\(I\) を線形形式の行列式の多項式的有理関数として表すことによって満たされる。条件 2) は、個々の行列式の代わりに、それらのある冪積を新しい変数として用いることによって満たされる。すると条件 3) は、\(I\) が量の行たちの多項式的有理対称関数になること、また逆に、そのような量の行たちの任意の対称関数が基本形式の不変式になることを述べるだけである。したがって、これらの量の行の基本対称関数が完全系を与える。

整係数不変式だけを考える場合も、線形形式の不変式について、より強い定理、すなわち\emph{二元線形形式のすべての整係数不変式は、それらの線形形式の行列式の整係数多項式関数である}ことが証明されれば、条件 1) は上とまったく同様に満たされる。私は § 3 で、これは実際に成り立つことを、これらの行列式の間の任意の素数を法とするすべての関係の全体が、すべての恒等関係の全体と一致する、という形で示す。すなわち、これらの関係に対応する加群は、任意の素数を法としても素加群であり続ける。この検証は、行列式の間の二次関係に対応する二次形式の加群に関して剰余類を正規化することによって行われる。その際、これらの二次形式は、任意の素数を法として加群の整基底をなすことが分かる。従来それらは、整性を無視したときにだけ加群基底として知られていたものである。

条件 2) を満たすためには変更は不要である。冪積の導入は整数係数に影響しないからである。条件 3) はいま、\(n_1!\cdots n_\mu! I\)（ここで \(n_1,\ldots,n_\mu\) は個々の基本形式の次数である）が、量の行たちの整係数対称関数になり、逆に、これらの行の任意の整係数対称関数が整係数不変式になることをいう。しかし容易に示されるように、これらの行の整係数対称関数は整基底をもち、その基底には基本関数だけではもはや十分でない。Hilbert の整加群基底に関する定理を適用すれば、したがって \(I\) 自身に対しても整基底が得られる。

§ 1 では、条件 1), 2), 3) に対応する特殊な有限性定理を集める。それらから § 2 で有限性の証明が従う。条件 3) に対応する特殊な有限性定理は、整数または代数的整数の置換からなる有限群の整係数不変式の有限性も与える。これを § 4 で簡潔に述べる。

\section*{§ 1. 特殊な有限性定理}

まず、導入で述べた不変性、重み、対称性の条件を満たすために用いる、次の特殊な有限性定理を述べる。

\textbf{I.} 二元線形形式の系の任意の整係数不変式は、これらの線形形式の行列式の、整数係数をもつ多項式的有理関数である。証明は § 3 で与える。

\textbf{II.} 非負指数をもつ \(n\) 個の変数の冪積からなる任意の系 \(\frS\) が、二つの冪積 \(f,g\) とともに、その商が変数に関する多項式である限りその商も含むとする。このとき \(\frS\) の関数からなる有限の乗法的基底 \(f_1,\ldots,f_k\) が存在し、\(\frS\) の任意の \(f\) に対して
\[
        f=f_1^{\lambda_1}\cdots f_k^{\lambda_k}
\]
という表示が、非負指数 \(\lambda_1,\ldots,\lambda_k\) によって成り立つ。

実際、Hilbert の加群基底定理により、\(\frS\) の関数の有限個 \(f_1,\ldots,f_k\) を、各々が実際に \(x\) を含むように選べて、\(\frS\) の任意の非定数 \(f\) はそれらで生成される加群に属する。
\[
        f\equiv0\pmod{(f_1,\ldots,f_k)}.
\]
したがって、任意の \(f=x_1^{i_1}\cdots x_n^{i_n}\) に対して、少なくとも一つの \(f_\nu=x_1^{\nu_1}\cdots x_n^{\nu_n}\) が存在し、\(\nu_1\le i_1,\ldots,\nu_n\le i_n\) を満たす。表示
\[
        f=x_1^{i_1-\nu_1}\cdots x_n^{i_n-\nu_n} f_\nu = A f_\nu
\]
において、因子 \(A\) は仮定により \(\frS\) に属するが、変数に関する次数は \(f\) より小さい。したがってこの議論を有限回繰り返すことで主張が従う。表示を \(f=1\) に対しても成り立たせるには、すべての指数 \(\lambda\) をゼロにすればよい。\footnote{この表示は直接にも証明でき、逆に Hilbert の定理はこの表示から従う。P. Gordan, Neuer Beweis des Hilbertschen Satzes über homogene Funktionen, Gött. Nachr. 1899 を参照。別の直接証明、およびそこからの定理 II の導出は、A. Ostrowski, Über die Existenz einer endlichen Basis bei gewissen Funktionensystemen, \emph{Math. Ann.} 78 (1916), § 2,1 および § 3,2 にある。§ 2 で用いる定理 II の特殊な場合、すなわち指数がディオファントス方程式系のすべての非負解を動く場合は、すでに Gordan-Kerschensteiner, \emph{Vorlesungen über Invariantentheorie}, vol. I, p. 199 に見られる。}

\textbf{III.} \(n\) 個の量の行に関するすべての多項式的整係数対称関数は、そのうち有限個の多項式的整係数関数である。

\(\tau_1(x),\ldots,\tau_n(x)\) を不定元 \(x_1,\ldots,x_n\) の基本対称関数とする。すると任意の指数 \(k\) について恒等式
\[
        x_i^k = G_0^{(k)}(\tau)+x_iG_1^{(k)}(\tau)+\cdots+x_i^{n-1}G_{n-1}^{(k)}(\tau),
        \qquad (i=1,2,\ldots,n),
\]
があり、ここで \(G_0^{(k)},\ldots,G_{n-1}^{(k)}\) は \(\tau(x)\) の整係数多項式関数である。\(x_i,y_i,\ldots,z_i\) が第 \(i\) 行の要素であるとき、これらの恒等式および \(y,\ldots,z\) についての対応する恒等式によって、すべての特殊な一行型対称関数
\[
        \sum_{i=1}^n x_i^{a}y_i^{b}\cdots z_i^{c}
\]
は、有限個の
\[
        \sum_{i=1}^n x_i^{a}y_i^{b}\cdots z_i^{c}
        \qquad
        (0\le a<n,\;0\le b<n,\;\ldots,\;0\le c<n)
\]
および基本対称関数 \(\tau(x),\tau(y),\ldots,\tau(z)\) の整係数多項式関数になる。\footnote{Hilbert, loc. cit. は、\(\tau(x),\ldots,\tau(z)\) の代わりに冪和を基底に入れる。そのため基底は一行型関数だけからなるが、表示の整性は失われる。}
\(n\) 行の最も一般の対称関数を考える場合（以下では生じない場合）も、同様にして、関数
\[
        \sum x_{i_1}^{a_1}y_{i_1}^{b_1}\cdots z_{i_1}^{c_1}
              x_{i_2}^{a_2}y_{i_2}^{b_2}\cdots z_{i_2}^{c_2}
              \cdots
              x_{i_n}^{a_n}y_{i_n}^{b_n}\cdots z_{i_n}^{c_n}
\]
\[
        (0\le a_\nu<n,\;0\le b_\nu<n,\;\ldots,\;0\le c_\nu<n),
\]
ここで \(i_1,i_2,\ldots,i_n\) はすべての置換を動くが、これらの関数は \(\tau(x),\tau(y),\ldots,\tau(z)\) とともに整基底をなすことが示される。

\textbf{IV.} \(F_1(x),\ldots,F_k(x)\) を \(x_1,\ldots,x_n\) の整係数多項式関数とし、\(a\) を固定された整数とする。このとき
\[
        \frac{G(F)}{a},
\]
ただし \(G\) は整係数多項式関数を表す、と書ける関数のうち、\(x\) に関して整になるものすべては、そのうち有限個の整係数多項式関数である。

実際、Hilbert の整加群基底に関する定理により、そのような各 \(G(F)/a\) について表示
\[
 \frac{G(F)}{a}
 =
 A_1(F_1\ldots F_k)\frac{G_1(F)}{a}
 +\cdots+
 A_\rho(F_1\ldots F_k)\frac{G_\rho(F)}{a},
\]
があり、\(A_1,\ldots,A_\rho\) は \(F\) の整係数多項式関数であり、
\[
        \frac{G_1(F)}{a},\ldots,\frac{G_\rho(F)}{a}
\]
はその系に属する。したがって関数
\[
        F_1=\frac{aF_1}{a},\ldots,F_k=\frac{aF_k}{a};
        \quad
        \frac{G_1(F)}{a},\ldots,\frac{G_\rho(F)}{a}
\]
が求める基底をなす。

\setcounter{equation}{0}
\section*{§ 2. 二元整係数不変式の有限性}

\(x,y\) を不定係数をもつ変換
\begin{equation}
        x=s_{11}x'+s_{12}y',\qquad y=s_{21}x'+s_{22}y',
\end{equation}
を受ける二元変数とし、\(f\) を係数が不定元である二元基本形式とする。このとき (1) により
\begin{equation}
        f=a_0x^n+a_1x^{n-1}y+\cdots+a_ny^n
        =a_0'x'^n+a_1'x'^{\,n-1}y'+\cdots+a_n'y'^n.
\end{equation}
\(f\) の整係数不変式とは、周知のように、整係数多項式関数 \(I(a)\) であって
\begin{equation}
        I(a')=|s_{ik}|^\rho I(a),\qquad \text{\(a,s_{ik}\) に関して恒等的に}
\end{equation}
を満たすものをいう。基本形式系の整係数不変式も対応して定義される。個々の基本形式の係数に関する同次性を定義に含めてもよい。なぜなら、残りの不変式は、そうした個別に同次な不変式から整線形に構成されるからである。

(2) に加えて、特別な二元形式 \(g\)、すなわち係数がやはり不定元である \(n\) 個の線形形式の積を考える。\footnote{不変式論でしばしば行われるように、\(f\) を二項係数付きで \(f=b_0x^n+\binom n1b_1x^{n-1}y+\cdots+b_ny^n\) と書くならば、すべての整係数不変式 \(I(b)\) の領域は \(I(a)\) の領域に比べて拡大し、ここで用いる有限性の議論は成り立たなくなる。\(I(b)\) はもはや根の整関数ではない。}
\begin{align}
        g&=(\alpha_1x+\beta_1y)(\alpha_2x+\beta_2y)\cdots(\alpha_nx+\beta_ny)\notag\\
         &=\sigma_0(\alpha,\beta)x^n+\sigma_1(\alpha,\beta)x^{n-1}y+\cdots+\sigma_n(\alpha,\beta)y^n\notag\\
         &=(\alpha_1'x'+\beta_1'y')(\alpha_2'x'+\beta_2'y')\cdots(\alpha_n'x'+\beta_n'y')\notag\\
         &=\sigma_0'x'^n+\sigma_1'x'^{\,n-1}y'+\cdots+\sigma_n'y'^n.
\end{align}
したがって \(\sigma_0(\alpha,\beta),\sigma_1(\alpha,\beta),\ldots,\sigma_n(\alpha,\beta)\) は同次化された基本対称関数を表し、変換後の係数 \(\sigma_i'\) は \(\sigma_i(\alpha',\beta')\) に等しくなる。

これらの基本関数の代数的独立性により、任意の関係
\begin{equation}
        G(\sigma_0,\sigma_1,\ldots,\sigma_n)=0
        \quad \text{\(\alpha,\beta\) に関して恒等的に}
\end{equation}
には、さらに関係
\begin{equation}
        G(\sigma_0,\sigma_1,\ldots,\sigma_n)=0
        \quad \text{\(\sigma_0,\ldots,\sigma_n\) に関して恒等的に}
\end{equation}
が対応し、したがって不定元 \(a_0,\ldots,a_n\) について
\begin{equation}
        G(a_0,a_1,\ldots,a_n)=0
        \quad \text{\(a_0,\ldots,a_n\) に関して恒等的に}
\end{equation}
が成り立つ。

いま \(f\) の任意の整係数不変式 \(I(a)\) は、置換 \(a_i=\sigma_i(\alpha,\beta)\) によって、\(g\) の整係数不変式 \(I(\sigma)=H(\alpha,\beta)\) を与える。このとき \(H\) は \(\alpha,\beta\) の整関数である。この不変式に対して関係 (3) は
\begin{equation}
        I(\sigma')=|s_{ik}|^\rho I(\sigma)
        \quad \text{\(\sigma,s_{ik}\) に関して恒等的に}
\end{equation}
となり、\(\sigma'=\sigma(\alpha',\beta')\) であるから、(8) は
\begin{equation}
        H(\alpha',\beta')=|s_{ik}|^\rho H(\alpha,\beta)
        \quad \text{\(\alpha,\beta;s_{ik}\) に関して恒等的に}
\end{equation}
を与える。したがって \(I(\sigma)=H(\alpha,\beta)\) は、\(n\) 個の線形形式 \((\alpha_i x+\beta_i y)\) の整係数不変式になる。逆に、\(H(\alpha,\beta)\) がこれら \(n\) 個の線形形式の整係数不変式、すなわち (9) を満たすものであり、同時に \(\sigma\) の整関数 \(I(\sigma)\) に等しいと仮定する。このとき
\[
        H(\alpha',\beta')=I(\sigma(\alpha',\beta'))=I(\sigma')
\]
であるから、式 (9) は
\begin{equation*}
        I(\sigma')=|s_{ik}|^\rho I(\sigma)
        \quad \text{\(\alpha,\beta;s_{ik}\) に関して恒等的に}.\tag{9a}
\end{equation*}
とも書ける。(5) と (6) から (8) が従い、(7) から (3) が従う。したがって各整係数不変式 \(I(a)\) には \(n\) 個の線形形式の整係数不変式 \(H(\alpha,\beta)=I(\sigma)\) が対応し、逆も成り立つ。さらに、\(I_1(a),\ldots,I_k(a)\) がすべての \(I(a)\) の整基底ならば、\(H_1(\alpha,\beta),\ldots,H_k(\alpha,\beta)\) は、同時に整関数 \(I(\sigma)\) であるすべての \(H(\alpha,\beta)\) の整基底である。ここでも逆が (5), (6), (7) により成り立つ。したがって、すべての \(I(a)\) に対する整基底の存在を証明するには、同時に整関数 \(I(\sigma)\) であるすべての整係数不変式 \(H(\alpha,\beta)\) についてそれを証明すれば十分であり、基底
\[
        H_1(\alpha,\beta)=I_1(\sigma),\ldots,H_k(\alpha,\beta)=I_k(\sigma)
\]
が \(I(a)\) に対する基底 \(I_1(a),\ldots,I_k(a)\) を与える。

\(H(\alpha,\beta)=I(\sigma)\) に (4) のすべての整係数不変式を動かせる。(3) により各不変式は \(\sigma\) に関して同次であり、(4) により各 \(\sigma\) は各線形形式の係数に関して同次かつ一次であるから、\(H(\alpha,\beta)\) は各行 \(\alpha_i,\beta_i\) に関して同次で、各行において同じ次元をもつ。さらにもちろん、\(n\) 行の対称関数である。逆に、対称関数に関する基本定理により、\(n\) 行の任意の整係数対称関数 \(H(\alpha,\beta)\) が、各行に関して別々に同次で、各行において同じ次元をもつならば、それは \(\sigma\) の整係数多項式関数である。したがって前の議論により、\(H(\alpha,\beta)\) が不変であれば、ただちに不変式 \(I(\sigma)\) である。よって (4) の整係数不変式の全体は、次の性質をもつ整関数 \(H(\alpha,\beta)\) の全体として特徴づけられる。
\[
\begin{array}{ll}
1) & \text{\(n\) 個の線形形式 \(\alpha_i x+\beta_i y\) の整係数不変式である;}\\
2) & \text{\(n\) 個の各行 \(\alpha_i,\beta_i\) に関して別々に同次で、同じ次元をもつ;}\\
3) & \text{\(n\) 個の行 \(\alpha_i,\beta_i\) に関して対称である。}
\end{array}
\]
特殊有限性定理 I により、\(H(\alpha,\beta)\) は 1) のために、行列式
\[
        (ik)=\alpha_i\beta_k-\beta_i\alpha_k;
\]
の整係数多項式関数になる。したがって
\[
        H(\alpha,\beta)=G((ik)).
\]
この表示において、対称性は \(n\) 行のすべての置換を \(G((ik))\) に適用することで形式的にも表せる。
\begin{equation}
        n!\,H(\alpha,\beta)
        =
        G^{(1)}((ik))+G^{(2)}((ik))+\cdots+G^{(n!)}((ik))
        =G^*((ik)).
\end{equation}
したがって \(G^*\) は、\((ik)\) の各冪積 \(P\) とともに、\(P\) のすべての置換にわたる和 \(Q=P^{(1)}+\cdots+P^{(n!)}\) も含む。実際、\(G^*\) は有限個の \(Q\) の整線形結合 \(G^*=\sum cQ\) になる。\(Q\) は条件 1), 2), 3) を満たし、したがってそれ自身不変式になるから、\(G^*((ik))\) の代わりに \(Q\) を考えれば十分である。

いま \(P\) を、\(Q\) に現れるすべての \((ik)\) の冪積、すなわち条件 2) を満たす冪積の系 \(\frS\) の中で動かす。二つの \(P\) の商が \((ik)\) に関して多項式ならば、それも条件 2) を満たし、したがって \(\frS\) に属する。したがって定理 II により、このような \(P\) は有限個 \(P_1,\ldots,P_\rho\) を用いて、任意の \(P\) が
\[
        P=P_1^{\lambda_1}\cdots P_\rho^{\lambda_\rho}
\]
という形に、非負指数 \(\lambda\) で表される。すべての置換を適用すると
\[
        Q=
        P_1^{(1)\lambda_1}\cdots P_\rho^{(1)\lambda_\rho}
        +\cdots+
        P_1^{(n!)\lambda_1}\cdots P_\rho^{(n!)\lambda_\rho}
\]
となる。こうして \(Q\) は、それぞれ \(\rho\) 個の要素をもつ \(n!\) 個の量の行
\[
        P_1^{(1)}\ldots P_\rho^{(1)};
        \quad
        P_1^{(2)}\ldots P_\rho^{(2)};
        \quad \ldots;\quad
        P_1^{(n!)}\ldots P_\rho^{(n!)};
\]
に関する整係数対称関数になる。よって定理 III により、それはこれらの行に関する有限個の対称関数の整係数多項式関数である。これらの基底関数は、それ自身 \(Q\) であるか、または \(n!\) 個の要素 \(P_i^{(1)},P_i^{(2)},\ldots,P_i^{(n!)}\) の基本対称関数であり、したがって条件 1), 2), 3) を満たし、整係数不変式
\begin{equation}
        H_1(\alpha,\beta)=I_1(\sigma);\ldots;\quad
        H_l(\alpha,\beta)=I_l(\sigma)
\end{equation}
になる。(10) と \(G^*=\sum cQ\) により、任意の \(H(\alpha,\beta)=I(\sigma)\) はしたがって
\[
        H(\alpha,\beta)=I(\sigma)=\frac1{n!}\Phi(I_1,\ldots,I_l)
\]
の形である。ここで \(\Phi\) は整係数多項式関数を表す。逆に、\(\frac1{n!}\Phi\) が \(\alpha,\beta\) に関して多項式であれば、それは整係数不変式である。したがって定理 IV により、さらに有限個の不変式 \(I_{l+1}(\sigma),\ldots,I_k(\sigma)\) を (11) に付け加えて、
\[
        H_1(\alpha,\beta)=I_1(\sigma);\ldots;H_l(\alpha,\beta)=I_l(\sigma);
\]
\[
        H_{l+1}(\alpha,\beta)=I_{l+1}(\sigma);\ldots;H_k(\alpha,\beta)=I_k(\sigma)
\]
がすべての不変式 \(H(\alpha,\beta)=I(\sigma)\) の整基底をなすようにできる。これで一つの基本形式の場合の定理が証明された。

基本形式の系を扱う場合には、先に述べたように、各基本形式の係数に関する不変式の同次性を仮定してよい。再び一般の基本形式を、それぞれ線形形式の積であるものに置き換える。その不変式は、1) これらの線形形式の不変式、2) 線形形式の係数からなる行に関して別々に同次で、各個別の基本形式に属する行において等しい次元をもつもの、3) 各個別の基本形式に属する行に関して対称なもの、になる。逆に、これらの条件を満たすすべての線形形式の整係数不変式は、再び基本形式系の整係数不変式に対応する。証明は今与えたものと完全に平行である。重み条件（定理 I と II）を満たす個々の行列式の冪積 \(P=P_1^{\lambda_1}\cdots P_\rho^{\lambda_\rho}\) ごとに、すべての
\[
        N=n_1!n_2!\cdots n_\mu!
\]
個の置換を適用する。ここで \(n_1,\ldots,n_\mu\) は基本形式の次数である。このようにして \(N\cdot I\) は、それぞれ \(\rho\) 個の要素をもつ \(N\) 個の行の整係数対称関数に変換される。これからすべての \(N\cdot I\) に対する整基底が得られ、したがって定理 IV により、すべての \(I\) に対する整基底が得られる。

\setcounter{equation}{0}
\section*{§ 3. 二元線形形式の整係数不変式の基底}

特殊有限性定理 I の証明を残している。Clebsch によって初めて証明された既知の定理により、\(n\) 個の線形形式 \(\alpha_i x+\beta_i y\) のすべての多項式的有理不変式は、行列式 \((ik)\) の多項式的有理関数である。したがって、とくに、これらの線形形式のすべての整係数不変式は、これらの行列式の多項式的有理関数になるが、必ずしも整係数関数になるとは限らない。

以下で示すのは、\((ik)\) の間に存在する恒等関係を考慮すると、この表示は常に整係数で可能であるということである。

式で、かつ加群論の言葉で、これは次のように表される。\(\xi_{ik}\), \(i<k\), を不定元とする。このとき、\(F((ik))\) が \(\alpha,\beta\) に関して恒等的に消えるという性質をもつすべての整係数多項式 \(F(\xi_{ik})\) の全体は加群 \(\frM\) をなす。そのような二つの \(F\) の和、およびそのような \(F\) と \(\xi_{ik}\) の任意の整係数多項式との積は、再びこの全体に属する。式では
\[
        F((ik))=0\quad \text{\(\alpha,\beta\) に関して恒等的に}
\]
ならば
\[
        F(\xi_{ik})\equiv0\pmod{\frM}\quad \text{\(\xi_{ik}\) に関して恒等的に},
\]
かつその逆も成り立つ。

いま次の完全に類似した主張を示す。
\[
        F((ik))\equiv0\pmod p\quad \text{\(\alpha,\beta\) に関して恒等的に}
\]
ならば
\begin{equation}
        \text{次が従う:}\qquad
        F(\xi_{ik})\equiv0\pmod{(\frM,p)}
        \quad \text{\(\xi_{ik}\) に関して恒等的に},
\end{equation}
かつその逆も、任意の素数 \(p\) について成り立つ。ゆえに \(\frM\) は、\((ik)\) の間の \(p\) を法とするすべての関係の全体に対応する加群 \(\frM_p\) とも同一である。この加群は、\(F((ik))\) が \(\alpha,\beta\) に関して恒等的に \(p\) を法として消えるような、すべての整係数多項式 \(F(\xi_{ik})\) からなる。

(1) が整係数表示可能性を述べていることは次のように分かる。それはまた、\(F((ik))=p\cdot H(\alpha,\beta)\) が \(\alpha,\beta\) に関して恒等的に成り立つならば、
\[
        F(\xi_{ik})-p\cdot G(\xi_{ik})\equiv0\pmod{\frM}
        \quad \text{\(\xi_{ik}\) に関して恒等的に}
\]
が従う、という形にも書ける。したがって \(\frM\) の定義により、
\[
        F((ik))=p\cdot G((ik))\quad \text{\(\alpha,\beta\) に関して恒等的に}.\footnote{この最後の結論は同時に (1) の逆も証明している。}
\]
ゆえに \(p\cdot H(\alpha,\beta)=p\cdot G((ik))\)、すなわち
\[
        H(\alpha,\beta)=G((ik)).
\]
したがって、もし表示
\[
        H(\alpha,\beta)
        =
        \frac1{p_1^{\lambda_1}\cdots p_\rho^{\lambda_\rho}}
        F((ik)),
\]
ただし \(p_1,\ldots,p_\rho\) は素数、が与えられているならば、同時に
\[
        H(\alpha,\beta)
        =
        \frac1{p_1^{\lambda_1-1}\cdots p_\rho^{\lambda_\rho}}
        G((ik))
\]
も従い、これを有限回繰り返せば表示の整性が得られる。

(1) の証明のため、
\[
        f_{ik,rs}=\xi_{ik}\xi_{rs}-\xi_{ir}\xi_{ks}+\xi_{is}\xi_{kr},
        \qquad i<k<r<s,
\]
で、\((ik)\) の間の二次関係
\[
        (ik)(rs)-(ir)(ks)+(is)(kr)=0
\]
に対応する \(\xi_{ik}\) の形式を表すものとし、これらから生成される整係数多項式の加群を \(\frN=(f_{ik,rs})\) とする。
\[
        F(\xi_{ik})\equiv0\pmod{\frN}
\]
からは
\[
        F((ik))=0\quad \text{\(\alpha,\beta\) に関して恒等的に}
\]
が従い、また
\begin{equation}
\begin{array}{rcl}
        \text{もし }F(\xi_{ik})&\equiv&0\pmod{(\frN,p)}\text{ ならば}\\
        F((ik))&\equiv&0\pmod p\quad\text{\(\alpha,\beta\) に関して恒等的に}
\end{array}
\end{equation}
が従う。以下では、\(\frN\) に関して剰余類を正規化することで、(2) の逆も成り立つことを示す。すると \(\frN=\frM=\frM_p\) となり、(1) が証明される。等式 \(\frN=\frM\) は、\(p=0\) の場合の \(\frN=\frM_p\) として得られる。以下の考察では、\(p=0\) の場合も常に含めておく。言い換えれば、\(f_{ik,rs}\) は \(\frM_p\) の整加群基底をなし、とくに \(\frM_0=\frM\) の基底でもある。

\subsection*{a) 二つまたは三つの行 \(\alpha_i,\beta_i\) の場合}

四つの異なる指数が現れないので、\(\frN\) は零加群である。したがって、
\begin{equation}
        F((ik))\equiv0\pmod p\quad \text{\(\alpha,\beta\) に関して恒等的に}
\end{equation}
から
\[
        F(\xi_{ik})\equiv0\pmod p\quad \text{\(\xi_{ik}\) に関して恒等的に}
\]
が従うことを示せばよい。これは \(\frM=\frM_p=0\) として (1) を証明する。

条件 \(F((ik))=p\cdot H(\alpha,\beta)\) は、個々の行に関して同次な成分ごとに別々に満たされなければならないので、以下では常に、\(F((ik))\) が各個別の行に関して同次、したがって \(F(\xi_{ik})\) が各指数に関して同次であると仮定して十分である。二つの行の場合には、行列式 \((12)\) と対応する不定元 \(\xi_{12}\) だけが存在する。仮定した同次性により
\[
        F=c\,\xi_{12}^\nu.
\]
\[
        c\,(12)^\nu\equiv0\pmod p\quad \text{\(\alpha,\beta\) に関して恒等的に}
\]
から、\((12)\not\equiv0\pmod p\) は \(\alpha,\beta\) に関して恒等的には成り立たないので、\(c\equiv0\pmod p\) が従う。ゆえに
\[
        F(\xi_{ik})\equiv0\pmod p\quad \text{\(\xi_{ik}\) に関して恒等的に}.
\]
これで (3)、したがって (1) が証明された。

三つの行の場合には、行列式 \((12),(13),(23)\) と対応する不定元 \(\xi_{12},\xi_{13},\xi_{23}\) がある。したがって
\[
        F=\sum c\,\xi_{12}^{\tau_{12}}\xi_{13}^{\tau_{13}}\xi_{23}^{\tau_{23}}.
\]
\(F\) の指数 \(1,2,3\) に関する次数をそれぞれ \(\sigma_1,\sigma_2,\sigma_3\) とする。このとき
\[
        \sigma_1=\tau_{12}+\tau_{13},\qquad
        \sigma_2=\tau_{12}+\tau_{23},\qquad
        \sigma_3=\tau_{13}+\tau_{23},
\]
であり、一意な逆写像
\[
        \tau_{12}=\frac12(\sigma_1+\sigma_2-\sigma_3),\quad
        \tau_{13}=\frac12(\sigma_1-\sigma_2+\sigma_3),\quad
        \tau_{23}=\frac12(-\sigma_1+\sigma_2+\sigma_3)
\]
をもつ。したがって各々の別々に同次な \(F\) は高々一項
\[
        F=c\,\xi_{12}^{\tau_{12}}\xi_{13}^{\tau_{13}}\xi_{23}^{\tau_{23}}
\]
を含むだけであり、そこから上と同様に (3)、したがって (1) が従う。

\subsection*{b) 四つの行 \(\alpha_i,\beta_i\) の場合}

六つの不定元 \(\xi_{12},\xi_{13},\xi_{14},\xi_{23},\xi_{24},\xi_{34}\) と対応する行列式が現れ、\(\frN=(f_{12,34})\) である。なぜなら
\begin{equation}
        \xi_{14}\xi_{23}\equiv \xi_{13}\xi_{24}-\xi_{12}\xi_{34}\pmod{\frN},
\end{equation}
であるから、任意の \(\xi_{ik}\) の冪積は、\(\xi_{14}\xi_{23}\) という積を含まない冪積の整線形結合に \(\frN\) を法として合同である。したがって任意の多項式 \(F(\xi_{ik})\) に対して
\begin{equation}
        F(\xi_{ik})\equiv F_0(\xi_{ik})\pmod{\frN},
\end{equation}
ただし \(F_0\) は \(\xi_{14}\xi_{23}\) という積を含まない。\(F_0\) を \(F\) の剰余類の正規化多項式と呼ぶ。(4) は別々に同次であるから、別々に同次な \(F\) には再び別々に同次な \(F_0\) が対応する。

いま
\begin{equation}
        F_0(\xi_{ik})=G(\xi_{ik})+\xi_{34}F_1(\xi_{ik}),
\end{equation}
とおく。ただし \(G\) は \(\xi_{34}\) で割れる冪積を含まない。さらに
\begin{equation}
        [G]_{\xi_{14}=\xi_{13}}=K(\xi_{12},\xi_{13},\xi_{23})
\end{equation}
とおく。
\[
        F((ik))\equiv0\pmod p\quad \text{\(\alpha,\beta\) に関して恒等的に}
\]
から、(2) を考慮した (5) により、
\begin{equation}
        F_0((ik))\equiv0\pmod p\quad \text{\(\alpha,\beta\) に関して恒等的に}
\end{equation}
が従う。\(\alpha_4=\alpha_3,\;\beta_4=\beta_3\) と特殊化すると、(6) と (7) は
\[
        K((ik))\equiv0\pmod p\quad \text{\(\alpha,\beta\) に関して恒等的に}
\]
を与える。したがって a) で証明した結果により、
\[
        K(\xi_{ik})\equiv0\pmod p\quad \text{\(\xi_{ik}\) に関して恒等的に}
\]
である。これから下で示すべき残りの主張
\[
        G(\xi_{ik})\equiv0\pmod p\quad \text{\(\xi_{ik}\) に関して恒等的に}
\]
が従う。したがって (6) は
\[
        F_0(\xi_{ik})\equiv \xi_{34}F_1(\xi_{ik})\pmod p
        \quad \text{\(\xi_{ik}\) に関して恒等的に}
\]
を与える。\((34)\not\equiv0\pmod p\) は \(\alpha,\beta\) に関して恒等的には成り立たないので、(8) から
\[
        F_1((ik))\equiv0\pmod p\quad \text{\(\alpha,\beta\) に関して恒等的に}
\]
が従う。ここで \(F_1(\xi_{ik})\) は再び正規化されており、\(\xi_{34}\) に関する次数は低くなっている。この手続きを有限回繰り返すと
\[
        F_0(\xi_{ik})\equiv\xi_{34}^{\lambda}F_\lambda(\xi_{ik})\pmod p,
\]
ただし \(F_\lambda\) は \(\xi_{34}\) に関して次数ゼロであり、したがって直前に示したことにより \(p\) を法として消える。よって
\begin{equation}
\begin{array}{rcl}
        F_0(\xi_{ik})&\equiv&0\pmod p\quad \text{\(\xi_{ik}\) に関して恒等的に},\quad\text{または}\\
        F(\xi_{ik})&\equiv&0\pmod{(\frN,p)}\quad \text{\(\xi_{ik}\) に関して恒等的に}
\end{array}
\end{equation}
を得る。これは (2) の逆を、したがって (1) を、同時に \(F_0(\xi_{ik})\) に関するより鋭い主張として証明している。\footnote{このより鋭い主張は、各剰余類がただ一つの正規化多項式を含むという意味である。}

あとは、\(K(\xi_{ik})\equiv0\pmod p\) から \(G(\xi_{ik})\equiv0\pmod p\) が従うことを付け加えるだけである。同値に、\(G(\xi_{ik})\not\equiv0\pmod p\) から \(K(\xi_{ik})\not\equiv0\pmod p\) が従う。正規化が必要なのは誘導のこの部分だけである。置換 \(\xi_{14}=\xi_{13}\) のもとで、\(G\) の個々の冪積は消えない。したがって \(K\) が消えるとすれば、それは \(G\) における異なる冪積が \(K\) では等しい冪積に対応するからでしかない。そしてこれは、それらの冪積が単に指数 3 と 4 の入れ替えによって異なる場合に限って起こり得る。たとえば、そのような正規化された冪積が
\[
        \xi_{12}^{\tau_{12}}\xi_{13}^{\tau_{13}}\xi_{14}^{\tau_{14}}\xi_{24}^{\tau_{24}}
\]
であるとしよう。各指数に関する別々の同次性の仮定により、一つの \(\xi_{ik}\) における 3 から 4 への置き換えは、別の \(\xi\) における 4 から 3 への置き換えを伴わなければならない。したがって関係する項は \(\xi_{13}\xi_{24}\) の代わりに \(\xi_{14}\xi_{23}\) という積を含むことになり、これは正規化によって排除されている。まったく同様に、
\[
        \xi_{12}^{\tau_{12}}\xi_{13}^{\tau_{13}}\xi_{14}^{\tau_{14}}\xi_{24}^{\tau_{24}}
\]
に対して相殺し得るのは、因子 \(\xi_{14}\xi_{23}\) を含む項だけであり、これも正規化により排除されている。ゆえに \(K(\xi_{ik})\equiv0\pmod p\) は \(G(\xi_{ik})\equiv0\pmod p\) を含意する。

\subsection*{c) \(n\) 個の行 \(\alpha_i,\beta_i\) の場合}

\(n\) 個の行の場合は、いま扱った四つの行の場合の厳密な一般化として、完全帰納法で解決される。\footnote{証明は \(n=3\) の場合にもなお適用される。}
まず再び、任意の多項式 \(F\) に対して合同
\begin{equation}
        F(\xi_{ik})\equiv F_0(\xi_{ik})\pmod{\frN}
        \quad \text{\(\xi_{ik}\) に関して恒等的に}
\end{equation}
が存在することを示す。ここで \(\frN=(f_{ik,rs})\) で、\(F_0(\xi_{ik})\) は正規化されている。正規化とは次のことを意味するものとする。\(i,k,r,s\) が四つの異なる指数で \(i<k<r<s\) のとき、\(F_0\) は積 \(\xi_{is}\xi_{kr}\) を含んではならない。すなわち、一つの \(\xi\) の二つの指数が、もう一つの \(\xi\) の二つの指数の間にともに入ってはならない。この正規化は、二つおよび三つの行の場合には何の制限も課さない。そこではすべての多項式がすでに正規化されているからである。四つの行の場合には、すでに示したように、各剰余類に正規化多項式が存在する。したがって \((n-1)\) 個の行については、各剰余類に正規化多項式が存在することを証明済みと仮定してよい。

\(F\) における任意の冪積を
\[
        \xi_{i_1n}\cdots\xi_{i_\nu n}\,P
\]
の形に書く。ただし \(P\) はもはや指数 \(n\) を含まない。誘導仮定により、\(P\) は \(\frN\) を法として整係数の正規化多項式 \(P_0\) に合同である。
\[
        \xi_{i_1n}\cdots\xi_{i_\nu n}P
        \equiv
        \xi_{i_1n}\cdots\xi_{i_\nu n}P_0
        \pmod{\frN}.
\]
もし \(P_0\) が、少なくとも一つの指数 \(i_\lambda\) に対して二つの指数 \(r,s\) が \(i_\lambda\) と \(n\) の間にある、すなわち \(i_\lambda<r<s<n\) となるような \(\xi_{rs}\) を含まないならば、すでに正規化に到達している。なぜなら、\(P_0\) から取った任意の二つの \(\xi\) の積について仮定が満たされるからである。そうでない場合、そのような指数 \(i_\lambda\), \(i_\lambda<r<s<n\), を取り、\(P_0^*\) を \(\xi_{rs}\) を含む \(P_0\) の冪積とする。
\[
        \xi_{i_\lambda n}\xi_{rs}\equiv
        \xi_{i_\lambda s}\xi_{rn}-\xi_{i_\lambda r}\xi_{sn}
        \pmod{\frN}
\]
であるから、
\begin{align*}
 &\xi_{i_1n}\cdots\xi_{i_\lambda n}\cdots\xi_{i_\nu n}P_0^*\\
 &\equiv
 \xi_{i_1n}\cdots\xi_{rn}\cdots\xi_{i_\nu n}Q
 +
 \xi_{i_1n}\cdots\xi_{sn}\cdots\xi_{i_\nu n}R
 \pmod{\frN}\\
 &\equiv
 \xi_{i_1n}\cdots\xi_{rn}\cdots\xi_{i_\nu n}Q_0
 +
 \xi_{i_1n}\cdots\xi_{sn}\cdots\xi_{i_\nu n}R_0
 \pmod{\frN},
\end{align*}
を得る。ここで \(Q_0\) と \(R_0\) は、\(Q\) と \(R\) の剰余類の正規化多項式を再び表す。これらは誘導仮定により存在する。\(Q\) と \(R\) は指数 \(n\) を含まないからであり、しかも \(Q\) と \(R\) は単なる冪積なので、それらは確かに整係数多項式である。\(i_\lambda<r<s<n\) であるから、さらに
\[
 ((n-i_1)+\cdots+(n-r)+\cdots+(n-i_\nu))
 <
 ((n-i_1)+\cdots+(n-i_\lambda)+\cdots+(n-i_\nu)),
\]
\[
 ((n-i_1)+\cdots+(n-s)+\cdots+(n-i_\nu))
 <
 ((n-i_1)+\cdots+(n-i_\lambda)+\cdots+(n-i_\nu))
\]
が成り立つ。したがって \(P_0^*\) から \(Q_0,R_0\) に移ると、この差の和は減少している。この和は必ず \(\ge\sigma\) であるから、この手続きは各冪積 \(P_0^*\) について有限回で停止しなければならない。ゆえに
\[
        F(\xi_{ik})
        =
        \sum \xi_{i_1n}\cdots\xi_{i_\nu n}P^{(1)}
        \equiv
        \sum \xi_{j_1n}\cdots\xi_{j_\sigma n}S_0^{(j)}
        =
        F_0(\xi_{ik})\pmod{\frN},
\]
ここで整係数多項式 \(F_0(\xi_{ik})\) は必ず正規化されている。そうでなければ差の和をさらに下げられるからである。これで (10) が証明された。\((n-1)\) 個の指数についての誘導仮定により、各々の別々に同次な \(F\) は別々に同次な正規化 \(F_0\) をもつので、いまの手続きはそれが \(n\) 個の指数についても成り立つことを示す。したがって以下では、別々に同次な多項式だけを考えればよい。

b) と同様に、
\begin{equation}
        F_0(\xi_{ik})=G(\xi_{ik})+\xi_{n-1,n}F_1(\xi_{ik}),\footnote{(10) と (11) に対応する恒等式
        \[
          F((ik))=F_0((ik))=G((ik))+(n-1,n)F_1((ik)),
        \]
        は、二つの行 \((n-1)\) と \((n)\) を変数の行 \(x,y\) で置き換え、したがって \((n-1,n)\) を \((xy)\) で置き換えるならば、Clebsch--Gordan 級数展開の類似物である。しかし Clebsch--Gordan が極形式による展開を含むのに対し、\(F_0\) の正規化は極形式の「初項」に対応し、それによって Clebsch--Gordan では成り立たない整性を強制する。整性を無視したとき \(f_{ik,rs}\) が \(\frM\) の基底をなすことの既知の証明は、まさに Clebsch--Gordan 級数展開によって行われる (Gordan-Kerschensteiner, \emph{Vorlesungen über Invariantentheorie}, p. 132)。}
\end{equation}
とおく。ここで \(G\) は \(\xi_{n-1,n}\) で割れる冪積を含まない。さらに
\begin{equation}
        [G]_{\xi_{in}=\xi_{i,n-1}}=K(\xi_{ik})
\end{equation}
とおく。正規化の定義が示すように、\(K\) も \(G\) と同じく正規化されている。

再び \(G\) と \(K\) の間には一対一対応がある。\(G\) の異なる冪積が \(K\) で同じ積に対応することがあり得るのは、それらが単に指数 \(n\) と \((n-1)\) の入れ替えによって異なる場合だけである。たとえば、\(G\) が指数 \((n-1)\) に関して次数 \(\rho\)、指数 \(n\) に関して次数 \(\sigma\) であり、
\[
        \chi_1\le\chi_2,\ldots,\le\chi_\rho\le\chi_{\rho+1},\ldots,\le\chi_{\rho+\sigma}
\]
が \(G\) の冪積の中で \((n-1)\) と \(n\) に結びつく指数であるとする。(11) により不定元 \(\xi_{n-1,n}\) は \(G\) に現れないので、その個数は \(\rho+\sigma\) でなければならない。正規化により、その冪積はしたがって
\[
        \xi_{\chi_1,n-1}\cdots\xi_{\chi_\rho,n-1}
        \xi_{\chi_{\rho+1},n}\cdots\xi_{\chi_{\rho+\sigma},n}\,P(\xi_{ik}),
\]
の形である。ここで \(P(\xi_{ik})\) はもはや指数 \(n\) と \((n-1)\) を含まない。\(n\) と \((n-1)\) を入れ替えると、\(i\ne j\) に対して因子 \(\xi_{in}\xi_{j,n-1}\)（\(i<j\)）が生じる。すると \(j,n-1\) は \(i\) と \(n\) の間にあり、その項は正規化されていない。したがって、固定された \(\chi_1,\chi_2,\ldots,\chi_{\rho+\sigma}\) と固定された \(P\) に対して、\(G\) は一つの冪積しか含まない。これにより、\(G\) の異なる冪積は \(K\) で異なる積に対応することが証明された。したがって
\begin{equation}
        K(\xi_{ik})\equiv0\pmod p\quad\Longrightarrow\quad
        G(\xi_{ik})\equiv0\pmod p.
\end{equation}

\[
        F((ik))\equiv0\pmod p\quad \text{\(\alpha,\beta\) に関して恒等的に}
\]
から、(2) を考慮した (10) により、
\begin{equation}
        F_0((ik))\equiv0\pmod p\quad \text{\(\alpha,\beta\) に関して恒等的に}
\end{equation}
が従う。\(\alpha_n=\alpha_{n-1},\;\beta_n=\beta_{n-1}\) と特殊化すると、(11) と (12) は
\begin{equation}
        K((ik))\equiv0\pmod p\quad \text{\(\alpha,\beta\) に関して恒等的に}
\end{equation}
を与える。\(K(\xi_{ik})\) は正規化されており、\((n-1)\) 個の指数だけを含むので、a) および b) の式 (9) で証明したことから
\[
        K(\xi_{ik})\equiv0\pmod p\quad \text{\(\xi_{ik}\) に関して恒等的に}
\]
が従う。すると (13) により
\[
        G(\xi_{ik})\equiv0\pmod p\quad \text{\(\xi_{ik}\) に関して恒等的に}
\]
である。したがって (11) は
\[
        F_0(\xi_{ik})\equiv \xi_{n-1,n}F_1(\xi_{ik})\pmod p
        \quad \text{\(\xi_{ik}\) に関して恒等的に}
\]
を与える。\((n,n-1)\not\equiv0\pmod p\) は \(\alpha,\beta\) に関して恒等的には成り立たないので、(14) から
\[
        F_1((ik))\equiv0\pmod p\quad \text{\(\alpha,\beta\) に関して恒等的に}
\]
が従う。ここで \(F_1(\xi_{ik})\) は再び正規化され、\(\xi_{n-1,n}\) に関する次数は \(F_0\) より低い。この過程を有限回繰り返すと
\[
        F_0(\xi_{ik})\equiv\xi_{n-1,n}^{\lambda}F_\lambda(\xi_{ik})\pmod p,
\]
ただし \(F_\lambda\) は \(\xi_{n-1,n}\) に関して次数ゼロであり、したがっていま証明したことによって \(p\) を法として消える。こうして
\begin{equation}
\begin{array}{rcl}
        F_0(\xi_{ik})&\equiv&0\pmod p\quad \text{\(\xi_{ik}\) に関して恒等的に},\quad\text{または}\\
        F(\xi_{ik})&\equiv&0\pmod{(\frN,p)}\quad \text{\(\xi_{ik}\) に関して恒等的に}
\end{array}
\end{equation}
を得る。これは (2) の逆を、したがって (1) を証明する。同時に、\(n\) 個の指数の場合、したがって一般に、誘導段階で用いた \(F_0(\xi_{ik})\) に対するより鋭い主張も証明している。

\setcounter{equation}{0}
\section*{§ 4. 有限群の整係数不変式の有限性}

私のノート「有限群の不変式に対する有限性定理」\footnote{\emph{Math. Ann.} 77, p. 89 (1915).}で与えた有限性の初等的証明は、§ 1 の特殊有限性定理 III と IV によって整性へと強めることができる。一方、Hilbert の加群基底定理に基づく通常の証明はこの整性を与えない。整加群基底を用いても、群の位数である整数 \(h\) の任意に高い冪が分母に現れる表示しか得られない。

以下では、いま引用したノートの記法を用いる。有限群 \(\frH\) は、非零行列式をもつ \(h\) 個の線形変換 \(A_1,\ldots,A_h\) からなるとする。\footnote{\(\frH\) が抽象的な有限群に単純に同型であると仮定するだけで十分であろう。すなわち、\(A\) のある変換行列式が消えるなら、\(\frH\) は群 \(\begin{pmatrix}\frS&0\\0&0\end{pmatrix}\) に相似であるべきであり、ここで \(\frS\) の置換の行列式は消えない。この場合、\(\frH\) の不変式は \(\frS\) の不変式と一致するので、行列式が消えないという条件は実際には一般性を制限しない。}ここで \(A_k\) は線形変換
\[
        x_1^{(k)}=\sum_{\nu=1}^{n}a_{1\nu}^{(k)}x_\nu,
        \quad\ldots,\quad
        x_n^{(k)}=\sum_{\nu=1}^{n}a_{n\nu}^{(k)}x_\nu,
\]
あるいは簡潔に \((x^{(k)})=A_k(x)\) を表す。したがって群 \(\frH\) は、要素 \(x_1,\ldots,x_n\) をもつ行 \((x)\) を、要素 \(x_1^{(k)},\ldots,x_n^{(k)}\) をもつ行 \((x^{(k)})\) へ送る。係数 \(a_{i\nu}^{(k)}\) は代数的整数であると仮定する。これは一般性を制限しない。J. Schur の定理\footnote{Über Gruppen linearer Substitutionen mit Koeffizienten aus einem algebraischen Zahlkörper. Satz VII. \emph{Math. Ann.} 71, p. 555 (1912).}により、任意の線形変換の有限群は、\(a_{i\nu}^{(k)}\) が代数的整数となるものに相似であるからである。

群の整絶対不変式とは、\(a_{i\nu}^{(k)}\) によって生成される数体 \(\frK\) の代数的整数を係数とする、\(x_1,\ldots,x_n\) の多項式的有理関数であって、\(A_1,\ldots,A_h\) の適用のもとで恒等的に変わらないものをいう。そのような不変式 \(f(x)\) に対しては
\begin{equation}
        f(x)=f(x^{(1)})=\cdots=f(x^{(h)})
        =\frac1h\sum_{k=1}^{h}f(x^{(k)}).
\end{equation}
逆に、量の \(h\) 個の行 \((x^{(k)})\) の任意の整係数対称関数も、群の整係数不変式になる。

\(\omega_1,\omega_2,\ldots,\omega_\rho\) を \(\frK\) のすべての整数の基底とする。すると任意の不変式は
\[
        f(x)=\omega_1 f_1(x)+\cdots+\omega_\rho f_\rho(x),
\]
と表せる。ここで \(f_1(x),\ldots,f_\rho(x)\) は有理整数係数をもつ。したがって (1) により
\[
        h\,f(x)
        =
        \omega_1\sum_{k=1}^{h}f_1(x^{(k)})
        +\cdots+
        \omega_\rho\sum_{k=1}^{h}f_\rho(x^{(k)}).
\]
ここで個々の和は、量の \(h\) 個の行 \((x^{(k)})\) の対称関数であり、しかも有理整数係数をもつ。したがって定理 III により、それらは、これらの行の有限個の整有理対称関数を用いて、有理整数係数をもつ多項式的かつ有理的な形で表される。前の議論により、これらの基底関数 \(I_1(x),\ldots,I_\sigma(x)\) は、一般には代数的整数係数をもつ群の整係数不変式になる。というのも、\(\frH\) は \(A\) とともにその代数的共役変換のすべてを含むとは限らないからである。

したがって表示
\begin{equation}
        h\cdot f(x)=\omega_1G_1(I)+\cdots+\omega_\rho G_\rho(I)
\end{equation}
がある。ここで \(G_1,\ldots,G_\rho\) は有理整数係数をもつ。

いま \(w_1,\ldots,w_\rho\) を不定元とし、
\[
        L=w_1G_1(I)+\cdots+w_\rho G_\rho(I)
\]
に、置換 \(w_i=\omega_i/h\) によって表示 (2) の不変式 \(f(x)\) に対応するすべての \(w\) の線形形式を動かせる。Hilbert の整加群基底に関する定理は
\begin{equation}
        L=A_1(I)L_1+\cdots+A_\tau(I)L_\tau,
\end{equation}
を与える。ここで \(A_i\) は有理整数係数をもち、またすべての \(L\) は \(w\) に関して線形なので、\(A_i\) はもはや \(w\) を含まない。置換 \(\omega_i=w_i/h\) が \(L_1,\ldots,L_\tau\) を不変式 \(\mathfrak I_1(x),\ldots,\mathfrak I_\tau(x)\) に送るなら、(3) は
\begin{equation}
        f(x)=A_1(I)\mathfrak I_1(x)+\cdots+A_\tau(I)\mathfrak I_\tau(x)
\end{equation}
となる。この表示は
\[
        I_1(x),\ldots,I_\sigma(x);\mathfrak I_1(x),\ldots,\mathfrak I_\tau(x)
\]
が群 \(\frH\) のすべての整係数不変式の基底をなし、各不変式がこれら有限個の不変式の、有理整数係数をもつ多項式的有理関数として表されることを示している。

とくに、群 \(\frH\) が、各変換 \(A\) とともに代数的共役な数係数をもつすべての \(A',A'',\ldots\) も含むという性質をもつなら、\(h\) 個の行 \((x^{(k)})\) の整有理対称関数は、有理整数係数をもつ \(x\) の関数になる。これは、とくに対称関数の基底関数についても成り立つ。したがって、有理整数係数をもつ不変式 \(f(x)\) に制限するなら、基底関数
\[
        I_1(x),\ldots,I_\sigma(x);\mathfrak I_1(x),\ldots,\mathfrak I_\tau(x)
\]
も有理整数係数をもつ。この特殊な仮定は一般の群 \(\frH\) では満たされないが、その仮定の下では、整係数有限性は有理整数係数をもつ不変式の部分領域についても成り立つ。


\end{document}
